% Options for packages loaded elsewhere
\PassOptionsToPackage{unicode}{hyperref}
\PassOptionsToPackage{hyphens}{url}
\PassOptionsToPackage{dvipsnames,svgnames,x11names}{xcolor}
\documentclass[
  ngerman,
  10pt,
  a4paper,
]{article}
\usepackage{xcolor}
\usepackage[margin=22mm]{geometry}
\usepackage{amsmath,amssymb}
\setcounter{secnumdepth}{-\maxdimen} % remove section numbering
\usepackage{iftex}
\ifPDFTeX
  \usepackage[T1]{fontenc}
  \usepackage[utf8]{inputenc}
  \usepackage{textcomp} % provide euro and other symbols
\else % if luatex or xetex
  \usepackage{unicode-math} % this also loads fontspec
  \defaultfontfeatures{Scale=MatchLowercase}
  \defaultfontfeatures[\rmfamily]{Ligatures=TeX,Scale=1}
\fi
\usepackage{lmodern}
\ifPDFTeX\else
  % xetex/luatex font selection
  \setmainfont[Path=/Users/stefanhamann/.cache/codex-runtimes/codex-primary-runtime/dependencies/native/libreoffice-headless/libreoffice/LibreOfficeDev.app/Contents/Resources/fonts/truetype/,Extension=.ttf,BoldFont=DejaVuSans-Bold,ItalicFont=DejaVuSans-Oblique,BoldItalicFont=DejaVuSans-BoldOblique]{DejaVuSans}
  \setmonofont[Path=/Users/stefanhamann/.cache/codex-runtimes/codex-primary-runtime/dependencies/native/libreoffice-headless/libreoffice/LibreOfficeDev.app/Contents/Resources/fonts/truetype/,Extension=.ttf,BoldFont=DejaVuSansMono-Bold,ItalicFont=DejaVuSansMono-Oblique,BoldItalicFont=DejaVuSansMono-BoldOblique]{DejaVuSansMono}
\fi
% Use upquote if available, for straight quotes in verbatim environments
\IfFileExists{upquote.sty}{\usepackage{upquote}}{}
\IfFileExists{microtype.sty}{% use microtype if available
  \usepackage[]{microtype}
  \UseMicrotypeSet[protrusion]{basicmath} % disable protrusion for tt fonts
}{}
\makeatletter
\@ifundefined{KOMAClassName}{% if non-KOMA class
  \IfFileExists{parskip.sty}{%
    \usepackage{parskip}
  }{% else
    \setlength{\parindent}{0pt}
    \setlength{\parskip}{6pt plus 2pt minus 1pt}}
}{% if KOMA class
  \KOMAoptions{parskip=half}}
\makeatother
\usepackage{longtable,booktabs,array}
\newcounter{none} % for unnumbered tables
\usepackage{calc} % for calculating minipage widths
% Correct order of tables after \paragraph or \subparagraph
\usepackage{etoolbox}
\makeatletter
\patchcmd\longtable{\par}{\if@noskipsec\mbox{}\fi\par}{}{}
\makeatother
% Allow footnotes in longtable head/foot
\IfFileExists{footnotehyper.sty}{\usepackage{footnotehyper}}{\usepackage{footnote}}
\makesavenoteenv{longtable}
\ifLuaTeX
  \usepackage{luacolor}
  \usepackage[soul]{lua-ul}
\else
  \usepackage{soul}
\fi
\ifLuaTeX
\usepackage[bidi=basic,shorthands=off]{babel}
\else
\usepackage[bidi=default,shorthands=off]{babel}
\fi
\ifPDFTeX
\else
\babelfont{rm}[Path=/Users/stefanhamann/.cache/codex-runtimes/codex-primary-runtime/dependencies/native/libreoffice-headless/libreoffice/LibreOfficeDev.app/Contents/Resources/fonts/truetype/,Extension=.ttf,BoldFont=DejaVuSans-Bold,ItalicFont=DejaVuSans-Oblique,BoldItalicFont=DejaVuSans-BoldOblique]{DejaVuSans}
\fi
\ifLuaTeX
  \usepackage{selnolig} % disable illegal ligatures
\fi
\setlength{\emergencystretch}{3em} % prevent overfull lines
\providecommand{\tightlist}{%
  \setlength{\itemsep}{0pt}\setlength{\parskip}{0pt}}
\usepackage{fvextra}
\RecustomVerbatimEnvironment{verbatim}{Verbatim}{breaklines=true,fontsize=\footnotesize}
\usepackage{xurl}
\usepackage{seqsplit}
\usepackage{adjustbox}
\usepackage{ragged2e}
\usepackage{titlesec}
\titleformat{\section}{\large\bfseries\raggedright}{}{0em}{}
\titleformat{\subsection}{\normalsize\bfseries\raggedright}{}{0em}{}
\DeclareRobustCommand{\codeword}[1]{\texttt{\seqsplit{#1}}}
\AtBeginEnvironment{longtable}{\small\RaggedRight}
\usepackage{newunicodechar}
\newunicodechar{𝒞}{\ensuremath{\mathcal C}}
\newunicodechar{𝔊}{\ensuremath{\mathfrak G}}
\newunicodechar{𝔤}{\ensuremath{\mathfrak g}}
\newunicodechar{𝓗}{\ensuremath{\mathcal H}}
\newunicodechar{ℋ}{\ensuremath{\mathcal H}}
\newunicodechar{𝔄}{\ensuremath{\mathfrak A}}
\newunicodechar{𝟙}{\ensuremath{\mathbf 1}}
\newunicodechar{𝕀}{\ensuremath{I}}
\newunicodechar{𝕋}{\ensuremath{\mathbb T}}
\newunicodechar{𝟘}{\ensuremath{\mathbf 0}}
\newunicodechar{⊞}{\ensuremath{\boxplus}}
\newunicodechar{∘}{\ensuremath{\circ}}
\newunicodechar{⟦}{\ensuremath{[\![}}
\newunicodechar{⟧}{\ensuremath{]\!]}}
\setlength{\emergencystretch}{4em}
\setlength{\parindent}{0pt}
\setlength{\parskip}{0.45em}
\AtBeginDocument{\hypersetup{colorlinks=true,linkcolor=black,urlcolor=blue}\pdfstringdefDisableCommands{\def\codeword#1{#1}}\RaggedRight}
\usepackage{fancyhdr}
\pagestyle{fancy}
\fancyhf{}
\fancyhead[L]{\small TFPT / Universalraum - v1.6.7}
\fancyfoot[R]{\thepage}
\setlength{\headheight}{15pt}
\usepackage{bookmark}
\IfFileExists{xurl.sty}{\usepackage{xurl}}{} % add URL line breaks if available
\urlstyle{same}
\hypersetup{
  pdflang={de-DE},
  colorlinks=true,
  linkcolor={Maroon},
  filecolor={Maroon},
  citecolor={Blue},
  urlcolor={Blue},
  pdfcreator={LaTeX via pandoc}}

\author{}
\date{}

\begin{document}

{
\setcounter{tocdepth}{2}
\tableofcontents
}
\section{TFPT / Universalraum: die richtige einfache
Kette}\label{tfpt-universalraum-die-richtige-einfache-kette}

\subsection{Vollständige Konsolidierung v1.6.7 - 15. September
2026}\label{vollstuxe4ndige-konsolidierung-v1.6.7---15.-september-2026}

\subsubsection{Ergebnis und
Leseschlüssel}\label{ergebnis-und-leseschluxfcssel}

Die zuletzt ergänzte Vogelperspektive steht im Teil B dieser Fassung:
\textbf{Gesucht wird der kleinste gemeinsame beobachtbare Prozess hinter
den verschiedenen Ansichten.} Ein nativer Rang-60-Kodierer ist bereits
vorhanden. Seine Informationsverluste, die minimalen
symmetrieverträglichen Dynamiken und der Unterschied zwischen
Korrelation und kausalem Eingriff wurden neu berechnet. Das verschiebt
die fundamentale Priorität von weiteren lokalen Näherungswerten zur
Herkunft der Operationen und zur gemeinsamen Schattenkarte.

Die neue Runde schließt zwei kleine Strukturfragen vollständig: Der
symmetrische Zustandsraum bei zwei Bosonen hat exakt vier Dimensionen;
der bisher unbestimmte Rest von 4035 bilinearen Komponenten zerfällt
vollständig in drei Darstellungen. Die vier Zustandsrichtungen und ihre
Kopplungen wurden zusätzlich am ursprünglichen Tensor berechnet. Daraus
folgt ein vollständiger sechs-dimensionaler Ausschnitt.

Die gewünschte einfache Kette existiert ebenfalls, aber anders als
vorgeschlagen: Man muss den \textbf{ganzen Hamiltonoperator}
wiederholen, nicht Bosonzahl und Iterationsnummer gleichsetzen. Die
ersten fünf Lanczos-Glieder sind jetzt exakt bekannt. Sie bestimmen zehn
Energiemomente des gefüllten Startzustands und verschärfen die obere
Grundzustandsgrenze geringfügig.

\textbf{Das ist keine vollständige TOE.} Insbesondere fehlen die
physische Auswahl des Modellvertrags, ausführbare ursprüngliche
Instrumente, derselbe geladene Transport im gemeinsamen Träger und die
relativistische Kontinuumskonstruktion. Dass die bisherigen Grunddaten
dafür noch nicht eindeutig genügen, wird unten an einer einfachen
Familie verschiedener Dynamiken nachgewiesen.

Kennzeichnungen: \textbf{exakt} bezeichnet eine bewiesene Aussage im
genannten Modell; \textbf{bedingt} eine Aussage mit zusätzlich gewährter
Operation oder bekanntem früherem Satz; \textbf{numerisch} eine
Näherung; \textbf{offen} einen fehlenden Nachweis. Komplexe
Algebradimensionen zählen weder physikalische Raumrichtungen noch
unabhängig hergeleitete Teilchen oder ausführbare Operationen.

Die Fassung baut auf dem vollständigen Hauptdokument v1.6.6 auf. Die
aktuelle Gesamtdarstellung steht voran; die frühere Fassung
einschließlich ihrer Herleitungen, Gegenbeispiele und älteren Anhänge
bleibt vollständig erhalten. Alle sechs in dieser Runde eingesandten
Texte werden unverändert dokumentiert und hier mit überprüften
Ergänzungen und ausdrücklich benannten Korrekturen gelesen.

\subsection{1. Das Gesamtbild: Grammatik, Prozess, Zustand,
Antwort}\label{das-gesamtbild-grammatik-prozess-zustand-antwort}

TFPTs Compiler beschreibt zunächst eine markierte algebraische
Grammatik: welche Bausteine zusammenpassen, welche Symmetrien gelten und
welche Zahlen innerhalb dieser Konstruktion folgen. Der Universalraum
soll diese Grammatik als einen ausführbaren Prozess realisieren. Der
entscheidende Unterschied ist der zwischen einem erlaubten Zusammenhang
und einer tatsächlich erzeugten physikalischen Entwicklung.

Im hier untersuchten Modell ist die Grundregel weiterhin sehr einfach:
\textbf{Zwei Fermionen können in einen Vermittler umgewandelt werden und
zurück.} Mit 64 Fermionmoden und 60 Bosonmoden lautet sie

\[\adjustbox{max width=\linewidth}{$\displaystyle 
H=\Delta N_b+g(T_++T_-),\qquad
T_+=\sum_A b_A^\dagger P_A,\quad T_-=T_+^\dagger,
\qquad P_A=\sum_{i<j}W_{A,ij}f_jf_i.
$}\]

Es gelten \(\Delta>0\), reelles \(g\) und die erhaltene Ladung
\(N=N_f+2N_b\). Der Tensor W besitzt 480 Einträge mit Werten +1 oder -1,
acht disjunkte Paare pro Zeile, und \(WW^\dagger=8I_{60}\). Seine innere
Darstellung ist \((16,4)\) für Fermionen und \((10,6)\) für Bosonen
unter Spin(10) mal SU(4).

Die heutige Beweiskette sieht so aus:

{\def\LTcaptype{none} % do not increment counter
\begin{longtable}[]{@{}
  >{\raggedright\arraybackslash}p{(\linewidth - 4\tabcolsep) * \real{0.3333}}
  >{\raggedright\arraybackslash}p{(\linewidth - 4\tabcolsep) * \real{0.3333}}
  >{\raggedright\arraybackslash}p{(\linewidth - 4\tabcolsep) * \real{0.3333}}@{}}
\toprule\noalign{}
\begin{minipage}[b]{\linewidth}\raggedright
Verbindung
\end{minipage} & \begin{minipage}[b]{\linewidth}\raggedright
Was bereits trägt
\end{minipage} & \begin{minipage}[b]{\linewidth}\raggedright
Was dadurch nicht automatisch folgt
\end{minipage} \\
\midrule\noalign{}
\endhead
\bottomrule\noalign{}
\endlastfoot
Compiler zu Paarregel & Konkreter Tensor, Symmetrie und Ladungserhaltung
& Physische Wahl von Energie, Kopplung und Instrumenten \\
Paarregel zu Zustand & Früherer eindeutiger Grundzustandssatz im
festgelegten Vertrag & Herleitung dieses Vertrags oder dessen
Präparation \\
Zustand zu geladener Antwort & Isolierte niedrige Fermion-Entnahmelinie
mit hohem Gewicht & Räumlich laufendes relativistisches Teilchen \\
Antwort zu Bewegung & Zusätzliche Links erlauben nachgewiesenen Transfer
& Dass derselbe Link aus dem ursprünglichen Operationssatz entsteht \\
Bewegung zu Raumzeit & Prüfprogramm und bedingte Modelle & Gemeinsame
3+1D-Welt, chirales Maß und dynamische Gravitation \\
\end{longtable}
}

Der endliche dokumentierte Clock gehört zur inneren Spin(10)-Symmetrie.
Er ist damit kein zusätzliches äußeres Symmetrieobjekt. Seine Periode
sechs ist aber auch noch keine Herleitung einer physikalischen
Zeiteinheit oder einer Raumzeit. Innere Spin(10)-Labels sind nicht von
selbst Lorentzspin.

\subsection{2. Was die beiden neuen Texte
beitragen}\label{was-die-beiden-neuen-texte-beitragen}

\subsubsection{2.1 Quellenstand und
Reproduktionsgrenzen}\label{quellenstand-und-reproduktionsgrenzen}

17 ausgewählte Eingaben wurden mit SHA-256 eingefroren: beide
Nutzertexte, die zugehörigen Programme und Ergebnisdateien, der native
Tensor sowie die vollständige frühere Fassung mit Prüfpaket. Keine
fremde Quelldatei wurde verändert. Die Herkunftsdateien gehören zu
parallel fortgeschriebenen Arbeitsständen; Dateipins sind deshalb
wichtiger als Versionsnamen allein.

Der neue Grundzustandslauf war beim ersten Blick noch ohne
Ergebnisdatei. Beim Einfrieren lag bereits sein fertiges Ergebnis vor:
\textbf{382 Guards, davon 379 exakt und drei numerisch}. Sein
Quellenhash stimmt mit der eingefrorenen Programmdatei überein. Er
findet ebenfalls die Multiplizitäten 1, 1 und 4. Der Zwischenhinweis
dieser Prüfung auf eine noch leere Datei war bei der späteren Übernahme
bereits überholt und wird hier korrigiert.

Das Operations-JSON enthält \textbf{306}, nicht die im Begleittext
genannten 307 Guards; sein interner Checkerhash passt nicht zur
gleichzeitig eingefrorenen Programmdatei. Das ist ein Versionskonflikt
des Berichts, kein Beweis, dass seine Mathematik falsch ist. Die
tragenden Aussagen werden deshalb unabhängig geprüft, nicht als exakter
Gesamtreplay dieses fremden Programms ausgegeben. Das Feld-JSON meldet
2161 Guards und besitzt einen passenden Quellenhash; es enthält dennoch
den unten korrigierten Normfehler und zu weit formulierte
Interpretationen.

Mit den weiteren Texten und zugehörigen Auditdateien sind es insgesamt
24 eingefrorene Eingaben. Unsere eigene neue Suite hat \textbf{822
exakte Prüfbedingungen pro Python-Modus}. Normaler und optimierter Lauf
erzeugen identische Ergebnisdateien; Warnungen werden als Fehler
behandelt. Zusätzlich wurde die ältere ganzzahlige Kontraktionsrechnung
frisch kompiliert und wiederholt, mit identischen Zahlen. Die
vier-Boson-Norm aus dem früheren Zertifikat bleibt ein gekennzeichneter
Eingang: Ihre vollständige Enumeration wurde in dieser Runde nicht
wiederholt. Auch die gesamten fremden Programme mit ihren Guards wurden
nicht alle erneut ausgeführt. Prüfzahl und Beweisumfang werden
ausdrücklich getrennt.

\subsubsection{2.2
Übernahmeentscheidung}\label{uxfcbernahmeentscheidung}

{\def\LTcaptype{none} % do not increment counter
\begin{longtable}[]{@{}
  >{\raggedright\arraybackslash}p{(\linewidth - 2\tabcolsep) * \real{0.5000}}
  >{\raggedright\arraybackslash}p{(\linewidth - 2\tabcolsep) * \real{0.5000}}@{}}
\toprule\noalign{}
\begin{minipage}[b]{\linewidth}\raggedright
Aussage aus den Eingaben
\end{minipage} & \begin{minipage}[b]{\linewidth}\raggedright
Entscheidung für v1.6.7
\end{minipage} \\
\midrule\noalign{}
\endhead
\bottomrule\noalign{}
\endlastfoot
Korrigierter Anfangskommutant 3829536 / 1444233216 & Bestätigt; keine
neue Größe gegenüber v1.6.6 \\
Nur SU(4) oder nur Spin(10) als zusätzliche Kontrollen & Übernommen;
echte Zwischenstufen statt binärer Verfügbarkeit \\
Volle Gruppe plus alle Modenbesetzungen erzeugt volle Sektoralgebra &
Unabhängig bestätigt, ausdrücklich bedingt \\
Clock außerhalb der zusammenhängenden Gruppe, weil nicht skalar &
Verworfen; das explizite Spin-Wort aus v1.6.6 widerlegt den Schluss \\
Vier Singuletts auf der zweiten Bosonstufe & Unabhängig bestätigt und
durch konkrete Zustandsprojektionen ergänzt \\
Eine Zustandsrichtung pro Bosonzahl könnte bis Stufe 32 genügen &
Verworfen; schon Stufe zwei hat vier Richtungen \\
Feldnorm 24 bei quadrierter Norm 192 & Korrigiert: Norm ist 8 mal Wurzel
3 in derselben Konvention \\
4035 bilineare Komponenten unbestimmt & Jetzt vollständig in drei innere
Darstellungen zerlegt \\
Jeder relativistische Abschluss müsse den Tensorfeldkanal verwenden &
Nur unter dem festgelegten gleichhändigen, ableitungsfreien Ansatz
gültig \\
Schwächere Sektorschranken trennen N=64 nicht & Richtig für diese Sonde;
widerlegt den stärkeren früheren Satz nicht \\
\end{longtable}
}

Die Entropie von ungefähr 1,66 Bit im ersten Text gehört zum dortigen
Ritz-Näherungszustand. Die Blockdiagonalität der reduzierten Bosondichte
rechtfertigt die Shannon-Untergrenze für diesen Zustand, nicht ungeprüft
denselben Zahlenwert für den wahren nativen Grundzustand. Ebenso ist
eine Kreuzung mit dem leeren Zustand bei einem bestimmten chemischen
Potential kein Beweis für die vollständige Sektor-Stabilitätszone.

\subsection{3. Vollständig gelöst: der erste verzweigte
Singulettbereich}\label{vollstuxe4ndig-geluxf6st-der-erste-verzweigte-singulettbereich}

\subsubsection{3.1 Warum vier und nicht
eins?}\label{warum-vier-und-nicht-eins}

Im Ladungssektor N=64 schreiben wir Zustände relativ zur gefüllten
Fermionreferenz \(F\). Bei k Bosonen gibt es 2k Löcher. Der
Singulettbereich auf dieser Stufe liegt in

\[\adjustbox{max width=\linewidth}{$\displaystyle 
\mathcal S_k=
\left[\Lambda^{2k}(\overline{16\otimes4})
\otimes\operatorname{Sym}^k(10\otimes6)\right]^{G},
\qquad G=\operatorname{Spin}(10)\times SU(4).
$}\]

Die ersten Dimensionen sind \textbf{1, 1, 4}. Dies ist eine Aussage über
alle Singuletts dieser Stufen, nicht über vier verschiedene
Grundzustände. Ein eindeutiger Grundzustand kann durchaus in einem
größeren Singulettbereich liegen. Die Zentrumsregel N gleich 0 modulo 4
ist nur eine notwendige Bedingung für Singuletts, kein Auswahlprinzip
für N=64.

Der Beweis für k=2 benötigt keinen riesigen Matrixkern. Die Bosonseite
zerfällt nach der symmetrischen Cauchy-Identität in

\[\adjustbox{max width=\linewidth}{$\displaystyle 
\operatorname{Sym}^2(10\otimes6)=
(1,1)\oplus(54,1)\oplus(1,20')\oplus(54,20')\oplus(45,15).
$}\]

Die Dimensionen 1, 54, 20, 1080 und 675 ergeben zusammen 1830. Die
vierte äußere Potenz der Fermionseite wird durch die fünf Partitionen
von vier beschrieben. Für die benötigten Überschneidungen reichen drei
Spinor-Schurfunktoren:

{\def\LTcaptype{none} % do not increment counter
\begin{longtable}[]{@{}
  >{\raggedright\arraybackslash}p{(\linewidth - 4\tabcolsep) * \real{0.3000}}
  >{\raggedright\arraybackslash}p{(\linewidth - 4\tabcolsep) * \real{0.3000}}
  >{\raggedleft\arraybackslash}p{(\linewidth - 4\tabcolsep) * \real{0.4000}}@{}}
\toprule\noalign{}
\begin{minipage}[b]{\linewidth}\raggedright
Bosontyp
\end{minipage} & \begin{minipage}[b]{\linewidth}\raggedright
Spinor-Schurfunktor auf der Fermionseite
\end{minipage} & \begin{minipage}[b]{\linewidth}\raggedleft
Anzahl gemeinsamer Typen
\end{minipage} \\
\midrule\noalign{}
\endhead
\bottomrule\noalign{}
\endlastfoot
(1,1) & Sym hoch 4 von 16 & 0 \\
(54,1) & Sym hoch 4 von 16 & 1 \\
(1,20') & S mit Partition (2,2) von 16 & 1 \\
(54,20') & S mit Partition (2,2) von 16 & 1 \\
(45,15) & S mit Partition (3,1) von 16 & 1 \\
\end{longtable}
}

Die Charaktere wurden aus den fünf Konjugationsklassen von S4 mit
ganzzahliger Arithmetik aufgebaut, durch eine unabhängige
Jacobi-Trudi-Rechnung kontrolliert und mit der Weyl-Alternierung
ausgewertet. Die volle äußere Potenz hat 635376 Dimensionen; der
gemeinsame Singulettbereich hier nur vier. Die Berechnung des anderen
Workers über den vollständigen Gewicht-Null-Raum kommt unabhängig auf
denselben Wert.

\subsubsection{3.2 Nicht nur gezählt: die vier Zustände und ihre
Kopplungen}\label{nicht-nur-gezuxe4hlt-die-vier-zustuxe4nde-und-ihre-kopplungen}

In der durch eine Bosonparität äquivalenten Lochdarstellung sei
\(v_k=(T_+)^kF\). Alle 108240 von null verschiedenen Koeffizienten von
\(v_2\) wurden neu berechnet, einschließlich der Boson-Fakultätsnormen.
Das ergibt \(\|v_2\|^2=439680\).

Sei \(P_R\) der orthogonale Projektor auf den jeweiligen inneren Typ der
Zwei-Boson-Seite. Dann sind die vier Vektoren \(P_Rv_2\) nicht null,
orthogonal und bilden wegen der bewiesenen Multiplizität eins die
gesamte Singulettbasis auf Stufe zwei. Ihre Normen und die Kopplungen
vom normierten Ein-Boson-Zustand lauten:

{\def\LTcaptype{none} % do not increment counter
\begin{longtable}[]{@{}
  >{\raggedright\arraybackslash}p{(\linewidth - 4\tabcolsep) * \real{0.2727}}
  >{\raggedleft\arraybackslash}p{(\linewidth - 4\tabcolsep) * \real{0.3636}}
  >{\raggedleft\arraybackslash}p{(\linewidth - 4\tabcolsep) * \real{0.3636}}@{}}
\toprule\noalign{}
\begin{minipage}[b]{\linewidth}\raggedright
Typ R
\end{minipage} & \begin{minipage}[b]{\linewidth}\raggedleft
Normquadrat von P\_R v2
\end{minipage} & \begin{minipage}[b]{\linewidth}\raggedleft
Quadrierte Kopplung geteilt durch g hoch 2
\end{minipage} \\
\midrule\noalign{}
\endhead
\bottomrule\noalign{}
\endlastfoot
(54,1) & 17280 & 36 \\
(1,20') & 7680 & 16 \\
(54,20') & 241920 & 504 \\
(45,15) & 172800 & 360 \\
Summe & 439680 & 916 \\
\end{longtable}
}

Diese Zahlen sind abgeleitet, nicht angepasst. Die Projektoren benutzen
nur Symmetrisierung sowie die beiden vorhandenen invarianten
Spurbildungen. Die gemeinsame skalare Projektion (1,1) verschwindet
identisch.

Nach Wahl der vier Basisphasen ist der Kopplungsvektor

\[\adjustbox{max width=\linewidth}{$\displaystyle 
g\,(6,\,4,\,6\sqrt{14},\,6\sqrt{10}).
$}\]

Alle vier Richtungen sind in v2 enthalten, aber von Stufe eins wird nur
ihre eine feste Linearkombination angeregt. Drei dazu orthogonale
Kombinationen sind \textbf{nach unten} dunkel. Sie sind nicht deshalb
auch gegenüber den höheren Stufen entkoppelt.

\subsubsection{3.3 Der vollständig bestimmte
Sechszustands-Ausschnitt}\label{der-vollstuxe4ndig-bestimmte-sechszustands-ausschnitt}

Die Kompression auf alle Singuletts mit k kleiner oder gleich zwei ist
exakt äquivalent zu

\[\adjustbox{max width=\linewidth}{$\displaystyle 
P_{\leq2}HP_{\leq2}\simeq
\begin{pmatrix}
0&g\sqrt{480}&0\\
g\sqrt{480}&\Delta&g\sqrt{916}\\
0&g\sqrt{916}&2\Delta
\end{pmatrix}
\oplus 2\Delta I_3.
$}\]

Damit ist dieser Ausschnitt vollständig gelöst. \textbf{Er ist nicht
invariant:} H führt von k=2 auch nach k=3. Seine Eigenwerte sind
folglich nicht automatisch Eigenwerte der ganzen Bank. Eine räumliche
Interpretation oder zusätzliche Wechselwirkung wurde für diese Reduktion
nicht eingeführt.

Explizit besitzt die Kompression dreimal den Wert 2 Delta und die drei
reellen Nullstellen des Polynoms

\[\adjustbox{max width=\linewidth}{$\displaystyle 
E^3-3\Delta E^2+(2\Delta^2-1396g^2)E+960\Delta g^2.
$}\]

\subsection{4. Die richtige einfache Kette: Wiederholung von
H}\label{die-richtige-einfache-kette-wiederholung-von-h}

\subsubsection{4.1 Der Fehler war die Gleichsetzung zweier verschiedener
Stufen}\label{der-fehler-war-die-gleichsetzung-zweier-verschiedener-stufen}

Die Vektoren \(v_k=T_+^kF\) sind nach Bosonzahl geordnet. Schon bekannt
ist

\[\adjustbox{max width=\linewidth}{$\displaystyle 
T_-v_3=\frac{299520}{229}v_2+w_2,\qquad
\langle v_2,w_2\rangle=0,\qquad
\|w_2\|^2=\frac{5001523200}{229}>0.
$}\]

Diese Rechnung wurde frisch reproduziert. Aus einer kleinen relativen
Norm dieses Seitenzweigs folgt keine Konvergenz der gesamten
Grundzustands- oder Spektralrechnung. Insbesondere fehlt dann immer noch
der große Übergang nach v4.

Eine Lanczos-Kette entsteht stattdessen aus \(F,HF,H^2F,\ldots\) durch
Orthogonalisierung. Für einen selbstadjungierten endlichen
Hamiltonoperator beschreibt sie dessen zyklischen Teilraum exakt. Das
ist ein zulässiges einfaches Rechenbild, aber keine Behauptung, dass
ihre Indizes physikalische Orte oder Bosonzahlen seien. Auch die gesamte
Bank muss nicht mit diesem einen zyklischen Teilraum übereinstimmen.

\subsubsection{4.2 Die ersten fünf Glieder sind jetzt exakt
bestimmt}\label{die-ersten-fuxfcnf-glieder-sind-jetzt-exakt-bestimmt}

Für die normierte Lanczos-Basis \(u_n\) schreiben wir

\[\adjustbox{max width=\linewidth}{$\displaystyle 
Hu_n=b_nu_{n-1}+a_nu_n+b_{n+1}u_{n+1},\qquad b_n>0.
$}\]

Die ersten Diagonalen lauten

\[\adjustbox{max width=\linewidth}{$\displaystyle 
\frac{(a_0,a_1,a_2,a_3,a_4)}{\Delta}
=\left(0,1,2,3,\frac{105168998}{26292551}\right).
$}\]

Die ersten quadrierten Nebendiagonalen sind

\[\adjustbox{max width=\linewidth}{$\displaystyle 
\frac{(b_1^2,b_2^2,b_3^2,b_4^2)}{g^2}
=\left(480,916,\frac{299520}{229},\frac{78877653}{47632}\right).
$}\]

Bis u3 stimmen die Vektoren, bis auf die Phase bei negativem g, mit den
normierten vk überein. Der nächste Vektor ist proportional zu
\(v_4+w_2\). Dabei liegen v4 und w2 auf verschiedenen Bosonstufen.
Deshalb gilt

\[\adjustbox{max width=\linewidth}{$\displaystyle 
\frac{b_4^2}{g^2}
=\frac{\|v_4\|^2+\|w_2\|^2}{\|v_3\|^2},
\qquad
\langle u_4,N_bu_4\rangle
=\frac{4\|v_4\|^2+2\|w_2\|^2}{\|v_4\|^2+\|w_2\|^2}.
$}\]

Die Bosonzahlvarianz dieses Vektors ist exakt

\[\adjustbox{max width=\linewidth}{$\displaystyle 
\operatorname{Var}_{u_4}(N_b)
=\frac{63416178576}{691298238087601}>0.
$}\]

Die frühere Quotientenfortsetzung mit nur \(\|v_4\|^2/\|v_3\|^2\) ließ
den positiven Zusatz \(1809/47632\) weg. Die korrigierte Kette bleibt
einfach, verliert aber gerade deshalb nicht die wirkliche
Mehrkanalstruktur. Ihre Länge ist nicht durch die 33 möglichen
Bosonzahlen festgelegt.

\subsubsection{4.3 Zehn vollständige Energiemomente der gefüllten
Referenz}\label{zehn-vollstuxe4ndige-energiemomente-der-gefuxfcllten-referenz}

Die Fünfer-Kompression bestimmt \(\mu_n=\langle F,H^nF\rangle\) für n
von 0 bis 9 exakt. Ein Weg, der den nächsten nicht enthaltenen
Lanczos-Vektor erreicht und nach F zurückkehrt, benötigt mindestens zehn
Schritte. Aus der Tridiagonalität folgt deshalb die angegebene
Momentengenauigkeit, ohne die übrige Kette zu erfinden.

Beispielsweise gilt

\[\adjustbox{max width=\linewidth}{$\displaystyle 
\mu_0=1,\quad \mu_1=0,\quad
\mu_2=480g^2,\quad \mu_3=480\Delta g^2,
\quad \mu_4=480\Delta^2g^2+670080g^4.
$}\]

Vollständig lauten die Koeffizienten in
\(\mu_n/\Delta^n=c_2x^2+c_4x^4+c_6x^6+c_8x^8\), mit \(x=g/\Delta\), für
n von 2 bis 9:

{\def\LTcaptype{none} % do not increment counter
\begin{longtable}[]{@{}lrrrr@{}}
\toprule\noalign{}
n & c2 & c4 & c6 & c8 \\
\midrule\noalign{}
\endhead
\bottomrule\noalign{}
\endlastfoot
2 & 480 & 0 & 0 & 0 \\
3 & 480 & 0 & 0 & 0 \\
4 & 480 & 670080 & 0 & 0 \\
5 & 480 & 2219520 & 0 & 0 \\
6 & 480 & 5527680 & 1510510080 & 0 \\
7 & 480 & 12353280 & 10437173760 & 0 \\
8 & 480 & 26213760 & 48253363200 & 4615972423680 \\
9 & 480 & 54144000 & 187581404160 & 54295325184000 \\
\end{longtable}
}

Diese Momente gehören \textbf{F}, nicht dem nativen Grundzustand Omega.
Sie sind auch nicht mit den geladenen Antwortmomenten aus der früheren
Fassung zu verwechseln. Sie machen eine endliche Modellrechnung genauer,
beweisen aber keine physische Präparation von F und keine
Vollständigkeit des Feldwörterbuchs.

\subsubsection{4.4 Kleine, rigorose Verschärfung der
Energiegrenze}\label{kleine-rigorose-verschuxe4rfung-der-energiegrenze}

Am Prüfpunkt g geteilt durch Delta gleich 1/20 hat die niedrigste
Eigenenergie der Fünfer-Kompression die durch rationale Sturm-Zählung
zertifizierte Lage

\[\adjustbox{max width=\linewidth}{$\displaystyle 
-1.12963813<\frac{E_{\mathrm{Ritz},5}}{\Delta}<-1.12963811.
$}\]

Rayleigh-Ritz liefert für den wahren Grundzustand daher

\[\adjustbox{max width=\linewidth}{$\displaystyle 
E_0<-1.12963811\Delta.
$}\]

Die frühere untere Schranke bleibt \(E_0>-1.158089\Delta\). Zusammen mit
dem früheren strengen N=63-Floor \(E_h>-1.121899\Delta\) folgt nun

\[\adjustbox{max width=\linewidth}{$\displaystyle 
\epsilon=E_h-E_0>0.00773911\Delta.
$}\]

Die Verbesserung gegenüber 0.007737 ist klein. Die obere Polgrenze
0.039079764 und das Gewicht über 88.007628 Prozent werden dadurch nicht
ungeprüft verändert. Auch die niedrigste Ritz-Energie ist kein
berechneter Zentralwert von E0.

\subsection{5. Vollständig gelöst: die innere bilineare
Zerlegung}\label{vollstuxe4ndig-geluxf6st-die-innere-bilineare-zerlegung}

Für den 64-dimensionalen Fermionträger gilt

\[\adjustbox{max width=\linewidth}{$\displaystyle 
\operatorname{End}(16\otimes4)
=(1,1)\oplus(45,1)\oplus(210,1)
\oplus(1,15)\oplus(45,15)\oplus(210,15).
$}\]

Die sechs Dimensionen sind 1, 45, 210, 15, 675 und 3150. Der offene Rest
ist also nicht mehr unbestimmt:

\[\adjustbox{max width=\linewidth}{$\displaystyle 
4035=210+675+3150.
$}\]

Eine kurze konkrete Konstruktion genügt. Auf dem chiralen 16er-Raum
liefern Cliffordprodukte der Grade 0, 2 und 4 genau 1, 45 und 210
unabhängige Matrizen. Ihr ganzzahliger Hilbert-Schmidt-Gram ist
\(16I_{256}\). Unter Spin(10)-Konjugation bleiben diese Gradräume
invariant. Auf der Farbseite zerfällt End(4) in die skalare Matrix und
die 15 spurlosen Matrizen. Die Tensorprodukte ergeben die gesamte
Zerlegung. Die Weyl-Multiplizitäten wurden zusätzlich überprüft.

Die verwendete Spinorzerlegung ist etablierte Darstellungstheorie, keine
neue Vorhersage von 210 physikalischen Teilchen. Der Literaturabgleich
bestätigt die drei Spin(10)-Kanäle 1, 45 und 210; unsere Rechnung ordnet
sie der tatsächlichen 16-mal-4-Bank zu. Siehe
\href{https://arxiv.org/abs/hep-th/0109116}{Nath und Syed, vollständige
Spinorkopplungen}.

Aus dieser Zerlegung folgt weder, dass alle Bilineare verfügbare
Kontrollen sind, noch, dass jede Darstellung ein propagierendes Feld
ist. Produkte einfacher Matrizen im Einteilchenraum dürfen insbesondere
nicht automatisch als dieselben Produkte ihrer
Vielteilchen-Hamiltonoperatoren gelesen werden.

\subsection{6. Welche Operationen wirklich verfügbar sein
müssten}\label{welche-operationen-wirklich-verfuxfcgbar-sein-muxfcssten}

\subsubsection{6.1 Die algebraische Frage ist präziser
beantwortet}\label{die-algebraische-frage-ist-pruxe4ziser-beantwortet}

Im N=3-Sektor bleiben für verschiedene \textbf{angenommene}
Operationsalphabete folgende komplexe Kommutantdimensionen:

{\def\LTcaptype{none} % do not increment counter
\begin{longtable}[]{@{}lr@{}}
\toprule\noalign{}
Angenommenes Alphabet & Kommutantdimension \\
\midrule\noalign{}
\endhead
\bottomrule\noalign{}
\endlastfoot
X und Nb & 1444233216 \\
Zusätzlich dokumentierter Clock & 240742144 \\
Stattdessen zusätzlich SU(4)-Generatoren & 2247168 \\
Stattdessen zusätzlich Spin(10)-Generatoren & 1648 \\
Zusätzlich beide vollständigen Gruppen & 7 \\
Volle Gruppen und alle Modenbesetzungen & 1 \\
\end{longtable}
}

Die physische Verfügbarkeit ist somit nicht auf eine binäre Alternative
zusammengeschrumpft. Es gibt geprüfte Zwischenstufen, darunter auch die
in v1.6.6 behandelten einzelnen Casimir-Auslesungen.

Die volle Algebra nach Hinzunahme aller Besetzungen wurde unabhängig mit
einem kleineren Beweis kontrolliert. Der Einteilchen-Supportgraph der
gewährten Gruppengeneratoren ist verbunden, ebenso der Bosongraph. Ein
verbundener Graph hat verbundene Besetzungsgraphen mit k
ununterscheidbaren Fermionen, solange 0 kleiner k kleiner 64. Der native
Tensor verbindet die Fermion- und Bosonseiten. Die Besetzungsoperatoren
trennen alle Basiszustände. Polynomiale Spektralprojektoren und die
nichtverschwindenden Verbindungseinträge erzeugen daher alle
Matrixeinheiten.

So folgt bedingt die volle Algebra in N=2 und N=3. Gewährt man
zusätzlich geladene f-Instrumente, verbinden sich auch die beiden
Sektoren. Das ist eine vollständige \textbf{assoziativ-algebraische}
Aussage, kein Nachweis beliebiger unitärer Steuerbarkeit, effizienter
Messung oder nativer Instrumentenherkunft.

\subsubsection{6.2 Warum bloßes Lesen noch keine fehlende Bewegung
erzeugt}\label{warum-blouxdfes-lesen-noch-keine-fehlende-bewegung-erzeugt}

Ohne die zusätzlich gewährten Gruppenkontrollen erhalten X, Nb, der
Clock und alle Modenbesetzungen weiterhin mindestens eine nichtskalare
Farb-Cartanladung. Jede Komposition dieser Operatoren behält diese
Erhaltung. Die alleinige Erlaubnis, mehr Moden zu unterscheiden, hebt
das Hindernis also nicht auf.

Das ist die präzise Grenze zwischen \textbf{Auslesen} und
\textbf{Eingreifen}. Ein mathematisch definierter Projektor ist nicht
schon ein physisch ausführbares Messinstrument. Dass das Modell unter
einer Transformation symmetrisch ist, bedeutet nicht, dass es diese
Transformation als kontrollierten Eingriff erzeugt. Die bedingten
Zwei-Banken- und Vierzustands-Transferbeispiele aus v1.6.5 und v1.6.6
bleiben nützlich, aber ihre zusätzlichen Links bleiben auch zusätzliche
Voraussetzungen.

\subsection{7. Feldwörterbuch: gesicherte Teile und korrigierte
Reichweite}\label{feldwuxf6rterbuch-gesicherte-teile-und-korrigierte-reichweite}

Für die unveränderte antisymmetrische Paarmatrix und gleichhändige
Weylfelder verschwindet die ableitungsfreie skalare Kontraktion. Die
drei symmetrischen Spinortensoren tragen einen nichtverschwindenden
Kanal des Typs (1,0), mit konjugiertem Typ (0,1). Das ist am gesamten
ursprünglichen Tensor neu geprüft.

In der Konvention des gelieferten Programms hat jeder der drei Kanäle
pro Zeile Normquadrat 64. Die Summe ist 192 und die direkte Summennorm
folglich \(8\sqrt3\), \textbf{nicht 24}. Die Summe von drei einzelnen
Normen darf nicht mit der Norm ihrer orthogonalen direkten Summe
verwechselt werden. Eine anders normierte Vertexdefinition ändert diese
Zahlen, nicht den Nullkanal.

Der Satz gilt unter den genannten Feld- und Ableitungsannahmen.
Zusätzliche unabhängige Komponenten oder Ableitungskopplungen sind
andere Modelle und nicht generell ausgeschlossen. Keine dieser
Möglichkeiten ist dadurch bereits aus der ursprünglichen Bank
hergeleitet.

Auch die Zerlegung eines Vermittler-mal-Weyl-Komposits in die
Lorentztypen (3/2,0) und (1/2,0) ist eine Darstellungsaussage. Der
Ranganteil 2/6 gleich 1/3 ist das Gewicht im maximal gemischten Zustand
dieses Sechserraums. Für allgemeine Zustände kann derselbe Projektor
Gewichte von null bis eins haben. Er liefert nicht automatisch ein
Drittel des nativen Polgewichts. Die frühere Ladungsunterscheidung
bleibt bestehen: Die EOM-Größe D trägt Ladung minus eins, das ältere
erzeugende Komposit chi-dagger plus drei.

\subsection{8. Die fundamentale Frage: was fehlt
wirklich?}\label{die-fundamentale-frage-was-fehlt-wirklich}

\subsubsection{8.1 Ein überprüfbarer
Unterbestimmtheitssatz}\label{ein-uxfcberpruxfcfbarer-unterbestimmtheitssatz}

Betrachte bei festem W dieselbe Familie \(H_{\Delta,g}=\Delta N_b+gX\).
Alle diese Hamiltonoperatoren besitzen die gleiche innere
Quellsymmetrie, dieselbe Ladungserhaltung und denselben kommutierenden
endlichen Clock. Dennoch unterscheiden sie sich dynamisch. Schon für
dieselbe gefüllte Referenz F ist die dimensionslose Momentenkombination

\[\adjustbox{max width=\linewidth}{$\displaystyle 
\frac{\mu_4\mu_2}{\mu_3^2}
=1+1396\left(\frac{g}{\Delta}\right)^2.
$}\]

Bei g geteilt durch Delta gleich 1/20 ist sie 4,49, bei 1/40 dagegen
1,8725. Beide Parameterpunkte liegen im früher abgesicherten
Grundzustandsintervall. Weil die Größe dimensionslos ist, handelt es
sich nicht lediglich um eine Änderung der Maßeinheit. Die gleiche
Grammatik und Symmetrie lässt also unterschiedliche Antworten auf
dieselbe Präparation zu.

\textbf{Damit ist die eindeutige physische Dynamik aus diesen Grunddaten
allein nicht ableitbar.} Das widerlegt nicht die Möglichkeit einer
einfachen zusätzlichen Ursprungsregel. Es identifiziert genau deren
Aufgabe: Sie muss den Prozess samt verfügbarem Operationssatz und
Zustand auswählen, statt nur seine Symmetrie zu benennen. Im fest
gesetzten Prüfvertrag sind g und Delta natürlich definiert; offen ist
ihre physische Herkunft.

\subsubsection{8.2 Der kleinste tragfähige
Forschungsansatz}\label{der-kleinste-tragfuxe4hige-forschungsansatz}

Für eine einzelne Antwortrechnung sollte der relevante Raum nicht von
vornherein der gesamte Fockraum sein. Der zyklische Raum der
tatsächlichen Operationen auf dem angegebenen Startzustand ist der
kleinere Arbeitsraum. Die H-Lanczos-Kette ist eine exakte Darstellung
eines solchen Raums. Mehrere Präparationen, Messungen oder
Compileroperationen können jedoch einen größeren gemeinsamen Prozessraum
verlangen. Ihre Kompatibilität muss dann nachgewiesen werden; getrennte
kleine Modelle dürfen nicht einfach zusammengefügt werden. Eine
vollständige Rekonstruktion des Grundzustands aus dieser Referenz
verlangt außerdem einen Nachweis ihrer nichtverschwindenden Überlappung
mit Omega; die Ritz-Obergrenze allein ersetzt ihn nicht.

Die nächste physikalische Schließung bleibt dieselbe konkrete Bedingung:
Auf einer gemeinsam hergeleiteten Quelle und demselben Zustand muss eine
verfügbare Operation eine geladene Anregung zwischen operational
bestimmten Teilen übertragen. Der Test muss die Mehrzeitantwort von
bloßem Überlapp unterscheiden und ein konsistentes relativistisches
Feldwörterbuch verwenden. Koordinatenwechsel, Gram-Überlapp oder ein
innerer Clock allein erfüllen diese Bedingung nicht.

\subsection{9. Aktueller Stand T1 bis
T8}\label{aktueller-stand-t1-bis-t8}

Die folgende Tabelle ordnet die Beiträge dieser Runde den bisherigen
offenen Beweispflichten zu; sie benennt keine neue Definition der
ursprünglichen Tore.

{\def\LTcaptype{none} % do not increment counter
\begin{longtable}[]{@{}
  >{\raggedright\arraybackslash}p{(\linewidth - 4\tabcolsep) * \real{0.3333}}
  >{\raggedright\arraybackslash}p{(\linewidth - 4\tabcolsep) * \real{0.3333}}
  >{\raggedright\arraybackslash}p{(\linewidth - 4\tabcolsep) * \real{0.3333}}@{}}
\toprule\noalign{}
\begin{minipage}[b]{\linewidth}\raggedright
Tor
\end{minipage} & \begin{minipage}[b]{\linewidth}\raggedright
Beitrag dieser Runde
\end{minipage} & \begin{minipage}[b]{\linewidth}\raggedright
Weiter fehlender entscheidender Nachweis
\end{minipage} \\
\midrule\noalign{}
\endhead
\bottomrule\noalign{}
\endlastfoot
T1 & Exakte bedingte Kontrollalgebra; Auslese- und Eingriffslücke
präzisiert & Tatsächlich erzeugbare Compileroperationen und
Instrumente \\
T2 & Bessere lokale Struktur und Antwortgrenze in der endlichen Bank &
Quellseitiges Half-Charge-Feld mit Energie- und Adjungiertenkontrolle
sowie E8-/Clock-Zuordnung \\
T3 & Kleiner gemeinsamer Rechenraum innerhalb einer Bank & Gemeinsamer
physischer 3+1D-Ursprung \\
T4 & Vollständige innere Bilineare; korrekte Feldkanäle & Vollständiges
chirales Maß mit dynamischen Eichfreiheitsgraden \\
T5 & Exakter lokaler Hamilton-Lanczos-Anfang & Kontrollierter
relativistischer Kontinuums- und Streuungsgrenzfall \\
T6 & Expliziter Nachweis, dass die bisherigen Grunddaten g/Delta nicht
auswählen & Physische Kopplungs-, Skalen- und Spektralzuordnung,
einschließlich offener Neutrinofragen \\
T7 & Keine neue Schließung & Aus derselben Quelle erzeugter dynamischer
Spin-2-Sektor und universelle Kopplung \\
T8 & Unterscheidung Startzustand, Grundzustand und Ausleseinstrument
verschärft & Gemeinsame physische Präparation, Aufzeichnung und
Zustandsauswahl \\
\end{longtable}
}

Alle acht Tore bleiben offen. Die heutigen Ergebnisse schließen
begrenzte mathematische Fragen innerhalb des Programms, nicht diese
physikalischen Gesamtnachweise. Frühere RH-, Faktorisierungs- oder
P-versus-NP-Grenzen werden durch die endlichen Rechnungen ebenfalls
nicht verändert. Der begrenzte RH-Registerabgleich und seine
Aktualitätsgrenze sind in Teil B dokumentiert; ein vollständiger
Neuaudit dieser Arbeitsfronten fand nicht statt.

\subsection{10. Die drei nächsten entscheidenden
Untersuchungen}\label{die-drei-nuxe4chsten-entscheidenden-untersuchungen}

\begin{enumerate}
\def\labelenumi{\arabic{enumi}.}
\tightlist
\item
  \textbf{Den minimalen Ursprungsvertrag auswählen.} Teil B
  klassifiziert die kleinen symmetrieverträglichen Terme und trennt
  echte Erweiterungen von einem bloßen Wechsel der Bosonvariablen. Der
  Compiler muss diese Auswahl und die relativen Parameter begründen.
  Symmetrie allein tut das noch nicht.
\item
  \textbf{Die gemeinsame Schattenkarte am vorhandenen Tensor prüfen.}
  Derselbe Zustand und dieselben Operationen müssen die verschiedenen
  Ausleseansichten tragen. Gesucht sind ergänzende statt nur duplizierte
  Informationen, ein kontrollierter gemeinsamer Kern und Vorhersagen
  ohne nachträgliche Anpassung.
\item
  \textbf{Einen echten Eingriffstransfer und denselben Feldadapter
  konstruieren.} Eine aus der Quelle abgeleitete lokale Operation muss
  später die unbedingte Statistik eines anderen operational bestimmten
  Teils verändern. Selbst zeitabhängige Kreuzkorrelation bei Anfangswert
  null reicht nicht; Teil B enthält ein exaktes Gegenbeispiel. Dazu
  müssen Ladung, Kinetik und relativistische Feldzuordnung passen, bevor
  räumliche Skalierung trägt.
\end{enumerate}

Die korrekte H-Kette mit Restschranke bleibt eine nützliche parallele
Rechenaufgabe. Die zehn bekannten Momente sind feste Gegenprüfungen. Sie
ersetzt aber nicht die jetzt vorrangige Ursprungs- und
Kompositionsfrage.

\subsection{11. Reproduktion und
Veröffentlichung}\label{reproduktion-und-veruxf6ffentlichung}

Das Prüfpaket enthält sechs neue Python-Prüfer, die eingefrorenen
Quellen, die alte und neu wiederholte Kontraktion, Ergebnisdateien,
Quellenmanifest und die vollständige v1.6.6 als historische Basis.
\codeword{replay.py} prüft alle Pins, beide Python-Modi und die
C++-Kontraktion. Die ausgelieferte ZIP wird nach dem Entpacken erneut
geprüft.

Die vollständige Hauptfassung, das kurze Änderungsdokument und die
einfache Erklärung werden gemeinsam versioniert direkt in Documents
abgelegt, jeweils als Markdown und PDF. Die PDF-Ausgaben werden
zusätzlich auf Darstellung, Textabdeckung und Seitenaufbau geprüft.
Frühere Fassungen bleiben erhalten. Es erfolgt keine Änderung der
zentralen TOE-Abnahmemarker, kein Commit, kein Push und keine externe
Web-Veröffentlichung.

Der Statusleitfaden beeinflusst die Trennung zwischen mathematischer
Struktur, angenommenen Ressourcen und nachgewiesener Fähigkeit. Der
Debugging-Leitfaden führte zur ausdrücklichen Sicherung
reell-ganzzahliger Tensorwerte vor der Typumwandlung; keine
Imaginärteile werden still verworfen.

\subsubsection{Schluss in einem Satz}\label{schluss-in-einem-satz}

\textbf{Die einfache Regel ist vorhanden; ihr erster verzweigter
Zustandsbereich und die richtige kleine Rechenkette sind jetzt genauer
gelöst. Was noch fehlt, ist die aus dem Ursprung folgende Auswahl und
Ausführung dieses Prozesses als dieselbe physikalische Welt.}

\begin{center}\rule{0.5\linewidth}{0.5pt}\end{center}

\section{Neue Vogelperspektive: derselbe Prozess in verschiedenen
Ansichten}\label{neue-vogelperspektive-derselbe-prozess-in-verschiedenen-ansichten}

\subsection{B1. Was sich durch die drei weiteren Texte
ändert}\label{b1.-was-sich-durch-die-drei-weiteren-texte-uxe4ndert}

Die nachgereichten Texte werden als Vorschläge geprüft, nicht als
bewiesene Ergebnisse oder neue Arbeitsanweisungen übernommen. Die
stärkste gemeinsame Idee trägt: \textbf{Innere Identität,
Zusammensetzung, Ort, Zeit und Auslese dürfen nicht stillschweigend
dasselbe Label erhalten.} Eine immer genauere Rechnung innerhalb einer
Bank ersetzt diese Verbindungen nicht. Die voranstehenden lokalen
Ergebnisse bleiben gültig; die fundamentale Priorität verschiebt sich
auf die gemeinsame Operations- und Kompositionsregel.

{\def\LTcaptype{none} % do not increment counter
\begin{longtable}[]{@{}
  >{\raggedright\arraybackslash}p{(\linewidth - 2\tabcolsep) * \real{0.5000}}
  >{\raggedright\arraybackslash}p{(\linewidth - 2\tabcolsep) * \real{0.5000}}@{}}
\toprule\noalign{}
\begin{minipage}[b]{\linewidth}\raggedright
Vorschlag
\end{minipage} & \begin{minipage}[b]{\linewidth}\raggedright
Prüfung und Konsequenz
\end{minipage} \\
\midrule\noalign{}
\endhead
\bottomrule\noalign{}
\endlastfoot
Ort in Multiplizitäten statt in den 64 inneren Labels suchen & Sinnvolle
Möglichkeit, aber Multiplizität beschreibt bereits chemische Umwandlung
ohne Raum \\
Der große Kommutant könnte verborgenen Raum enthalten & Hamilton- und
Symmetriekommutant sind verschieden; die Größe allein trägt diesen
Schluss nicht \\
Den ganzen Z4-Ursprung statt nur den ladungserhaltenden Ausschnitt
betrachten & Konkrete zusätzliche Terme existieren am selben W; ihre
Verfügbarkeit ist noch nicht hergeleitet \\
Ladung +3 und Ladung -1 sind modulo vier gleich & Ein zusätzlicher Term
kann sie dynamisch verbinden; die Operatoren sind nicht identisch \\
Zeitabhängige Kreuzantwort zeigt Transport & Im Allgemeinen falsch;
unten steht ein exaktes Gegenbeispiel sogar bei Anfangskorrelation
null \\
Mehrere Schatten könnten einen gemeinsamen Ursprung bestimmen & Möglich
bei gemeinsamem Vertrag und ausreichender Unterscheidbarkeit;
Übereinstimmung allein genügt nicht \\
Holografie könnte die fehlende Verbindung sein & Ein nativer
isometrischer Kodierungsbaustein ist vorhanden; Raumgrenze, Schutz und
gemeinsame Raumzeit fehlen \\
\end{longtable}
}

Die zusätzlich eingefrorenen Texte heißen im Prüfpaket
\codeword{multiplicity\_z4.txt}, \codeword{fundamental\_reduction.txt}
und \codeword{runtime\_synthesis.txt}. Sie werden im Hauptdokument
vollständig angehängt. Die anschließende Schatten- und Holografiefrage
des Nutzers wurde durch eigene Rechnungen weiterverfolgt.

\subsection{B2. Noethers Frage: Welche Unterschiede kann die Dynamik
überhaupt
verändern?}\label{b2.-noethers-frage-welche-unterschiede-kann-die-dynamik-uxfcberhaupt-veruxe4ndern}

Bei ungebrochener innerer Gruppe G zerfällt ein Zustandsraum als

\[\adjustbox{max width=\linewidth}{$\displaystyle 
\mathcal H=\bigoplus_\lambda V_\lambda\otimes M_\lambda,
\qquad H=\bigoplus_\lambda I_{V_\lambda}\otimes h_\lambda,
\qquad \operatorname{End}_G(\mathcal H)
=\bigoplus_\lambda I_{V_\lambda}\otimes\operatorname{End}(M_\lambda).
$}\]

G-invariante Entwicklung darf die Kopien eines Typs verändern, nicht die
inneren Komponenten dieses Typs willkürlich auseinanderziehen. Das macht
Multiplizitäten zu einem möglichen Träger weiterer Dynamik. Es macht sie
noch nicht zu Orten. Schon die drei hellen N=3-Typen enthalten zwei
Kopien: einmal drei Fermionen, einmal einen Fermion plus einen Boson.
Ihre jeweilige Hamiltonmatrix lautet

\[\adjustbox{max width=\linewidth}{$\displaystyle 
h_\lambda=\begin{pmatrix}0&g\sqrt\lambda\\g\sqrt\lambda&\Delta\end{pmatrix},
\qquad \lambda=7,10,12.
$}\]

Am Prüfpunkt g/Delta gleich 1/20 sind die maximalen
Umwandlungswahrscheinlichkeiten 7/107, 1/11 und 3/28. Das ist echte
symmetrieverträgliche Dynamik, zunächst aber eine Veränderung der
Zusammensetzung und kein Ortswechsel.

Wichtig ist die Richtung des Kommutanten: Im N=3-Sektor hat der
Kommutant der vollen inneren Gruppe die Dimension 16. Die viel größere
Zahl 1444233216 betrifft den kleineren Kontrollsatz aus X und Bosonzahl.
Auf dem isolierten irreduziblen 64er-Pol ist der G-Kommutant sogar nur
skalar, während der Kommutant der dort skalaren Energie alle
64-mal-64-Matrizen enthält. Die große Energiedegeneration ist kein
Nachweis vieler Orte.

Für einen G-Singulett-Grundzustand und die ursprünglichen irreduziblen
Fermionoperatoren folgt aus Schurs Lemma

\[\adjustbox{max width=\linewidth}{$\displaystyle 
C_{rs}(t)=\delta_{rs}\,c(t).
$}\]

Feste lineare Umbenennungen liefern lediglich einen Gramfaktor mal
derselben Zeitfunktion. Auch eine perfekte Kenntnis aller lokalen
Spektrallinien würde aus den 64 inneren Komponenten deshalb noch keine
räumliche Anordnung machen. Ein stationärer Grundzustand kann
zeitabhängige Korrelationsfunktionen haben; Nichtstationarität ist keine
notwendige Voraussetzung dafür. Physische Uhrenablesung und Zeitrichtung
sind weitere Fragen.

\subsection{B3. Korrelation ist noch kein Eingriff: ein exakt gelöstes
Gegenbeispiel}\label{b3.-korrelation-ist-noch-kein-eingriff-ein-exakt-geluxf6stes-gegenbeispiel}

Zwei nicht miteinander wechselwirkende Qubits genügen:

\[\adjustbox{max width=\linewidth}{$\displaystyle 
H=\tfrac12(Z_A+Z_B),\qquad
|\Omega\rangle=(|01\rangle+|10\rangle)/\sqrt2,
\qquad A=X_A,\quad B=Y_B.
$}\]

Der Zustand ist stationär. Trotzdem ist

\[\adjustbox{max width=\linewidth}{$\displaystyle 
\langle\Omega|B e^{-itH}A|\Omega\rangle=\sin t.
$}\]

Die Kreuzkorrelation startet bei null und verändert sich mit der Zeit.
Es gibt dennoch keinerlei Signal von A nach B: Alle A-lokalen
Operationen kommutieren mit jedem zeitentwickelten B-lokalen
Observablen. Für jedes spurtreue lokale Quanteninstrument bleibt daher
die unbedingte B-Statistik unverändert. Eine Konditionierung auf ein
Messergebnis in A wäre etwas anderes.

Der belastbare Übertragungstest ist stattdessen

\[\adjustbox{max width=\linewidth}{$\displaystyle 
\delta_{A\to B}(t)=
\operatorname{tr}\!\left[B U_t\mathcal E_A(\rho)U_t^\dagger\right]
-\operatorname{tr}\!\left[B U_t\rho U_t^\dagger\right].
$}\]

A und B müssen zuerst als unabhängig adressierbare Teile definiert sein;
das Instrument muss verfügbar sein, ohne durch seine Definition schon B
anzusteuern. Erfolg bedeutet eine nichtverschwindende unbedingte
Änderung nach dem A-Eingriff. Die lineare Antwort misst entsprechend
einen retardierten Kommutator, nicht bloß eine Zweipunktkorrelation.
Lieb-Robinson-Abschätzungen können anschließend die Ausbreitung in einem
bereits lokal gekoppelten Modell begrenzen; sie erzeugen diese Lokalität
oder drei Raumdimensionen nicht von selbst. Siehe die Primärarbeit von
\href{https://arxiv.org/abs/quant-ph/0603121}{Bravyi, Hastings und
Verstraete}.

Ein weiterer exakter Befund verhindert ein zu starkes Paritäts-No-go:
Der native Graph der 480 Fermionpaare verbindet alle 64 Moden. Eine
Vorzeichenumkehr auf einer Fermionteilmenge, bei unveränderten Bosonen,
erhält jede Kante genau dann, wenn sie trivial oder die globale
Fermionparität ist. Die binäre Inzidenzmatrix hat Rang 63. Zusätzliche
unabhängige Bankparitäten entstehen beim Kopieren von Banken; sie sind
nicht als solche ein Hindernis innerhalb der einen ursprünglichen Bank.
Werden auch Bosonen umgezeichnet, sind andere gemeinsame Gradierungen
möglich.

\subsection{B4. Die einfachste fehlende Auswahl: Warum genau dieser
Hamiltonoperator?}\label{b4.-die-einfachste-fehlende-auswahl-warum-genau-dieser-hamiltonoperator}

Die Z4-Idee besitzt einen präzisen Kern. Ein ganzzahliger, auf jedem
Grad konstanter Ladungslift mit neutralem Grad null müsste bei den
entsprechenden nichtverschwindenden Klammern erfüllen

\[\adjustbox{max width=\linewidth}{$\displaystyle 
2q_1=q_2,\qquad q_1+q_3=0,\qquad 2q_3=q_2.
$}\]

Die Koeffizientenmatrix hat Determinante 4 und Rang drei. Über den
ganzen Zahlen bleibt nur q1=q2=q3=0; modulo vier gilt dagegen die
gewohnte Belegung 1, 2, 3. Das verbietet nicht andere kontinuierliche
Cartanladungen und beweist nicht, dass der aktuelle Operator-Hamiltonian
schon alle E8-Klammern realisiert.

Eine wörtliche Identifikation der Grad-eins-Labels mit primitiven CAR-
Vernichtern scheitert bereits an einem kleinen Beispiel: Die native
W-Spalte des Paares (0,1) ist null, aber \([f_0,f_1]=2f_0f_1\ne0\). Die
innere Lieklammer und die gewöhnliche Operatoralgebra der Fermionfelder
dürfen nicht gleichgesetzt werden. Eine andere, zusammengesetzte
E8-Realisierung wird dadurch nicht ausgeschlossen.

\subsubsection{B4.1 Zwei zusätzliche Terme am unveränderten
Tensor}\label{b4.1-zwei-zusuxe4tzliche-terme-am-unveruxe4nderten-tensor}

Der Bosonträger (10,6) besitzt eine symmetrische reelle invariante
Paarung eta. In der nativen Gewichtsbasis koppelt sie entgegengesetzte
Spin- Vektorgewichte und komplementäre Farbpaare. Es gilt eta Quadrat
gleich eins. Damit sind bereits

\[\adjustbox{max width=\linewidth}{$\displaystyle 
B_+=\tfrac12\sum_{AB}b_A^\dagger\eta_{AB}b_B^\dagger,
\qquad
R_+=\sum_{AB}b_A^\dagger\eta_{AB}P_B^\dagger
$}\]

G-invariant. Bplus erzeugt zwei Bosonen, Rplus einen Boson und zwei
Fermionen. Beide ändern N um vier; beide erhalten N modulo vier. Die
vollständige Lie-Invarianz wurde an allen 45 plus 15 Generatoren
überprüft, nicht nur an den Cartanladungen. Rplus wirkt auf dem leeren
Zustand nicht trivial: Das Normquadrat beträgt exakt 480. Der
quadratische Bosonpaarterm ist sogar einfacher als der vorgeschlagene
zusätzliche kubische Kanal.

Auf dem unveränderten Träger lautet die allgemeine G-invariante,
fermionparitätsgerade, normalgeordnete Hamiltonfamilie vom Grad
höchstens drei

\[\adjustbox{max width=\linewidth}{$\displaystyle 
H_{\rm ext}=c+\varepsilon N_f+\Delta N_b
+\kappa B_++\bar\kappa B_-
+g\sum_A b_A^\dagger P_A+\bar g\sum_A P_A^\dagger b_A
+\lambda R_++\bar\lambda R_-.
$}\]

Die Vollständigkeit gilt nur innerhalb dieser ausdrücklichen
Polynomialklasse: End(F) und End(B) haben jeweils einen Skalar, Sym hoch
zwei von B einen Skalar, und B tritt in Lambda hoch zwei von F genau
einmal auf. Konkret zerfällt letzterer Raum in (10,6), (126,6) und
(120,10). Die SU(4)-Zentrumsregel schließt die verbleibenden
bosonisch-kubischen und gemischten Besetzungs-Boson-Terme aus. Höhere
Polynome, räumliche Kopien und neue Feldträger sind damit nicht
klassifiziert.

Für das native Modell hat das lineare System aller diagonalen
Modenphasen exakt neun unabhängige kontinuierliche Erhaltungen: die acht
inneren Cartans und N. Bei zusätzlichem nichtverschwindendem Rplus-Kanal
oder Bosonpaar-Kanal bleiben exakt acht. Die Ränge sind 115
beziehungsweise 116 auf 124 Moden; explizite ganzzahlige Kerne und
modulare Rangzertifikate beweisen sie. Dies ist keine Behauptung, die
volle nichtabelsche Gruppe schrumpfe auf acht Dimensionen. Sie bleibt
erhalten. Außerdem ist N modulo vier bereits die Wirkung des
SU(4)-Zentrums und keine neu hinzugewonnene unabhängige Symmetrie.

Für reelle Kopplungen und epsilon gleich null enthält die
Fermiongleichung nun zwei verschiedene Beiträge:

\[\adjustbox{max width=\linewidth}{$\displaystyle 
[H_{\rm ext},f_r]=-gD_r-\lambda\widetilde\chi_r^\dagger,
\qquad D_r\sim bf^\dagger,\quad
\widetilde\chi_r^\dagger\sim b^\dagger\eta f^\dagger.
$}\]

Sie haben N-Ladungen -1 und +3. Eine Dynamik kann diese Kanäle bei
bloßer Z4-Erhaltung verbinden, ohne die Operatoren oder ihre
Referenzzustände fälschlich gleichzusetzen. \textbf{Zulässig ist aber
noch nicht aus dem Compiler abgeleitet.} Die Grundzustands- und
Polschranken des unveränderten Modells gelten für diese Erweiterung
nicht ungeprüft weiter.

Auch ein mu-N-Term bleibt G- und Z4-invariant. Bereits der
Zwei-Zustandsblock aus leerem Zustand und einer normierten
Quartetterzeugung im kleinen Modell zeigt

\[\adjustbox{max width=\linewidth}{$\displaystyle 
H+\mu N=\begin{pmatrix}0&\lambda\\\lambda&\Delta+4\mu\end{pmatrix}.
$}\]

Seine niedrigere Eigenenergie hängt von mu ab. Der Übergang zu Z4 allein
wählt also weder ein chemisches Potential noch einen eindeutigen
physikalischen Grundzustand aus. Für Delta größer als Betrag kappa ist
die bosonische Quadratik positiv; die linearen Bosonkopplungen an
beschränkte Fermionoperatoren lassen sich relativ dazu abschätzen. Das
gibt eine untere Energieschranke, aber noch keinen neuen eindeutigen
Grundzustandssatz.

\subsubsection{B4.2 Eine weitere Vereinfachung: neue Dynamik oder andere
Teilchenvariablen?}\label{b4.2-eine-weitere-vereinfachung-neue-dynamik-oder-andere-teilchenvariablen}

Manche scheinbaren Erweiterungen beschreiben nur andere Bosonvariablen.
Auf der reellen, phasengleichen Teilfamilie setze man

\[\adjustbox{max width=\linewidth}{$\displaystyle 
b=c\cosh r+\eta c^\dagger\sinh r.
$}\]

Dies ist eine G-verträgliche kanonische Bogoliubov-Transformation. Sie
ergibt

\[\adjustbox{max width=\linewidth}{$\displaystyle 
\Delta'=\Delta\cosh2r+\kappa\sinh2r,\quad
\kappa'=\Delta\sinh2r+\kappa\cosh2r,
$}\] \[\adjustbox{max width=\linewidth}{$\displaystyle 
g'=g\cosh r+\lambda\sinh r,\quad
\lambda'=g\sinh r+\lambda\cosh r.
$}\]

Die additive Konstante ist 30 mal (Delta prime minus Delta). Bei g
Quadrat größer lambda Quadrat verschwinden kappa prime und lambda prime
gemeinsam genau dann, wenn

\[\adjustbox{max width=\linewidth}{$\displaystyle 
\mathcal I=\kappa(g^2+\lambda^2)-2\Delta g\lambda=0.
$}\]

Beispiel: Delta=1, kappa=4/5, g=2, lambda=1 und tanh r=-1/2 ergeben
Delta prime=3/5, g prime=Wurzel 3, kappa prime=lambda prime=0. Obwohl
die alte Zahl N nicht erhalten ist, gibt es dann die verborgene
kontinuierliche Erhaltung Nf plus zweimal Nc. Daher wäre auch die
Folgerung „N ist gebrochen, also ist das Fundament nur Z4`` ohne Prüfung
zu schnell. Für I ungleich null ist nur diese uniforme kanonische
Reduktion ausgeschlossen, nicht jede denkbare nichtlineare verborgene
Symmetrie.

Der sehr einfache reelle Quadraturansatz lambda=g bei kappa=0 ist
dagegen nicht durch eine endliche solche Transformation auf den nativen
Ansatz reduzierbar. Er ist ein prüfbarer Kandidat, kein hergeleitetes
Naturgesetz. Ein weiterer Herkunftstest ist besonders schlicht: Bei
freiem H0=epsilon Nf+Delta Nb haben die beiden kubischen Terme die
Frequenzen Delta minus 2 epsilon und Delta plus 2 epsilon. Am
verwendeten epsilon=0 gibt es keine schnelle gegen langsame Frequenz,
die das Weglassen des zweiten Kanals durch eine Rotating-Wave-Näherung
rechtfertigen würde. Ein exaktes U(1)-Prinzip könnte das dennoch
rechtfertigen; dieses Prinzip wäre dann auszuweisen und abzuleiten.

\subsection{B5. Holografie und Schatten: die konkrete native
Verbindung}\label{b5.-holografie-und-schatten-die-konkrete-native-verbindung}

\subsubsection{B5.1 Ein Kodierer ist bereits
vorhanden}\label{b5.1-ein-kodierer-ist-bereits-vorhanden}

Der vorhandene, frisch geprüfte Tensor erfüllt WW dagger gleich 8 I60.
Daher ist

\[\adjustbox{max width=\linewidth}{$\displaystyle 
V=W^\dagger/\sqrt8:\ \mathbb C^{60}\longrightarrow\Lambda^2\mathbb C^{64},
\qquad V^\dagger V=I_{60},\quad
\Pi=VV^\dagger=W^\dagger W/8
$}\]

eine exakte Isometrie mit einem Rang-60-Projektor Pi. Ein
60-dimensionaler logischer Zustand kann verlustfrei als kohärente
Überlagerung von Fermionpaaren dargestellt werden. Die andere Richtung
sieht aber nicht den ganzen 2016-dimensionalen Paarraum: Ihr Kern hat
Dimension 1956. Diese „dunklen`` Richtungen sind nicht deshalb generell
unphysikalisch; sie sind für genau diese Abbildung unsichtbar.

Die beiden Ansichten sind durch W also bereits verbunden. Die native
N=2-Dynamik im hellen Paar-plus-Boson-Bereich ist sogar exakt

\[\adjustbox{max width=\linewidth}{$\displaystyle 
H_{N=2,\mathrm{hell}}=
\begin{pmatrix}0&\sqrt8 g\\\sqrt8 g&\Delta\end{pmatrix}\otimes I_{60}.
$}\]

Sie wandelt die Kodierungsform um; am Prüfpunkt erreicht die
entsprechende Übergangswahrscheinlichkeit höchstens 2/27. Der nur in die
Paarseite eingebettete Code ist bei g ungleich null kein invarianter
H-Unterraum. Dieser kleine gemeinsame dynamische Träger ist ein
positiver Befund. Er ist N=2, nicht der native N=64-Grundzustand und
nicht dessen geladener Pol. Insbesondere ist er noch kein Transport
durch physikalischen Raum.

\subsubsection{B5.2 Warum das noch kein holografischer
Fehlerkorrekturcode
ist}\label{b5.2-warum-das-noch-kein-holografischer-fehlerkorrekturcode-ist}

Für jeden der 64 Fermionmoden wurde die logische Kompression seiner
Besetzungsmessung berechnet:

\[\adjustbox{max width=\linewidth}{$\displaystyle 
V^\dagger n_rV\quad\hbox{hat Spektrum}\quad
0\ (45\text{-fach}),\quad 1/8\ (15\text{-fach}).
$}\]

Das Resultat ist nicht skalar. Die Umgebung kann daher schon durch
Ablesen einer einzigen verlorenen Mode etwas über den logischen Zustand
erfahren. Die notwendige Fehlerkorrekturbedingung für die vollständige
60er-Codemenge ist verletzt: Der Code korrigiert nicht die beliebige
Löschung dieser einen Mode. Fermionparität ändert dieses
Besetzungsargument nicht, denn n\_r ist gerade. Kleinere geeignete
Untercodes oder ein aus nativen Operationen erzeugtes größeres
Kodierungsnetz sind dadurch nicht ausgeschlossen, aber noch nicht
konstruiert. Bloßes Aneinanderhängen desselben Tensors garantiert keinen
Code.

In etablierten holografischen Spielzeugmodellen sind isometrische
Kodierung, Rekonstruktion logischer Operationen aus verschiedenen
Randteilen und Fehlerkorrektur genau kontrollierte Eigenschaften. Das
ist ein hilfreiches Prüfschema, keine direkte Identifikation mit TFPT.
Siehe \href{https://arxiv.org/abs/1503.06237}{Pastawski, Yoshida, Harlow
und Preskill}. Die Rekonstruktion aus einem Teilrand ist auch in der
AdS/CFT-Arbeit von \href{https://arxiv.org/abs/1601.05416}{Dong, Harlow
und Wall} an konkrete Quanteninformations- und geometrische
Voraussetzungen gebunden. Diese Voraussetzungen wurden für TFPT nicht
nachgewiesen.

\subsubsection{B5.3 Was mehrere Schatten gemeinsam eindeutig machen
können}\label{b5.3-was-mehrere-schatten-gemeinsam-eindeutig-machen-kuxf6nnen}

Seien Ri festgelegte Mess- oder Reduktionsabbildungen eines gemeinsamen
Zustands rho. Für eine endliche, nicht eingeschränkte Zustandsklasse
sind alle Zustände aus den Schatten Ri(rho) genau dann unterscheidbar,
wenn auf den hermiteschen spurlosen Differenzen gilt

\[\adjustbox{max width=\linewidth}{$\displaystyle 
\bigcap_i\ker R_i=\{0\}.
$}\]

Das ist das einfache gemeinsame-Kern-Kriterium. Anschaulich: Was eine
Ansicht nicht sieht, muss eine andere sehen. Zwei Qubitansichten auf
(x,z) und (y,z) genügen gemeinsam zur Bloch-Rekonstruktion, jede für
sich nicht. Zweimal dieselbe Ansicht auf (x,z) genügt weiterhin nicht.
Das gilt ebenso für zwei Theorieberichte, deren Übereinstimmung bereits
aus denselben eingesetzten Annahmen stammt: Sie liefern keine zweite
unabhängige Messung.

Drei exakt nachgerechnete Grenzen sind wesentlich:

\begin{enumerate}
\def\labelenumi{\arabic{enumi}.}
\tightlist
\item
  Die orthogonalen Zustände (000 plus 111)/Wurzel 2 und (000 minus
  111)/Wurzel 2 haben dieselben reduzierten Zustände auf sämtlichen
  echten Teilmengen der drei Qubits. Eine gemeinsame phasensensitive
  XXX-Messung unterscheidet sie mit Erwartungswert +1 oder -1. Auch sehr
  viele lokale Schatten können also eine globale Phase übersehen.
\item
  Paarweise passende Überschneidungen garantieren nicht einmal einen
  gemeinsamen Ursprung. Drei binäre Variablen können nicht paarweise
  ausnahmslos verschieden sein. Jede einzelne perfekte Antikorrelation
  besitzt aber dieselben gleichverteilten Einzelmarginalen. Gemeint ist
  hier ein gemeinsames klassisches Tripel, nicht eine Behauptung gegen
  kontextuelle Quantenexperimente.
\item
  Selbst die vollständige Antwort einer Referenz bestimmt keinen völlig
  entkoppelten dunklen Zusatzsektor. H und H direkt plus K besitzen vom
  Startvektor (F,0) aus dieselben Energiemomente und dieselbe
  Zeitantwort. Ein exaktes kleines Matrixbeispiel prüft diesen
  Sachverhalt.
\end{enumerate}

Das erreichbare Ziel ist daher zunächst der \textbf{kleinste gemeinsame
beobachtbare Prozess}, nicht ein aus endlichen Schatten bewiesener
einzigartiger ontologischer Innenraum. Zwei Historien gelten operational
als gleich, wenn alle verfügbaren künftigen Eingriffs- und Auslesefolgen
dieselben Wahrscheinlichkeiten liefern. Das ist eine eindeutige
Unterscheidungsregel für gegebenes Verhalten; sie garantiert weder eine
eindeutige Hilbertraumdarstellung noch von selbst eine einfache
Geometrie.

\subsubsection{B5.4 Die drei Wände müssen unterschiedlich geprüft
werden}\label{b5.4-die-drei-wuxe4nde-muxfcssen-unterschiedlich-gepruxfcft-werden}

{\def\LTcaptype{none} % do not increment counter
\begin{longtable}[]{@{}
  >{\raggedright\arraybackslash}p{(\linewidth - 4\tabcolsep) * \real{0.3333}}
  >{\raggedright\arraybackslash}p{(\linewidth - 4\tabcolsep) * \real{0.3333}}
  >{\raggedright\arraybackslash}p{(\linewidth - 4\tabcolsep) * \real{0.3333}}@{}}
\toprule\noalign{}
\begin{minipage}[b]{\linewidth}\raggedright
Gemeinte Wand
\end{minipage} & \begin{minipage}[b]{\linewidth}\raggedright
Was sie gegenwärtig bedeutet
\end{minipage} & \begin{minipage}[b]{\linewidth}\raggedright
Welcher Anschluss fehlt
\end{minipage} \\
\midrule\noalign{}
\endhead
\bottomrule\noalign{}
\endlastfoot
Universalraum & Grenze zwischen gemeinsamem Operatorobjekt und seinen
reduzierten Ansichten & Ein gemeinsamer verfügbarer Operationssatz,
Zustand und Rekonstruktionsvertrag \\
TFPT & Grenze zwischen interner Compilerstruktur und physikalischer
Realisierung & Herleitung der Ausführung, relativer Parameter und
operationaler Teilung \\
Beobachtete Realität & Endlicher Zugang über Messungen, Präparationen
und Korrelationen & Ein quantitatives, nicht nachträglich angepasstes
Feld- und Messwörterbuch \\
\end{longtable}
}

Diese Erkenntnisgrenzen sind nicht schon drei bewiesene physikalische
Holografieschirme. Ein kosmologischer Horizont, ein Code-Rand und eine
offene mathematische Beweispflicht sind verschiedene Dinge. Ebenso sind
Universalraum und TFPT Kandidatenbeschreibungen der Realität, nicht
bereits drei experimentell etablierte wechselwirkende Welten.

Eine stärkere holografische Deutung müsste eine einzige Kodierung mit
Zuständen, Operationen und Zeitentwicklung verbinden. Auf einem
invarianten Codesektor wäre etwa V dagger V=I und Hrand V=V Hin eine
passende Intertwining-Bedingung. Bei nichtinvarianten reduzierten
Ansichten muss die entstehende Erinnerung mitgeführt werden; autonome
reduzierte Dynamik darf nicht einfach vorausgesetzt werden. Zusätzlich
braucht es eine Bedeutung von Randteilen, kontrollierte Rekonstruktion
und eine skalierende Geometrie.

\subsection{B6. Die einfachste neue Arbeitsrichtung und ihre
Abnahmekriterien}\label{b6.-die-einfachste-neue-arbeitsrichtung-und-ihre-abnahmekriterien}

Der plausible übersehene Punkt ist kein weiterer großer Tensor:
\textbf{Wir brauchen die gemeinsame Regel, die festlegt, welche
Operationen, Teilchenvariablen und Ansichten tatsächlich dieselbe
Ausführung beschreiben.} Noethers Symmetrieprüfung sagt, was erhalten
bleibt. Die Schattenprüfung sagt, was überhaupt unterscheidbar ist. Der
Eingriffstest sagt, was etwas anderes beeinflussen kann. Zusammen sind
diese drei kleinen Fragen strenger und informativer als weitere
isolierte Zahlenübereinstimmungen.

\begin{enumerate}
\def\labelenumi{\arabic{enumi}.}
\tightlist
\item
  \textbf{Den minimalen Ursprungsvertrag auswählen.} Den tatsächlichen
  Compiler auf den nativen kubischen Term, den zweiten kubischen Kanal
  und den quadratischen Paarterm prüfen. Die vollständige kleine Familie
  und der Bogoliubov-Test verhindern, dass derselbe Prozess mehrfach
  gezählt wird. Erfolg ist eine Ableitung oder ein präziser Ausschluss
  samt relativen Parametern, nicht nur die Feststellung erlaubter
  Symmetrie.
\item
  \textbf{Die Schattenkarte am vorhandenen W aufbauen.} Für denselben
  Vertrag Präparationen, Aufzeichnungen und kontrollierte
  Mehrzeitantworten in den verschiedenen Ansichten auflisten. Gemeinsam
  unsichtbare Richtungen und unverträgliche Überlappungen berechnen. Aus
  einem Teil der Daten eine Aussage bestimmen, die eine andere Ansicht
  ohne Nachjustierung testet. Der jetzt bewiesene Rang-60-Code und sein
  Einmoden-Leck sind Ausgangsdaten, kein Anlass, Raum oder perfekte
  Fehlerkorrektur bereits anzunehmen.
\item
  \textbf{Daraus einen gemeinsamen kausalen Teilträger gewinnen.} Eine
  verfügbare Operation muss die unbedingte Statistik eines anderen
  operational bestimmten Teils verändern. Derselbe Träger muss das
  relativistische Feldwörterbuch tragen. Erst dann sind Skalierung,
  chirales Maß und ein Spin-2-Sektor belastbare nächste Rechnungen.
\end{enumerate}

Die lokale Lanczos-Fortsetzung bleibt als kontrolliertes
Diagnoseinstrument nützlich, steht aber nicht mehr an erster Stelle der
fundamentalen Suche. Die drei Schritte ersetzen keine einzelnen
T1-T8-Nachweise. Auch eine erfolgreiche endliche Kodierung würde RH,
effiziente Faktorisierung oder P-versus-NP nicht automatisch lösen.

Der RH-Absatz der Eingabe enthält keine neue vollständige arithmetische
Beweiskonstruktion. Das vorhandene Forschungsregister wurde begrenzt
abgeglichen; sein Aktualitätslauf scheitert derzeit an einer nicht
verfügbaren historischen Quelle beziehungsweise Quellen-/Review-Drift.
Bestehende endliche Determinantsonden sind ausdrücklich Diagnosemodelle.
Es wurde kein neuer RH-Beweis, kein umfassender RH-Neulauf und keine
Erneuerung der entsprechenden Vertrauensmarker behauptet.

\textbf{Bilanz:} Ein gemeinsamer Kodierungsbaustein ist schon da;
außerdem sind die minimalen symmetrieverträglichen Dynamiken und zwei
wichtige Verwechslungen jetzt wesentlich schärfer bestimmt. Die
vollständige universelle physikalische Lösung bleibt offen. Der nächste
Schritt ist kleiner und konkreter geworden: dieselbe Ausführung durch
mehrere kalibrierte, dynamisch konsistente Ansichten rekonstruieren und
testen.

\subsection{B7. Abgleich des zuletzt eingesandten Berichts
v1.6.3}\label{b7.-abgleich-des-zuletzt-eingesandten-berichts-v1.6.3}

Der sechste Text dieser Runde, „Operationssatz, nativer Grundzustand und
Feldwörterbuch``, enthält überwiegend bereits berücksichtigte Befunde:
die Kommutantenleiter, das Lochbild, die Normen der ersten Stufen, den
nichtverschwindenden Seitenzweig und den verschwindenden
Weyl-Skalarkanal. Seine zweite vorgeschlagene Folgeaufgabe wurde
inzwischen wesentlich weitergeführt: Alle vier Zwei-Boson-Singuletts
sind konstruiert und der richtige H-Lanczos-Anfang ist fünf Glieder weit
bestimmt.

Vier Präzisierungen verhindern eine Übernahme veralteter Schlüsse:

\begin{enumerate}
\def\labelenumi{\arabic{enumi}.}
\tightlist
\item
  \textbf{Sieben und sechzehn sind verschiedene Kommutanten.} Sieben
  gilt für volle innere Gruppe zusammen mit X und Nb; sechzehn für die
  Gruppe allein. Die drei hellen inneren Typen kommen auf beiden Seiten
  der Umwandlung vor. Einzelne Tabellenzeilen sind multiplizitätsfrei,
  der ganze N=3-Sektor aber nicht. Der Schluss vom nichtskalaren Clock
  auf eine äußere, nicht zusammenhängende Symmetrie ist weiterhin
  falsch.
\item
  \textbf{Eine kleine Zweignorm beweist keine globale Konvergenz.} Der
  genannte Defekt 2,90 mal 10 hoch -5 betrifft eine bestimmte
  Kontraktion. Die höhere Bosonstufe ist weiterhin gekoppelt. „Numerisch
  fast allein durch die Normen lösbar`` und „vernachlässigbar`` sind
  ohne Restschranke zu stark. Auch Zh plus Zadd gleich eins ist eine
  Summenregel, keine Bestimmung sämtlicher Spektrallinien.
  Ritz-Erwartungswerte bleiben Erwartungswerte des Ritz-Zustands, nicht
  automatisch von Omega.
\item
  \textbf{Der Feldschluss ist ansatzabhängig.} Der angegebene Fünfzyklus
  (16,45,0,30,37) wurde frisch am W überprüft. Damit ist eine einzige
  binäre Chiralitätsbelegung der primitiven Labels mit entgegengesetzter
  Händigkeit auf jeder Kante unmöglich. Das verbietet kein allgemeines
  Dirac-Wörterbuch mit erweiterten Trägern, Ableitungen oder
  zusammengesetzten Feldern. Der Tensorfeldkanal ist unter dem
  angegebenen engen Ansatz sinnvoll, nicht schon der universell einzige
  physikalische Feldtyp.
\item
  \textbf{Der gespeicherte Rechnerstand ist älter als der Text.} Die
  lokal zugehörige Datei \codeword{native\_ground\_state.json} hat einen
  anderen Checkerhash als das aktuelle Programm. Ihr gespeichertes
  w2-Normquadrat ist um 229 Quadrat gegenüber der physikalischen Norm
  skaliert. Das aktuelle Programm enthält bereits die richtige
  Entskalierung; die alte Datei ist also kein aktueller Replaynachweis.
  Sie meldet außerdem \codeword{norms\_by\_traces\_used=false} und
  Ritz-Tiefe drei. Die im Text erwähnte Datei
  \codeword{norms\_by\_traces.json} und das dortige Replaymanifest waren
  bei dieser Prüfung nicht vorhanden. Diese Beobachtung schließt ein
  Resultat an anderem Ort nicht aus, liefert hier aber keines.
\end{enumerate}

Daher wird insbesondere die Behauptung, nu4 sei in dieser Quelle auf
zwei unabhängigen Wegen vollständig reproduziert, nicht übernommen. Der
Text selbst bezeichnet nu4 zugleich als „nur Spurnetzwerk`` und
beschreibt den Brute-Force-Abgleich nur bis nu3. Der frühere gepinnte
exakte nu4-Eingang unserer eigenen Rechnung bleibt davon getrennt und
wurde nicht als neue Enumeration ausgegeben. Ein fertig berechnetes nu5
wird aus dem vorhandenen Programm mit \codeword{MAX\_N=5} nicht
abgeleitet; eine Einstellung ist kein Ergebnis.

\textbf{Die interessante verbleibende Richtung ist die
Spurnetzwerk-Methode.} Sie kann Normen durch Kontraktionen des
vorhandenen Tensors berechnen, statt alle Vielteilchenamplituden zu
speichern. Zur fundamentalen Vereinfachung wird sie erst dann, wenn sie
auch gemischte H-Momente, Seitenzweige und kontrollierte
Mehrzeitantworten liefert. Der konkrete nächste Test ist daher kein
bloßer weiterer Normwert: Eine zweite kleine Kontraktionsrechnung soll
die schon exakt bekannten zehn H-Momente reproduzieren und dann
mindestens ein noch nicht bekanntes gemischtes Matrixelement samt
Fehlerkontrolle bestimmen. Diese Methode lässt sich mit der
Schattenkarte verbinden, ersetzt aber keine native Operation.

Das Prüfpaket bewahrt neben dem Text die beiden widersprechenden
Programmstempel und das vorhandene Spurnetzwerkprogramm. Diese fremden
Dateien wurden nicht verändert oder als vollständig neu ausgeführt
ausgegeben.

\begin{center}\rule{0.5\linewidth}{0.5pt}\end{center}

\section{Historischer Anhang: vollständige Fassung
v1.6.6}\label{historischer-anhang-vollstuxe4ndige-fassung-v1.6.6}

Die folgende vollständige Fassung wird unverändert bewahrt. Maßgeblich
für den aktuellen Status sind die vorangestellten Kapitel v1.6.7. Ältere
Vorhaben, Modellannahmen und Zahlen erhalten dadurch keinen neuen
Beweisstatus; die kleinen Energieverschärfungen und die neuen Singulett-
und Bilinearergebnisse stehen oben.

\section{TFPT / Universalraum: Clock, gemeinsame Quelle und
tatsächlicher
Transport}\label{tfpt-universalraum-clock-gemeinsame-quelle-und-tatsuxe4chlicher-transport}

\subsection{Konsolidierung und eigene Forschungsfortsetzung v1.6.6 - 15.
September
2026}\label{konsolidierung-und-eigene-forschungsfortsetzung-v1.6.6---15.-september-2026}

\textbf{Ergebnis:} Der dokumentierte endliche Clock ist jetzt
ausdrücklich als Element der vorhandenen inneren Spin(10)-Symmetrie
konstruiert. Dadurch lässt sich die bislang offene gemeinsame
Clock-/Casimir-Auslesung exakt bestimmen. Daneben wurde ein kleiner,
echter Vierzustandskanal der ursprünglichen Paarwechselwirkung gefunden,
der nach Zulassung eines Bosonmischers eine Fermionenmarke überträgt.
Die Herkunft dieses Mischers ist nicht bewiesen.

Der neue Universalraum-Text enthält eine sinnvolle Suchrichtung -
gemeinsame Quelle statt vorab unabhängiger Banken -, aber Überlappung
allein erzeugt weder Transport noch Raumzeit, Eichkrümmung oder die
arithmetische Spur. Mehrere seiner vermeintlich automatischen Übergänge
werden hier präzisiert.

Diese Revision integriert \textbf{drei neue Nutzertexte}, die
unabhängigen parallelen Prüfungen und eigene neue Herleitungen. Die
vollständige frühere Konsolidierung v1.6.5 einschließlich v1.6.4 bleibt
im historischen Anhang erhalten. Das ist eine aktualisierte
Forschungsfassung in Markdown mit kurzem Änderungsdokument und einfacher
Erklärung; keine behauptete neue PDF-/Web-Publikation oder vollständige
Theory of Everything.

\subsubsection{Leseschlüssel}\label{leseschluxfcssel}

\begin{itemize}
\tightlist
\item
  \textbf{Exakt:} Identität oder mathematische Folgerung im angegebenen
  Modell.
\item
  \textbf{Numerisch:} berechneter Näherungswert, ohne zertifizierte
  Einschließung.
\item
  \textbf{Bedingt:} exakter Satz erst nach ausdrücklich zusätzlicher
  Ressource.
\item
  \textbf{Offen:} fehlende Herleitung oder physikalische Identifikation.
\end{itemize}

Eine grüne Prüfung ersetzt weder die Voraussetzungen eines Satzes noch
eine ausführbare Operation. Assoziative Algebradimension, dynamische
Kontrollierbarkeit und physische Verfügbarkeit werden getrennt geführt.

\subsection{1. Die unveränderte gemeinsame
Grundlage}\label{die-unveruxe4nderte-gemeinsame-grundlage}

Mit 64 Fermionmoden, 60 Bosonmoden und demselben gepinnten Tensor gilt

\[\adjustbox{max width=\linewidth}{$\displaystyle 
H_{\rm nat}=\Delta N_b+g\sum_A(b_A^\dagger P_A+P_A^\dagger b_A),
\qquad P_A=\sum_{i<j}W_{A,ij}f_jf_i,
\qquad N=N_f+2N_b.
$}\]

W besitzt 480 Einträge ±1, acht disjunkte Paare je Zeile, und
\codeword{WW†=8 I60}. Die Fermionmarkierung ist \codeword{16⊗4}, die
Bosonmarkierung \codeword{10⊗6}. Spin(10)×SU(4) ist zunächst eine
\textbf{innere} Quellsymmetrie, keine bereits hergeleitete
Raumzeit-Lorentzgruppe.

Der frühere native Satz bleibt unverändert maßgeblich: Für
Δ\textgreater0 und \codeword{0<|g|/Δ≤1/20} hat dieser μ=0-Vertrag einen
eindeutigen globalen Grundzustand Ω bei N=64. Am Prüfpunkt
\codeword{g/Δ=1/20} gelten unter anderem

\[\adjustbox{max width=\linewidth}{$\displaystyle 
-1.158089\Delta<E_0<-1.129636\Delta,
\qquad 0.842846<\langle N_b\rangle_\Omega<1.245656.
$}\]

Die isolierte niedrige Entnahmelinie im N=63-Sektor hat weiterhin die
bewiesenen Schranken aus v1.6.4/v1.6.5:

\[\adjustbox{max width=\linewidth}{$\displaystyle 
0.007737\Delta<\epsilon<0.039079764\Delta,
\qquad
Z_{\rm low}>\frac{40912436089}{46487375000}>0.88007628.
$}\]

Z bezieht sich auf das gesamte normierte Fermionspektralgewicht pro
Mode. Die gesamte Entnahmenorm ist dagegen
\codeword{ν=<f\_r†f\_r>=1-<Nb>/32}. Die beiden Größen sind nicht
austauschbar. Der Pol ist eine wohldefinierte isolierte Eigenlinie
dieses endlichen Modenmodells; seine relativistische Teilchen-, Massen-
oder räumliche Propagationsdeutung ist eine weitere Aufgabe.

Der vollständige bedingte Zwei-Banken-Satz aus v1.6.5 bleibt erhalten:
mit deklariertem Fermionlink und Rotorvertrag gelingt der reine niedrige
Transfer mit Wahrscheinlichkeit \textgreater99,267 \%, die unfiltrierte
ursprüngliche Entnahme mit \textgreater89,726 \%. Der Link war und ist
\textbf{zusätzlich vorausgesetzt}. Die neuen Vierzustandszahlen unten
gehören zu einem anderen Vertrag und ersetzen diese Ergebnisse nicht.

\subsection{2. Quellenprüfung der beiden neuen
Ergebnisrunden}\label{quellenpruxfcfung-der-beiden-neuen-ergebnisrunden}

\subsubsection{2.1 Bestätigte Zerlegung, korrigierte
Operationsaussagen}\label{bestuxe4tigte-zerlegung-korrigierte-operationsaussagen}

Die 944, 31 und 152 mitgelieferten Prüfbedingungen der Symmetrie-,
Verfügbarkeits- und Lorentzprogramme sind normal und optimiert
byteidentisch reproduziert. Darunter befinden sich neun wörtliche
\codeword{need(True, ...)}-Guards mit theoretischen Folgerungen. Sie
zählen als ausgeführte Guards, nicht als neun unabhängige Beweise. Die
Clock-Normalisatorbehauptung gehörte dazu; Abschnitt 3 ersetzt sie durch
eine explizite Konstruktion.

Im gesamten N=3-Raum gilt die Zerlegung:

{\def\LTcaptype{none} % do not increment counter
\begin{longtable}[]{@{}
  >{\raggedright\arraybackslash}p{(\linewidth - 8\tabcolsep) * \real{0.1579}}
  >{\raggedleft\arraybackslash}p{(\linewidth - 8\tabcolsep) * \real{0.2105}}
  >{\raggedleft\arraybackslash}p{(\linewidth - 8\tabcolsep) * \real{0.2105}}
  >{\raggedleft\arraybackslash}p{(\linewidth - 8\tabcolsep) * \real{0.2105}}
  >{\raggedleft\arraybackslash}p{(\linewidth - 8\tabcolsep) * \real{0.2105}}@{}}
\toprule\noalign{}
\begin{minipage}[b]{\linewidth}\raggedright
Typ
\end{minipage} & \begin{minipage}[b]{\linewidth}\raggedleft
irreduzible Dimension
\end{minipage} & \begin{minipage}[b]{\linewidth}\raggedleft
Multiplizität
\end{minipage} & \begin{minipage}[b]{\linewidth}\raggedleft
Spin-Casimir
\end{minipage} & \begin{minipage}[b]{\linewidth}\raggedleft
Farb-Casimir
\end{minipage} \\
\midrule\noalign{}
\endhead
\bottomrule\noalign{}
\endlastfoot
(1200,20) & 24000 & 1 & 141/4 & 39/4 \\
(560,20′) & 11200 & 1 & \textbf{117/4} & 63/4 \\
(672,4̄) & 2688 & 1 & 165/4 & 15/4 \\
(144,20) & 2880 & 2 & 85/4 & 39/4 \\
(144,4̄) & 576 & 2 & 85/4 & 15/4 \\
(16̄,20) & 320 & 2 & 45/4 & 39/4 \\
(16̄,4̄) & 64 & 1 & 45/4 & 15/4 \\
\end{longtable}
}

Die gewichtete Gesamtdimension ist 45504. Nur die getrennten Teile Λ³64
und 60⊗64 sind jeweils multiplizitätsfrei, nicht ihre direkte Summe. Die
drei dunklen Typen besitzen alle Gesamtcasimir 45, werden aber bereits
durch \textbf{einen} getrennten Casimir unterschieden.

Eine unabhängige rationale Rechnung auf den kleinen Multiplizitätsräumen
liefert für die jeweils \textbf{gewährten} Operationen:

{\def\LTcaptype{none} % do not increment counter
\begin{longtable}[]{@{}lr@{}}
\toprule\noalign{}
Operationsvertrag & komplexe assoziative Algebradimension \\
\midrule\noalign{}
\endhead
\bottomrule\noalign{}
\endlastfoot
Ein fixes H=Nb+X/20 & 8 \\
X und Nb unabhängig & 14 \\
Zusätzlich ein dunkler Isotypieprojektor & 15 \\
Zusätzlich ein getrennter Casimir & 16 \\
Zusätzlich nur Gesamtcasimir & 14 \\
\end{longtable}
}

Somit gibt es keine universelle binäre Frage „alle Symmetriegriffe oder
gar kein Fortschritt``. Die Dimension 15 ist ein konkretes
Gegenbeispiel. Die vollständige Gruppen-Generatoralgebra hat mit X,Nb
die Dimension 743583744 und Kommutantdimension 7. Sie ist unter
Gruppenkonjugation \textbf{stabil}, aber nicht punktweise invariant.
Sieben ist kein universeller Boden für beliebige symmetriestabile
Algebren: End(H) wäre ebenfalls stabil und hätte einen eindimensionalen
Kommutanten. Keine dieser Zahlen beweist, dass die zugehörigen
Operationen im Compiler ausführbar sind.

\subsubsection{2.2 Der Grundzustandsbericht enthält einen
Versionskonflikt}\label{der-grundzustandsbericht-enthuxe4lt-einen-versionskonflikt}

Der neue Text hat die Nichtschlussnorm bereits korrigiert, das
danebenliegende \codeword{native\_ground\_state.json} stammt jedoch von
einer älteren Codefassung. Sein als PASS markierter Wert ist um
\codeword{229²} zu groß; die dortige relative orthogonale Norm ist sogar
größer als eins. Der daraus berechnete Ritzwert −1,18826198 wird
\textbf{nicht} übernommen.

Eine neue direkte ganzzahlige Kontraktion umgeht die 15,25 Millionen
Komponenten von v₃. Sie berechnet \codeword{u=T T†v₂} auf 293280
nichtverschwindenden Komponenten und bestätigt exakt

\[\adjustbox{max width=\linewidth}{$\displaystyle 
\|u\|^2=752194252800,\quad
\langle v_2,u\rangle=575078400,\quad
w_2=u-\frac{299520}{229}v_2,
\quad\|w_2\|^2=\frac{5001523200}{229}>0.
$}\]

Der eine Vektor je Bosonlage schließt daher nicht. w₂ ist ein zweiter
Singulettvektor derselben Lage. Für die kleine vier- bzw.
fünfdimensionale Ritzmatrix ergeben sich bei g/Δ=1/20 \textbf{numerisch}

\[\adjustbox{max width=\linewidth}{$\displaystyle 
E_{R,4}/\Delta=-1.0942308394409612,
\qquad E_{R,5}/\Delta=-1.0942318710142458.
$}\]

Die zusätzliche w₂-Richtung ändert diese Näherung tatsächlich erst um
etwa 10⁻⁶. Das bedeutet \textbf{keine} Konvergenz zum wahren
Grundzustand: Beide Werte liegen noch über dessen bereits bewiesener
Obergrenze.

Die neue einfache Restformel zeigt die Ursache. Für den normierten
Vier-Vektor-Ritzeigenzustand mit letzter Komponente c₃ gilt

\[\adjustbox{max width=\linewidth}{$\displaystyle 
\|(H-E_R)\psi_3\|^2
=g^2|c_3|^2\frac{\nu_4+\|w_2\|^2}{\nu_3}.
$}\]

Mit dem früher geprüften vierten Moment ist die Restnorm numerisch rund
0,373063 Δ. Der kleine w₂-Zweig trägt weniger als 23 Millionstel des
quadrierten Restes; die ausgelassene v₄-Richtung dominiert. Die neuen
Zahlen \codeword{Zh≈0,9737} und \codeword{εh≈0,031Δ} gehören ebenfalls
zur Ritz-Näherung: Gesamtentnahmenorm und erstes Moment, nicht neu
bestimmte Polgrößen auf Ω.

Die Suche nach \codeword{E<−1,2Δ} oder \codeword{E<−1,1625Δ} wäre sogar
mit der bestehenden unteren Schranke unvereinbar. Die schwächere neue
Beweismethode öffnet den bereits geschlossenen nativen Grundzustandssatz
nicht wieder.

\subsection{3. Neuer exakter Anschluss: Der Clock liegt bereits in
Spin(10)}\label{neuer-exakter-anschluss-der-clock-liegt-bereits-in-spin10}

\subsubsection{3.1 Ausgangspunkt ist der tatsächliche
Quell-Clock}\label{ausgangspunkt-ist-der-tatsuxe4chliche-quell-clock}

Der gepinnte endliche Compilerpräfix liefert auf fünf Hilfsachsen

\[\adjustbox{max width=\linewidth}{$\displaystyle 
p=(0\mapsto2,\ 1\mapsto0,\ 2\mapsto1,\ 3\mapsto4,\ 4\mapsto3),
\qquad\det p=-1.
$}\]

Seine vorzeichenrichtige Exteriorhebung auf dem geraden
5-Moden-Hilfsfockraum ist Γ(p)\textbar even. Auf den nativen Fermionen
ist \codeword{GF=Γ(p)|even⊗I4}; auf den Bosonen ist
\codeword{GB=(-P10)⊗I6}, wobei P10 beide Fünfergruppen permutiert. Diese
Darstellung wird aus dem vorhandenen Clockpräfix erneut extrahiert,
nicht aufgrund passender Eigenwerte ausgewählt.

\subsubsection{3.2 Ein ausdrückliches
Spinwort}\label{ein-ausdruxfcckliches-spinwort}

Auf dem 32-dimensionalen Hilfsfockraum seien

\[\adjustbox{max width=\linewidth}{$\displaystyle 
\gamma_j=a_j+a_j^\dagger,\qquad
\gamma_{j+5}=i(a_j^\dagger-a_j),\qquad 0\le j<5,
$}\]

und Π dessen Fermionparität. Setze

\[\adjustbox{max width=\linewidth}{$\displaystyle 
\Omega_{10}=\gamma_0\gamma_1\cdots\gamma_9=i\Pi,
\qquad
R_{ab}=\frac{(\gamma_a-\gamma_b)(\gamma_{a+5}-\gamma_{b+5})}{2}.
$}\]

R\_ab ist das Produkt zweier reeller Einheitsvektoren der
Cliffordalgebra, also ein Spin(10)-Element. Direkt gilt
\codeword{R\_ab=i Γ((ab))}. Da \codeword{p=(01)(02)(34)} mit rechts
zuerst angewandter Permutation, folgt

\[\adjustbox{max width=\linewidth}{$\displaystyle 
\boxed{S=\Omega_{10}R_{01}R_{02}R_{34}=\Pi\Gamma(p)\in\mathrm{Spin}(10).}
$}\]

S ist ein Produkt von 16 reellen Clifford-Einheitsvektoren. Die
Matrixprüfung bestätigt ohne Rundungstoleranzen

\[\adjustbox{max width=\linewidth}{$\displaystyle 
S_{even}=\Gamma(p)_{even},\qquad S_{odd}=-\Gamma(p)_{odd},
\qquad S\gamma_jS^\dagger=-\gamma_{p(j)}
$}\]

mit entsprechender Fortsetzung auf die zweite Fünfergruppe. Damit
stimmen \textbf{beide nativen Darstellungen} und die ursprüngliche
Kopplung überein:

\[\adjustbox{max width=\linewidth}{$\displaystyle 
G_F=S_{even}\otimes I_4,\qquad
G_B=(-P_{10})\otimes I_6,\qquad
W\Lambda^2G_F=G_BW.
$}\]

S hat Ordnung sechs. Die Konjugation aller 45 Spin-Erzeuger und die
Kommutation mit allen 15 Farb-Erzeugern wurden zusätzlich direkt
geprüft. Die Aussage ist stärker und präziser als ein lediglich
behaupteter Normalisatorsatz.

\textbf{Bedeutung:} Der dokumentierte Clock und die innere
Spin-Symmetrie brauchen hier keine getrennten zusätzlichen
mathematischen Träger. Es handelt sich um einen konkreten endlichen
Schritt derselben Darstellung. Daraus folgen aber weder kontinuierliche
Spin-Kontrollen noch die Identifikation des Clocks mit Hamiltonzeit. Die
Hilfs-Cliffordachsen sind nicht bereits fünf Raumrichtungen.

\subsection{4. Clock und getrennte Casimire lassen sich gemeinsam
auslesen}\label{clock-und-getrennte-casimire-lassen-sich-gemeinsam-auslesen}

Weil G aus Spin(10) stammt und auf SU(4) trivial wirkt, kommutiert es
mit beiden getrennten Casimiren, auch auf ihren Fockhebungen. Die Frage
aus dem neuen Text ist damit auf algebraischer Ebene beantwortet.

Die exakten Charakterfolgen von S\_even, S\_odd und dem Vektor sind

\[\adjustbox{max width=\linewidth}{$\displaystyle 
(16,0,4,0,4,0),\quad(16,0,4,0,4,0),\quad(10,0,4,-6,4,0).
$}\]

Die schon unabhängig nachgerechneten Zerlegungen
\codeword{Sym³16=672+144}, \codeword{Λ³16=560},
\codeword{S21(16)=1200+144+16̄} und \codeword{10⊗16=144+16̄} bestimmen die
Charaktere der sieben Typen. Die diskrete Fourierinversion erfolgt exakt
in Z{[}ζ6{]}, nicht durch gerundete komplexe Eigenwerte. Für die drei
zuvor gemeinsam dunklen Räume ergibt sich:

{\def\LTcaptype{none} % do not increment counter
\begin{longtable}[]{@{}
  >{\raggedright\arraybackslash}p{(\linewidth - 12\tabcolsep) * \real{0.1111}}
  >{\raggedleft\arraybackslash}p{(\linewidth - 12\tabcolsep) * \real{0.1481}}
  >{\raggedleft\arraybackslash}p{(\linewidth - 12\tabcolsep) * \real{0.1481}}
  >{\raggedleft\arraybackslash}p{(\linewidth - 12\tabcolsep) * \real{0.1481}}
  >{\raggedleft\arraybackslash}p{(\linewidth - 12\tabcolsep) * \real{0.1481}}
  >{\raggedleft\arraybackslash}p{(\linewidth - 12\tabcolsep) * \real{0.1481}}
  >{\raggedleft\arraybackslash}p{(\linewidth - 12\tabcolsep) * \real{0.1481}}@{}}
\toprule\noalign{}
\begin{minipage}[b]{\linewidth}\raggedright
Dunkler Typ
\end{minipage} & \begin{minipage}[b]{\linewidth}\raggedleft
Phase 0
\end{minipage} & \begin{minipage}[b]{\linewidth}\raggedleft
1
\end{minipage} & \begin{minipage}[b]{\linewidth}\raggedleft
2
\end{minipage} & \begin{minipage}[b]{\linewidth}\raggedleft
3
\end{minipage} & \begin{minipage}[b]{\linewidth}\raggedleft
4
\end{minipage} & \begin{minipage}[b]{\linewidth}\raggedleft
5
\end{minipage} \\
\midrule\noalign{}
\endhead
\bottomrule\noalign{}
\endlastfoot
(1200,20) & 4000 & 4000 & 4000 & 4000 & 4000 & 4000 \\
(560,20′) & 1920 & 1840 & 1840 & 1920 & 1840 & 1840 \\
(672,4̄) & 464 & 440 & 440 & 464 & 440 & 440 \\
Summe, bisherige Clock-Auslesung & 6384 & 6280 & 6280 & 6384 & 6280 &
6280 \\
\end{longtable}
}

Ein separater Casimir trennt jede der sechs dunklen Clockzellen in genau
drei Teile. Die hellen Zweier-Multiplizitätsräume bleiben erhalten.
Daraus folgt im N=3-Sektor

\[\adjustbox{max width=\linewidth}{$\displaystyle 
\boxed{\dim\operatorname{Alg}^*(X,N_b,G,C_S)=96,}
\qquad
\boxed{\dim\operatorname{Alg}^*(X,N_b,G,C_S)'=119597824.}
$}\]

Zum Vergleich: ohne C\_S sind es 84 und 240742144. C\_C kann C\_S
ersetzen. Die Lagrangeprojektoren auf den drei dunklen C\_S-Werten
bestätigen die Trennung ausdrücklich; beide Casimire zugleich werden
nicht benötigt.

\textbf{Es handelt sich um eine gelöste Kombinationsfrage, nicht um eine
neue Kontrollfreigabe.} Der benötigte Casimir ist noch kein vom Compiler
bereitgestelltes Messinstrument. Der große verbleibende Kommutant zeigt
weiterhin viele nicht aufgelöste Freiheitsgrade. Die Dimensionszahlen
gelten für N=3, nicht als vollständige Klassifikation aller
Ladungssektoren.

\subsection{5. Neuer kleiner Transportzeuge mit tatsächlichen
W-Kanälen}\label{neuer-kleiner-transportzeuge-mit-tatsuxe4chlichen-w-kanuxe4len}

Ein gemeinsamer Fermion-/Boson-Dreierstern hätte die einfache Identität

\[\adjustbox{max width=\linewidth}{$\displaystyle 
Q_x=b^\dagger s f_x,\quad Q_y=b^\dagger s f_y,
\qquad [Q_y^\dagger,Q_x]=(N_b+n_s)f_y^\dagger f_x.
$}\]

Auf festem \codeword{K=Nb+n\_s≥1} ist \codeword{c†=b†s/√K} exakt eine
CAR-Mode. Das ergibt eine gewöhnliche Dreimoden-Transferkette.
\textbf{Aber dieser Stern ist nicht in einer nativen W-Zeile enthalten:}
Jede Zeile ist ein Matching aus acht disjunkten Paaren. Die Analogie
allein wäre daher eine falsche Herkunftsangabe.

Im echten W gibt es stattdessen die Paare \codeword{(4,57)} in Kanal 0
und \codeword{(4,58)} in Kanal 1. Mit den sieben besetzten
Pauli-Blockern

\[\adjustbox{max width=\linewidth}{$\displaystyle 
S_b=\{8,16,28,32,44,52,56\}
$}\]

und einem \textbf{zusätzlichen} Mischer \codeword{J(b1†b0+b0†b1)}
schließt die volle native Wechselwirkung exakt auf den vier Zuständen

\[\adjustbox{max width=\linewidth}{$\displaystyle 
S_b+\{4,57\}\ \longleftrightarrow\ S_b+b_0
\ \longleftrightarrow\ S_b+b_1
\ \longleftrightarrow\ S_b+\{4,58\}.
$}\]

Alle haben N=9. Die vollständige Wirkung sämtlicher 480 signierter
Paarterme bleibt in diesem Raum, mit Matrix

\[\adjustbox{max width=\linewidth}{$\displaystyle 
H_4=\begin{pmatrix}
0&g&0&0\\g&\Delta&J&0\\0&J&\Delta&g\\0&0&g&0
\end{pmatrix}.
$}\]

Für \codeword{g/Δ=1/20}, \codeword{J/Δ=1-1/(10√3)} und
\codeword{tΔ=20π√3} gilt analytisch

\[\adjustbox{max width=\linewidth}{$\displaystyle 
P(57\to58)>\frac{2009992727}{2022609600}
=0.9937620819\ldots>99.3\%.
$}\]

Die Eigenfrequenzen Δ±J bleiben positiv. Dies ist keine unkontrollierte
Vierzustandsabschneidung: Zuerst ist die Invarianz unter der
vollständigen Wechselwirkung bewiesen, danach wird deren exakte
Einschränkung gelöst. Ohne Mischer zerfällt sie in zwei getrennte Blöcke
und der Transfer ist null.

Der reine Einzelmischer kommutiert nicht mit dem ursprünglichen Clock.
Die Clock-Orbitsumme der Mischer \codeword{(0,1)+(12,13)+(6,7)} ist
dagegen Clock-invariant und hat auf diesem Vierzustandsraum dieselbe
Wirkung. Somit wäre „jeder solche Transfer muss den Clock brechen``
falsch. Beide Mischer sind jedoch noch zusätzliche Operationen;
Clock-Invarianz allein beweist keine Erzeugbarkeit. Ein gemeinsamer
SU(4)-Cartan Q kommutiert mit H, Nb, Clock und beiden getrennten
Casimiren. Auf den vier Zuständen besitzt er aber die Werte
\codeword{(7,7,9,9)}. Die beiden Mischer ändern diese Quantenzahl:
\codeword{||[Q\_B,M01]||²\_F=8}, \codeword{||[Q\_B,Morb]||²\_F=24}.
\textbf{Auch die Clock-invariante Orbitsumme ist damit aus diesem
Operationssatz ausgeschlossen.}

Die noch fehlende Ressource ist hier konkret: eine Operation, die diese
innere Cartanladung verändert, oder eine global hergeleitete Kopplung an
einen ausdrücklich mitgerechneten Ladungsausgleich. Mehr Wörter aus dem
unveränderten erhaltenden Alphabet können beide Mischer nicht liefern.

Dieser Zeuge ist \textbf{keine} Propagation der
N=64-Grundzustands-Lochanregung, keine räumliche Zwei-Banken-Übertragung
und keine Konstruktion ihrer Präparation. Er zeigt konstruktiv, wie nah
eine kleine Paarumwandlung an einem Transfermechanismus liegen kann, und
benennt die fehlende Ressource.

\subsection{6. Das Feldwörterbuch: eine Basisänderung allein reicht
nicht}\label{das-feldwuxf6rterbuch-eine-basisuxe4nderung-allein-reicht-nicht}

Der bestätigte Grundfehler bleibt: \codeword{M\_A⊗ε\_Lorentz} ist
symmetrisch und verschwindet als Grassmann-Bilinear. Das betrifft die
lokale ableitungsfreie Zuordnung einer gleichhändigen Weylkopie je
ursprünglichem Label bei unverändertem antisymmetrischem W. Der
nichtverschwindende gleichhändige Kanal trägt \codeword{(1,0)}; daraus
folgt noch keine gesunde Kinetik.

Der neue Fünf-Zyklus im Trägergraphen widerlegt eine rein diagonale
Zuweisung entgegengesetzter Händigkeiten an jedes Paar. Ein neuer
eigener Satz geht darüber hinaus. Für irgendeine hermitesche interne
Händigkeit Γ müsste im unveränderten Ein-Kopie-Vektoransatz gelten

\[\adjustbox{max width=\linewidth}{$\displaystyle 
\Gamma^TM_A+M_A\Gamma=0\quad\forall A.
$}\]

Adjungieren und Einsetzen liefert \codeword{[Γ,M\_A†M\_B]=0}. Der
ursprüngliche Tensor faktorisiert exakt als

\[\adjustbox{max width=\linewidth}{$\displaystyle 
M_{ka}=S_k\otimes C_a,\quad
\sum_kS_k^\dagger S_k=5I_{16},\quad
\sum_aC_a^\dagger C_a=3I_4.
$}\]

Die S\_k†S\_l erzeugen M16; die C\_a†C\_b erzeugen M4. Dafür wurden
vollständige Rangzertifikate modulo 101 mit 256 bzw. 16 unabhängigen
ganzzahligen Matrixwörtern konstruiert. Ein nichtverschwindender
Determinant modulo 101 beweist Nichtverschwindung über C; keine
Rangdefizienz wird übertragen. Damit ist Γ skalar, und die ursprüngliche
Gleichung erzwingt Γ=0.

\textbf{Keine hermitesche Involution Γ²=I auf demselben Träger repariert
diesen Ein-Kopie-Vektoransatz.} Der Ausschluss ist nun basisunabhängig,
aber weiterhin vertragsgebunden.

Nicht ausgeschlossen sind zusätzliche Dirackomponenten, ein unabhängiger
Zweierfaktor oder Ableitungsterme. Insbesondere ist

\[\adjustbox{max width=\linewidth}{$\displaystyle 
M_A\otimes\epsilon_{aux}\otimes\epsilon_{Lorentz}
$}\]

antisymmetrisch und nicht null. Ein ableitungshaltiger gleichgeladener
Vektorstrom \codeword{M\_IJ ε\_ab ψ\_Ia ↔∂\_μ ψ\_Jb} ist ebenfalls nicht
null; der pauschale Vektorausschluss durch Ladung gilt nicht außerhalb
der ableitungsfreien Klasse.

Die Zahlen 180/128 oder 60/256 zählen Komponenten eines bestimmten
Feldansatzes, nicht automatisch unabhängige propagierende Oszillatoren.
Schur fixiert interne Intertwiner, nicht alle Lorentz-/Ableitungsterme.
Der behauptete Zwang „zuerst räumlich skalieren, danach Feldtyp prüfen``
folgt nicht daraus. Zuerst bleibt ein kleiner konsistenter
Typ-/Kinetiktest sinnvoll; seine reale räumliche Herkunft muss
anschließend gemeinsam geprüft werden.

\subsection{7. Was der nachgereichte Universalraum-Text
trägt}\label{was-der-nachgereichte-universalraum-text-truxe4gt}

\subsubsection{7.1 Eine sinnvolle Hypothese, kein bereits
identifiziertes universelles
Objekt}\label{eine-sinnvolle-hypothese-kein-bereits-identifiziertes-universelles-objekt}

Nichtorthogonale oder überlappende lokale Einbettungen in eine
gemeinsame Algebra sind eine ernstzunehmende Alternative zu unabhängigen
Banken. Sie können Voraussetzungen der lokalen Paritätsschranke ändern.
Das ist ein Forschungsauftrag, kein Widerspruch zur bisherigen Schranke.

Bei Einbettungen \codeword{J\_A:V\_A→V} gilt für die CAR jedoch

\[\adjustbox{max width=\linewidth}{$\displaystyle 
\{f_A(u),f_B(v)^\dagger\}=\langle J_Au,J_Bv\rangle.
$}\]

Die Überlappungs-Grammatrix gehört also zwingend zum Modell. Die alten
Produktzustände \codeword{Ω\_A⊗Ω\_B}, ihre beiden unabhängigen
Ladungssektoren und der Zwei-Banken-Polraum dürfen nicht ungeprüft in
einen überlappenden Träger übernommen werden. Der gemeinsame
Hamiltonoperator und sein Zustand müssen dort neu beziehungsweise durch
einen bewiesenen Adapter konstruiert werden.

\subsubsection{7.2 Der vorgeschlagene Kill-Test muss korrigiert
werden}\label{der-vorgeschlagene-kill-test-muss-korrigiert-werden}

Für unabhängige orthogonale Banken kann Π\_AΠ\_B deren gemeinsame
Parität sein. Für überlappende Unterräume gilt das nicht allgemein; ihre
lokalen Paritäten müssen nicht einmal kommutieren. Der Test
\codeword{[U,Π\_AΠ\_B]=0} ist deshalb im neuen Bild nicht der richtige
universelle Test. Stattdessen muss die tatsächliche globale Parität
Π\_global separat definiert und erhalten werden. Ein
nichtverschwindender Kommutator mit einer lokalen Parität belegt zudem
allein noch keinen gerichteten Transfer.

\subsubsection{7.3 Der einfachste False-Positive: Ein Signal ohne
Bewegung}\label{der-einfachste-false-positive-ein-signal-ohne-bewegung}

Seien \codeword{a=c1}, \codeword{b=(c1+c2)/√2} auf zwei globalen
CAR-Moden und \codeword{H=ωN}. Dann haben die beiden
Einteilchen-Chartzustände bereits Überlappung \codeword{1/√2}. In ihrer
nichtorthogonalen Basis ist \codeword{K=ωS} mit einem
nichtverschwindenden Offdiagonalelement. Trotzdem bleibt die betreffende
Wahrscheinlichkeit zu jeder Zeit genau 1/2. Es wurde nichts
transportiert.

Ein belastbarer gemeinsamer Test muss daher mindestens

\[\adjustbox{max width=\linewidth}{$\displaystyle 
S_{AB}=\langle\psi_A,\psi_B\rangle,\qquad
K_{AB}=\langle\psi_A,H\psi_B\rangle
$}\]

und die tatsächliche Zeitantwort gemeinsam bestimmen. Auf einem positiv
definiten Gramträger ist
\codeword{S\textasciicircum{}\{-1/2\}KS\textasciicircum{}\{-1/2\}} der
korrekte orthonormierte Operator; bei singulärem S ist zunächst der
Nullraum zu quotientieren. Im Gegenbeispiel ergibt dies nur ωI. Die
Offdiagonale von K war kein Hop.

Das trifft sogar unmittelbar den vorhandenen nativen Pol. Dessen
64-dimensionaler Raum besitzt genau \textbf{eine} Energie E\_h, also
\codeword{P\_h H P\_h=E\_h P\_h}. Sind A und B nur neue Charts innerhalb
dieses selben Polraums, gilt zwangsläufig \codeword{K=E\_h S}. Bloßes
Umbenennen der 64 entarteten Polrichtungen erzeugt daher keine
Ausbreitung durch H. Eine tatsächliche globale Dynamik beziehungsweise
eine verfügbare aktive Operation muss hinzukommen und aus der Quelle
ausgewiesen werden.

Für die eigentliche TFPT-Lochantwort wäre der gemeinsame Gegenstand

\[\adjustbox{max width=\linewidth}{$\displaystyle 
G_{AB}(t)=\langle\Omega,
 f_B^\dagger e^{-it(H-E_0)}f_A\Omega\rangle,
$}\]

mit einer einzigen Quelle, einem H und einem Ω. \codeword{G(0)} bezahlt
die schon vorhandene Überlappung; die dynamische Änderung muss davon
unterschieden werden. Das ist die geschärfte, kleinere nächste
Forschungsaufgabe.

\subsubsection{7.4 Was nicht automatisch
folgt}\label{was-nicht-automatisch-folgt}

\begin{itemize}
\tightlist
\item
  \textbf{Zeit:} Für einen Grundzustand ist
  \codeword{e\textasciicircum{}\{-itH\}Ω=e\textasciicircum{}\{-itE0\}Ω};
  sein physischer Zustand bleibt gleich. Das ist kein eigener Zeitpfeil.
  Nichtstationäre Präparation, Korrelationen und Records müssen
  ausgewiesen werden.
\item
  \textbf{Eichfeld:} Für bloße Basiswechsel \codeword{Uxy=Vx†Vy}
  teleskopiert jede geschlossene Holonomie zu I. Echte Krümmung braucht
  zusätzliche Verbindungs-/Transportdaten oder einen nachgewiesenen
  Mechanismus.
\item
  \textbf{Spin:} Zwei Orientierungsmarken liefern nicht von selbst einen
  SL(2,C)-Spinor, dessen Transformationsgesetz und Kinetik. Der
  zusätzliche antisymmetrische Zweierfaktor muss unabhängig konstruiert
  sein.
\item
  \textbf{Raum:} Minimale Operationskosten sind ohne Reversibilität
  zunächst gerichtete Kosten, nicht notwendig eine symmetrische Metrik.
  Kubisches Volumenwachstum allein beweist keine 3D-Mannigfaltigkeit,
  keine Lorentzstruktur und keinen gemeinsamen Lichtkegel.
\item
  \textbf{Gravitation:} Dynamische Beziehungen sind ein möglicher
  Träger, aber ein masseloser Spin-2-Sektor, zwei Helizitäten,
  Energiepositivität und universelle konsistente Kopplung bleiben zu
  beweisen.
\item
  \textbf{Arithmetik:} Primitive Schleifen besitzen generisch neue
  gemischte primitive Zyklen. Sie sind nicht automatisch Primzahlen. Die
  Xi-Determinante bleibt eine unbewiesene vollständige
  Operatoridentität. Der getrennte RH-Anhang dokumentiert Vorarbeiten,
  Beispiele und Quellenstatus.
\end{itemize}

Der Text wird deshalb als \textbf{präzisierter globaler
Forschungsansatz} integriert, nicht als Nachweis, dass derselbe
unbekannte Baustein schon alle T1-T8-, RH-, Faktorisierungs-, P/NP- und
Hylæan-Fragen beantwortet.

\subsection{8. T1-T8: gemeinsamer Fortschritt ohne
Statusinflation}\label{t1-t8-gemeinsamer-fortschritt-ohne-statusinflation}

{\def\LTcaptype{none} % do not increment counter
\begin{longtable}[]{@{}
  >{\raggedright\arraybackslash}p{(\linewidth - 4\tabcolsep) * \real{0.3333}}
  >{\raggedright\arraybackslash}p{(\linewidth - 4\tabcolsep) * \real{0.3333}}
  >{\raggedright\arraybackslash}p{(\linewidth - 4\tabcolsep) * \real{0.3333}}@{}}
\toprule\noalign{}
\begin{minipage}[b]{\linewidth}\raggedright
Front
\end{minipage} & \begin{minipage}[b]{\linewidth}\raggedright
In dieser Revision weitergekommen
\end{minipage} & \begin{minipage}[b]{\linewidth}\raggedright
Noch erforderlicher Nachweis
\end{minipage} \\
\midrule\noalign{}
\endhead
\bottomrule\noalign{}
\endlastfoot
T1 - Ursprung und Auswahl & Explizite Identität zwischen Quell-Clock und
innerer Spinwirkung; Operationsverträge genauer & Primitive Auswahl von
P1/P2, Träger, erlaubten Instrumenten und Zustand \\
T2 - Half-Charge/E8-Feld & Bestehender geladener Pol bleibt abgesichert,
innere Darstellung präzisiert & Renormiertes Half-Charge-Feld mit
Energie-/Adjungiertenkontrolle und E8-/Clock-Adapter \\
T3 - gemeinsamer 3+1D-Parent & Kleiner exakter Paartransferzeuge und
korrigierter Überlappungstest & Quellenseitig ausgewählter räumlicher
Parent auf einem gemeinsamen Hilbertraum und Zustand \\
T4 - chirale Materie & Basisunabhängiger Ausschluss eines konkreten
Ein-Kopie-Vektoradapters & Konsistenter chiraler Feldadapter,
Maß/Anomalien/Index und Spiegelentkopplung \\
T5 - Kontinuum und Dynamik & Exakte endliche Clock-/Casimir-Kombination
und bedingter invarianten Transferblock & Kontrollierter
wechselwirkender Grenzübergang, Lorentzverhalten, Clustering und
Streuung \\
T6 - Kopplungen und Texturen & Keine neue Herleitung; neue
Mischparameter sichtbar zusätzlich & Herkunft aller Kopplungen und
vollständiger Neutrinostruktur/-skala \\
T7 - Gravitation & Der falsche pauschale Tensor-Ausschluss bleibt
zurückgenommen & Masseloser dynamischer Spin zwei, beide Helizitäten,
konsistente universelle Kopplung im selben Parent \\
T8 - Präparation und Auslesung & Gram-/Antworttest gemeinsam formuliert;
spezielle Blockerressource konkret & Natives geladenes Instrument,
Quellzustandswahl, Records und ein gemeinsames Auslesefunktional \\
\end{longtable}
}

Keiner der acht vollständigen Abschlussverträge wird in dieser Revision
auf „gelöst`` gesetzt. Gelöst wurden ausdrücklich benannte Teilfragen.

\subsection{9. Die nächsten drei Arbeiten - klein, gemeinsam und
entscheidbar}\label{die-nuxe4chsten-drei-arbeiten---klein-gemeinsam-und-entscheidbar}

\textbf{A. Den gemeinsamen Quellenvertrag wirklich hinschreiben.} Zwei
lokale Einbettungen, ihr CAR-Gram, globale Parität und Ladung,
ursprüngliche Operationen und ein einziges H angeben. Abbruch für eine
Kandidatenfassung, wenn sie nur zwei Koordinatenansichten desselben
unveränderten Signals umbenennt. Die gemeinsamen Generatoren müssen
einen nachweislich neuen Antwortprozess ermöglichen, nicht nur eine
andere Beschriftung.

\textbf{B. Die fehlende Zwischenoperation ausweisen.} Der
Vierzustandszeuge gibt eine konkrete positive Zieloperation vor. Zu
prüfen ist, ob der Compiler ein geeignetes Mischinstrument oder eine
andere geteilte Paarstruktur wirklich liefert. Clock und Casimire
alleine sind kein Herkunftsbeweis. Zeigt ein erhaltener gemeinsamer
Generator die Unmöglichkeit für einen festen Operationssatz, endet
dessen weitere Kommutator-Suche; dann muss ein anderer bereits
vorhandener Quellbestandteil benannt werden.

\textbf{C. Dieselbe Anregung und denselben Zustand behalten.} Auf dem
Kandidaten \codeword{G\_AB(0)} und \codeword{G\_AB(t)} einschließlich
vollständiger Restantwort bestimmen. Der N=9-Zeuge darf nicht still zum
N=64-Lochpol umgedeutet werden. Ein minimaler Lorentz-/Kinetikadapter
muss die nachgerechneten Tensorbedingungen erfüllen. Erst mit diesem
gemeinsamen Anschluss trägt eine großräumige Skalierungsrechnung.

Der stärkste Vereinfachungsgewinn dieser Runde ist damit konkret:
\textbf{Clock und innere Symmetrie sind nicht zwei voneinander
unabhängige Mechanismen; reine Überlappung und echte Bewegung sind
dagegen zwei verschiedene Dinge.} Die weitere Suche sollte diese erste
Identität nutzen und die zweite Unterscheidung im selben kleinen Modell
erzwingen.

\subsection{10. Belege und
Reproduzierbarkeit}\label{belege-und-reproduzierbarkeit}

Die drei Eingangstexte bleiben unverändert im Prüfpaket. Ein Teil-Audit
reproduziert die 1127 gelieferten Guards; eigene Programme behandeln
Clock/Casimir, Quellenreichweite, direkte Krylov-Kontraktion,
basisunabhängige Feldbedingungen, den tatsächlichen Vierzustandskanal
und die Überlappungsgegenkontrollen. Exakte und numerische Teile sind in
den jeweiligen Ergebnisdateien getrennt markiert. Die abschließenden
Zählungen und Bytevergleiche stehen im Auslieferungs-/Replayprotokoll;
interne erneute Aufrufe der 944 Guards werden nicht mehrfach als neue
Tests gezählt.

Die drei ausführlichen parallelen Berichte und der
Überlappungs-/RH-Nachtrag folgen im vollständigen Dokument. Ihr Status
ersetzt widersprechende historische Deutungen, nicht unveränderte ältere
Beweise. Das native frühere Archiv bleibt zusätzlich unverändert im
Paket; dessen große frühere Enumeration wird nicht als in dieser
Revision erneut durchgeführt ausgegeben.

\begin{center}\rule{0.5\linewidth}{0.5pt}\end{center}

\section{Anhang A: Quellenprüfung Symmetrie und
Verfügbarkeit}\label{anhang-a-quellenpruxfcfung-symmetrie-und-verfuxfcgbarkeit}

\section{Quellenprüfung: Quellsymmetrie, Verfügbarkeit und
Lorentztypen}\label{quellenpruxfcfung-quellsymmetrie-verfuxfcgbarkeit-und-lorentztypen}

Stand: 15. September 2026. Eigenständiger Teil-Audit der ersten neuen
Anlage und dreier Programme. Keine Änderung der gelieferten Programme,
Ergebnisdateien oder Dokumente. Keine Aussage, dass T1-T8 geschlossen
seien.

\subsection{1. Reproduktion}\label{reproduktion}

Alle drei Programme wurden vollständig gelesen und anschließend in
isolierten Kopien ausgeführt. Die relevanten Tensoren wurden an ihre von
den unveränderten Programmen erwarteten relativen Orte kopiert. Normaler
und optimierter Lauf reproduzieren jeweils die gelieferten
Ergebnisdateien byteidentisch, ohne Standardfehlerausgabe:

{\def\LTcaptype{none} % do not increment counter
\begin{longtable}[]{@{}lrl@{}}
\toprule\noalign{}
Programm & Gelieferte Prüfbedingungen & Reproduktion \\
\midrule\noalign{}
\endhead
\bottomrule\noalign{}
\endlastfoot
\codeword{operation\_symmetry.py} & 944 & PASS, identisch \\
\codeword{symmetry\_availability.py} & 31 & PASS, identisch \\
\codeword{lorentz\_types.py} & 152 & PASS, identisch \\
Summe & 1127 & normal und \codeword{-OO} \\
\end{longtable}
}

Die Zusatzprüfung \codeword{check\_scope.py} rechnet kleine exakte
Gegenprüfungen zur Reichweite der Aussagen. Das ausführbare
Reproduktionsprogramm ist \codeword{replay.py}; Quellenpins, Kopien,
Originalberichte und Ausgaben liegen in diesem Auditordner.
\codeword{replay\_receipt.json} dokumentiert alle Pins und die
Unverändertheit der Eingaben.

Programm-Pins:

{\def\LTcaptype{none} % do not increment counter
\begin{longtable}[]{@{}
  >{\raggedright\arraybackslash}p{(\linewidth - 2\tabcolsep) * \real{0.5000}}
  >{\raggedright\arraybackslash}p{(\linewidth - 2\tabcolsep) * \real{0.5000}}@{}}
\toprule\noalign{}
\begin{minipage}[b]{\linewidth}\raggedright
Quelle
\end{minipage} & \begin{minipage}[b]{\linewidth}\raggedright
SHA-256
\end{minipage} \\
\midrule\noalign{}
\endhead
\bottomrule\noalign{}
\endlastfoot
\codeword{operation\_symmetry.py} &
\codeword{ec93f759d6582f42542a597bf916ee2a62a51c510240aa777ae439e90fe28382} \\
\codeword{symmetry\_availability.py} &
\codeword{88cc3e3361a198fe7b1b32a7174f3dd45f341106978d8abd8d5722e0cb631967} \\
\codeword{lorentz\_types.py} &
\codeword{43777c81ea8ba2d2ae12b711321a54abd7f7152a326611a296598d5c5dd463da} \\
Nativer Tensor &
\codeword{3f00a0892157a691e4ef26920c3702772c7bf76a250fad98ff04a079bb64b763} \\
\end{longtable}
}

Die Erzählung über zwei unabhängig entstandene, vollständig
übereinstimmende Programme ist durch diese Reproduktion \textbf{nicht}
geprüft: Vorgelegen hat die erhaltene Fassung mit 944 Guards, nicht die
vorherige Fassung mit 528 Guards.

\subsection{2. Bestätigte Mathematik}\label{bestuxe4tigte-mathematik}

Der N=3-Raum zerfällt als Darstellung von Spin(10) × SU(4) in sieben
Isotypien. Drei haben Multiplizität zwei, die anderen vier Multiplizität
eins. Die zwei Fockteile Λ³(C⁶⁴) und C⁶⁰⊗C⁶⁴ sind jeweils
multiplizitätsfrei; der gesamte N=3-Raum ist es ausdrücklich
\textbf{nicht}.

Die irreduziblen Dimensionen sind 24000, 11200, 2688, 2880, 576, 320,
64; die mit Multiplizitäten gewichtete Summe ist 45504. Die
Spin(10)-Darstellung mit Dimension 16 in BF und Λ³ ist in den
angegebenen Gewichten die konjugierte Spinordarstellung; die
Dimensionskurzschreibweise \codeword{16} darf diese Information nicht
ersetzen.

Bestätigt sind:

\begin{itemize}
\tightlist
\item
  die exakte Zerlegung über Gewichtsmultiplizitäten und
  Racah-Alternation;
\item
  \codeword{C(560)=117/4} in der erklärten Casimir-Normierung;
\item
  die Quellenintertwiner für alle 60 Lie-Erzeuger;
\item
  \codeword{rank C3=3776}, \codeword{dim ker C3=37888};
\item
  Spektrum von \codeword{C3 C3†}: 0,7,10,12 mit Multiplizitäten
  64,2880,576,320;
\item
  die drei dunklen Typen mit Gesamtcasimir 45, aber getrennten
  Spin-Casimiren 141/4,117/4,165/4;
\item
  für den \textbf{gewährten} Satz X,Nb eine assoziative Algebra der
  komplexen Dimension 14;
\item
  für alle punktweise symmetrieinvarianten Operatoren die Dimension 16;
\item
  nach zusätzlicher Gewährung aller 60 Gruppenerzeuger die assoziative
  Dimension 743583744 und Kommutantdimension 7.
\end{itemize}

Die Eigenwertzuordnung im Code prüft pro S-Eigenwert einen Vektor, nicht
eine vollständige Basis dieses Eigenraums. Zusammen mit der unabhängig
geprüften Multiplizitätsfreiheit, der Kommutation mit allen Erzeugern
und der eindeutigen Dimensionszuordnung der BF-Teilsummen ist die
Zuordnung trotzdem abgesichert. Die Guard-Beschriftung »entire
eigenspace« ist als direkte Beschreibung dieser einen Rechnung zu weit.

\subsection{3. Invarianz und Kovarianz
auseinanderhalten}\label{invarianz-und-kovarianz-auseinanderhalten}

Mit \codeword{A0=Alg*(X,Nb)} und Gruppenwirkung ρ bezeichne

\codeword{C = End\_G(H3) = ⊕\allowbreak{}\_i End(C\textasciicircum{}\{m\_i\}) ⊗ I\_\{d\_i\}}.

Dann gilt \codeword{A0 ⊂ C}, \codeword{dim A0=14} und
\codeword{dim C=Σ m\_i²=16}. Jedes Wort aus X und Nb kommutiert mit
ρ(G). Nichtzentrale Lie-Erzeuger können deshalb aus diesem Alphabet
nicht entstehen. Dieser Ausschluss ist korrekt und bleibt richtig, wenn
nur ein fixes H statt zweier unabhängiger Kontrollen gewährt ist.

Nach Gewährung der Gruppenerzeuger entsteht dagegen

\codeword{Afull = ⊕\allowbreak{}\_i End(C\textasciicircum{}\{m\_i\} ⊗ V\_i)}.

Diese Algebra ist \textbf{unter Gruppenkonjugation stabil}, aber ihre
Operatoren sind nicht sämtlich invariant. In
\codeword{symmetry\_availability.py} ist die Kennzeichnung
\codeword{symmetry\_invariant: true} für diese Zeile daher falsch. Der
Kommutant von Afull ist die sieben-dimensionale skalare Blockmitte.

Sieben ist ein Boden, solange alle zusätzlichen Operationen diese sieben
Isotypie-Projektoren erhalten; insbesondere senken zusätzliche
punktweise G-invariante Operationen den bereits erreichten Kommutanten
nicht weiter. Sieben ist aber \textbf{kein} universeller Boden für unter
G stabile Algebren oder alle denkbaren nativen Verträge.
\codeword{End(H3)} ist unter jeder Gruppenkonjugation stabil und hat nur
den skalaren Kommutanten. Ein exaktes 2×2-Gegenmodell in der
Zusatzprüfung macht die Unterscheidung direkt sichtbar.

Die Zahlen sind zudem Dimensionen \textbf{komplexer assoziativer
Sternalgebren}, keine nachgewiesenen Dimensionen einer dynamischen
Lie-Algebra oder erreichbaren Unitärgruppe. Burnside liefert nicht
automatisch beliebige ausführbare Gatter auf jedem irreduziblen Block.

\subsection{4. Verfügbarkeit ist keine universelle
Ja/Nein-Frage}\label{verfuxfcgbarkeit-ist-keine-universelle-janein-frage}

Ein vorgegebener Hamiltonoperator allein gewährt nicht bereits das
unabhängige Schalten von X und Nb. Am erklärten Punkt g/Δ=1/20 besitzt
ein fixes H im N=3-Raum acht verschiedene Energien, also nur eine
acht-dimensionale kommutative Spektralalgebra. Die großzügigere Annahme
unabhängiger Kontrollen X,Nb ergibt 14.

Eine treue, auf die Multiplizitätsräume reduzierte rationale
Matrixrechnung ergibt:

{\def\LTcaptype{none} % do not increment counter
\begin{longtable}[]{@{}lr@{}}
\toprule\noalign{}
Zusätzlich gewährter Satz & Assoziative Dimension \\
\midrule\noalign{}
\endhead
\bottomrule\noalign{}
\endlastfoot
Nur fixes H=Nb+X/20 & 8 \\
X,Nb & 14 \\
X,Nb plus Projektor auf genau einen dunklen Typ & 15 \\
X,Nb plus Spin(10)-Casimir & 16 \\
X,Nb plus SU(4)-Casimir & 16 \\
X,Nb plus Gesamtcasimir & 14 \\
\end{longtable}
}

Bereits \textbf{einer} der getrennten Casimire reicht zur vollständigen
Trennung der drei dunklen Typen. Die zwei fehlenden Algebradimensionen
sind kein Beweis, dass genau zwei physische Griffe fehlen. Die Dimension
15 zeigt eine mögliche invariante Zwischenstufe. Die behauptete
universelle Dichotomie »getrennte Griffe ja oder unverändert 14« und das
\codeword{IF AND ONLY IF} zur Verfügbarkeit der vollen Symmetrie sind
daher zu stark.

Die zusätzlichen Projektoren und Casimirkontrollen werden hier nur als
algebraische Gegenbeispiele verwendet. Ihre Implementierung aus dem
Compiler wird nicht behauptet.

\subsection{5. Der Clock-Satz ist im gelieferten Checker nicht
geprüft}\label{der-clock-satz-ist-im-gelieferten-checker-nicht-gepruxfcft}

\codeword{operation\_symmetry.py:478} schreibt den Normalisator- und
Nichts-hinzufügen-Satz als \codeword{need(True, ...)}. Insgesamt enthält
diese Datei neun solche als Prüfbedingungen gezählten theoretischen
Folgerungen. Einige sind durch die vorausgehende Mathematik gut
begründet; der bloße grüne Guard zertifiziert sie jedoch nicht
unabhängig.

Die direkte Clock-Definition liegt in
\codeword{universalraum-native-operations-ground-response-20260915/common.py},
Funktion \codeword{clock\_lift}, ab Zeile 187: eine Permutation der fünf
Oszillatorachsen wird mit Exterior-Vorzeichen auf den geraden Spinorraum
gehoben; \codeword{GF=G16⊗I4}. Der Bosonlift wird entsprechend gebaut
und die Tensorintertwining-Gleichung geprüft. Eine neue direkte Prüfung
der Clock-/Casimir-Beziehungen erfolgt im parallelen Hauptstrang, nicht
in diesem Audit.

\subsection{6. Feldtyp: was tatsächlich erzwungen
ist}\label{feldtyp-was-tatsuxe4chlich-erzwungen-ist}

Für ein \textbf{lokales, ableitungsfreies Bilinear zweier gleichhändiger
Weylfelder}, das alle 64 inneren Quellenmarken unverändert als
unabhängige innere Komponenten und genau den gegebenen antisymmetrischen
Tensor W verwendet, gilt

\codeword{(1/2,0) ⊗ (1/2,0) = (1,0) ⊕\allowbreak{} (0,0)}.

Die skalare Epsilon-Kontraktion zusammen mit antisymmetrischem W ist im
gemeinsamen Index symmetrisch und verschwindet wegen
Grassmann-Antikommutation. Der symmetrische Spinortensor, also der
(1,0)-Kanal, bleibt nichtverschwindend. Im genannten Vertrag ist dies
korrekt. Ein zusätzliches unabhängiges gleichgeladenes Hilfsdublett
erlaubt stattdessen wieder die skalare Kontraktion.

Die Zahlen 180/128 beziehungsweise 60/256 zählen \textbf{komplexe lokale
Feldkomponenten im jeweiligen direkten Tensorprodukt-Ansatz}. Sie zählen
nicht automatisch zusätzliche unabhängige physische Oszillatoren oder
propagierende Freiheitsgrade. Diese hängen von Kinetik,
Nebenbedingungen, positiver Energie, Teilchen-/Antiteilchenstruktur und
dem tatsächlichen Feldadapter ab. Deshalb folgt aus der Multiplikation
mit zwei oder drei kein allgemeiner Ausschluss jeder relativistischen
Lesart mit den nativen Modenzahlen.

Vor allem folgt daraus keine vorgeschriebene Reihenfolge »erst räumliche
Skalierung, dann Feldwörterbuch«. Lokale Darstellungstypen und mögliche
Kinetik können vor einer Skalierungsrechnung geprüft werden. Welche
Feldkomponenten tatsächlich aus dem räumlichen Quellprozess entstehen,
muss anschließend gemeinsam mit dem Adapter untersucht werden.

\subsubsection{Konkreter Gegencheck jenseits der ableitungsfreien
Klasse}\label{konkreter-gegencheck-jenseits-der-ableitungsfreien-klasse}

Der allgemeine Satz »ein Vektor verlangt ψ†ψ und scheitert deshalb an
der Ladung« ist falsch, sobald Ableitungen zugelassen werden.
Beispielsweise ist

\codeword{J\textasciicircum{}A\_mu = Σ\_IJab M\textasciicircum{}A\_IJ ε\_ab ψ\_Ia ↔∂\_mu ψ\_Jb}

ein Lorentzvektor mit Ladung −2. Die Ableitungsantisymmetrisierung macht
ihn bei antisymmetrischem M nichtverschwindend. Für M=ε und zwei inneren
Marken liefert die exakte Grassmann-Jetrechnung vier nichtverschwindende
Monome mit Koeffizienten +2,−2,−2,+2. Ein entsprechend geladener
Vektorvermittler könnte algebraisch durch
\codeword{b†\_mu J\textasciicircum{}mu + h.c.} koppeln, ohne
Ladungsverletzung.

Das ist \textbf{nur ein Gegenbeispiel gegen den überbreiten Ausschluss},
keine native Lösung: Es fügt eine Ableitung und einen Feldtyp hinzu; bei
kanonischen 4D-Felddimensionen ist es ein höherdimensionaler
Kopplungsterm. Gesunde Kinetik, Quellenherkunft, Skalierung und T1-T8
werden dadurch nicht bewiesen.

\subsection{7. Konsequenz}\label{konsequenz}

Die neuen Quellen verkleinern und präzisieren den endlichen
N=3-Operationsraum. Sie schließen weder die physische Verfügbarkeit der
Operationen noch den Lorentzadapter. Die sachlich tragfähige Integration
übernimmt die Zerlegung, Casimire, 14/16/7 im jeweils erklärten Vertrag
sowie den eingeschränkten Feldtypsatz; sie ersetzt die überbreiten
Ausschlüsse und die behauptete Programmumkehr durch explizite
Bedingungen.

\begin{center}\rule{0.5\linewidth}{0.5pt}\end{center}

\section{Anhang B: unabhängige Krylov- und
Feldprüfung}\label{anhang-b-unabhuxe4ngige-krylov--und-feldpruxfcfung}

\section{Unabhängige Prüfung: Krylov-Verzweigung und
Feldwörterbuch}\label{unabhuxe4ngige-pruxfcfung-krylov-verzweigung-und-feldwuxf6rterbuch}

Arbeitsstand 2026-09-15. Geprüft wird die Anlage
\codeword{7380c383-f364-4ef8-8066-24aa4d180d85} gegen den Vertrag
\codeword{universalraum-native-operations-ground-response-20260915}.
Fremde Quellen und Ausgaben wurden nicht verändert. Die hier angegebenen
Matrix- und Graphbefunde beziehen sich auf denselben gepinnten Tensor W
mit SHA-256
\codeword{3f00a0892157a691e4ef26920c3702772c7bf76a250fad98ff04a079bb64b763}.

\subsection{1. Der neue Nichtschluss ist exakt
bestätigt}\label{der-neue-nichtschluss-ist-exakt-bestuxe4tigt}

Für \codeword{v\_n=(T†)\textasciicircum{}n F} habe ich
\codeword{u=T T† v\_2=T v\_3} unmittelbar aus den 480 nativen Paaren neu
kontrahiert. Das materialisiert nicht die 15.252.960 Komponenten von v₃.
Es arbeitet mit 108.240 Komponenten von v₂ und 293.280
nichtverschwindenden Komponenten des Ergebnisses u. Es wurden 45.771.360
zulässige Erzeugungsübergänge und 139.289.760 anschließende
Entnahmebeiträge exakt mit ganzzahliger Fermion-Parität und
Boson-Multiplizität zusammengeführt.

Die unabhängige Rechnung liefert

\[\adjustbox{max width=\linewidth}{$\displaystyle 
\|v_2\|^2=439680,\quad
\langle v_2,u\rangle=575078400,\quad
\|u\|^2=752194252800.
$}\]

Damit folgen rein rational

\[\adjustbox{max width=\linewidth}{$\displaystyle 
\alpha=\frac{299520}{229},\qquad
w_2=u-\alpha v_2,\qquad
\|w_2\|^2=\frac{5001523200}{229}>0,
$}\]

und

\[\adjustbox{max width=\linewidth}{$\displaystyle 
\frac{\|w_2\|^2}{\nu_3}=\frac{1809}{47632}.
$}\]

Der Ein-Vektor-pro-Bosonlage-Ansatz ist somit nicht invariant. Da T die
native Symmetrie erhält und v₂ und u Singuletts sind, liefert w₂ einen
zweiten unabhängigen Singulettvektor in derselben Bosonlage. Das ist ein
Strukturresultat, kein physikalischer Mehrteilchen- oder
Raumzeitnachweis.

\subsection{2. Das mitgelieferte Ground-JSON ist veraltet und teilweise
falsch}\label{das-mitgelieferte-ground-json-ist-veraltet-und-teilweise-falsch}

Die aktuelle Pythonquelle hat SHA-256
\codeword{8ae0f2f103ba0cb0a5cb832681a33130b6756fe627395368b6be635587a2eb02}.
Das vorhandene \codeword{native\_ground\_state.json} nennt dagegen
\codeword{046ffa3f41d9dd365e2c9a199e938734bf735d762f23373a3f85b541c3da30fe}
und enthält noch

\[\adjustbox{max width=\linewidth}{$\displaystyle 
\|w_2\|^2_{\rm alt}=1145348812800
=229^2\|w_2\|^2.
$}\]

Sein \codeword{closure\_defect=1.5226768996658837} ist schon als
relative orthogonale Norm unmöglich. Das daraus abgeleitete
Ritz-Ergebnis −1,18826198 ist ungültig. Der spätere Anlagentext
korrigiert den Normfehler, aber die JSON-Belegkette wurde noch nicht
entsprechend neu erzeugt. Das grüne \codeword{status: PASS} dieser
historischen Ausgabe darf deshalb nicht mit dem korrigierten Text
zusammengezogen werden.

Die korrigierte Norm und die folgenden kleinen Matrizen wurden hier
unabhängig reproduziert. Die ganze fremde große Ausführung wurde
ausdrücklich nicht erneut gestartet.

\subsection{3. Die sechste Dezimalstelle ist richtig - aber nur für die
eine zusätzliche
Richtung}\label{die-sechste-dezimalstelle-ist-richtig---aber-nur-fuxfcr-die-eine-zusuxe4tzliche-richtung}

Die neue numerische Diagonalisierung ergibt:

{\def\LTcaptype{none} % do not increment counter
\begin{longtable}[]{@{}
  >{\raggedright\arraybackslash}p{(\linewidth - 6\tabcolsep) * \real{0.2000}}
  >{\raggedleft\arraybackslash}p{(\linewidth - 6\tabcolsep) * \real{0.2667}}
  >{\raggedleft\arraybackslash}p{(\linewidth - 6\tabcolsep) * \real{0.2667}}
  >{\raggedleft\arraybackslash}p{(\linewidth - 6\tabcolsep) * \real{0.2667}}@{}}
\toprule\noalign{}
\begin{minipage}[b]{\linewidth}\raggedright
g/Δ
\end{minipage} & \begin{minipage}[b]{\linewidth}\raggedleft
vierdimensionale Ritz-Matrix
\end{minipage} & \begin{minipage}[b]{\linewidth}\raggedleft
plus w₂
\end{minipage} & \begin{minipage}[b]{\linewidth}\raggedleft
Änderung
\end{minipage} \\
\midrule\noalign{}
\endhead
\bottomrule\noalign{}
\endlastfoot
1/20 & −1,0942308394409612 & −1,0942318710142458 & −1,0315733·10⁻⁶ \\
1/40 & −0,29535718178032033 & −0,29535720496683315 & −2,3186513·10⁻⁸ \\
1/100 & −0,04789424798601009 & −0,04789424801311775 &
−2,7107663·10⁻¹¹ \\
\end{longtable}
}

Das sind numerische Eigenwerte exakter kleiner Ritz-Matrizen, keine nach
außen gerundeten vollständigen Spektraleinschließungen. Durch
Eliminieren der einzigen neuen w₂-Koordinate bekommt die v₃-Diagonale
die exakt bestimmte energieabhängige Korrektur

\[\adjustbox{max width=\linewidth}{$\displaystyle 
-\frac{g^2(1809/47632)}{2\Delta-E}.
$}\]

Das erklärt die kleine Änderung durch diesen einen Zweig. Es sagt nichts
über die ausgelassenen höheren Bosonlagen.

\subsubsection{Eigene einfache Folgeherleitung: Das vollständige
Ritz-Residuum}\label{eigene-einfache-folgeherleitung-das-vollstuxe4ndige-ritz-residuum}

Sei ψ₃ der normierte Ritz-Eigenvektor in \codeword{span(v₀,…,v₃)}, und
c₃ seine Komponente entlang \codeword{v₃/√ν₃}. Außerhalb dieses Raums
führt Hψ₃ genau in zwei orthogonale Richtungen: v₄ und w₂. Deshalb gilt
exakt

\[\adjustbox{max width=\linewidth}{$\displaystyle 
\|(H-E_{\rm Ritz})\psi_3\|^2
=g^2|c_3|^2\frac{\nu_4+\|w_2\|^2}{\nu_3}.
$}\]

Hier wird \codeword{ν₄=952296652800} aus dem früher vollständig
geprüften Momentbeleg übernommen; die vierte Ordnung wurde in dieser
Runde nicht erneut enumeriert.

Bei g/Δ=1/20 folgt numerisch ein Residuum von rund \textbf{0,373063 Δ}.
Der w₂-Zweig trägt weniger als \textbf{23 Millionstel} des quadrierten
Residuums. Die große ausgelassene Richtung ist v₄, nicht w₂. Aus dem
winzigen Nichtschlussanteil darf daher keine praktisch vollständige
Konvergenz abgeleitet werden.

Außerdem liegt die neue Fünf-Zustands-Ritzenergie mindestens
\textbf{0,0354 Δ über} der bereits bewiesenen Obergrenze des wahren
Grundzustands. Das ist eine nachweislich noch relevante Lücke,
unabhängig von einer Konvergenzschätzung.

\subsection{4. Der bereits bewiesene native Grundsatz bleibt
gültig}\label{der-bereits-bewiesene-native-grundsatz-bleibt-guxfcltig}

Der ältere, gepinnte Beweis lautet: Für Δ\textgreater0 und
0\textless\textbar g\textbar/Δ≤1/20 besitzt der unveränderte
μ=0-Hamiltonoperator einen eindeutigen globalen Grundzustand in N=64; er
ist ein Spin(10)×SU(4)-Singulett. Bei 1/20 bestehen außerdem positive
Lücke und die bekannten Energie-, Besetzungs- und Polschranken.

Die neue Ausführung benutzt schwächere Sektor-Untergrenzen und weniger
Ritz-Vektoren. Dass diese schwächere Methode bei 1/20 nicht genügt, ist
weder eine Widerlegung noch ein Wiederöffnen des früheren Satzes.
Insbesondere sind die vorgeschlagenen Ziele \codeword{E<−1,2Δ} bzw.
\codeword{E<−1,1625Δ} unmöglich, weil bereits \codeword{E₀>−1,158089Δ}
bewiesen ist. Bessere Schranken für die Konkurrenzsektoren sind
erforderlich, nicht eine unrealistisch tiefere Energie.

Die neuen Zahlen \codeword{Nb≈0,84211266}, \codeword{Zh≈0,97368398} und
\codeword{εh≈0,03107309Δ} gehören zum nichtstationären K₃-Ritz-Zustand.
Nb liegt sogar knapp unter der früher bewiesenen Schranke
\codeword{Nb(Ω)>0,842846}. Zh bezeichnet die gesamte Entnahmenorm dieser
Näherung, nicht das Gewicht der isolierten niedrigen Linie. Ein erstes
Moment auf einem Ritz-Zustand ist keine neu bestimmte Polenergie auf Ω.

Offen bleiben der vollständige Ω-Vektor, die physische Auswahl des
μ=0-Vertrags und die Herkunft ausführbarer räumlicher Operationen.

\subsection{5. Feldwörterbuch: bestätigte Algebra und notwendige
Einschränkungen}\label{feldwuxf6rterbuch-bestuxe4tigte-algebra-und-notwendige-einschruxe4nkungen}

Frisch exakt geprüft wurden:

\begin{itemize}
\tightlist
\item
  Der Trägergraph ist zusammenhängend, 15-regulär und dreiecksfrei.
\item
  Der Fünf-Zyklus \codeword{16→45→0→30→37→16} liegt wirklich im
  Trägergraphen.
\item
  Für jeden der 60 nativen Kanäle verschwindet der gleichhändige
  Ein-Kopie-Weyl-Skalaransatz \codeword{M\_A⊗ε} als Grassmann-Bilinear.
\item
  Bei Dirac-Kernen sind C, Cγ⁵ und alle vier Cγ\^{}μγ⁵ antisymmetrisch;
  alle vier Cγ\^{}μ und alle sechs Cσ\^{}\{μν\} sind symmetrisch.
\item
  Die 43 getesteten nichttrivialen Gradierungen ergeben genau die
  angegebenen Kanalzahlen 24/36, 40/20 bzw. 32/28. Die
  Stabilisatordimensionen 36, 52 und 28 wurden diesmal durch exakt
  diagonale reelle Gram-Matrizen der Kommutatorbedingungen bestätigt,
  nicht durch einen SVD-Schwellwert.
\end{itemize}

Der Fünf-Zyklus allein schließt nur eine diagonale
±-Händigkeitszuweisung aus, die an jedem Paar entgegengesetzte
Händigkeiten verlangt. Er allein ist kein basisunabhängiger Ausschluss
und kein Ausschluss zusätzlicher Feldkomponenten.

\subsubsection{Eigener stärkerer Satz: auch ein Basiswechsel repariert
den Ein-Kopie-Vektoransatz
nicht}\label{eigener-stuxe4rkerer-satz-auch-ein-basiswechsel-repariert-den-ein-kopie-vektoransatz-nicht}

Sei Γ eine hermitesche Händigkeit auf dem ursprünglichen
64-dimensionalen internen Raum. Für einen ausschließlich
entgegengesetzt-händigen, unveränderten Paartensor müsste

\[\adjustbox{max width=\linewidth}{$\displaystyle 
\Gamma^T M_A+M_A\Gamma=0\quad\text{für alle }A
$}\]

gelten. Adjungieren und Einsetzen zeigt

\[\adjustbox{max width=\linewidth}{$\displaystyle 
[\Gamma,M_A^\dagger M_B]=0\quad\text{für alle }A,B.
$}\]

Die native Faktorisierung wurde exakt rekonstruiert:

\[\adjustbox{max width=\linewidth}{$\displaystyle 
M_{k,a}=S_k\otimes C_a,\quad
\sum_k S_k^\dagger S_k=5I_{16},\quad
\sum_a C_a^\dagger C_a=3I_4.
$}\]

Daher muss Γ sowohl mit allen \codeword{S\_k†S\_l⊗I₄} als auch mit allen
\codeword{I₁₆⊗C\_a†C\_b} kommutieren. Die von den ersten Produktmatrizen
erzeugte Algebra ist die volle M₁₆; die zweite ist M₄. Hierfür enthält
die Prüfung vollständige Rangzertifikate von ganzzahligen Matrixwörtern
modulo 101: \textbf{256 unabhängige Spinorwörter bis Länge zwei und 16
unabhängige Farbprodukte bis Länge eins}. Voller Rang modulo einer
Primzahl beweist vollen Rang über den komplexen Zahlen; es wird
ausdrücklich keine Rangdefizienz von einem endlichen Körper übertragen.

Folglich ist Γ skalar. Die ursprüngliche Gleichung erzwingt dann Γ=0.
Insbesondere gibt es \textbf{keine hermitesche Involution Γ²=I} mit der
verlangten Eigenschaft. Diese Schlusskette schließt den
Ein-Weyl-pro-Label-Vektoradapter jetzt unabhängig von einer Wahl der
internen Basis aus.

Der Satz betrifft unveränderte W-Kanäle und unveränderten
64-dimensionalen Träger. Er schließt \textbf{nicht} Diracfelder pro
Label, zusätzliche Kopien, Ableitungsadapter oder andere ausdrücklich
erweiterte Modelle aus.

\subsection{6. Die minimale Erweiterung darf nicht aus Versehen mit
ausgeschlossen
werden}\label{die-minimale-erweiterung-darf-nicht-aus-versehen-mit-ausgeschlossen-werden}

Zwei explizite Gegenkontrollen wurden auf allen 60 Kanälen neu
gerechnet:

\begin{enumerate}
\def\labelenumi{\arabic{enumi}.}
\tightlist
\item
  \textbf{Vierkomponenten-Diracfelder pro Label:} \codeword{M\_A⊗Cγ⁰}
  ist nichtverschwindend und antisymmetrisch im gesamten
  Grassmann-Index. Es enthält zusätzliche linke und rechte
  Freiheitsgrade pro ursprünglichem Label. Deshalb widerspricht es dem
  Ein-Kopie-Satz nicht. Eine chirale Standardmodell-Entstehung folgt
  daraus nicht.
\item
  \textbf{Unabhängige Zweier-Erweiterung:}
  \codeword{M\_A⊗ε\_aux⊗ε\_Lorentz} ist nichtverschwindend und
  antisymmetrisch. Der skalare Kanal wird möglich, weil die zusätzlich
  antisymmetrische interne Zweierform die Symmetrie des internen
  Paarfaktors umkehrt. Das ist die bereits bekannte minimale Reparatur
  im deklarierten Tensorprodukt-Ansatz, keine aus der Quelle
  hergeleitete Verdopplung.
\end{enumerate}

Eine solche Zweier-Erweiterung macht auch den Trägergraphen durch
\codeword{M\_A⊗σ\_x} bipartit: \codeword{Γ=I₆₄⊗diag(1,−1)}
antikommutiert exakt mit allen so erweiterten Matrizen. Damit ist
sichtbar, welche zusätzliche Ressource den ursprünglichen Ausschluss
überwindet.

``Der Vermittler kann nie Skalar sein'' muss folglich auf den
unveränderten Ein-Kopie-Ansatz begrenzt werden. Ebenso fixiert Schur nur
die interne Matrixstruktur innerhalb eines irreduziblen Blocks; es
beweist nicht die Einzigartigkeit aller möglichen Lorentz- und
Ableitungsterme. Eine gesunde Kinetik, Nebenbedingungen, Eichherkunft
und ein dynamischer Spin-2-Sektor sind hier nicht konstruiert.

\subsection{Reproduktion und Status}\label{reproduktion-und-status}

\codeword{contracted\_krylov.cpp} ist der unabhängige ganzzahlige
Kontraktionskern. \codeword{verify.py} prüft diesen Kern, die
korrigierte Pythagoras-Rechnung, kleine Ritz-Matrizen, sämtliche
angeführten Graph-, Kern- und Gradierungsidentitäten sowie die
Matrixalgebra-Zertifikate. Die Ergebnisse stehen in
\codeword{audit.json} und \codeword{audit\_optimized.json}.

Es bestehen \textbf{376 Prüfbedingungen: 371 exakte und fünf
numerische}. Für eine eigenständige Reproduktion zuerst den kleinen
C++-Kern mit C++17 und Optimierung übersetzen; anschließend
\codeword{verify.py} normal und mit \codeword{-OO} starten. Sämtliche
Eingaben liegen eingefroren unter \codeword{sources/}; die
Ergebnisdateien enthalten ihre SHA-256-Werte. Fremde Pythonprogramme
werden hierbei nicht ausgeführt. Die numerischen Ritz-Prüfungen sind als
solche markiert, alle Algebra-Zertifikate und die vollständige direkte
Kontraktion verwenden exakte Ganzzahl-/Bruchrechnung.

Beide abschließenden Ausführungen bestanden und lieferten byteidentische
JSON-Dateien mit SHA-256
\codeword{148047028236a887de1b28a0a9c12d84037a7c633892abc51dd9047f0b7a7baf}.

Keine T1-T8-Abschlussbehauptung. Keine Änderungen an fremden Quellen,
Hauptpaper, Webseite, Ledger oder Git-Historie.

\begin{center}\rule{0.5\linewidth}{0.5pt}\end{center}

\section{Anhang C: minimaler Transfer und seine
Quellen-Grenze}\label{anhang-c-minimaler-transfer-und-seine-quellen-grenze}

\section{Minimaler Ursprungsstrang: geteilte Moden und ein echter
nativer
Vierzustandskanal}\label{minimaler-ursprungsstrang-geteilte-moden-und-ein-echter-nativer-vierzustandskanal}

Stand: 15. September 2026. Unabhängiger, begrenzter Forschungsstrang.

\subsection{Ergebnis und Reichweite}\label{ergebnis-und-reichweite}

Es gibt zwei verschiedene Ergebnisse, die ausdrücklich nicht verwechselt
werden dürfen:

\begin{enumerate}
\def\labelenumi{\arabic{enumi}.}
\tightlist
\item
  \textbf{Ein exakt lösbarer Überlappungsbaustein:} Eine gemeinsame
  Fermionmode und ein gemeinsamer Bosonkanal werden in jedem festen
  positiven Ressourcensektor zu einer gewöhnlichen effektiven
  Fermionmode. Dafür ist keine direkte Endpunkt-Hopping-Wechselwirkung
  nötig. Dieser Dreierstern ist jedoch \textbf{kein Ausschnitt einer
  tatsächlichen nativen W-Zeile}.
\item
  \textbf{Ein tatsächlicher W-basierter Vierzustandskanal:} Im
  unveränderten 64-Fermion/60-Boson-Paartensor existiert ein exakt
  invariantes Vierzustandsystem, wenn ein \textbf{zusätzlicher reiner
  Bosonmischer} zwischen Kanal 0 und 1 zugelassen und ein bestimmter
  Zustand mit Gesamtladung 9 präpariert wird. Der vollständige native
  Hamiltonoperator plus dieser Mischer überträgt dort eine innere
  Fermionmarke mit einer streng abgesicherten Wahrscheinlichkeit von
  \textbf{mehr als 99,3 \%}.
\end{enumerate}

Das zweite Ergebnis benötigt keine Projektion, die die übrigen nativen
Zustände künstlich wegschneidet: Der Vierzustandsraum ist unter allen 60
ursprünglichen Paaroperatoren plus dem Mischer exakt invariant. Es ist
aber \textbf{weder eine Untersuchung der Entnahmeantwort des
N=64-Grundzustands noch ein Transport zwischen zwei unabhängigen
Banken}. Die Herkunft des Bosonmischers, die spezielle Präparation und
die räumliche Bedeutung der Marken bleiben offen.

\subsection{1. Warum eine Clockphase die bisherige Paritätsschranke
nicht
aufhebt}\label{warum-eine-clockphase-die-bisherige-parituxe4tsschranke-nicht-aufhebt}

Seien zwei operational unabhängig definierte Banken mit lokalen
Paritäten

\[\adjustbox{max width=\linewidth}{$\displaystyle 
\Pi_x=(-1)^{N_{f,x}},\qquad \Pi_y=(-1)^{N_{f,y}}
$}\]

gegeben. Kommutiert jede ausführbare lokale Operation und jeder einzelne
realisierte Messzweig mit beiden Paritäten, gilt dies auch für beliebige
Produkte, Summen, Adjungierte und starke beschränkte Grenzwerte. Eine
Clockoperation, die lediglich innerhalb jeder Bank Fermionmoden
zahlenerhaltend permutiert, bleibt in dieser Algebra.

Ein echter Einfermiontransfer erfüllt dagegen

\[\adjustbox{max width=\linewidth}{$\displaystyle 
\Pi_x f_y^\dagger f_x=-f_y^\dagger f_x\Pi_x,
\qquad
\Pi_y f_y^\dagger f_x=-f_y^\dagger f_x\Pi_y.
$}\]

Er kann daher nicht allein durch mehr Kompositionen jener Operationen
entstehen. Eine bloß paritätskovariante CP-Abbildung ist nicht die
gleiche Voraussetzung: Sie kann ungerade Krauszweige besitzen. Die
Aussage gilt nur für den ausdrücklich geraden ausführbaren
Operationssatz.

Bei geteilten Fermionmoden sind die beiden lokalen CAR-Algebren von
Anfang an nicht unabhängig. Liegt dieselbe Mode s in beiden, kann sie
nicht zugleich als zwei unabhängige antikommutierende Kopien behandelt
werden: \(\{s,s^\dagger\}=1\). Auch lokale Ladungen, die s doppelt
mitzählen, sind keine unabhängigen additiven Bankladungen. Überlappung
widerlegt somit den Paritätssatz nicht, sondern ändert seine
Voraussetzungen.

\subsection{2. Exakter CAR-Baustein aus einem geteilten Boson und
Fermion}\label{exakter-car-baustein-aus-einem-geteilten-boson-und-fermion}

Betrachtet werden drei CAR-Moden \(f_x,s,f_y\), ein CCR-Boson b und

\[\adjustbox{max width=\linewidth}{$\displaystyle 
Q_x=b^\dagger s f_x,\qquad Q_y=b^\dagger s f_y,
\qquad
H_\star=\Delta N_b+g(Q_x+Q_y+Q_x^\dagger+Q_y^\dagger).
$}\]

Der Ressourcenzähler

\[\adjustbox{max width=\linewidth}{$\displaystyle 
K=N_b+n_s
$}\]

kommutiert mit allen diesen Operationen. Auf jedem exakten K-Sektor mit
ganzzahligem \(K\ge1\) setze

\[\adjustbox{max width=\linewidth}{$\displaystyle 
c^\dagger=\frac{b^\dagger s}{\sqrt K},\qquad
c=\frac{s^\dagger b}{\sqrt K}.
$}\]

Direkt aus CCR und CAR folgen

\[\adjustbox{max width=\linewidth}{$\displaystyle 
(c^\dagger)^2=c^2=0,\qquad
\{c,c^\dagger\}=\frac{N_b+n_s}{K}=1.
$}\]

Die Mode c antikommutiert mit beiden Endpunktmoden. Ferner

\[\adjustbox{max width=\linewidth}{$\displaystyle 
N_b=K-1+n_c,
\qquad
N_f+2N_b=2K-1+n_x+n_c+n_y.
$}\]

Damit gilt exakt

\[\adjustbox{max width=\linewidth}{$\displaystyle 
\boxed{
H_\star=\Delta(K-1)+\Delta n_c
 +g\sqrt K\big(c^\dagger f_x+c^\dagger f_y+\mathrm{h.c.}\big).
}
$}\]

Die gewöhnliche Drei-Moden-Transferkette wurde hier aus einer
überlappenden Paarumwandlung erhalten. Ein nützlicher Operatorausdruck
derselben Tatsache ist

\[\adjustbox{max width=\linewidth}{$\displaystyle 
\boxed{[Q_y^\dagger,Q_x]=K f_y^\dagger f_x.}
$}\]

Das ist zunächst eine Operatoridentität. Die Verfügbarkeit getrennter
Kommutator-Kontrollsequenzen folgt daraus nicht automatisch. Schon der
konstante Summen-Hamiltonoperator zeigt jedoch Transfer. Der minimale
positive Ressourcensektor K=1 benötigt am Anfang eine besetzte
gemeinsame Fermionmode s und ein leeres Boson.

\subsubsection{Vollständige Transferrechnung im
Vergleichsmodell}\label{vollstuxe4ndige-transferrechnung-im-vergleichsmodell}

Im effektiven Einteilchenraum und bei K=1 lautet die Matrix

\[\adjustbox{max width=\linewidth}{$\displaystyle 
\begin{pmatrix}0&g&0\\g&\Delta&g\\0&g&0\end{pmatrix}.
$}\]

Der antisymmetrische Endpunktzustand ist dunkel. Der symmetrische
koppelt mit \(\sqrt2g\) an c.~Für \(g/\Delta=1/20\), \(m=101\) und

\[\adjustbox{max width=\linewidth}{$\displaystyle 
t=\frac{202\pi}{\Delta\sqrt{51/50}}
$}\]

verschwindet die mittlere Besetzung wieder exakt. Die
Zielwahrscheinlichkeit ist

\[\adjustbox{max width=\linewidth}{$\displaystyle 
p_\star=\sin^2\left[\frac\pi2\,101\left(1-\sqrt{50/51}\right)\right].
$}\]

Durch Quadrieren rationaler Schranken erhält man

\[\adjustbox{max width=\linewidth}{$\displaystyle 
\frac{199}{200}<101\left(1-\sqrt{50/51}\right)<1.
$}\]

Mit \(|\sin u|\le|u|\) und \(\pi<355/113\) folgt

\[\adjustbox{max width=\linewidth}{$\displaystyle 
p_\star>1-\left(\frac{355/113}{400}\right)^2>0.9999.
$}\]

Diese Zahl gehört \textbf{nur zum Vergleichsmodell mit gemeinsamem b}.

\subsubsection{Die direkte
Quellen-Grenze}\label{die-direkte-quellen-grenze}

Die tatsächliche native W-Zeile enthält jeweils acht disjunkte
Fermionpaare, also 16 verschiedene Marken. Eine einzelne Zeile enthält
deshalb niemals zugleich \(s f_x\) und \(s f_y\) mit \(x\ne y\). Das ist
an allen 60 Zeilen exakt geprüft.

Ein gemeinsamer-b-Dreierstern darf somit nicht allein wegen der Form
\(b^\dagger f f\) als nativer Baustein bezeichnet werden. Insbesondere
wäre ein Wechsel von zwei Bosonkanälen zu einer gemeinsamen Mode ohne
Kontrolle des orthogonalen Kanals eine zusätzliche Modellannahme.

\subsection{3. Exakter Vierzustandskanal mit dem vollständigen nativen
W}\label{exakter-vierzustandskanal-mit-dem-vollstuxe4ndigen-nativen-w}

Die gepinnte Quelle ist

\codeword{universalraum-native-ground-response-20260915/ground\_replay/outputs/simple\_core/spinor\_tensors.npz},

SHA-256
\codeword{3f00a0892157a691e4ef26920c3702772c7bf76a250fad98ff04a079bb64b763}.

Die zwei relevanten Quellenzeilen sind

\[\adjustbox{max width=\linewidth}{$\displaystyle 
\begin{aligned}
P_0={}&-f_{57}f_4+f_{56}f_5+f_{53}f_8-f_{52}f_9
       -f_{45}f_{16}+f_{44}f_{17}+f_{33}f_{28}-f_{32}f_{29},\\
P_1={}&-f_{58}f_4+f_{56}f_6+f_{54}f_8-f_{52}f_{10}
       -f_{46}f_{16}+f_{44}f_{18}+f_{34}f_{28}-f_{32}f_{30}.
\end{aligned}
$}\]

Die Konvention ist \(P_A=\sum_{i<j}W_{A,ij}f_jf_i\). Das gemeinsame
aktive Fermion ist s=4; die Endpunktmarken sind 57 und 58.

Präpariert werden sieben unveränderte Pauli-Blocker

\[\adjustbox{max width=\linewidth}{$\displaystyle 
S=\{8,16,28,32,44,52,56\}.
$}\]

Sei \(|S\rangle\) der geordnet erzeugte Fockzustand dieser Moden. Die
vier Basiszustände sind

\[\adjustbox{max width=\linewidth}{$\displaystyle 
\begin{array}{c|c|c}
\text{Zustand}&\text{besetzte Fermionmoden}&\text{Bosonen}\\\hline
L&S\cup\{4,57\}&0\\
B_0&S&1\text{ in Kanal }0\\
B_1&S&1\text{ in Kanal }1\\
R&S\cup\{4,58\}&0
\end{array}
$}\]

Jeder Zustand besitzt die native Gesamtladung \(N=N_f+2N_b=9\).

Zugelassen wird nun genau eine zusätzliche Hamiltonoperation,

\[\adjustbox{max width=\linewidth}{$\displaystyle 
H_J=J(b_1^\dagger b_0+b_0^\dagger b_1).
$}\]

Alle ursprünglichen 60 nativen Paaroperatoren bleiben erhalten:

\[\adjustbox{max width=\linewidth}{$\displaystyle 
H=\Delta N_b+g\sum_{A=0}^{59}(b_A^\dagger P_A+P_A^\dagger b_A)+H_J.
$}\]

Der Prüfer wendet sämtliche 480 signierten Paarkanäle einschließlich
aller CAR-Vorzeichen und der Bosonfaktoren auf jeden der vier Zustände
an. Es entstehen \textbf{keine Zustände außerhalb ihres linearen
Spanns}. Die Pauli-Blocker verhindern insbesondere die sieben
unerwünschten Rückpaarungen jedes aktiven Bosonkanals.

Es ergibt sich exakt

\[\adjustbox{max width=\linewidth}{$\displaystyle 
\boxed{
H_4=\begin{pmatrix}
0&g&0&0\\
g&\Delta&J&0\\
0&J&\Delta&g\\
0&0&g&0
\end{pmatrix}_{(L,B_0,B_1,R)}.
}
$}\]

Ohne Mischer, J=0, zerfällt die Matrix in zwei getrennte 2×2-Blöcke.
Dann ist \(\langle R|e^{-itH}|L\rangle\) zu jeder Zeit exakt null. Mit J
ungleich null beginnt die Transferamplitude bei dritter Ordnung; die
Wahrscheinlichkeit hat den führenden Term \(g^4J^2t^6/36\).

\textbf{Die Quelle enthält hier also die beiden Endstücke und eine exakt
funktionierende Pauli-Blockierung, aber nicht die geprüfte
Bosonverbindung als bereits verfügbare Operation.}

\subsection{4. Strenge vollständige Transfergrenze im tatsächlichen
W-Zeugen}\label{strenge-vollstuxe4ndige-transfergrenze-im-tatsuxe4chlichen-w-zeugen}

Unter Spiegelung L↔R und B₀↔B₁ zerfällt H₄ in

\[\adjustbox{max width=\linewidth}{$\displaystyle 
H_+=\begin{pmatrix}0&g\\g&\Delta+J\end{pmatrix},\qquad
H_-=\begin{pmatrix}0&g\\g&\Delta-J\end{pmatrix}.
$}\]

Wähle den vorhandenen Prüfquotienten \(g/\Delta=1/20\) und die
ausdrücklich neu gesetzte Mischerstärke

\[\adjustbox{max width=\linewidth}{$\displaystyle 
J=\Delta-\frac{2g}{\sqrt3}
 =\Delta\left(1-\frac1{10\sqrt3}\right),
\qquad t=\frac{\pi\sqrt3}{g}=\frac{20\pi\sqrt3}{\Delta}.
$}\]

Da \(0<J<\Delta\), bleiben die Eigenfrequenzen des Bosonmischers
\(\Delta\pm J\) positiv. Es wurde \textbf{keine Nullfrequenz durch
exakte Aufhebung von \(\Delta N_b\)} erzeugt. Auch auf dem ganzen
Fockraum bleibt die hinzugefügte Hamiltonfamilie nach unten
beschränkbar: Das Bosonquadrat ist strikt positiv und die endlichen
Fermionoperatoren können durch quadratische Ergänzung kontrolliert
werden. Das bedeutet nicht, dass ihr Grundzustand noch der unveränderte
native Grundzustand ist.

Für H₋ ist \(\Delta-J=2g/\sqrt3\). Zu der gewählten Zeit kehrt dessen
Endpunktkomponente exakt mit Phase −1 zurück, ohne mittlere
Restbesetzung.

Für H₊ setze

\[\adjustbox{max width=\linewidth}{$\displaystyle 
a=\Delta-\frac g{\sqrt3},\quad
\omega=\sqrt{a^2+g^2},\quad
w_- =\frac{1-a/\omega}{2}.
$}\]

Die Endpunktamplitude lautet

\[\adjustbox{max width=\linewidth}{$\displaystyle 
A_+=(1-w_-)e^{i(\omega-a)t}+w_-e^{-i(\omega+a)t}.
$}\]

Mit \(\omega-a\le g^2/(2a)\), \(w_-\le g^2/(4a^2)\),
\(\cos u\ge1-u^2/2\) und \(|A_+|^2\ge1-g^2/a^2\) folgt

\[\adjustbox{max width=\linewidth}{$\displaystyle 
\begin{aligned}
p_{L\to R}
 &=\left|\frac{A_++1}{2}\right|^2\\
 &\ge1-\frac{g^4t^2}{16a^2}-\frac{g^2}{2a^2}\\
 &=1-\frac{3\pi^2+8}{6400(a/\Delta)^2}.
\end{aligned}
$}\]

Aus \(\sqrt3>17/10\) folgt \(a/\Delta>33/34\). Somit ergibt sich rein
rational

\[\adjustbox{max width=\linewidth}{$\displaystyle 
\boxed{
p_{L\to R}>
1-\frac{3(355/113)^2+8}{6400(33/34)^2}
=\frac{2009992727}{2022609600}
>0.9937620819>0.993.
}
$}\]

Die Aussage ist eine \textbf{analytische Schranke der vollständigen
Dynamik auf einem exakt invarianten Quellen-Unterraum}. Es gibt hier
keinen numerisch abgeschnittenen Restzustandsraum. Die
64-Fermion/60-Boson-Quelle ist nicht durch vier frei erfundene
Matrixeinträge ersetzt worden: Ihre exakte Einschränkung wurde zuerst
vollständig geprüft.

Die Zeit in Einheiten \(\hbar=1\) ist ungefähr \(108.83/\Delta\). Diese
numerische Orientierung und die gesetzte Mischerstärke sind keine
vorhergesagten Naturkonstanten.

\subsection{5. Was dies für die fundamentale Suche
ändert}\label{was-dies-fuxfcr-die-fundamentale-suche-uxe4ndert}

Der konstruktive Anschluss ist klein: \textbf{Paarumwandlung → geteilte
Zwischenressource → Paar-Rückumwandlung} kann einen
Einfermion-Endpunktwechsel bewirken. Für die tatsächliche W-Quelle
braucht man dazu weder einen neuen direkten Fermion-Hopping-Term noch
die Kontrolle jedes einzelnen der 480 Paarmonome. Ein einzelner
Bosonmischer plus eine spezielle Pauli-blockierte Präparation reicht im
belegten internen Zeugen.

Offen bleiben aber genau die folgenden Herkunftsfragen:

\begin{enumerate}
\def\labelenumi{\arabic{enumi}.}
\tightlist
\item
  \textbf{Operationssatz:} Ist der reine Mischer
  \(b_1^\dagger b_0+\mathrm{h.c.}\) tatsächlich verfügbar? Eine
  simultane Fermion-und-Boson-Symmetrie ist nicht automatisch eine
  unabhängige Bosonoperation.
\item
  \textbf{Präparation:} Wie wird der N=9-Blockerzustand mit dem
  zugelassenen Operationssatz erzeugt? Ladungserhaltende native
  Evolution präpariert ihn nicht aus dem N=64-Grundzustand.
\item
  \textbf{Operationaler Raum:} Sind die Endpunktmarken 57 und 58
  überhaupt verschiedene räumliche Teile oder nur interne Marken
  derselben Bank? Die Rechnung allein liefert keine Raumposition.
\item
  \textbf{Gemeinsamer Grundzustand:} Besteht ein entsprechender
  Mechanismus für die ursprüngliche Entnahmeantwort des nativen
  N=64-Grundzustands, statt für diesen besonders präparierten Zeugen?
\item
  \textbf{Zwei Banken:} Eine Bosonverbindung zwischen wirklich
  unabhängigen Banken erhält deren lokale Fermionparitäten. Dieser
  interne Zeuge hebt jene Schranke nicht auf. Dafür wäre eine
  ursprünglich geteilte Fermionstruktur oder eine andere insgesamt
  gerade, lokal ungerade Quelloperation zu konstruieren.
\end{enumerate}

T1-T8, das relativistische Feldwörterbuch und die Herkunft der Raumzeit
sind damit nicht geschlossen. Das Ergebnis verkleinert eine konkrete
Suche: Statt eines völlig beliebigen neuen Fermionlinks kann jetzt ein
bestimmter Bosonmischer mitsamt seinem Quellen- und Präparationsvertrag
geprüft werden. Die nachfolgende Quellenprüfung weist diesen Mischer
jedoch ausdrücklich \textbf{nicht} als Synthese der bisher gewährten
Kontrollen aus.

\subsection{6. Anschluss des neuen Chart-/Glue-Vorschlags und exakter
Quellen-No-go}\label{anschluss-des-neuen-chart-glue-vorschlags-und-exakter-quellen-no-go}

Der neue Nutzeranhang vom 15. September 2026 wurde vollständig gelesen:

\codeword{/Users/stefanhamann/.codex/attachments/59dc0059-8914-48ca-953d-85933f66e00b/pasted-text.txt},

SHA-256
\codeword{1a75e84b28868dd682e4b2337d6546f275f49321e489b3b23af4754a844d1ac4}.

Die Hypothese, Banken zunächst als überlappende lokale Beschreibungen
einer gemeinsamen CAR-Struktur zu untersuchen, passt zum hier
konstruierten gemeinsamen Fermion-/Boson-Zwischenraum. \textbf{Sie löst
den Übergang von Überlappung zu Dynamik aber nicht automatisch.} Bereits
zwei nichtorthogonale lokale Moden \(f_A=f_1\),
\(f_B=\cos\theta f_1+\sin\theta f_2\) bei H=0 besitzen ein
nichtverschwindendes Kreuz-Antikommutator
\(\{f_A,f_B^\dagger\}=\cos\theta\), obwohl überhaupt keine
Zustandsentwicklung stattfindet. Eine Änderung der Beschreibung und ein
dynamischer Transfer müssen getrennt geprüft werden.

\subsubsection{6.1 Eine nötige Korrektur des vorgeschlagenen
Paritätstests}\label{eine-nuxf6tige-korrektur-des-vorgeschlagenen-parituxe4tstests}

Der Anhang fordert für überlappende Charts einen Intertwiner, der lokal
Paritäten ändert, aber mit \(\Pi_A\Pi_B\) kommutiert. \textbf{Bei
Überlappung ist dieses Produkt nicht automatisch die Gesamtparität.} Im
kleinsten gemeinsamen-Moden-Zeugen

\[\adjustbox{max width=\linewidth}{$\displaystyle 
A=\{x,s\},\qquad B=\{s,y\}
$}\]

gilt

\[\adjustbox{max width=\linewidth}{$\displaystyle 
\Pi_A\Pi_B=(-1)^{n_x+n_y},
\qquad
\Pi_{\rm global}=(-1)^{n_x+n_s+n_y}.
$}\]

Die gemeinsame Mode s wurde im Produkt zweimal gezählt und fällt heraus.
Die Paarumwandlung \(b^\dagger s f_x\) ist bezüglich
\(\Pi_{\rm global}\) gerade, kommutiert aber nicht mit \(\Pi_A\Pi_B\).
Der vorgeschlagene Kill-Test würde hier einen legitimen global geraden
Überlappungsmechanismus fälschlich aussortieren. Alle vier Identitäten
wurden auf dem vollständigen Drei-Fermion-CAR-Raum exakt geprüft.

\textbf{Der korrigierte Test muss die globale Parität aus der
gemeinsamen CAR-Darstellung selbst verwenden}, nicht aus einem
ungeprüften Produkt lokaler Paritäten. Zusätzlich benötigt er eine
definierte Anfangspräparation, dynamisch unterschiedliche
Endpunktbeobachtungen und einen aus der Quelle stammenden Generator.
Passive Chartwechsel allein reichen nicht.

\subsubsection{6.2 Tatsächlicher endlicher Clock statt unterstellter
voller
Gruppe}\label{tatsuxe4chlicher-endlicher-clock-statt-unterstellter-voller-gruppe}

Die Prüfung lädt den gepinnten ursprünglichen Clock-Konstruktor über

\codeword{sources/repo/experiments/theory-contracts/universalraum-native-operations-ground-response-20260915/common.py},

SHA-256
\codeword{2cc97522457ecc7e6774d96a25b2581b5d67cd051ae35e28a26f0e98e2adb994}.

Der Konstruktor kontrolliert seine ursprünglichen Quellenpins und die
exakte W-Kovarianz. Er liefert tatsächlich

\[\adjustbox{max width=\linewidth}{$\displaystyle 
p=(2,0,1,4,3),\qquad s_{\rm Clock}=-1,
$}\]

und auf Bosonmoden \(A=6k+c\)

\[\adjustbox{max width=\linewidth}{$\displaystyle 
G_B e_{6k+c}=-e_{6(p(k\bmod5)+5\lfloor k/5\rfloor)+c}.
$}\]

Es wird \textbf{nicht} vorausgesetzt, dass der Clock der volle
Spin(10)×SU(4)-Operationssatz ist. Sein endlicher, direkt reproduzierter
Lift reicht für den folgenden Ausschluss.

Schreibe \(M_{ij}=|i\rangle\langle j|+|j\rangle\langle i|\) auf dem
Boson-Einteilchenraum. Exakt gilt

\[\adjustbox{max width=\linewidth}{$\displaystyle 
G_BM_{01}G_B^\dagger=M_{12,13},\quad
G_BM_{12,13}G_B^\dagger=M_{6,7},\quad
G_BM_{6,7}G_B^\dagger=M_{01}.
$}\]

Insbesondere

\[\adjustbox{max width=\linewidth}{$\displaystyle 
\|[G_B,M_{01}]\|_F^2=4\ne0.
$}\]

Die native W-Kovarianz impliziert \([G,H]=[G,N_b]=0\) für den
gemeinsamen Fock-Lift G. Auch ein gewährter getrennter Gesamt-Casimir
von Spin(10) oder SU(4) kommutiert mit diesem tatsächlichen G. Daher
liegt

\[\adjustbox{max width=\linewidth}{$\displaystyle 
\operatorname{Alg}(H,N_b,G,C_{\rm Spin},C_{\rm Color})\subset\{G\}'
$}\]

und der spezifische Einzelmischer \(b_1^\dagger b_0+\mathrm{h.c.}\)
liegt \textbf{nicht} in dieser erzeugten Algebra. Auf den endlichen
Ladungssektoren gibt es dabei keine Domänenprobleme; für die volle
Hilbert-Darstellung formuliert man dieselbe Aussage mit den erzeugten
beschränkten Zeitentwicklungen und Spektraloperationen.

\subsubsection{6.3 Clock-Mittelung ist möglich, aber der Farb-Cartan
sperrt sie
weiterhin}\label{clock-mittelung-ist-muxf6glich-aber-der-farb-cartan-sperrt-sie-weiterhin}

Der Clock-Ausschluss allein wäre zu schwach. Die Orbitsumme

\[\adjustbox{max width=\linewidth}{$\displaystyle 
M_{\rm orb}=M_{01}+M_{12,13}+M_{6,7}
$}\]

kommutiert exakt mit \(G_B\). Ihre beiden zusätzlichen Linkblöcke wirken
auf den vier Zuständen aus Abschnitt 3 null. Ein neu zugelassener
globaler Bosonmischer \(J b^\dagger M_{\rm orb}b\) erzeugt deshalb
\textbf{denselben} exakt invarianten Vierzustandskanal und dieselbe
Transfergrenze. Ein neuer globaler Zusammenhang muss also nicht zwingend
den endlichen Clock brechen.

Es gibt jedoch einen zweiten, stärkeren gemeinsamen Erhaltungssatz. Aus
dem tatsächlichen Quell-Gewichtswörterbuch wählen wir die zweite
SU(4)-Cartankomponente und definieren den gemeinsamen Ladungsoperator

\[\adjustbox{max width=\linewidth}{$\displaystyle 
Q=\sum_r q_r n_{f,r}+\sum_A q_A n_{b,A}.
$}\]

Jeder ursprüngliche Paarterm erfüllt exakt \(q_i+q_j=q_A\); dies wird
für alle 480 Quellenpaare geprüft. Der tatsächliche Clock wirkt auf dem
Spinindex und lässt die SU(4)-Farbkomponente unverändert. Daher

\[\adjustbox{max width=\linewidth}{$\displaystyle 
[Q,H]=[Q,N_b]=[Q,G]=0.
$}\]

Die getrennten Gesamt-Casimire kommutieren ebenfalls mit diesem Cartan.
Diese Aussage erfordert \textbf{keine} Verfügbarkeit aktiver
kontinuierlicher SU(4)-Kontrollen.

Auf den beiden Bosonkanälen gilt dagegen

\[\adjustbox{max width=\linewidth}{$\displaystyle 
q_0=0,\qquad q_1=2,
$}\]

und auf dem ganzen Boson-Einteilchenraum

\[\adjustbox{max width=\linewidth}{$\displaystyle 
\|[Q_B,M_{01}]\|_F^2=8,\qquad
\|[Q_B,M_{\rm orb}]\|_F^2=24.
$}\]

\textbf{Damit sind sowohl der Einzelmischer als auch seine
Clock-invariante Orbitsumme aus dem genannten Operationssatz
ausgeschlossen.} Die Mittelung beseitigt das Clock-Hindernis, nicht den
separaten Farb-Cartan-Erhaltungssatz.

Die tatsächlichen Cartanladungen im Vierzustandsraum sind

\[\adjustbox{max width=\linewidth}{$\displaystyle 
(Q_L,Q_{B_0},Q_{B_1},Q_R)=(7,7,9,9).
$}\]

Dieser Zeuge verändert also eine innere SU(4)-Quantenzahl. Er ist noch
deutlicher als ein bloßer Markenwechsel vom Transport \textbf{derselben}
niedrigenergetischen Fermionmode zwischen zwei räumlichen Banken zu
unterscheiden.

\subsubsection{6.4 Kleinste noch zu suchende
Quellressource}\label{kleinste-noch-zu-suchende-quellressource}

Für den hier exakt spezifizierten Zeugen fehlt mindestens eine
Operation, die nicht mit dem betrachteten Farb-Cartan des ursprünglichen
Systems kommutiert, oder eine aus einer erweiterten globalen Quelle
abgeleitete Kopplung mit einer \textbf{expliziten Ausgleichsressource
für diese Cartanladung}. Ein zusätzlicher aktiver SU(4)-Generator wäre
ein neuer Operationsvertrag; außerdem bewegt sein simultaner
Fermion-/Boson-Lift im Allgemeinen auch die Blocker und ist nicht mit
dem reinen Bosonmischer gleichzusetzen.

Die Befunde schließen nicht alle globalen Glue-Modelle aus. Sie zeigen
präzise, warum der konkret schon gelöste Vierzustandstransfer noch keine
Herkunftslösung ist und welche neue algebraische Eigenschaft eine echte
Quellenfortsetzung besitzen müsste. Für eine überwiegend kinematische
Chart-Interpretation muss zusätzlich gezeigt werden, dass diese
Erweiterung eine tatsächliche Dynamik und nicht bloß eine
Basisumbenennung erzeugt.

\subsection{Reproduktion und Status}\label{reproduktion-und-status-1}

\codeword{verify.py} prüft \textbf{307 Bedingungen}, verwendet keine
Python-Assertions und reproduziert normal sowie mit \codeword{-OO}
byteidentische JSON-Ergebnisse. Die Prüfung enthält:

\begin{itemize}
\tightlist
\item
  CAR und Ladungsbuchhaltung für den geteilten Baustein in K=1,2,3;
\item
  die rationale Transfergrenze des klar getrennten Vergleichssterns;
\item
  die Matching-Eigenschaft aller 60 nativen W-Zeilen;
\item
  die volle symbolische Wirkung aller 480 nativen Paarterme auf alle
  vier Zeugen-Zustände;
\item
  die zusätzliche Mischerwirkung, exakte Invarianz und vollständige
  Spiegelungszerlegung;
\item
  die rationale Schranke größer 99,3 \% bei strikt positiver
  Bosonfrequenz.
\item
  den tatsächlich konstruierten Clock, den Einzelmischer-Ausschluss und
  die Clock-invariante Orbitsumme;
\item
  die gemeinsame Farb-Cartan-Erhaltung aller Quellenpaare und den
  präzisen Ausschluss beider Mischinstrumente;
\item
  das Gegenbeispiel gegen die Gleichsetzung lokaler Paritätsprodukte mit
  der globalen Parität bei überlappenden Charts.
\end{itemize}

Die analytischen Ungleichungen sind im Text hergeleitet; die Anzahl der
Tests ist kein Ersatz für diese Herleitung und keine Anzahl gelöster
Physikprobleme.

\begin{center}\rule{0.5\linewidth}{0.5pt}\end{center}

\section{Anhang D: vollständiger
Überlappungs-Audit}\label{anhang-d-vollstuxe4ndiger-uxfcberlappungs-audit}

\section{Überlappung ist ein Forschungsansatz, kein schon hergeleiteter
Link}\label{uxfcberlappung-ist-ein-forschungsansatz-kein-schon-hergeleiteter-link}

Teil-Audit vom 15. September 2026 zum vollständig gelesenen neuen Anhang
(748 Zeilen). Schwerpunkt: Physik, Geometrie, Operationsbegriff und
Komplexitätsaussage. Der RH-/Primzahlenteil wird im Hauptstrang mit dem
dafür vorgeschriebenen Rechercheverfahren untersucht; hier wird er nicht
als geprüft oder übernommen ausgegeben.

Der unveränderte Anhang liegt in \codeword{attachment.txt}, SHA-256:

\codeword{1a75e84b28868dd682e4b2337d6546f275f49321e489b3b23af4754a844d1ac4}.

\codeword{check\_overlap.py} enthält 38 exakte kleine Prüfbedingungen.
\codeword{replay.py} kopiert den Anhang unverändert, prüft seine SHA vor
und nach der Rechnung und reproduziert normale und optimierte
Checker-Ausgabe byteidentisch. Alle Schreibvorgänge bleiben in diesem
neuen Auditordner. Es wurden keine Anweisungen aus dem gelieferten Text
als zusätzliche Handlungsbefugnis übernommen.

\subsection{1. Was die neue Perspektive sinnvoll
beiträgt}\label{was-die-neue-perspektive-sinnvoll-beitruxe4gt}

Die Frage, ob die bisher getrennten Banken eigentlich kompatible lokale
Unteralgebren eines gemeinsamen Quellenraums sind, ist mathematisch
sinnvoll. Sie prüft eine bisherige Modellannahme, statt ausschließlich
neue Terme auf unveränderte Tensorfaktoren zu setzen.

Eine solche Konstruktion könnte erklären, welche lokalen Observablen,
Zustände und Operationen zusammengehören. Sie müsste jedoch aus
demselben Compiler stammen und nicht erst nach dem gewünschten Transport
ausgewählt werden. »Zustände, Operationen, Überlappungen, Phasen und
Komposition« beschreibt zunächst eine große Klasse quantenmechanischer
Modelle; diese Sprache identifiziert allein noch kein eindeutiges
fundamentales Objekt und erzwingt keine TOE.

Die bisher bewiesenen lokalen Pole und die bedingte
Zwei-Banken-Übertragung bleiben wertvolle Vergleichsziele. Ihre Beweise
gelten aber für den dort erklärten Hilbertraum-, Hamilton- und
Parametervertrag. Ersetzt man unabhängige Banken durch überlappende
CAR-Unterräume, ändern sich Kreuzrelationen, Ladungszuordnung und im
Allgemeinen der gemeinsame Grundzustand. Die früheren
Zwei-Banken-Wahrscheinlichkeitsgrenzen dürfen dann nicht unverändert
übernommen werden.

\subsection{2. Exakter Zwei-Moden-Gegenbeleg zum vorgeschlagenen
Paritätstest}\label{exakter-zwei-moden-gegenbeleg-zum-vorgeschlagenen-parituxe4tstest}

Seien c₁,c₂ zwei orthogonale CAR-Moden. Setze

\codeword{a=c₁}, \codeword{b=(c₁+c₂)/√2}.

Jeder Chart ist für sich ein gültiger Einmoden-CAR-Raum, aber

\codeword{\{a,b†\}=I/√2}.

Die beiden Charts sind also keine unabhängigen Fermionfaktoren.
Definiere ihre lokalen Paritäten

\codeword{Π\_A=I−2a†a}, \codeword{Π\_B=I−2b†b}.

In der Besetzungsbasis \codeword{|00>,|10>,|01>,|11>} ist die wirkliche
globale Parität

\codeword{Π\_global=diag(1,−1,−1,1)}.

Das Produkt der lokalen Paritäten lautet dagegen

\begin{verbatim}
Π_A Π_B = [[1, 0, 0, 0],
           [0, 0, 1, 0],
           [0,-1, 0, 0],
           [0, 0, 0, 1]].
\end{verbatim}

Es ist weder die globale Parität noch hermitesch, und sein Quadrat ist
nicht die Identität. Auch \codeword{[Π\_A,Π\_B]≠0}.

Der vorgeschlagene Test \codeword{[U\_AB,Π\_AΠ\_B]=0} charakterisiert
deshalb bei solchen Überlappungen \textbf{keine globale Geradheit}. Das
lässt sich noch direkter widerlegen: Der Hamiltonoperator

\codeword{H=a†a+b†b}

ist global gerade, kommutiert mit keiner der beiden lokalen Paritäten
und auch nicht mit ihrem Produkt. Sein Cayley-Transform

\codeword{U=(I−iH)(I+iH)⁻¹}

ist eine exakt unitäre, global gerade Operation mit denselben drei
Nichtkommutationen. Der im Anhang vorgeschlagene Kill-Test würde diese
zulässige globale Operation fälschlich ausschließen.

Die richtige Reihenfolge ist daher: zuerst die gemeinsame CAR-Einbettung
und ihre echte globale Graduierung bestimmen. Erst danach lässt sich
entscheiden, welche lokalen Paritätskriterien anwendbar sind. Bei einer
disjunkten orthogonalen Zerlegung gilt die gewohnte Produktformel; bei
überlappenden Charts im Allgemeinen nicht.

\subsection{3. Noch wichtiger: Ein Überlappungsmatrixelement ist kein
Transport}\label{noch-wichtiger-ein-uxfcberlappungsmatrixelement-ist-kein-transport}

Für dieselben zwei Moden seien

\codeword{|A>=a†|0>}, \codeword{|B>=b†|0>}.

Dann ist ihre Gram-Matrix

\codeword{S=[[1,1/√2],[1/√2,1]]}.

Wähle den vollkommen trivialen Hamiltonoperator \codeword{H₀=ωN}, mit
\codeword{N=c₁†c₁+c₂†c₂}. Die Matrix der Hamilton-Matrixelemente in
diesen beiden Chartvektoren ist

\codeword{K\_ij=<i|H₀|j>=ω S\_ij}.

K hat also eine nichtverschwindende Offdiagonale, obwohl auf dem
gesamten Einteilchenraum nur dieselbe Phase
\codeword{e\textasciicircum{}(−iωt)} entsteht. Das korrekte Eigenproblem
ist \codeword{Kv=E Sv}, nicht \codeword{Kv=Ev}. Es liefert
ausschließlich die entartete Energie ω. Die Wahrscheinlichkeit

\codeword{|<B|exp(−itH₀)|A>|²=1/2}

ist für jede Zeit gleich groß. Es hat keine Übertragung stattgefunden;
die Hälfte war schon bei t=0 als statische Überlappung vorhanden.

Das ist unmittelbar für den bisherigen nativen TFPT-Pol relevant. Dieser
besitzt im erklärten N=63-Vertrag eine einzelne Energie E\_h auf einem
64-dimensionalen Polraum:

\codeword{P\_h H P\_h = E\_h P\_h}.

Sind A und B lediglich zwei Beschreibungen oder Sonden innerhalb
desselben Polraums, gilt wiederum \codeword{K=E\_h S}: ein gemeinsamer
Phasenfaktor, keine durch H verursachte Ausbreitung zwischen den Marken.
Diese Aussage braucht keine große Diagonalisierung. Sie folgt allein aus
der bereits bewiesenen Entartung.

Ein passiver Chartwechsel verändert nur die Beschreibung. Ein aktiv
implementierter Rotationsoperator kann Zustände verändern; dann müssen
jedoch genau dieser Operator, seine physische Verfügbarkeit, Zeit- oder
Ressourcenskala und sein Instrument aus der Quelle nachgewiesen werden.
Das ist nicht durch den Namen »Intertwiner« erledigt.

Die mögliche globale Überlappungskonstruktion ist dadurch nicht
ausgeschlossen. Sie müsste tatsächliche zusätzliche globale Dynamik
beziehungsweise eine Bandaufspaltung aus der Quelle zeigen und die
lokalen Polbeweise darin erneut absichern.

\subsection{4. Chartwechsel, Verbindungen und Krümmung sind verschiedene
Daten}\label{chartwechsel-verbindungen-und-kruxfcmmung-sind-verschiedene-daten}

Wenn mehrere Charts lediglich vollständige orthonormale Rahmen R\_x
desselben festen Vektorraums sind, lauten die Basiswechsel

\codeword{U\_xy=R\_x† R\_y}.

Dann teleskopiert jedes Dreiecksprodukt:

\codeword{U\_12 U\_23 U\_31=I}.

Das gilt auch bei nichtkommutierenden R\_x; der Checker prüft ein
konkretes solches Beispiel. Durch bloße lokale Basiswahl entsteht daher
nicht automatisch ein physisch gekrümmtes Eichfeld. Auf einem Bündel
beschreiben Übergangsfunktionen seine Verklebung; eine Verbindung
enthält zusätzliche Paralleltransportdaten. Eine nichttriviale
Bündeltopologie ist ebenfalls nicht dasselbe wie bereits gewählte lokale
Krümmung.

Bei tatsächlich variierenden \textbf{echten Unterräumen} ist die
Situation interessanter. Überlappungen \codeword{ι\_x†ι\_y} müssen dann
nicht unitär sein. Drei Strahlen \codeword{(1,0)}, \codeword{(1,1)/√2},
\codeword{(1,i)/√2} besitzen das nichtreelle Schleifenprodukt
\codeword{(1+i)/4}. Das liefert eine geometrische Phase, aber sein
Betragsquadrat ist 1/8 und nicht eins. Eine solche Unterraumgeometrie
kann Ansatzpunkt für eine geometrische Verbindung sein; sie ist kein
automatisch ausgeführter normerhaltender Transport und keine schon
gewonnene Eichfeldkinetik.

Zudem legt eine interne Eichverbindung alleine keine Raumzeitmetrik
fest. Verschiedene Linkphasen können auf demselben Graphen bei denselben
Operationskosten existieren. »Die U\_xy ändern sich« bedeutet ohne
weiteren Nachweis nicht bereits »die Raumzeitgeometrie ändert sich«.

\subsection{5. Zeit: Ein Grundzustand bewegt sich unter seinem H nicht
beobachtbar}\label{zeit-ein-grundzustand-bewegt-sich-unter-seinem-h-nicht-beobachtbar}

Die im Anhang skizzierte Entwicklung \codeword{Ω -> exp(−itH)Ω} erzeugt
bei einem Energieeigenzustand lediglich \codeword{exp(−itE₀)Ω}. Der
Dichteoperator bleibt exakt gleich. Der Checker bestätigt das an einem
kleinen Fockzustand.

Damit wird weder die Zeitordnung noch ein Zeitpfeil erzeugt.
Nichtstationäre Zustände und mehrzeitige Korrelationsfunktionen können
natürlich eine Dynamik anzeigen; eine gerichtete Record- oder
Präparationsgeschichte braucht ihre eigenen Bedingungen. Auch eine
relationale Zeit aus bedingten Zuständen wäre als Forschungsansatz
möglich, verlangte aber eine explizite Uhr, ihren gemeinsamen Zustand
mit dem System und eine Konditionierungsregel. Der bloße globale
Eigenzustandsphasenfaktor reicht dafür nicht.

Die nun geprüfte endliche Clock kann eine aktive Operation auf inneren
Marken darstellen, falls sie physisch gewährt ist; dass sie zur
Quellsymmetrie gehört, macht sie nicht automatisch identisch mit der
Hamiltonzeit.

\subsection{6. Der binäre Index muss als Freiheitsgrad bewiesen
werden}\label{der-binuxe4re-index-muss-als-freiheitsgrad-bewiesen-werden}

Die skalare Hilfsdublett-Reparatur braucht zwei unabhängige,
gleichgeladene Komponenten, auf denen eine nichtentartete
antisymmetrische Form wirkt. Eine Orientierung der Überlappungsgeometrie
\textbf{könnte} so etwas liefern, tut es aber nicht allein durch das
Vorhandensein zweier Bezeichnungen.

Die drei einfachen Verwechslungen sind:

\begin{itemize}
\tightlist
\item
  Zwei Namen oder ±-Vorzeichen für dieselbe Mode bilden nur eine
  eindimensionale, redundante Beschreibung. Der Rückzug der
  antisymmetrischen Zweiform auf diesen Raum ist null.
\item
  Ist Rückwärts die Adjungierte der Vorwärtsoperation, haben f und f†
  entgegengesetzte U(1)-Ladungen. Das ist kein gleichgeladenes
  Hilfsdublett. Ihr bilinearer Kanal ist neutral, nicht von Ladung −2.
\item
  Bedeutet Vorwärts/Rückwärts links-/rechtshändige Weylfelder, wurden
  zwei verschiedene Lorentzdarstellungen eingeführt. Das ist nicht die
  zusätzliche gleichartige Hilfskopie des bereits geprüften
  Tensorvertrags.
\end{itemize}

Eine geometrisch hergeleitete Zweifachheit könnte also eine gute
Erklärung des Hilfsindex sein. Nötig wären zwei tatsächliche unabhängige
Kanäle, ihre CAR-/Ladungsrelationen, die Spin-Lorentz-Wirkung und die
nichtverschwindende Kopplung auf demselben Quellraum. Ein
Orientierungsbit liefert weder automatisch Weylspin noch die
Transformationsregel unter einer vollen Lorentzdrehung.

\subsection{7. Operationsgeometrie benötigt mehr als
Erreichbarkeit}\label{operationsgeometrie-benuxf6tigt-mehr-als-erreichbarkeit}

Minimale Kosten erfüllen unter geeigneter Komposition eine
Dreiecksungleichung. Ohne reversible Operationen mit symmetrischen
Kosten ist die resultierende Funktion jedoch nur eine gerichtete
Distanz; der Checker liefert den gerichteten Dreierzyklus mit
\codeword{d(A,B)=1}, \codeword{d(B,A)=2}. Unerreichbarkeit ergibt
unendliche Distanzen. Kostenfreie Chartwechsel können verschiedene
Beschreibungen auf Distanz null setzen; dann muss zunächst nach
physischer Äquivalenz quotiert werden.

Wachstum \codeword{V(R)\textasciitilde{}R³} ist ein Nachweis einer
dreidimensionalen \textbf{Wachstumsdimension im gewählten Kostenmaß},
nicht automatisch einer glatten dreidimensionalen Mannigfaltigkeit,
einer 3+1D-Lorentzmetrik oder einer universellen Lichtgeschwindigkeit.
Schon ein kubisches Gitter mit Laplace-Hamiltonoperator hat dieses
Volumenwachstum, aber bei kleinen Impulsen

\codeword{E(k)=Σ\_j(2−2cos k\_j) \textasciitilde{} |k|²},

nicht eine lineare relativistische Dispersion. Der eindimensionale
Taylor-Koeffizient dieser separierbaren Gegenkonstruktion wird exakt
geprüft.

Eine lokale Spektrallücke ist ebenso noch keine relativistische
Teilchenmasse. Dazu braucht es die gemeinsame Energie-/Impulsdeutung,
den entsprechenden Dispersionszweig und den Ladungs-/Referenzvertrag.
»Kopieren« sollte bei Teilchenpropagation nur bildlich verwendet werden:
kohärenter Transfer erzeugt nicht zwei unabhängige Kopien eines
unbekannten Quantenzustands.

\subsection{8. Gravitation bleibt eine dynamische
Verpflichtung}\label{gravitation-bleibt-eine-dynamische-verpflichtung}

Die linearen Fluktuationen einer aus der Quelle bestimmten globalen
Geometrie zu prüfen, ist ein sinnvoller Forschungsauftrag. Eine
transversale spurfreie Tensorzerlegung allein garantiert aber weder
einen masselosen physikalischen Spin-2-Pol noch positive Norm, genau
zwei Helizitäten oder universelle Kopplung. Selbst ein gesunder linearer
Spin-2-Sektor ersetzt nicht den Nachweis seiner nichtlinearen
Eich-/Zwangsstruktur und konsistenten Materiekopplung.

Das vorgeschlagene wechselseitige Schema Materie -\textgreater{}
bevorzugte Überlappungen -\textgreater{} Geometrie ist daher ein
mögliches Wirkungsprinzip, noch keine aus dem vorhandenen nativen H
gewonnene Gleichung.

\subsection{9. P versus NP: eine konkrete
Korrektur}\label{p-versus-np-eine-konkrete-korrektur}

Der Anhang sagt, P≠NP würde einen intrinsisch exponentiellen Suchaufwand
bedeuten. Das ist falsch. P≠NP würde ausschließen, dass \textbf{jedes}
NP-Entscheidungsproblem deterministisch in Polynomialzeit gelöst werden
kann. Es folgt daraus keine exponentielle untere Schranke: Eine
hypothetische Laufzeit wie \codeword{2\textasciicircum{}(√n)} ist
superpolynomiell und zugleich subexponentiell. Stärkere
Exponentialzeitaussagen benötigen zusätzliche Sätze beziehungsweise
Hypothesen.

NP bezieht sich zudem auf polynomial lange, polynomial prüfbare
Zertifikate in der Länge einer wohldefinierten Eingabekodierung.
Beliebige Erreichbarkeit in einem knapp beschriebenen Graphen mit
möglicherweise exponentiell langen Wegen ist nicht allein deshalb ein
NP-Problem. Eine geometrische Umformulierung muss Eingabelänge,
Zertifikatlänge, erlaubte Operationen und deren Kosten erhalten. Eine
Pfadmetapher löst diese Verpflichtungen nicht.

\subsection{10. Ein korrigierter, wirklich entscheidbarer
Anschluss-Test}\label{ein-korrigierter-wirklich-entscheidbarer-anschluss-test}

Ein sinnvoller nächster Versuch besteht nicht nur aus einem
unbeschrifteten \codeword{P U\_AB P}. Er sollte zusammen liefern:

\begin{enumerate}
\def\labelenumi{\arabic{enumi}.}
\tightlist
\item
  \textbf{Globale Quelle und Einbettungen:} konkrete primitive Algebra
  und zwei CAR-/Tensor-erhaltende Abbildungen; alle Kreuzrelationen und
  die tatsächliche globale Graduierung.
\item
  \textbf{Gemeinsamer Zustand:} derselbe globale H und Zustand;
  Nachweis, welche bisherigen lokalen Pole und Lücken darin noch gelten.
  Kein stiller Rückgriff auf den Produktgrundzustand unabhängiger
  Banken.
\item
  \textbf{Physische Operation:} aus einem erklärten Quellenwort
  erzeugter aktiver Operator oder Generator, mit Zeit- beziehungsweise
  Ressourcenskala. Passive Chartwechsel bleiben getrennt.
\item
  \textbf{Gram-bereinigte Dynamik:} \codeword{S\_ij=<h\_i|h\_j>} und
  \codeword{K\_ij=<h\_i|H|h\_j>} berechnen; das generalisierte
  Eigenproblem beziehungsweise eine orthonormalisierte Darstellung
  verwenden. Prüfen, ob mehr als \codeword{K=E\_h S} entsteht.
\item
  \textbf{Operativer Transfer:} dieselbe Anfangspräparation und
  Zielmessung; zeitabhängige Änderung gegenüber der statischen
  Überlappung und Fehlerkontrolle gegenüber dem übrigen Zustandsraum.
  Erst dann mit dem bekannten bedingten Zwei-Banken-Transfer
  vergleichen.
\end{enumerate}

Das nimmt die mögliche gemeinsame Herkunft ernst, korrigiert aber die
falschen Automatismen. Der stärkste sofortige Erkenntnisgewinn dieses
Audits ist negativ und präzise: \textbf{Allein durch Überlappung oder
Umbenennung eines entarteten nativen Polraums entsteht keine
Transportdynamik.} Ob der Compiler darüber hinaus genau die passende
globale Dynamik besitzt, ist die konkrete offene Frage.

\begin{center}\rule{0.5\linewidth}{0.5pt}\end{center}

\section{Anhang E: arithmetischer
Schleifenvorschlag}\label{anhang-e-arithmetischer-schleifenvorschlag}

\section{Zusatzprüfung: Schleifen, Primzahlen und der
Universalraum-Vorschlag}\label{zusatzpruxfcfung-schleifen-primzahlen-und-der-universalraum-vorschlag}

\begin{enumerate}
\def\labelenumi{\arabic{enumi}.}
\setcounter{enumi}{14}
\tightlist
\item
  September 2026, v1.6.6. Quellenprüfung des nachgereichten Textes
  \codeword{59dc0059-8914-48ca-953d-85933f66e00b}, kein neuer RH-Beweis
  und kein neuer arithmetischer Quellenoperator. Seine Aussagen über
  gemeinsame lokale Beschreibungen werden getrennt im Überlappungsaudit
  behandelt.
\end{enumerate}

\subsection{Urteil}\label{urteil}

\textbf{Ein globaler Wegraum ist eine sinnvolle Suchklasse; ein
Eulerprodukt über primitive Wege ist noch nicht das Eulerprodukt der
Riemannschen Zetafunktion.} Die im Text vorgeschlagene Xi-Determinante
bleibt ein hinreichendes Ziel unter expliziten Operatorannahmen. Das
Umbenennen ihres noch fehlenden Operators in „Generator der primitiven
Schleifen`` konstruiert ihn nicht.

\subsection{1. Was die endliche Bank tatsächlich
ausschließt}\label{was-die-endliche-bank-tatsuxe4chlich-ausschlieuxdft}

Der beibehaltene Beweis aus v1.6.5 gilt für den ursprünglichen
Hamiltonoperator mit 64 Fermion- und 60 Bosonmoden, Δ\textgreater0. Der
Hilbertraum ist wegen der Bosonen \textbf{nicht endlichdimensional}.
Endlich ist die Zahl der Moden. Mit \codeword{C=960g²/Δ} folgt aus
quadratischem Ergänzen

\[\adjustbox{max width=\linewidth}{$\displaystyle 
H\ge\frac\Delta2N_b-C,
\qquad
N_H(E)\le2^{64}\binom{\lfloor2(E+C)/\Delta\rfloor+60}{60}.
$}\]

Dies widerspricht einem vollständigen energieerhaltenden Spektrum
\codeword{E\_* log n + E\_off}, E\_*\textgreater0, weil dessen
Zählfunktion exponentiell wächst. Es ist kein Ausschluss beliebiger
arithmetischer Untersektoren oder eines anderen Operators mit
nichtlinearer Energiezuordnung.

Ein großer Wegraum kann eine andere Zählfunktion besitzen. Die
Gleichsetzung „mehr Wege = mehr orthogonale Zustände unter derselben
Energie`` muss aber bewiesen werden. Verschiedene Wörter können
denselben Operator oder Zustand darstellen; ein Quantenüberlagerungsraum
ist nicht ohne Weiteres der freie Wortraum. Bei unendlich vielen
identischen Zellen kann stattdessen die globale Wärmespur divergieren.
Der Text entfernt diese Fragen nicht durch Vergrößerung.

\subsection{2. Der kleinste Gegencheck zur
Primzahl-Automatik}\label{der-kleinste-gegencheck-zur-primzahl-automatik}

Nehmen wir einen gerichteten Knoten mit zwei erlaubten Schleifen a und
b. Neben a und b ist auch ab ein primitiver periodischer Weg: Es ist
keine Potenz eines kürzeren Wortes. Ebenso entstehen weitere gemischte
primitive Wörter. Setzen wir nachträglich L(a)=log 2, L(b)=log 3, so hat
ab die Länge log 6. Die 6 ist keine Primzahl. „Primitiver Weg`` bedeutet
nicht „arithmetische Primzahl``.

Schon formal unterscheiden sich die zugehörigen erzeugenden Funktionen:

\[\adjustbox{max width=\linewidth}{$\displaystyle 
Z_{\rm Wege}(x,y)=\frac1{1-x-y},\qquad
Z_{\rm zwei\ Primarten}(x,y)=\frac1{(1-x)(1-y)}.
$}\]

Der Koeffizient von xy ist links 2 und rechts 1. Links werden die beiden
Wörter ab und ba gezählt; rechts eine kommutative Besetzung. Im
periodischen Eulerprodukt bilden ab und ba eine zyklische Klasse, aber
diese ist ein \textbf{neuer primitiver Faktor}. Die Abweichung
verschwindet damit nicht.

Eine kommutative freie Halbgruppe über vorgegebenen Primarten
reproduziert eindeutige Faktorzerlegung. Sie leitet aber weder die
natürliche arithmetische Markierung noch die logarithmischen Längen oder
die passende Spur her. Eine gekoppeltere Geometrie müsste gemischte
Bahnen mit den richtigen Amplituden, Relationen oder nachgewiesenen
Auslöschungen behandeln. Dies ist kein allgemeiner Ausschluss von
Schleifenmodellen mit Interferenz.

Der Unterschied ist in den Originalarbeiten sichtbar: Kuipers, Hummel
und Richter konstruieren Quantengraphen mit dem passenden oszillierenden
Anteil der Nullstellendichte, weisen aber auf den anderen glatten Anteil
und damit das andere Spektrum hin. Das ist ein nützlicher Vorläufer,
kein RH-Operator. \href{https://arxiv.org/abs/1307.6055}{Originalarbeit,
Phys. Rev.~Lett. 112, 070406}

Graphische Eulerprodukte besitzen ihre eigene Theorie. Auch eine aus
einer Graph-Zetafunktion bestimmbare Größe ist nicht automatisch
effizient auslesbar. Storm trennt in seiner Arbeit ausdrücklich
berechenbare Zeta-Darstellung und teure Auswertung bestimmter
Graphinformationen.
\href{https://arxiv.org/abs/0708.1923}{Originalarbeit zu
Edge-Zetafunktionen}

\subsection{3. Die richtige
RH-Zielbedingung}\label{die-richtige-rh-zielbedingung}

Sei A strikt positiv und selbstadjungiert, mit kompakter Inverser und
\codeword{A\textasciicircum{}\{-2\}} von Spurklasse. Dann ist der
Fredholm-Ausdruck

\[\adjustbox{max width=\linewidth}{$\displaystyle 
D(z)=\det(I-z^2A^{-2})
$}\]

ganz, und seine Nullstellen liegen bei den reellen Zahlen ±λ\_j(A), mit
ihren Multiplizitäten. Würde unabhängig und auf ganz C bewiesen

\[\adjustbox{max width=\linewidth}{$\displaystyle 
\Xi(z)/\Xi(0)=D(z),\qquad \Xi(z)=\xi(1/2+iz),
$}\]

so folgte RH. Das ist die korrekte bedingte Implikation. Nicht
ausreichend sind ein paar passende Eigenwerte, die imaginären Teile
bereits eingesetzter Nullstellen, ein Eulerprodukt nur für
Re(s)\textgreater1 oder eine Formaldeterminante ohne Domäne, Spurklasse
und vollständige archimedische Faktoren.

Der vorgelegte Text liefert A nicht, bestimmt keine Domäne und leitet
weder die vollständige Spurformel noch die ganze Funktionsidentität her.
Die allgemeine Selbstadjungiertheit eines anderen Operators hilft dabei
nicht.

\subsection{4. Vorarbeiten und aktuelle Grenze der
Quellenprüfung}\label{vorarbeiten-und-aktuelle-grenze-der-quellenpruxfcfung}

Der installierte Leitfaden \codeword{rh-graph-research} und die
projektspezifische Korpussuche wurden angewandt. Die aktuelle
Konsistenzprüfung und der anschließende Wiederaufbauversuch brachen mit

\codeword{SOURCE\_UNAVAILABLE: /Users/stefanhamann/Documents/Codex/2026-09-04/scha-2/research}

ab. Die gesonderte Paper-Abfrage meldete
\codeword{PAPER\_SOURCE\_OR\_REVIEW\_DRIFT}. Keine dieser Prüfungen
wurde durch Ersetzen von Pins grün gemacht. Der Forschungsindex darf
hier \textbf{nicht als global aktuell geprüft} bezeichnet werden.
Vorhandene Suchresultate dienten lediglich zum Auffinden real
verfügbarer Originaldateien; die aktuelle Queue ist keine Beweisquelle.

Gezielt gelesen wurden die relevanten Abschnitte von
\codeword{rh/catalog/analysis/geometry\_audit.md} (Quantengraphen) und
\codeword{rh/catalog/analysis/event\_log\_function.md} (arithmetische
Ereignisse), sowie die vollständige maßgebliche
Zählungs-/Determinantenpassage aus v1.6.5. Die Dateien werden für diese
Revision eingefroren. Ihre historischen Korpus-Abwesenheitsangaben
werden nicht als aktueller Vollständigkeitsbeweis übernommen.

Im bestehenden Index berühren \codeword{r618 / STRUCTURAL\_MISMATCH} die
Verwechslung RH-neutraler E8-Daten mit einer Xi-Identität und
\codeword{ledger:E8.COXETER.EULER.COMPLETION.01 / NO\_BRIDGE} die
fehlende globale Vervollständigung. Solche Treffer allein widerlegen
kein neues Objekt; die hier tragende Prüfung ist der explizite
Unterschied zwischen gemischten primitiven Wegen und arithmetischen
Primarten. Es wird kein neuer globaler RH-Beweisstatus registriert oder
freigegeben.

\subsection{5. Tragfähiger Anschluss, ohne acht Versprechen an ein
unbekanntes
Objekt}\label{tragfuxe4higer-anschluss-ohne-acht-versprechen-an-ein-unbekanntes-objekt}

Der nächste arithmetische Test wäre erst nach Definition einer
tatsächlichen Quell-Wegstruktur sinnvoll: primitive Bahnen, ihre Längen,
Amplituden und die Spur gemeinsam bestimmen, ohne Primlisten oder
Nullstellen einzusetzen. Ein positiver Kontrolltest muss insbesondere
gemischte Zyklen erklären und den archimedischen Anteil auf demselben
Träger liefern. Gelingt nur ein generisches Graph-Eulerprodukt, ist das
ein Graphresultat und bleibt von RH getrennt.

Faktorisierung und P versus NP werden dadurch nicht gelöst. Außerdem ist
„P≠NP bedeutet notwendigerweise exponentielle Suchkosten`` zu stark: Aus
einer fehlenden polynomialen Zeitgrenze folgt nicht allein eine
exponentielle untere Schranke. Für Hylæan enthält der Text ein mögliches
Operationsbild, aber keinen neu überprüften Lern- oder
Gedächtnisnachweis.

\begin{center}\rule{0.5\linewidth}{0.5pt}\end{center}

\section{Historischer vollständiger Herleitungsstand v1.6.5
einschließlich
v1.6.4}\label{historischer-vollstuxe4ndiger-herleitungsstand-v1.6.5-einschlieuxdflich-v1.6.4}

Der folgende Text wird unverändert bewahrt. Neue Gesamtstatusaussagen
und Reichweitenkorrekturen stehen in der vorangestellten Revision
v1.6.6. Historische Arbeitsaufträge oder Verfügbarkeitsannahmen werden
dadurch nicht erneut zu aktuellen Beweisen erklärt.

\section{TFPT / Universalraum: kontrollierter
Zwei-Banken-Transfer}\label{tfpt-universalraum-kontrollierter-zwei-banken-transfer}

Forschungsfortsetzung und Quellenaudit · 15. September 2026 · v1.6.5

\subsection{Ergebnis und
Geltungsbereich}\label{ergebnis-und-geltungsbereich}

Eine konkrete bisher offene Rechnung lässt sich schließen: \textbf{Im
ausdrücklich um einen Rotor-Link erweiterten Modell zweier nativer
Fermion-Boson-Banken ist der Transfer der isolierten Lochanregung mit
einer Fehlergrenze gegenüber dem vollständigen physikalischen
Zustandsraum kontrollierbar.} Dafür ist keine Diagonalisierung dieses
enorm großen Raums nötig. Die vorhandene Casimiridentität,
Ladungserhaltung und eine Spektrallücke reichen aus.

Am unten vollständig angegebenen, bewusst sehr schwachen Kopplungspunkt
gilt:

{\def\LTcaptype{none} % do not increment counter
\begin{longtable}[]{@{}
  >{\raggedright\arraybackslash}p{(\linewidth - 4\tabcolsep) * \real{0.3000}}
  >{\raggedleft\arraybackslash}p{(\linewidth - 4\tabcolsep) * \real{0.4000}}
  >{\raggedright\arraybackslash}p{(\linewidth - 4\tabcolsep) * \real{0.3000}}@{}}
\toprule\noalign{}
\begin{minipage}[b]{\linewidth}\raggedright
Aussage
\end{minipage} & \begin{minipage}[b]{\linewidth}\raggedleft
Strenge Schranke
\end{minipage} & \begin{minipage}[b]{\linewidth}\raggedright
Status
\end{minipage} \\
\midrule\noalign{}
\endhead
\bottomrule\noalign{}
\endlastfoot
Zielwahrscheinlichkeit aus dem normierten Zustand der isolierten
Lochlinie & \textgreater{} 99,2677105 \% & Analytisch mit rationalen
Zertifikaten, im zusätzlichen Linkmodell \\
Zielwahrscheinlichkeit aus der normierten ursprünglichen Entnahme
\(f_r\Omega/\sqrt\nu\), \textbf{ohne anfänglichen Energiefilter} &
\textgreater{} 89,7260353 \% & Gleicher vollständiger Hamiltonoperator,
Ziel ist die niedrige Lochlinie rechts \\
Verlassen des ungestörten niedrigen Bandes, aus einem Zustand dieses
Bandes & \textless{} \(1{,}8\cdot10^{-11}\) zu jeder Zeit & Voller
erlaubter Ladungssektor, keine Bosonen- oder Flussabschneidung \\
Herkunft des Links, der Bankzerlegung und der Präparation & nicht
hergeleitet & Offen \\
Relativistische Felder, gemeinsamer 3+1D-Ursprung, T1-T8 & nicht
konstruiert & Offen \\
\end{longtable}
}

Die Zahlen sind \textbf{untere beziehungsweise obere Einschließungen},
keine berechneten Zentralwerte. Sie beschreiben kein durchgeführtes
physisches Experiment. „Vollständig`` bezieht sich hier auf die
Fehlerkontrolle im deklarierten Zwei-Banken-Ladungssektor, nicht auf die
TOE.

Die gelieferten Arbeiten wurden mit 1.149 beziehungsweise 753
Prüfbedingungen normal und optimiert reproduziert, jeweils byteidentisch
mit den mitgelieferten Ergebnissen. Unsere neue Transferrechnung hat
4.516 Prüfbedingungen; eine gesonderte Reichweitenprüfung der
währenddessen geänderten Symmetrieaussagen hat zehn. Die zwei zuletzt
eingegangenen Anlagen wurden mit 62 weiteren Bedingungen nachgeprüft.
Insgesamt sind das 6.490 Bedingungen, jeweils in beiden Laufarten. Diese
Zahl ist keine Zahl unabhängiger Theoreme. Insbesondere ersetzen die
Zertifikate nicht die nachstehenden analytischen Argumente und sind kein
Lean-Beweis.

\subsection{1. Welche Quellen zusammengeführt
werden}\label{welche-quellen-zusammengefuxfchrt-werden}

\begin{enumerate}
\def\labelenumi{\arabic{enumi}.}
\tightlist
\item
  Die Benutzeranlage
  \codeword{addbf22d-4764-4685-bdfe-f422bfc563d0/pasted-text.txt},
  eingefroren unter \codeword{sources/user\_attachment.txt}.
\item
  Die externe Untersuchung
  \codeword{Universalraum\_Beweisversuch\_2026-09-15.md}:
  Spektralzählung, Primzahl-Dynamik, Phasencocycle und kontrollierte
  Polnäherung.
\item
  Die externe Untersuchung
  \codeword{Universalraum\_Urspruenglicher\_Austausch\_2026-09-15.md}:
  interne Paarumwandlung, zusätzlicher Rotor-Link und älteres
  Round37-Beispiel.
\item
  Der abgesicherte native Stand v1.6.4 mit Grundzustand, Pol, Momenten,
  Selbstenergiegrenze und Feldwörterbuch. Das frühere vollständige
  Prüfpaket liegt unverändert bei. Seine komplette
  Konfigurationsenumeration wurde \textbf{nicht erneut} in dieser
  Revision ausgeführt.
\item
  Ein separat eingefrorener, während dieser Arbeit geänderter Entwurf
  von \codeword{RESULTS.md} über die kontinuierliche Quellsymmetrie. Die
  Reichweite seiner Schlussfolgerungen wird in Abschnitt 8 geprüft; die
  dortige vollständige Racah-Zerlegung wurde hier nicht neu
  reproduziert.
\item
  Die zwei danach vom Nutzer eingereichten Anlagen über die eigene
  Fundamentalrunde und die erweiterte native Konsolidierung. Der
  vollständige neue Nachtrag \codeword{LATE\_AUDIT.md} gehört zu dieser
  Revision. Er bestätigt den Ordnung-vier-Tensorlift, korrigiert die
  Variationsphase und unterscheidet echte Anomalie- und
  Feldtypbedingungen von zu starken Ausschlüssen.
\end{enumerate}

Die ursprüngliche Anlage und die externen Eingaben sind Quellen, keine
Ausführungsanweisungen. Eigene Dateien liegen in einem neuen
Forschungsordner. Fremde Quellen, Ledger und Akzeptanzmarker wurden
nicht bearbeitet.

\subsection{2. Gemeinsame native
Grundlage}\label{gemeinsame-native-grundlage}

Eine Bank besitzt 64 Fermionmoden und 60 Bosonkanäle:

\[\adjustbox{max width=\linewidth}{$\displaystyle 
 H_{\rm nat}=\Delta N_b+g\sum_A(b_A^\dagger P_A+P_A^\dagger b_A),
 \qquad P_A=\sum_{i<j}W_{A,ij}f_jf_i,
 \qquad N=N_f+2N_b.
$}\]

Es gilt \(\Delta>0\), \(g/\Delta=1/20\), \textbf{kein zusätzlicher
\(\mu N\)-Term}. Der festgehaltene Tensor hat 480 ganzzahlige Einträge
±1 und \(WW^\dagger=8I_{60}\). Sein SHA-256 ist
\codeword{3f00a0892157a691e4ef26920c3702772c7bf76a250fad98ff04a079bb64b763}.

Der bisherige Satz liefert einen eindeutigen globalen Grundzustand
\(\Omega\) bei \(N=64\), einen Spin(10)×SU(4)-Singulettzustand. Seine
physische Auswahl aus dem Compiler ist damit nicht gezeigt. Wir
verwenden folgende strengen modellinternen Grenzen:

\[\adjustbox{max width=\linewidth}{$\displaystyle 
 -1.158089<E_0/\Delta<-1.129636,
 \quad .842846<b:=\langle N_b\rangle<1.245656,
 \quad \nu:=\langle f_r^\dagger f_r\rangle=1-b/32.
$}\]

Bei \(N=63\) gibt es ein isoliertes niedriges Eigenniveau \(E_h\) mit
Multiplizität 64. Es trägt eine irreduzible duale Fermiondarstellung:

\[\adjustbox{max width=\linewidth}{$\displaystyle 
 -1.121899<E_h/\Delta<-1.095812,
 \quad .007737<(E_h-E_0)/\Delta<.039079763822.
$}\]

Für seinen Spektralprojektor \(P_h\) ist

\[\adjustbox{max width=\linewidth}{$\displaystyle 
 Z=\|P_hf_r\Omega\|^2,
 \quad Z_{\rm lo}:=\frac{40912436089}{46487375000}<Z
 <\frac{15578577}{16000000}=:Z_{\rm hi}.
$}\]

Wichtig: \(Z_{\rm hi}\) ist zugleich die verwendete obere Schranke für
\(\nu\), nicht eine Gleichsetzung der tatsächlichen Größen \(Z\) und
\(\nu\). Das niedrige Gewicht ist \textgreater88,0076 \% des gesamten
normierten CAR-Spektralmaßes. Innerhalb der \textbf{Entnahmeantwort} ist
sein Anteil \(w=Z/\nu>90,38836\,\%\).

Alle weiteren Eigenwerte von \(H_{63}\) sind
\textgreater{}\(-.75\Delta\); alle angeregten Eigenwerte von \(H_{64}\)
ebenfalls. Die Entnahme-Restantwort beginnt relativ zu \(E_0\) oberhalb
\(.379636\Delta\), die Additionsantwort oberhalb \(.329636\Delta\).

Diese Schranken gehören zu demselben Hamiltonoperator und demselben
Zustand. Die ursprüngliche Größe \(\chi_r^\dagger\Omega\) bei Ladung 67
wird nicht mit \(f_r\Omega\) bei Ladung 63 identifiziert.

\subsection{3. Prüfung der gelieferten neuen
Resultate}\label{pruxfcfung-der-gelieferten-neuen-resultate}

\subsubsection{3.1 Die interne Paarumwandlung ist wirklich
nativ}\label{die-interne-paarumwandlung-ist-wirklich-nativ}

Bei Gesamtladung 2 zerfällt der 2.076-dimensionale Raum in 60 helle
Zweizustandsblöcke und 1.956 dunkle Zustände. In jedem hellen Block gilt

\[\adjustbox{max width=\linewidth}{$\displaystyle 
 H_A=\begin{pmatrix}0&\sqrt8g\\\sqrt8g&\Delta\end{pmatrix},
 \qquad p_{\max}=\frac{32g^2}{\Delta^2+32g^2}=\frac2{27}.
$}\]

Das wurde aus dem unveränderten Tensor reproduziert. Es ist die
Umwandlung zweier Fermionen in einen Bosonkanal, \textbf{kein
abgeleiteter Ortswechsel einer Ladung eins}, und der Ladung-2-Zustand
ist nicht der native Grundzustand. Die acht internen Pfade pro Kanal
müssen kohärent bleiben; getrennte Pfadaufzeichnungen verändern den
Gramoperator und damit die Dynamik.

Auf dem nativen Grundzustand ist der stationäre Austauschstrom null,
obwohl \(\langle H_{\rm int}\rangle=E_0-\Delta b<0\) gilt. Eine
vollständige \(N_b\)-Dephasierung erhöht die Energie um
\(\Delta b-E_0\in(1.972482,2.403745)\Delta\). Statische Kohärenz ist
nicht dasselbe wie ein gerichteter Strom.

\subsubsection{3.2 Der neue Link ist konsistent, aber
zusätzlich}\label{der-neue-link-ist-konsistent-aber-zusuxe4tzlich}

Die externe Arbeit definiert ein Rotorpaar \(E,U\) auf
\(\ell^2(\mathbb Z)\): \(E|e\rangle=e|e\rangle\),
\(U|e\rangle=|e+1\rangle\), \([E,U]=U\). Mit zwei \textbf{gradiert}
zusammengesetzten nativen Banken lautet der Zusatz

\[\adjustbox{max width=\linewidth}{$\displaystyle 
 B=\sum_{r=1}^{64}f_{y,r}^\dagger U f_{x,r},
 \qquad V=\tau B+\bar\tau B^\dagger,
 \qquad H=H_x+H_y+\frac\kappa2E^2+V.
$}\]

Dieser Operator erhält die Gesamtladung und die passend definierten
lokalen Gaussgrößen. Er ist beschränkt mit \(\|V\|\le64|\tau|\). Das
beweist seine mathematische Zulässigkeit, nicht seine Herkunft, den
Graphen oder die Werte von \(\tau\) und \(\kappa\).

Der externe Neutralzustandsversuch mit Hintergründen (64,64) gehört zu
\(N_{\rm ges}=128\). Dort sind die erzeugten Loch-Additions-Paare eine
andere Antwort als eine bereits vorhandene einzelne Lochanregung. Die
berechnete Gewichtsnorm \(64\nu(1-\nu)=2b-b^2/16\), das Faltungsmaß und
die Schwelle \(>.337373\Delta+\kappa/2\) sind damit vereinbar. Sie
werden nicht als Einloch-Transferwahrscheinlichkeit ausgegeben.

\subsubsection{3.3 Das ältere 99,8718-\%-Beispiel ist ein anderer
Träger}\label{das-uxe4ltere-998718--beispiel-ist-ein-anderer-truxe4ger}

Die mitgelieferte Round37-Rechnung wurde exakt reproduziert. Ihre sechs
Zustände sind der vollständige Gausssektor eines
\textbf{Vierfermion-Modells} mit einem Link. Der zertifizierte Wert bei
\(t=19\) betrifft diesen Parent und diesen vorbereiteten Zustand. Er ist
weder eine Zwei-64-Moden-Banken-Simulation noch ein Beweis für deren
ursprüngliche Auswahl. Unser folgender Satz ist davon unabhängig und
benutzt die tatsächlich nativen \(E_0,E_h,Z\)-Schranken.

\subsection{4. Neuer Satz: vollständige Kontrolle des
Einloch-Transfers}\label{neuer-satz-vollstuxe4ndige-kontrolle-des-einloch-transfers}

\subsubsection{4.1 Den richtigen physikalischen Sektor
festlegen}\label{den-richtigen-physikalischen-sektor-festlegen}

Für genau ein Loch ist eine Ladungsreferenz notwendig. Wir deklarieren
Hintergründe \((q_x,q_y)=(63,64)\) und fordern

\[\adjustbox{max width=\linewidth}{$\displaystyle 
 G_x=N_x-63+E=0,\qquad G_y=N_y-64-E=0.
$}\]

Dann ist \(N_x+N_y=127\) und \textbf{\(E=63-N_x\) exakt festgelegt}. Auf
diesem Ein-Kanten-Baum verbleibt kein unabhängiger unbeschränkter
Flussindex. \(0\le N_x,N_y\le127\) beschränkt auch die Bosonenzahlen.
Der gesamte physikalische Sektor ist endlichdimensional; keine
Flussabschneidung und kein willkürlicher Bosonen-Cutoff werden
eingeführt.

Definiere \(h_r=P_hf_r\Omega/\sqrt Z\). Die 64 Vektoren sind
orthonormal. Das niedrige Band \(P\) wird von

\[\adjustbox{max width=\linewidth}{$\displaystyle 
 |L,r\rangle=h_{x,r}\otimes\Omega_y\otimes|0\rangle,
 \qquad |R,r\rangle=\Omega_x\otimes h_{y,r}\otimes|-1\rangle
$}\]

aufgespannt und besitzt Dimension 128. Für den ungekoppelten Operator
\(H_0=H_x+H_y+\kappa E^2/2\) sind die Energien \(e_*=E_h+E_0\) und
\(e_*+\kappa/2\).

Nach einheitlicher Wahl der relativen Fermionphase wirkt die exakte
Kompression \(PHP\), abzüglich \(e_*\), wie

\[\adjustbox{max width=\linewidth}{$\displaystyle 
 \begin{pmatrix}0&-\bar\tau Z\\-\tau Z&\kappa/2\end{pmatrix}
 \otimes I_{64}.
$}\]

Das Vorzeichen ist für die Population unerheblich. Das \(I_{64}\) zeigt,
dass die Rechnung für jeden normierten internen Überlagerungszustand
gilt. Die projizierten Fermionoperatoren sind Hubbard-Übergänge und
keine vollständige CAR-Algebra auf diesem 128-dimensionalen Band.

\subsubsection{4.2 Die Lücke zum gesamten übrigen
Raum}\label{die-luxfccke-zum-gesamten-uxfcbrigen-raum}

Die native, nicht nur perturbative Casimiridentität lautet

\[\adjustbox{max width=\linewidth}{$\displaystyle 
 \sum_AP_A^\dagger P_A=\tfrac12(15N_f-C_{\rm Spin(10)}-C_{\rm SU(4)}).
$}\]

In einem Block \(N=n,N_b=b\) ist die rechte Seite höchstens
\(\tfrac{15}2(n-2b-\eta)\) mit \(\eta=n\bmod2\). Bei ungerader
Fermionzahl kommt der Casimir-Mindestwert 15 hinzu. Die Norm des
Boson-Erzeugungszeilenoperators auf dem Zielblock ist \(\sqrt{b+1}\).
Damit liefert die Blocknorm-Abschätzung eine skalare Jacobi-Untergrenze
mit Diagonale \(b\) und quadrierten Nebendiagonalen

\[\adjustbox{max width=\linewidth}{$\displaystyle 
 a_b^2=\frac{15(b+1)(n-2b-\eta)}{800}
$}\]

in Einheiten \(\Delta=1\). Man nimmt negative Nebendiagonalen; die
Quadratformabschätzung benutzt die Normen der vollen
Bosonzahlkomponenten eines beliebigen Zustands, nicht einzelne
Testvektoren.

Für \textbf{jedes \(n=65,\ldots,127\)} wurden sämtliche LDL-Pivots der
Vergleichsmatrix oberhalb \(-4/5\) exakt positiv nachgewiesen. Der erste
Bosonindex ist \(\lceil(n-64)/2\rceil\), der letzte
\(\lfloor n/2\rfloor\). Somit gilt \(H_n>-.8\Delta\) in allen 63 hohen
Ladungssektoren.

Für alle \(n\le62\) liefert die Minimierung der Cauchy-Schwarz-Schranke

\[\adjustbox{max width=\linewidth}{$\displaystyle 
 H_n/\Delta\ge-\frac{n-\eta}{4}(\sqrt{23/20}-1)>-1.121899.
$}\]

Jede andere Ladungsverteilung als (63,64) oder (64,63) enthält eine Bank
mit \(n\le62\) und eine mit \(n\ge65\). Ihre Gesamtenergie liegt also
oberhalb \(-1.921899\Delta\). Die elektrische Energie ist nichtnegativ.
Das obere Bandende von \(P\) ist kleiner als
\(-2.225448\Delta+\kappa/2\).

Für die beiden zentralen Ladungsverteilungen wurden außerdem die
Vergleichsmatrizen ab \(b=1\) bei \(n=63,64\) oberhalb \(-3/4\) erneut
rational geprüft. Durch Minmax mit Kodimension 64 beziehungsweise eins
ergeben sich die bekannten angeregten Energieböden; es wird nicht
behauptet, dass der exakte Spektralprojektor mit dem
Null-Bosonenprojektor identisch ist.

Zusammen ergibt sich eine Lücke zwischen \(P\) und \(Q=I-P\) von
mindestens

\[\adjustbox{max width=\linewidth}{$\displaystyle 
 \boxed{\delta=.303549\Delta-\kappa/2>0.}
$}\]

Die anderen beiden Vergleichslücken sind \(.345812\Delta-\kappa/2\) und
\(.379636\Delta-\kappa/2\), also größer. Diese Schranke umfasst
\textbf{alle} Ladungsaufteilungen, Bosonzustände und internen Moden des
deklarierten physikalischen Sektors.

\subsubsection{4.3 Gleichmäßige Auslaufkontrolle und endliche
Transferzeit}\label{gleichmuxe4uxdfige-auslaufkontrolle-und-endliche-transferzeit}

Setze \(v=64|\tau|\) und fordere \(2v<\delta\). Sei \(a\) das obere
Eigenwertende von \(H_0|_P\). Minmax liefert ein Spektralband \(P'\) von
\(H\) derselben Dimension 128, dessen oberer Rand ≤\(a+v\) ist. Die
Kompression \(QHQ\) liegt ≥\(a+\delta-v\).

Mit \(Y=QP'\), als Abbildung aus \(\operatorname{Ran}P'\), gilt die
Sylvestergleichung

\[\adjustbox{max width=\linewidth}{$\displaystyle 
 (QHQ)Y-Y(H|_{P'})=-QVP\,PP'.
$}\]

Die beiden Spektren sind geordnet und mindestens \(\delta-2v\) getrennt.
Die Lösung über das konvergente Exponentialintegral liefert
\(\|Y\|\le v/(\delta-2v)\): nach einem gemeinsamen Skalarshift ist der
Integrand durch \(v e^{-s(\delta-2v)}\) beschränkt. Wegen gleicher
endlicher Ränge gilt dieselbe Schranke für \(\|P'-P\|\).

Da \(P'\) mit \(H\) kommutiert, folgt zu \textbf{jeder} Zeit

\[\adjustbox{max width=\linewidth}{$\displaystyle 
 \boxed{\|Qe^{-itH}P\|\le L:=\frac{2v}{\delta-2v}.}
$}\]

Aus der projizierten Schrödingergleichung und Duhamel folgt weiter

\[\adjustbox{max width=\linewidth}{$\displaystyle 
 \|Pe^{-itH}\psi-e^{-itPHP}\psi\|
 \le |t|vL=:\eta(t),\qquad \psi\in P,\ \|\psi\|=1.
$}\]

Für den Zielprojektor \(P_R\subset P\) weichen die Wahrscheinlichkeiten
damit höchstens um \(2\eta(t)\) ab. Dies ist eine nichtperturbative
Fehlergrenze für die volle Entwicklung; nur der ausgewählte
Kopplungsbereich ist klein. Es handelt sich nicht um das Weglassen eines
unbekannten höheren Störungsterms.

\subsubsection{4.4 Ein vollständig numerisch spezifizierter, rational
zertifizierter
Punkt}\label{ein-vollstuxe4ndig-numerisch-spezifizierter-rational-zertifizierter-punkt}

Mit \(\hbar=1\) wähle

\[\adjustbox{max width=\linewidth}{$\displaystyle 
 \tau/\Delta=10^{-8},\quad \kappa/\Delta=10^{-10},
 \quad T:=t\Delta=169500000.
$}\]

Diese Zahlen sind ein konservativer Existenzpunkt, \textbf{keine aus
TFPT abgeleiteten Konstanten und keine Behauptung schneller oder
optimaler Übertragung}. Die physikalische Sekundenskala ist nicht
bestimmt.

Die genaue komprimierte Rabi-Wahrscheinlichkeit ist

\[\adjustbox{max width=\linewidth}{$\displaystyle 
 p_P(t)=\frac{|\tau|^2Z^2}{|\tau|^2Z^2+(\kappa/4)^2}
 \sin^2\!\left(t\sqrt{|\tau|^2Z^2+(\kappa/4)^2}\right).
$}\]

Für \textbf{jedes} zulässige unbekannte \(Z\) liegt der Winkel weniger
als 0,08 von \(\pi/2\) entfernt. Das folgt aus
\(Z_{\rm lo},Z_{\rm hi}\), \(\sqrt{x^2+y^2}\le x+y^2/(2x)\) und
rationalen Pi-Grenzen. Letztere wurden zusätzlich mit der
Machin-Identität und endlichen alternierenden Arctan-Reihen
eingeschlossen; Fließkomma-Pi ist keine Beweisvoraussetzung.

Mit \(\sin^2(\pi/2+u)\ge1-u^2\) folgt

\[\adjustbox{max width=\linewidth}{$\displaystyle 
 p_P(T)>\left(1-\frac{\kappa^2}{16\tau^2Z_{\rm lo}^2}\right)
 (1-.08^2)>.993591982278.
$}\]

Die vollständigen Fehlergrenzen sind

\[\adjustbox{max width=\linewidth}{$\displaystyle 
 L=\frac{25600}{6070954399}<4.216800\cdot10^{-6},
 \quad L^2<1.8\cdot10^{-11},
 \quad\eta(T)=\frac{2777088}{6070954399}<.000457438455.
$}\]

Also

\[\adjustbox{max width=\linewidth}{$\displaystyle 
 \boxed{p_{\rm voll}(T)>p_P^{\rm lo}-2\eta(T)
 =\frac{25218291754002904578202303151726961}
 {25404324948782158166703488466197500}>.992677105368.}
$}\]

Die unbekannten \(E_0,E_h\) treten nur als gemeinsame Phase auf und
müssen für diese Schranke nicht genau ausgerechnet werden. Eine arbiträr
hohe Treue ist innerhalb des zusätzlichen Modells prinzipiell
erreichbar: bei bekanntem \(Z\), \(\kappa/|\tau|\to0\),
\(|\tau|/\Delta\to0\) und \(t\sim\pi/(2|\tau|Z)\) verschwinden sowohl
Detuning als auch der Fehler \(tv^2/\delta\). Das ist eine bedingte
Grenzaussage mit wachsender Laufzeit, kein ausführbares natives
Optimierungsverfahren.

\subsubsection{4.5 Dieselbe ursprüngliche Fermionantwort ohne
anfänglichen
Filter}\label{dieselbe-urspruxfcngliche-fermionantwort-ohne-anfuxe4nglichen-filter}

Für \(\psi_f=(f_{x,r}\Omega_x/\sqrt\nu)\otimes\Omega_y\otimes|0\rangle\)
gilt

\[\adjustbox{max width=\linewidth}{$\displaystyle 
 \psi_f=\sqrt w\,|L,r\rangle+\sqrt{1-w}\,q,
 \quad q\in Q,\quad w=Z/\nu\ge Z_{\rm lo}/Z_{\rm hi}.
$}\]

Die Auslaufnorm in umgekehrter Richtung erfüllt ebenfalls
\(\|Pe^{-itH}Q\|\le L\), indem man die obige Schranke adjungiert und die
Zeit umkehrt. Der unerwünschte Anteil kann also nicht beliebig stark
destruktiv in das niedrige Zielband einstreuen. Mit der
Dreiecksungleichung, \(2\sqrt{w(1-w)}\le1\) und der gesondert geprüften
Positivität vor dem Quadrieren folgt

\[\adjustbox{max width=\linewidth}{$\displaystyle 
 \boxed{\|P_Re^{-itH}\psi_f\|^2
 \ge w\,p_{\rm voll}^{\rm lo}-L
 >.897260353455>.897.}
$}\]

Damit ist \textbf{kein vorbereitender Projektor auf die isolierte Linie}
nötig, um eine starke Übertragungsaussage über die ursprüngliche
Entnahmeantwort zu erhalten. Noch benötigt werden die deklarierte
geladene Präparation und der zusätzliche Link. Die Aussage betrifft das
niedrige Zielband; sie behauptet nicht, dass die vollständige
hochenergetische Restantwort ebenfalls formtreu übertragen wird. Für
diesen ungefilterten Anfangszustand gilt die extrem kleine
Auslaufwahrscheinlichkeit aus 4.3 nicht: Er startet bereits teilweise
außerhalb von \(P\).

\subsection{5. Warum die Herkunftslücke nicht durch längeres Rechnen
verschwindet}\label{warum-die-herkunftsluxfccke-nicht-durch-luxe4ngeres-rechnen-verschwindet}

Sei \(\Pi_x=(-1)^{N_{f,x}}\) die lokale Fermionparität. Die bisher
angegebenen bankinternen Paarumwandlungen, Zahloperationen,
zahlenerhaltenden Symmetrielifts sowie reine Bosonverbindungen
kommutieren mit jeder \(\Pi_x\). Jede endliche Komposition, lineare
Kombination, Adjungierung und jeder durch solche Operatoren realisierte
einzelne Messzweig tut das ebenfalls. Starke beschränkte Grenzwerte
bleiben im Kommutanten der Parität.

Der benötigte Link erfüllt dagegen \(\Pi_xB=-B\Pi_x\) und
\(\Pi_yB=-B\Pi_y\), während er die Gesamtparität erhält. \textbf{Er
liegt nicht in dieser bekannten lokalen geraden Algebra.} Auch adaptives
Wiederholen und Nachselektieren gerader Krauszweige kann die
Sektorgrenze nicht überwinden.

Die Voraussetzung „jeder realisierte Zweig ist gerade`` ist wesentlich.
Eine bloß paritätskovariante CP-Abbildung kann ungerade Krausoperatoren
besitzen. Der im Prüfer enthaltene Reset-Kanal auf einer Fermionmode ist
ein exaktes Gegenbeispiel zur unzulässigen stärkeren Behauptung.

Das beweist keinen universellen Unmöglichkeitssatz für den Compiler. Es
beweist, \textbf{welche neue Quelleneigenschaft gesucht werden muss}:
eine gerade Gesamtoperation, die bezüglich der zwei operational
definierten Teilbereiche ungerade ist, oder eine ursprüngliche
Überlappung, durch die die angenommene Zerlegung in unabhängig gerade
Banken falsch war. Eine andere bloße Clockphase oder weitere
Boson-Paarumwandlung reicht nicht.

\subsection{6. Was von der arithmetischen Untersuchung
trägt}\label{was-von-der-arithmetischen-untersuchung-truxe4gt}

\subsubsection{6.1 Eine einzelne endliche Bank ist nicht der volle
Logarithmusgenerator}\label{eine-einzelne-endliche-bank-ist-nicht-der-volle-logarithmusgenerator}

Quadratisches Ergänzen liefert mit der nativen Casimirnorm
\(\|\sum P_A^\dagger P_A\|\le480\)

\[\adjustbox{max width=\linewidth}{$\displaystyle 
 H_{\rm nat}\ge\tfrac\Delta2N_b-C,
 \qquad C=960g^2/\Delta.
$}\]

Die externe Rechnung verwendet teilweise die schwächere Dreiecksgrenze
\(C=7680g^2/\Delta\), die ebenfalls gültig ist. Über Minmax,
\textbf{nicht} eine allgemeine Operator-Monotonie der
Exponentialfunktion, folgen

\[\adjustbox{max width=\linewidth}{$\displaystyle 
 N_H(E)\le2^{64}\binom{\lfloor2(E+C)/\Delta\rfloor+60}{60},
 \qquad \operatorname{Tr}e^{-\beta H}
 \le\frac{2^{64}e^{\beta C}}{(1-e^{-\beta\Delta/2})^{60}}.
$}\]

Das Zustandswachstum ist polynomial, während ein vollständiges Spektrum
\(E_*\log n+E_{\rm off}\), \(E_*>0\), exponentiell viele Zustände unter
wachsender Energie verlangt. Die ursprüngliche Bank kann dieses
komplette Spektrum daher nicht energieerhaltend mit nur affiner
Skalenänderung tragen. Diese Schranke widerlegt nicht jeden
arithmetischen Untersektor, jede nichtlineare Umparametrisierung oder
jeden Nullstellenoperator.

Endlich viele solche Banken und endlich viele energetisch kontrollierte
Rotoren ändern den qualitativen Gegensatz nicht. Unendlich viele
identische Banken lösen ihn nicht automatisch: Bei gleichmäßig
beschränkten lokalen Anregungskosten divergiert bereits die globale
Wärmespur. Ein geeigneter relativer oder lokaler Spurbegriff wäre ein
zusätzlicher Gegenstand.

\subsubsection{6.2 Primzahlen können eine gewählte Dynamik organisieren,
erzwingen sie aber
nicht}\label{primzahlen-kuxf6nnen-eine-gewuxe4hlte-dynamik-organisieren-erzwingen-sie-aber-nicht}

Auf \(\ell^2(\mathbb N)\) kann man
\(H_{\rm ar}|n\rangle=\log n|n\rangle\) definieren. Die additive
Primfaktorbesetzung erklärt dann die Eulerstruktur. Der
Definitionsbereich lautet \(\sum_n(\log n)^2|\psi_n|^2<\infty\). Dass
quantenstatistische Systeme mit Zeta-Zustandssumme existieren, ist
bereits ein Ergebnis von
\href{https://alainconnes.org/wp-content/uploads/bostconnesscan.pdf}{Bost
und Connes}. Das ist ein relevanter Anschluss, nicht die fehlende native
Herleitung.

Der gelieferte Cocycle \(c(m,n)=(-1)^{v_2(m)v_3(n)}\) ist assoziativ,
weil Bewertungen unter Multiplikation additiv sind. Für
\(T_m|n\rangle=c(m,n)|mn\rangle\) gilt \(T_lT_m=c(l,m)T_{lm}\),
insbesondere \(T_2T_3=-T_3T_2\). Die unkontrollierten CP-Kanäle
verlieren diese globale Phase und erfüllen
\(\mathcal E_l\mathcal E_m=\mathcal E_{lm}\). In kohärent kontrollierten
Wegen bleibt die relative Phase messbar. Das ist eine exakte
Phasenstruktur, aber kein Beweis einer Quellenauswahl dieses Cocycle
oder von RH.

Strikt multiplikative Schritte \(m>1\) ergeben auch keine geschlossenen
Bahnen: \(n\mapsto mn>n\). Die Frequenz \(\log p\) eines Operators, eine
Bahnlänge \(r\log p\) und der Schwingungszeitraum \(2\pi/\log p\) müssen
auseinandergehalten werden. Hier hat der RH-Graph-Skill die Prüfung auf
Generator, Spur und Phase getrennt und eine bloße Analogieschließung
verhindert.

\subsubsection{6.3 Der RH-Zielsatz bleibt
unbewiesen}\label{der-rh-zielsatz-bleibt-unbewiesen}

Die vorgeschlagene Identität

\[\adjustbox{max width=\linewidth}{$\displaystyle 
 \frac{\Xi(z)}{\Xi(0)}=\det(I-z^2A^{-2}),\quad
 \Xi(z)=\xi(\tfrac12+iz),
$}\]

für einen strikt positiven selbstadjungierten Operator \(A\) mit
kompakter Inverser und \(A^{-2}\) von Spurklasse wäre tatsächlich
hinreichend: Die gesamte rechte Funktion hat ihre Nullstellen nur bei
reellen \(\pm\lambda_j(A)\), einschließlich Multiplizitäten. Aber dieser
Operator und die gesamte Funktionsidentität sind \textbf{nicht
konstruiert}. Ein Operator mit lediglich den Imaginärteilen der
Nullstellen wäre unzureichend.

Weder diese Prüfung noch die im externen Code bereits verwendete
Faktorisierung kleiner Testzahlen liefert einen neuen effizienten
Faktorisierungsalgorithmus, eine native Quantenoperation dafür oder eine
Lösung von P versus NP. Hylæan wurde in dieser Revision nicht neu
geprüft.

\subsection{7. Die lokale Polnäherung wird präziser, das relativistische
Feld bleibt
offen}\label{die-lokale-polnuxe4herung-wird-pruxe4ziser-das-relativistische-feld-bleibt-offen}

Die externe Resolventenabschätzung ist korrekt: Im Kreis mit Radius
\(.1\Delta\) um den \textbf{tatsächlichen}, unbekannten negativen Pol
\(-\epsilon\) gilt

\[\adjustbox{max width=\linewidth}{$\displaystyle 
 G(z)=\frac{Z}{z+\epsilon}+R(z),\qquad
 |R(z)|<.505213/\Delta.
$}\]

Der übrige Spektralträger ist mindestens \(.337373\Delta\) vom Pol
entfernt. Die relative Abweichung vom echten Polterm beträgt im Kreis
unter 5,75 \%. Die Abschätzung setzt keine Intervallmitte anstelle von
\(Z,\epsilon\) ein und ist keine globale Vernachlässigbarkeit der
Selbstenergie.

Aus v1.6.4 bleibt bestehen: Die exakte Zwei-Feld-Resolvente besitzt
einen nichtverschwindenden Rest \(\Sigma_2(z)\). Seine komplette
Funktion ist noch nicht bestimmt. Das neue kontrollierte Transportlemma
benötigt sie nicht vollständig, weil es den Rest durch eine
Spektrallücke kontrolliert.

Ebenso bleibt der Feldtypentest bestehen: Der native antisymmetrische
Paartensor zusammen mit zwei gleichhändigen Weylfeldern und einem
skalaren Vermittler ergibt den verschwindenden Kanal. Ein symmetrischer
Lorentz- Spinortensor oder eine zusätzliche gleichgeladene
zweidimensionale Antisymmetriemarke eröffnet einen nichtverschwindenden
Kanal, führt aber zusätzliche Struktur ein. Ein internes
Spin(10)-Spinorlabel ist kein Lorentz-Weylindex. Der neue Link behebt
diese Typfrage nicht. Die Konventionen des bisherigen Tests sind mit der
Primärübersicht von \href{https://arxiv.org/abs/0812.1594}{Dreiner,
Haber und Martin} verbunden; ein neuer vollständiger relativistischer
Adapter wurde hier nicht konstruiert.

Die spätere Feldtypenprüfung in \codeword{LATE\_AUDIT.md} zeigt außerdem
einen Lorentz-erlaubten Zwei-Ableitungs-Kandidaten für den symmetrischen
Spinortensor und widerlegt einen pauschalen (1,1)-Ausschluss: Der
Energie-Impuls-Tensor ist schon bei Dimension vier zulässig. Weder ein
gesunder nativer Feldadapter noch ein dynamischer Gravitonpol folgt
daraus.

\subsection{8. Während der Arbeit entdeckte Änderung: „Kommutant 7``
richtig
einordnen}\label{wuxe4hrend-der-arbeit-entdeckte-uxe4nderung-kommutant-7-richtig-einordnen}

Die strikte Originalquellenprüfung stoppte korrekt, als eine andere
Arbeit den früheren nativen Bericht ergänzte. Der alte Stand mit Hash
\codeword{1062e13fe0fb79cce015657a7b4d40a61e9639922ff95fa60c689c0068c38ffc}
blieb eingefroren. Der gesonderte geänderte Entwurf hat Hash
\codeword{ab745e74b1232b1ddf35dd0e0d0da4ace0176c797d769e5db8bedd0010bcba6a}.
Die in dieser Revision verwendeten Energie- und Polzertifikate änderten
sich bei dieser Prüfung nicht. Der Fehlerbeleg wurde aufbewahrt.

Der neue Entwurf listet drei dunkle Darstellungstypen mit Dimensionen
24.000, 11.200, 2.688; drei helle mit 2.880, 576, 320; und einen χ-Typ
mit 64. \textbf{Unter Voraussetzung dieser Zerlegung} lässt sich seine
Folgerung exakt nachprüfen:

\[\adjustbox{max width=\linewidth}{$\displaystyle 
 \mathcal H_{N=3}\simeq
 \bigoplus_{d\in D}V_d\ \oplus\
 \bigoplus_{b\in B}(\mathbb C^2\otimes V_b)\ \oplus V_\chi.
$}\]

Die Summe ist \(37888+2\cdot3776+64=45504\). Die drei hellen Typen
kommen im \textbf{vollen} Raum jeweils zweimal vor. Daraus folgen:

\begin{itemize}
\tightlist
\item
  Nur die Gruppenwirkung hat einen Kommutanten der Dimension
  \(3+3\cdot4+1=16\), nicht sieben.
\item
  Werden zusätzlich die nichtverschwindenden hellen Paarmischungen und
  \(N_b\) als Operationen zugelassen, erzeugen sie die vollständigen
  \(M_2\)-Multiplizitätsalgebren. Dann bleiben sieben zentrale Skalare.
\item
  Sieben ist \textbf{keine absolute Untergrenze für beliebige weitere
  Operationen}. Ein erlaubter Operator, der diese sieben Blöcke
  verbunden mischt, zwingt die sieben Skalarwerte zur Gleichheit und
  lässt nur einen übrig. Der endliche Inzidenzmatrix-Test hat Rang sechs
  und beweist dieses Gegenargument.
\item
  Die Zerlegung allein beweist weder operative Verfügbarkeit der
  kontinuierlichen Symmetrie noch einen Einzelfermion-Link. Selbst alle
  bankinternen Symmetrieoperationen erhalten die lokale Fermionparität.
\end{itemize}

Der interessante Wert sieben wird dadurch nicht verworfen. Korrigiert
werden die Aussagen „voller N=3-Raum multiplizitätsfrei``, „absolute
Grenze für jede Operation`` und eine unbewiesene Gleichsetzung von
Symmetrie und verfügbarem Instrument. Die genaue Darstellungszerlegung
selbst bleibt in diesem Zusatz eine externe Voraussetzung, kein frisch
reproduzierter Satz.

Der Debugging-Skill führte hier zur Versionsprüfung statt zum Ersetzen
eines Pins. Das neue Manifest trennt deshalb einen erfolgreichen
\textbf{eingefrorenen Quellenreplay} von
\codeword{originals\_unchanged=false}.

\subsection{9. T1-T8: was dieses Ergebnis beiträgt und nicht
beiträgt}\label{t1-t8-was-dieses-ergebnis-beitruxe4gt-und-nicht-beitruxe4gt}

Die Begriffe folgen der aktuellen offenen Problemliste und ihrer
Forschungszuordnung. Keine Akzeptanzmarkierung wurde heraufgesetzt.

{\def\LTcaptype{none} % do not increment counter
\begin{longtable}[]{@{}
  >{\raggedright\arraybackslash}p{(\linewidth - 4\tabcolsep) * \real{0.3333}}
  >{\raggedright\arraybackslash}p{(\linewidth - 4\tabcolsep) * \real{0.3333}}
  >{\raggedright\arraybackslash}p{(\linewidth - 4\tabcolsep) * \real{0.3333}}@{}}
\toprule\noalign{}
\begin{minipage}[b]{\linewidth}\raggedright
Tor
\end{minipage} & \begin{minipage}[b]{\linewidth}\raggedright
Fehlender Nachweis
\end{minipage} & \begin{minipage}[b]{\linewidth}\raggedright
Beitrag dieser Revision
\end{minipage} \\
\midrule\noalign{}
\endhead
\bottomrule\noalign{}
\endlastfoot
T1 & Ursprung von P1/P2, Dimension, Compiler- und Zustandswahl &
Bekannte gerade Operationsalgebra ist für Einzeltransport unzureichend;
ursprüngliche Auswahl weiter offen \\
T2 & Tatsächlich halbgeladenes, markiertes E8-Feld samt Renormierung,
Energie- und Adjungiertenkontrolle & Der Ladung-eins-Lochzustand ist
kontrolliert, aber nicht mit dem fehlenden Half-Charge-Feld
identifiziert \\
T3 & Ein gemeinsamer, aus TFPT ausgewählter lokaler unitärer 3+1D-Parent
& Vollständiger bedingter Zwei-Banken-Transfertest; kein ausgewählter
räumlicher Parent \\
T4 & Chirales SM-Maß, Anomalien/Index, gleichmäßige Spiegelentkopplung &
Feldtyp-Widerspruch weiter sichtbar; kein chirales Maß \\
T5 & Wechselwirkender Kontinuumsübergang, Lorentzverhalten, Clustering,
Confinement und Streuung & Endliche-Sektor-Fehlerkontrolle ist ein
Baustein, aber kein räumlicher Grenzübergang \\
T6 & Interne Herleitung aller Eichkopplungen und vollständiger
Neutrinotextur/-skala & Keine neue Herleitung; \(\tau,\kappa\) sind
deklarierte Testparameter \\
T7 & Masseloser dynamischer Spin zwei, beide Helizitäten, universelle
Kopplung im selben Parent & Offen; der neue pauschale Ausschluss eines
Dimension-vier-(1,1)-Tensors wurde korrigiert \\
T8 & Physische Präparation/Anfangszustand und ein gemeinsames
Quellfunktional aller Auslesungen & Anfänglicher Polfilter wird für
\textgreater89,7 \% unnötig; geladene Präparation und Messinstrument
weiter offen \\
\end{longtable}
}

\subsection{10. Nächste entscheidende Arbeiten mit klaren
Abbruchkriterien}\label{nuxe4chste-entscheidende-arbeiten-mit-klaren-abbruchkriterien}

\textbf{A. Die Quelle des geraden Gesamtlinks finden.} Nicht noch eine
riesige Kontrollmatrix bauen, sondern für jeden tatsächlich
ursprünglichen Generator \(O\) prüfen, ob \([O,\Pi_x]\ne0\) bei
erhaltener Gesamtparität möglich ist. Eine positive Antwort muss
Generator, Teilbereichsdefinition, Hilbertraum, Gaussreferenz und
Matrixelement liefern. Bleiben alle Generatoren lokal gerade, ist die
Suche nach dem Link durch deren weitere Komposition beendet. Dann muss
die ursprüngliche Teilraumdefinition oder der Quellenvertrag explizit
geändert werden; man darf keinen Hopterm als Ergebnis ausgeben.

\textbf{B. Die ursprüngliche Entnahme als gemeinsames Instrument
realisieren.} Eine Referenzmode könnte Ladung erhalten, während
\(f_r\Omega\) präpariert und ausgelesen wird. Zu prüfen sind ihre
Herkunft, Energiekosten, Erfolgshäufigkeit und Erhaltung der internen
Kohärenz. Der neue Satz entfernt bereits die zusätzliche Pflicht eines
perfekten anfänglichen Energiefilters. Er entfernt nicht die Pflicht,
ein geladenes Instrument auszuweisen.

\textbf{C. Den Feldtyp vor räumlichen Rechnungen entscheiden.} Entweder
folgt ein passender Lorentzträger wirklich aus derselben Quelle, oder
der skalare Weyl-Ansatz wird in seiner jetzigen Form verworfen. Eine
bloße zusätzliche Verdopplung ohne Herkunft ist ein Modellvorschlag,
keine Ableitung.

\textbf{D. Erst dann viele Banken.} Der nächste räumliche Test benötigt
dieselbe Quelle, einen bestimmten Graphen, Lokalisierung der Operationen
und mit dem Volumen verträgliche Fehler- und Energiebounds. Ein einziges
neues Linklemma liefert weder Raumdimension noch Lichtkegel, chirale
Materie oder Gravitation.

\textbf{E. Arithmetik getrennt scharf halten.} Eine echte Verbindung
müsste eine gemeinsame Quellabbildung mit Generator,
Zustands-/Spurstruktur, Phasen und Domänen liefern. Die volle signierte
Weilform oder die ganze Determinanten- identität bleibt das Beweisziel.
Der bloße Auftritt von Primzahlen erfüllt keines davon. Der globale
RH-Index ließ sich wegen einer fehlenden historischen Quelle und
veränderter Review-Quellen nicht vollständig aktualisieren; die gezielte
Originalquellenprüfung ist enger als ein frischer Gesamtindex. Es wurde
kein RH-Abschlusskandidat registriert.

\subsection{11. Reproduktion und ehrliche
Liefergrenze}\label{reproduktion-und-ehrliche-liefergrenze}

\codeword{replay.py --frozen-only} führt die fünf Prüfer normal und mit
\codeword{-OO} aus, vergleicht beide JSON-Ergebnisse und die zwei
externen Originalberichte, prüft alle eingefrorenen Quellen und
protokolliert Änderungen der Originalorte gesondert. Ein Fehler führt zu
FAIL. Der zuvor aufgetretene Live-Quellenfehler bleibt in
\codeword{source\_drift\_failure\_receipt.json} sichtbar.

Die abgesicherten Zahlen stehen vollständig rational in
\codeword{two\_bank\_transfer\_normal.json}; Dezimalwerte dienen nur der
Lesbarkeit. Die analytische Beweiskette steht in Abschnitt 4. Das Paket
umfasst die unveränderten externen Arbeiten, deren Programme, die native
v1.6.4-Abhängigkeit und alle neuen Berichte. NumPy, SymPy und für das
ältere native Paket dessen separat dokumentierte Voraussetzungen sind zu
unterscheiden.

Die versionierte ausführliche Markdown-Lieferung enthält zusätzlich den
vollständigen eingefrorenen nativen Herleitungsbericht v1.6.4 als
historischen Anhang. Daneben gibt es ein kurzes Update und eine
bildliche Erklärung. \textbf{Die großen TFPT-Haupt-PDFs und die Webseite
werden in dieser Revision nicht als aktualisiert ausgegeben. Es gab
keinen Commit oder Push.}

\begin{center}\rule{0.5\linewidth}{0.5pt}\end{center}

\section{Nachprüfung der beiden zuletzt eingegangenen
Untersuchungen}\label{nachpruxfcfung-der-beiden-zuletzt-eingegangenen-untersuchungen}

Integrierter Nachtrag zu v1.6.5 · 15. September 2026

Quellen: die Anlagen \codeword{8e110ae1-19e2-4d47-bfc8-ebe4fe9ab81b} und
\codeword{b928b72a-6669-4c67-a75d-124095791c6f}, unverändert
eingefroren. Zusätzlich wurden die vier genannten Programme
\codeword{t1\_fixed.py}, \codeword{stabilizer.py},
\codeword{t5\_twobank.py}, \codeword{t5\_varresp.py} vollständig gelesen
und archiviert. Eine eigene Nachrechnung prüft die kritischen
Folgerungen; sie ist \textbf{kein vollständiger Replay aller sechs
Sonden} des fremden Arbeitsordners.

\subsection{1. Übersicht: übernehmen, korrigieren oder
offenlassen}\label{uxfcbersicht-uxfcbernehmen-korrigieren-oder-offenlassen}

{\def\LTcaptype{none} % do not increment counter
\begin{longtable}[]{@{}
  >{\raggedright\arraybackslash}p{(\linewidth - 2\tabcolsep) * \real{0.5000}}
  >{\raggedright\arraybackslash}p{(\linewidth - 2\tabcolsep) * \real{0.5000}}@{}}
\toprule\noalign{}
\begin{minipage}[b]{\linewidth}\raggedright
Neue Aussage
\end{minipage} & \begin{minipage}[b]{\linewidth}\raggedright
Ergebnis der Prüfung
\end{minipage} \\
\midrule\noalign{}
\endhead
\bottomrule\noalign{}
\endlastfoot
Innerer Viererzyklus erhält alle 60 W-Kanäle & Bestätigt, jetzt
einschließlich der richtigen Fermion-Vorzeichen und explizitem
Bosonlift \\
Grundzustandseindeutigkeit/Singulett/Gaplücke weiterhin unbekannt &
Überholter Stand: unter dem festgelegten nativen Vertrag bereits
bewiesen; voller Vektor weiter unbekannt \\
Leerer Grundzustand für alle μ≥0 & Falsch; bei μ=0 liegt der eindeutige
N=64-Grundzustand unter −1,129636Δ \\
Kommutant 7 als absolute Untergrenze & Nur im benannten blocktreuen
erweiterten Operationsvertrag; siehe Hauptbericht Abschnitt 8 \\
Beliebige Einteilchenmatrix ist damit als natives Fock-Wort verfügbar &
Nicht gezeigt; assoziative Matrixalgebra, Lie-Kontrolle und Fock-Lift
sind verschiedene Fragen \\
Kanal-Variationszustand hat E=−0,0196Δ & Die angegebene Phase im Code
ergibt gegen den gepinnten P=f\_jf\_i-Vertrag +0,05736Δ; ein relatives
Vorzeichen korrigiert es \\
99,892-\%-Transfer bei t≈2484 & Numerisches Maximum eines
Vierzustands-Paarmodells im untersuchten Zeitfenster; kein
erster/globaler exakter Maximalsatz und kein nativer Einloch-Transfer \\
Gleiche Zweigableitungen genau bei ε=0 & Algebraisch richtig; die
Eigenwertlücke bleibt dort √32· \\
SU(4)³-Anomalie =16 & Richtig unter A(4)=1 und der zusätzlichen
Interpretation aller (16,4) als gleichhändige Weylfelder; Eichung ist
eine weitere Voraussetzung \\
Gravitationsanomalie =64 & Keine reine perturbative Gravitationsanomalie
in 3+1D; 64 kann bei zusätzlich deklarierter gleichgeladener U(1) die
gemischte U(1)-Gravitationsanomalie zählen \\
Mit einer Ableitung kein (1,1)-Tensor bei Dimension ≤4 & Falsch: der
Energie-Impuls-Tensor ist ein Gegenbeispiel; daraus folgt aber noch kein
dynamisches Graviton \\
Kein lokaler kovarianter kinetischer Term für (1,0) & Zu pauschal: der
bilineare Ein-Ableitungs-Term fehlt, ein Zwei-Ableitungs-Skalar
existiert; Positivität/Constraints/Herkunft bleiben zu prüfen \\
\end{longtable}
}

„Negativ geschlossene Route`` bedeutet nicht „T2, T3, T4 oder T7
geschlossen``. Ein ausgeschlossener Ansatz beantwortet nicht die
jeweilige physikalische Existenzfrage. Die beiden gelieferten Texte
widersprechen sich vor allem beim bereits bewiesenen modellinternen
Grundzustand; sie dürfen nicht als gleichzeitig aktueller einheitlicher
Status zitiert werden.

\subsection{2. Positiver neuer Anschluss: der korrekt angehobene
Viererzyklus}\label{positiver-neuer-anschluss-der-korrekt-angehobene-viererzyklus}

Die Abbildung der Fermionmarken lautet \(p(4s+a)=4s+(a+1\bmod4)\). Auf
Paaren ist zwingend das Exteriorvorzeichen zu beachten: Falls
\(p(i)>p(j)\), erhält das sortierte Paar ein Minuszeichen. Der
gelieferte \codeword{stabilizer.py} lässt dieses Zeichen weg. Das wäre
im Allgemeinen falsch; für den untersuchten Viererzyklus überlebt das
positive Resultat jedoch auch den korrekten Test.

Alle 60 W-Zeilen werden mit ihren Vorzeichen wieder auf W-Zeilen
abgebildet. Wir haben den zugehörigen signierten Bosonoperator \(R_b\)
explizit gebildet und exakt geprüft:

\[\adjustbox{max width=\linewidth}{$\displaystyle 
 W'=R_bW,\qquad R_bR_b^T=I_{60},\qquad R_b^4=I_{60},
 \qquad R_b^2\ne I_{60}.
$}\]

Damit existiert ein \textbf{gemeinsamer nativer Tensor-Automorphismus}
auf Fermionen und Bosonen, nicht nur eine Übereinstimmung von
Zeilenzahlen. Die gleichzeitige Transformation erhält Paarwechselwirkung
und Bosonzahl. Sie beweist eine Ordnung-vier-Symmetrie, nicht deren
operative Verfügbarkeit.

Dieser Permutationszyklus ist nicht die zentrale Matrix \(iI_4\) von
SU(4). Seine Determinante auf dem Viererfaktor ist −1. Mit einer
zusätzlichen Ladungsphase kann man geeignete Gruppenlifts vergleichen;
die Permutation, die zentrale Phase, der Clock und die geometrische
Glue-Markierung sind dadurch aber nicht schon identifiziert. Genau diese
Intertwiner-Frage ist ein sinnvoller neuer Anschluss, statt noch einmal
nur vier zu zählen.

\subsection{3. Warum der Einteilchen-Abschluss den Fock-Operationssatz
nicht
schließt}\label{warum-der-einteilchen-abschluss-den-fock-operationssatz-nicht-schlieuxdft}

Die Irreduzibilität der Einteilchendarstellung kann ihre assoziative
Matrixalgebra zu \(B(\mathbb C^{64})\) machen. Daraus folgt nicht, dass
jede Matrix als ein zulässiges physisches Wort verfügbar ist, und auch
nicht, dass ihre zweite Quantisierung bereits erzeugt wird.

Ein kleinstes exaktes Gegenbeispiel zur falschen Liftregel:
\(A=|1\rangle\langle1|\), \(B=|2\rangle\langle2|\) auf zwei Moden. Dann
\(AB=0\), also \(d\Gamma(AB)=0\), aber
\(d\Gamma(A)d\Gamma(B)=n_1n_2\ne0\). Der zweite Quantisierungsschritt
ist ein Lie-, nicht ein assoziativer Algebra-Homomorphismus dieser Art.

Auch \(\operatorname{diag}(i,1,1,1)\) und
\(\operatorname{diag}(i,i,1,1)\) haben Determinanten i beziehungsweise
−1. Ihre Zugehörigkeit zur \textbf{linearen Matrixalgebra} aus
SU(4)-Generatoren und Identität ist kein exaktes SU(4)-Gruppenwort.
Projektive Gleichheit oder eine zusätzliche U(1)-Phase kann helfen, muss
aber mit der Bosonwirkung, Ladung und Referenz gemeinsam ausgewiesen
werden.

Der fremde Code folgert außerdem die Irreduzibilität des
37.888-dimensionalen dunklen Dreifermionraums aus der vollen
Einteilchenmatrixalgebra. Diese Schlussregel ist nicht begründet.
Bereits die im anderen neuen Text angegebenen drei dunklen irreduziblen
Typen widersprechen der pauschalen Eins-Block-Lesart unter der bloßen
Quellsymmetrie. Die Zahl 8.732.673 mag als sehr schwache obere Schranke
anderweitig verträglich sein; die im Code angegebene Herleitung belegt
sie nicht. Die SVD-/Toleranzprüfungen des Codes sind zudem numerisch,
nicht allein wegen „PASS`` exakte Lie-Beweise.

\subsection{4. Zustandswahl: keine Rückkehr vor den abgesicherten
Grundsatz}\label{zustandswahl-keine-ruxfcckkehr-vor-den-abgesicherten-grundsatz}

Der vorhandene native Satz beweist Eindeutigkeit, N=64, Singulett und
positive Lücke bei μ=0 im angegebenen Kopplungsbereich. Eine komplette
Clebsch-Gordan-Serie oder Voll-Diagonalisierung ist dafür nicht nötig;
Vergleichsungleichungen und Minmax waren gerade der einfache Ausweg. Die
physische Auswahl von H beziehungsweise μ=0 bleibt eine andere Frage.

Das grobe Sandwich \([-1.2,-.0196]\Delta\) ist nach einer
Phasenkorrektur zwar verträglich, aber wesentlich schwächer als
\((-1.158089,-1.129636)\Delta\). Es ist keine Verschärfung.

Im tatsächlich angegebenen Variationscode werden die Fermionen in der
Reihenfolge j, dann i gelöscht. Das liefert \(f_if_jF=-f_jf_iF\),
während der native Vertrag \(P_A=\sum W_{A,ij}f_jf_i\) verwendet. Bei
unverändertem \(g=+\Delta/20\) und den angegebenen
Variationskoeffizienten folgt deshalb

\[\adjustbox{max width=\linewidth}{$\displaystyle 
 E_{\rm Code}/\Delta=\frac12-\frac{23\sqrt3}{90}
 =+.057364793621\ldots,
$}\]

nicht der behauptete negative Wert. Mit korrigierter relativer Phase
lautet er \(1/2-3\sqrt3/10=-.019615242271\ldots\). Alternativ wäre eine
konsequente andere P/g-Phasenkonvention möglich; die Quelle muss sie
dann überall führen. Die gleichzeitige Einteilchendichte des einfachen
Variationszustands bleibt von dieser Phasenreparatur unberührt, weil
seine Bosonzahlkomponenten orthogonal sind. Sie ist keine dynamische
Greenfunktion auf dem nativen Grundzustand und ersetzt dessen schon
kontrollierte Antwort nicht.

\subsection{5. Feldtheorie: falsche Verbote entfernen, echte Bedingungen
behalten}\label{feldtheorie-falsche-verbote-entfernen-echte-bedingungen-behalten}

\subsubsection{5.1 Anomalien}\label{anomalien}

Für die zusätzlich angenommene 3+1D-Weylinterpretation lautet das lokale
Anomaliepolynom bis auf Konventionsvorzeichen

\[\adjustbox{max width=\linewidth}{$\displaystyle 
 I_6=[\widehat A(T)\,\mathrm{ch}_R(F)]_6
 =\mathrm{ch}_3(F)-\frac{p_1(T)}{24}\,\mathrm{ch}_1(F).
$}\]

Der reine gravitative Anteil \([\widehat A]_6\) verschwindet; seine
Formgrade sind Vielfache von vier. Das ist die bekannte dimensionale
Unterscheidung bei
\href{https://collaborate.princeton.edu/en/publications/gravitational-anomalies/}{Álvarez-Gaumé
und Witten}. Eine Spur-/Weylanomalie ist nochmals ein anderer Begriff.

Auf (16,4) ist der SU(4)-Kubikkoeffizient 16 bei Normierung A(4)=1. Der
exakte Test \(t=\operatorname{diag}(1,1,1,-3)\) hat
\(\operatorname{tr}t=0\), \(\operatorname{tr}t^3=-24\). Ist SU(4)
\textbf{dynamisch geeicht}, muss die vollständige Theorie diese
Eichanomalie kompensieren. Als globale Symmetrie kann sie eine
't-Hooft-Anomalie tragen; dann folgt keine pauschale Spiegelpflicht. Ein
vollständiger konjugierter Spiegel ist ein möglicher Ausgleich, aber
hier nicht als einzige oder allgemein minimale Lösung bewiesen.

„Gravitativ 64`` kann sinnvoll eine \textbf{gemischte
U(1)-Gravitationsanomalie} meinen, wenn alle 64 linksgetragenen
Weylfelder explizit U(1)-Ladung eins erhalten. Das muss samt
Eich-/Globalstatus gesagt werden. Für SU(4) allein ist der gemischte
lineare Spurkoeffizient null. Die Quellenmarke N darf nicht ohne Beweis
in eine dynamisch geeichte chirale U(1) umgedeutet werden.

\subsubsection{5.2 Der behauptete Spin-2-Ausschluss ist
falsch}\label{der-behauptete-spin-2-ausschluss-ist-falsch}

Die Lorentzdarstellungen liefern exakt

\[\adjustbox{max width=\linewidth}{$\displaystyle 
 (\tfrac12,\tfrac12)\otimes(\tfrac12,\tfrac12)
 =(0,0)\oplus(1,0)\oplus(0,1)\oplus(1,1).
$}\]

Eine Ableitung des Vektorbilinears \(\psi^\dagger\bar\sigma_\mu\psi\)
kann daher einen symmetrischen spurfreien (1,1)-Tensor bilden.
Insbesondere hat der Energie-Impuls-Tensor
\(T_{\mu\nu}\sim i\psi^\dagger\bar\sigma_{(\mu}
\overleftrightarrow\partial_{\nu)}\psi\), mit passender Spurbehandlung,
Dimension \(3/2+3/2+1=4\). Ein Energie-Impuls-Tensor kann bereits in
einer festen Hintergrundraumzeit definiert werden; dynamische
Diffeomorphismen müssen nicht vorausgesetzt werden, um dieses
Gegenbeispiel zu formulieren.

Das schließt \textbf{T7 nicht}. Ein vorhandener Tensor ist kein
masseloser Spin-2-Pol. Das sinnvolle Ziel ist später sein transversaler,
spurfreier Korrelator auf demselben räumlichen Parent, mit positiver
Norm, beiden Helizitäten und universeller Kopplung. Wir entfernen hier
eine falsche Darstellungsobstruktion, nicht die dynamischen
Nachweispflichten.

\subsubsection{5.3 (1,0)-Kinetik ist nicht generell
verboten}\label{kinetik-ist-nicht-generell-verboten}

Der bilineare Ein-Ableitungs-Term mit einem (1,0)-Feld und seinem
Adjungierten besitzt keinen Lorentzskalar. Mit zwei Ableitungen
existiert aber etwa

\[\adjustbox{max width=\linewidth}{$\displaystyle 
 \mathcal L_2\propto
 (\partial^{\alpha\dot\alpha}\Phi_{\alpha\beta})
 (\partial^{\beta\dot\beta}\bar\Phi_{\dot\alpha\dot\beta})
$}\]

für symmetrisches \(\Phi\). Alle Spinorindizes sind kontrahiert. Bei
rein zeitartigem Impuls ist sein Symbol proportional zu
\(\omega^2(|\Phi_{11}|^2+2|\Phi_{12}|^2+|\Phi_{22}|^2)\), also nicht
identisch null. Das ist ein Gegenbeispiel zum uneingeschränkten Satz
„kein lokaler kovarianter kinetischer Term``. Es beweist \textbf{nicht}
Positivität des vollen Hamiltonoperators, korrekte Constraints,
gewünschte Freiheitsgrade oder eine native Quelle. Diese bleiben die
entscheidenden Tests.

\subsection{6. Die Vierzustandsrechnung richtig
lesen}\label{die-vierzustandsrechnung-richtig-lesen}

Das Modell mit \(s=\sqrt8g\), Bosonverbindung η und Diagonalen (0,Δ,Δ,0)
erlaubt kohärenten \textbf{Paartransport}. Die exakte Identität
\((H^3)_{41}=8g^2\eta\) stimmt. Unsere unabhängige Fließkomma-Abtastung
reproduziert auf \([0,3000]\) mit Schrittweite .05 den größten
gefundenen Wert .998920021869 bei 2484.05. Der erste
Gitter-Lokalhöchstwert liegt bereits bei 6.05. Weder „erster
Maximalzeitpunkt`` noch ein globales exaktes Maximum ist damit bewiesen.
Diese kleine Rechnung wird nicht mit dem vollständigen Einloch-Satz des
Hauptberichts verwechselt.

Für \(F_\pm(\epsilon)=(\epsilon\pm\sqrt{\epsilon^2+32g^2})/2\) gilt
\(F'_+-F'_-=\epsilon/\sqrt{\epsilon^2+32g^2}\). Bei ε=0 stimmen die
Ableitungen überein, aber die \textbf{Eigenwertlücke} beträgt
\(\sqrt{32}|g|>0\). Ein verschwindender nackter Vermittlerparameter und
eine verschwindende wechselwirkende Anregungslücke sind verschieden. Ein
allgemeiner relativistischer Kegelsatz oder ein zwingendes neues
Renormierungsprogramm folgt aus dieser Ableitungsidentität nicht.

\subsection{7. Konsequenz für die nächste
Forschungsrevision}\label{konsequenz-fuxfcr-die-nuxe4chste-forschungsrevision}

Die sinnvollen positiven Anschlüsse sind jetzt schärfer:

\begin{enumerate}
\def\labelenumi{\arabic{enumi}.}
\tightlist
\item
  Den expliziten Ordnung-vier-Lift mit dem tatsächlich markierten Clock
  und der Herkunft des A3-Index vergleichen: vollständige Intertwiner,
  Ladungsphase und Bosonwirkung, nicht nur Gruppennamen.
\item
  Einen Quellenoperator finden, der die lokale Parität beider
  operational bestimmten Teile gleichzeitig ändert. Die native
  Z4-Symmetrie selbst tut das nicht. Der kontrollierte Einloch-Transfer
  ist bereits als Akzeptanztest für einen solchen Operator verfügbar.
\item
  Den Zwei-Ableitungs-Feldkandidaten auf Positivität und Constraints
  prüfen, bevor er als relativistischer Adapter zählt. Ein bloßes
  Nichtnull-Symbol ist noch keine gesunde Feldtheorie.
\item
  Den Energie-Impuls-Tensor als zulässigen \textbf{Kandidaten} behalten,
  aber erst auf dem gemeinsam konstruierten räumlichen Träger nach einem
  dynamischen Spin-2-Pol suchen. Kein vorhandener Spin-2-No-go
  rechtfertigt hier das Aufgeben dieses Anschlusses.
\end{enumerate}

62 zusätzliche Prüfbedingungen bestehen normal und optimiert identisch.
Exakte Identitäten, bedingte Darstellungsargumente und die numerische
Zeitfenstersuche sind im Ergebnisbericht getrennt. Der externe Status
„alle Restfragen geschlossen oder auf genau ein Stück reduziert`` wird
nicht übernommen. Die vollständige TOE bleibt offen.

\begin{center}\rule{0.5\linewidth}{0.5pt}\end{center}

\section{Historischer Beweisanhang: vollständiger nativer
Herleitungsstand
v1.6.4}\label{historischer-beweisanhang-vollstuxe4ndiger-nativer-herleitungsstand-v1.6.4}

Dieser Anhang bewahrt die frühere Herleitung vollständig. Seine
damaligen Statussätze werden durch die aktuelle Konsolidierung und den
Nachtrag oben ergänzt; er ist kein ungeprüft übernommener neuer
Gesamtstatus. Die externe laufende Symmetrieergänzung wurde separat
geprüft und ist nicht unbemerkt in diesen eingefrorenen Anhang
eingegangen.

\section{TFPT / Universalraum: nativer Pol, Bewegungsgleichung und
minimale
Feldtypen}\label{tfpt-universalraum-nativer-pol-bewegungsgleichung-und-minimale-feldtypen}

\textbf{Konsolidierte Forschungsfortsetzung v1.6.4 · 15. September 2026}

Diese Revision verbindet den während der Arbeit neu eingegangenen
Polsatz v1.6.3 mit einer unabhängig begonnenen Untersuchung der nativen
Bewegungsgleichung. Sie ergänzt die bisherigen Hauptdokumente; sie ist
keine verkürzte Neufassung des vollständigen Hauptbuchs. Die vorhandenen
Haupt- und Update-PDFs wurden in dieser Runde nicht verändert.

\subsection{1. Ergebnis in einem Satz}\label{ergebnis-in-einem-satz}

Auf dem abgesicherten Grundzustand des festgelegten nativen Fockmodells
existiert eine isolierte ursprüngliche Fermion-Entnahmelinie; nach
Kombination beider Rechnungen trägt sie \textbf{mehr als 88,007628 \%
des gesamten normierten Spektralgewichts pro Mode}. Die zusätzliche
Antwort ist jedoch nicht exakt auf eine einzige weitere Linie
reduzierbar. Ihr erster Rückwirkungsschritt wird direkt von derselben
ursprünglichen Wechselwirkung bestimmt.

Am Prüfpunkt \(g/\Delta=1/20\), ausdrücklich ohne \(\mu N\)-Zusatz:

{\def\LTcaptype{none} % do not increment counter
\begin{longtable}[]{@{}
  >{\raggedright\arraybackslash}p{(\linewidth - 4\tabcolsep) * \real{0.2727}}
  >{\raggedleft\arraybackslash}p{(\linewidth - 4\tabcolsep) * \real{0.3636}}
  >{\raggedleft\arraybackslash}p{(\linewidth - 4\tabcolsep) * \real{0.3636}}@{}}
\toprule\noalign{}
\begin{minipage}[b]{\linewidth}\raggedright
Größe
\end{minipage} & \begin{minipage}[b]{\linewidth}\raggedleft
Eingegangene v1.6.3
\end{minipage} & \begin{minipage}[b]{\linewidth}\raggedleft
Konsolidierte strengere Grenze
\end{minipage} \\
\midrule\noalign{}
\endhead
\bottomrule\noalign{}
\endlastfoot
Mittlere Bosonenzahl & \(0.77<\bar b<1.45\) &
\(0.842846<\bar b<1.245656\) \\
Energie der niedrigen Entnahmelinie &
\(0.007737<\epsilon/\Delta<0.062277\) &
\(0.007737<\epsilon/\Delta<0.039079764\) \\
Gewicht dieser Linie im gesamten CAR-Maß &
\(Z_{\rm low}>0.864972353\ldots\) &
\(Z_{\rm low}>0.880076280689\ldots\) \\
Obere Gewichtsgrenze & \(Z_{\rm low}<0.9759375\) &
\(Z_{\rm low}<0.9736610625\) \\
Übrige Entnahmeenergien & \(>0.379636\Delta\) & unverändert \\
Sämtliche Additionsenergien & \(>0.329636\Delta\) & unverändert \\
\end{longtable}
}

Die Dezimalzahlen sind gerundete Darstellungen rationaler Schranken. Es
sind weder exakte Polpositionen noch angepasste Zentralwerte. Der Pol
steht im retardierten Spektrum bei negativer Frequenz \(-\epsilon\). Das
niedrige Niveau im N=63-Hilbertraum ist 64-fach entartet; jede diagonale
Modenantwort sieht dieselbe Linie.

\textbf{Nicht bewiesen:} die physische Herleitung des Hamiltonoperators,
sein vollständiger Operationssatz, native Präparation, räumliche
Ausbreitung, ein vollständiges relativistisches Feldwörterbuch oder eine
TOE. T1-T8 bleiben als vollständige Aufgaben offen.

\subsection{2. Ein unveränderter
Modellvertrag}\label{ein-unveruxe4nderter-modellvertrag}

Es werden keine Hopping-, Massen-, Ladungs- oder Projektionsglieder zu H
hinzugefügt:

\[\adjustbox{max width=\linewidth}{$\displaystyle 
H=\Delta N_b+g(Q_++Q_-),\quad Q_+=\sum_A b_A^\dagger P_A,
\quad Q_-=Q_+^\dagger,
$}\] \[\adjustbox{max width=\linewidth}{$\displaystyle 
P_A=\sum_{i<j}W_{A,ij}f_jf_i,\quad N=N_f+2N_b,
\quad WW^\dagger=8I_{60}.
$}\]

Es gibt 64 CAR-Fermionmoden, 60 CCR-Bosonmoden und 480 von null
verschiedene reelle W-Einträge mit ihren ursprünglichen Vorzeichen.
\(\Delta>0\), g ist reell; die Zahlen beziehen sich auf
\(g/\Delta=1/20\). Die Tensorquelldatei ist durch SHA-256
\codeword{3f00a0892157a691e4ef26920c3702772c7bf76a250fad98ff04a079bb64b763}
fixiert.

Die Fockrealisierung, H, der sektorenübergreifende Energievergleich und
der Prüfparameter sind weiterhin Modellvoraussetzungen. Die geometrische
oder arithmetische Herkunft eines W-Tensors leitet diese Voraussetzungen
nicht allein her. Insbesondere sind die fünf elementaren CAR in einer
Spinor-Matrixkonstruktion nicht mit den hier verwendeten 64 Fockmoden
gleichzusetzen.

\subsection{3. Was tatsächlich erneut geprüft
wurde}\label{was-tatsuxe4chlich-erneut-gepruxfcft-wurde}

\subsubsection{3.1 Grundzustand: frische Berechnung statt bloßer
Statusübernahme}\label{grundzustand-frische-berechnung-statt-blouxdfer-statusuxfcbernahme}

Die ursprünglichen Programme wurden mit fixierten Dateihashes in eine
eigene Arbeitskopie übernommen. Die Paar-Konfigurationen wurden frisch
mit neu kompiliertem Quellcode berechnet:

{\def\LTcaptype{none} % do not increment counter
\begin{longtable}[]{@{}
  >{\raggedright\arraybackslash}p{(\linewidth - 4\tabcolsep) * \real{0.2727}}
  >{\raggedleft\arraybackslash}p{(\linewidth - 4\tabcolsep) * \real{0.3636}}
  >{\raggedleft\arraybackslash}p{(\linewidth - 4\tabcolsep) * \real{0.3636}}@{}}
\toprule\noalign{}
\begin{minipage}[b]{\linewidth}\raggedright
Ordnung
\end{minipage} & \begin{minipage}[b]{\linewidth}\raggedleft
Vollständig enumerierte Bosonkonfigurationen
\end{minipage} & \begin{minipage}[b]{\linewidth}\raggedleft
Normquadrat von \(Q_+^kF\)
\end{minipage} \\
\midrule\noalign{}
\endhead
\bottomrule\noalign{}
\endlastfoot
0 & analytischer Ausgangszustand & 1 \\
1 & 60 Kanäle / 480 Paare & 480 \\
2 & 1830 & 439680 \\
3 & 37820 & 575078400 \\
4 & 595665 & 952296652800 \\
\end{longtable}
}

Zusätzlich wurden unabhängige Wortentwicklungen, Überlaufgrenzen, die
vollständige Zweikörper-Casimiridentität, rationale Sektorvergleiche und
die unabhängige Symmetriespurrechnung wiederholt. Die vier ausgewählten
ursprünglichen Zertifikate bestehen normal und unter \codeword{-OO} mit
identischen JSON-Bytes. Dies ist eine gezielte erneute
Grundzustandsprüfung, nicht die Behauptung, alle historischen
Repository-Tests erneut ausgeführt zu haben. Der unveränderte alte
Normprüfer gibt unter dem aktuellen NumPy eine ComplexWarning bei der
Ganzzahlkonversion aus. Die neuen Prüfer bestätigen vor jeder Konversion
exakt verschwindende Imaginärteile und ganzzahlige Realteile des
gepinnten W. Die Warnung ist gespeichert und wird nicht als verlorener
Tensoranteil oder als unterdrückter Prüffehler ausgegeben.

Der modellinterne Satz bleibt bestehen: Für \(0<|g|/\Delta\le1/20\) ist
der globale Grundzustand \(\Omega\) eindeutig, hat N=64 und ist
Spin(10)×SU(4)-invariant. Am Zahlenprüfpunkt:

\[\adjustbox{max width=\linewidth}{$\displaystyle 
-1.158089\Delta<E_0<-1.129636\Delta,
\qquad \operatorname{gap}(H)>0.007737\Delta.
$}\]

Der volle Grundzustandsvektor ist damit nicht ausgerechnet. Die fünf
Versuchsvektoren sind keine behauptete invariante Fünferbasis für H.

\subsubsection{3.2 Das neu eingegangene
Polergebnis}\label{das-neu-eingegangene-polergebnis}

Der vollständige technische Bericht v1.6.3 und sein 248-zeiliger Prüfer
wurden gelesen; Quellen, Bericht, Ergebnismatrix und Prüfpaket wurden
unverändert archiviert. Der Prüfer wurde normal und optimiert
wiederholt: \textbf{317 Prüfbedingungen, identische JSON-Bytes auch zum
gelieferten Bericht}.

Der Herkunftspin des Prüfers ist
\codeword{fdbcabd244450c302182086d67c68284634c994e3fee026200c42c508f731654}.
Er verwendet denselben ursprünglichen
\codeword{verify\_hole.py}-Quellstand wie die bisherige Fortsetzung. Die
neuen eigenen Sektorvergleiche wurden zusätzlich unabhängig mit
rationalen LDL-Pivots berechnet.

Die Prüfzahlen zählen auch Komponenten und Wiederholungen; sie sind
keine Zahl unabhängiger Entdeckungen. Die analytischen Argumente sind
ausgeschrieben, aber nicht vollständig in einem Beweisassistenten
formalisiert.

\subsection{4. Der einfache native Anschluss: Bewegungsgleichung statt
Umdeutung}\label{der-einfache-native-anschluss-bewegungsgleichung-statt-umdeutung}

Erweitere jede W-Zeile zur antisymmetrischen Matrix \(M_A\) mit
\((M_A)_{ij}=W_{A,ij}\) für i\textless j. Definiere

\[\adjustbox{max width=\linewidth}{$\displaystyle 
D_r=\sum_{A,j}(M_A)_{rj}b_Af_j^\dagger,
\qquad \phi_r=D_r/\sqrt{15}.
$}\]

Direkte CAR/CCR-Normalordnung liefert \textbf{Operatoridentitäten auf
dem endlichen Teilchenkern}, nicht nur Gleichheiten auf Testzuständen:

\[\adjustbox{max width=\linewidth}{$\displaystyle 
\boxed{[H,f_r]=-gD_r,\qquad [H,f_r^\dagger]=gD_r^\dagger,}
$}\] \[\adjustbox{max width=\linewidth}{$\displaystyle 
[N,D_r]=-D_r,\qquad \{f_r,D_s^\dagger\}=0.
$}\]

Der neue Vergleichsoperator hat also dieselbe Ladung −1 wie f. Er ist
nicht das früher auf dem leeren Referenzzustand betrachtete
\(\chi^\dagger\sim b^\dagger f^\dagger\) mit Ladung +3. Auf dem
Grundzustand liegen \(f\Omega\) und \(D\Omega\) in N=63,
\(\chi^\dagger\Omega\) dagegen in N=67. Neutrale Zeitentwicklung hebt
diese Unterscheidung nicht auf.

Die Normalordnungsprüfung erfasst alle 64 ursprünglichen f-Operatoren
und die tatsächlichen W-Vorzeichen. Für die Kompositnorm gilt exakt

\[\adjustbox{max width=\linewidth}{$\displaystyle 
\{D_r,D_s^\dagger\}=
\sum_{A,B,j,k}(M_A)_{rj}(M_B)_{sk}
\left(\delta_{AB}f_j^\dagger f_k+\delta_{jk}b_B^\dagger b_A\right).
$}\]

Mit \(\bar b=\langle N_b\rangle\) und der Grundzustandssymmetrie:

\[\adjustbox{max width=\linewidth}{$\displaystyle 
\langle\{D_r,D_s^\dagger\}\rangle=\delta_{rs}S,
\qquad S=15-\frac7{32}\bar b.
$}\]

\(\phi\) hat entsprechend Norm \(1-7\bar b/480\), nicht globale
kanonische CAR. Der normierte Zustandserwartungswert ersetzt keine
Operatorrelation.

\subsection{5. Dieselbe Antwort auf demselben
Grundzustand}\label{dieselbe-antwort-auf-demselben-grundzustand}

Für Im z\textgreater0 und \(H_n=H|_{N=n}\):

\[\adjustbox{max width=\linewidth}{$\displaystyle 
G_{rs}(z)=\langle\Omega|f_r(z+E_0-H_{65})^{-1}f_s^\dagger|\Omega\rangle
+\langle\Omega|f_s^\dagger(z-E_0+H_{63})^{-1}f_r|\Omega\rangle.
$}\]

Die innere Symmetrie macht G diagonal und alle Diagonalelemente gleich.
Schreibe mit positiven Entnahme- und Additionsmaßen auf \(\epsilon>0\)

\[\adjustbox{max width=\linewidth}{$\displaystyle 
G(z)=\int\frac{d\nu_+(\epsilon)}{z-\epsilon}
+\int\frac{d\nu_-(\epsilon)}{z+\epsilon}.
$}\]

Dann gelten exakt

\[\adjustbox{max width=\linewidth}{$\displaystyle 
Z_+=\bar b/32,\quad Z_-=1-\bar b/32,
\quad a:=\int\epsilon\,d\nu_+=\int\epsilon\,d\nu_-
=\frac{\Delta\bar b-E_0}{64}.
$}\]

Für das gesamte signierte Spektralmaß:

\[\adjustbox{max width=\linewidth}{$\displaystyle 
\boxed{m_0=1,\quad m_1=0,\quad m_2=g^2S,\quad
m_3=g^2(\Delta S+7a).}
$}\]

Die ersten beiden Momente und die Kanalgewichte stimmen mit der
eingegangenen unabhängigen Rechnung überein. \textbf{Das dritte Moment
ist der zusätzliche eigene Schritt dieser Revision.}

\subsubsection{5.1 Herleitung des dritten
Moments}\label{herleitung-des-dritten-moments}

Setze \(\mathcal L A=[A,H]\). Stationarität macht diesen Operator
symmetrisch in der positiven, nach Nullvektoren quotientierten Metrik
\((A,B)=\langle\{A^\dagger,B\}\rangle\). Direkte Normalordnung ergibt

\[\adjustbox{max width=\linewidth}{$\displaystyle 
\sum_r\{[D_r,X],D_r^\dagger\}=-14Q_+.
$}\]

Die beiden Beiträge sind \(16Q_+\) und \(-30Q_+\). Ihre Koeffizienten
folgen aus den vollständig geprüften Kontraktionen
\(\sum_{rj}M_{A,rj}M_{B,rj}=16\delta_{AB}\) und
\(\sum_{Ar}M_{A,rj}M_{A,rk}=15\delta_{jk}\). Mit \([D,N_b]=D\) und
\(\langle Q_+\rangle=(E_0-\Delta\bar b)/(2g)\) folgt die Formel.

Als unabhängige Kontrolle wurden sämtliche Formeln einschließlich m3 an
der früher exakt geschlossenen N=4/N=5-Antwort symbolisch geprüft. Diese
Kontrolle verwendet den früheren Referenzzustand nur als Test der
Identitäten, nicht als Ersatz für den nativen Grundzustand.

Die allgemeine Methode, Spektralmomente aus Bewegungsgleichungen zu
gewinnen, ist etabliert; siehe
\href{https://arxiv.org/abs/0907.1284}{Freericks und Turkowski, Phys.
Rev.~B 80, 115119}. Die hier angegebenen W-Kontraktionen und Konstanten
wurden eigenständig für dieses Modell berechnet; die zitierte Arbeit
beweist keine TFPT-Aussage.

\subsubsection{5.2 Neue engere Dichte- und
Gewichtsgrenzen}\label{neue-engere-dichte--und-gewichtsgrenzen}

Positivität der beiden Antwort-Grammatrizen beziehungsweise zweimal
Cauchy-Schwarz liefern

\[\adjustbox{max width=\linewidth}{$\displaystyle 
m_2\ge a^2\left(\frac1{Z_-}+\frac1{Z_+}\right),
$}\] \[\adjustbox{max width=\linewidth}{$\displaystyle 
\boxed{(\Delta\bar b-E_0)^2\le
60g^2\bar b(32-\bar b)(1-7\bar b/480).}
$}\]

Mit der bereits bewiesenen Energieobergrenze und \(g/\Delta=1/20\) muss
das folgende rationale Polynom positiv sein:

\[\adjustbox{max width=\linewidth}{$\displaystyle 
P(b)=\frac7{3200}b^3-\frac{61}{50}b^2
+\frac{317591}{125000}b-\frac{79754843281}{62500000000}>0.
$}\]

Seine beiden im alten zulässigen Bereich liegenden Nullstellen liegen
bei ungefähr 0.842846697 und 1.245655664. Exakte rationale
Wurzelisolation ergibt die nach außen gerundete strenge Schranke

\[\adjustbox{max width=\linewidth}{$\displaystyle 
\boxed{0.842846<\bar b<1.245656.}
$}\]

Damit:

{\def\LTcaptype{none} % do not increment counter
\begin{longtable}[]{@{}
  >{\raggedright\arraybackslash}p{(\linewidth - 2\tabcolsep) * \real{0.4286}}
  >{\raggedleft\arraybackslash}p{(\linewidth - 2\tabcolsep) * \real{0.5714}}@{}}
\toprule\noalign{}
\begin{minipage}[b]{\linewidth}\raggedright
Größe pro Mode
\end{minipage} & \begin{minipage}[b]{\linewidth}\raggedleft
Strenges offenes Intervall
\end{minipage} \\
\midrule\noalign{}
\endhead
\bottomrule\noalign{}
\endlastfoot
Gesamtes Additionsgewicht & (0.0263389375, 0.03892675) \\
Gesamtes Entnahmegewicht & (0.96107325, 0.9736610625) \\
\(\langle\{\phi,\phi^\dagger\}\rangle\) & (0.981834183333\ldots,
0.987708495833\ldots) \\
Frühere \(\chi\)-Kompositnorm, nicht \(\phi\) & (0.040386370833\ldots,
0.059687683333\ldots) \\
\end{longtable}
}

Das sind Erwartungswerte und integrierte Spektralgewichte, keine
direkten Nachweise von Produktionsraten oder von bereits verfügbaren
Messinstrumenten.

\subsection{6. Der Polsatz und seine zusätzliche
Verschärfung}\label{der-polsatz-und-seine-zusuxe4tzliche-verschuxe4rfung}

Der eingegangene Beweis verwendet \(u_{k,r}=Q_+^kf_rF=f_rQ_+^kF\).
Invarianz und Besetzung ergeben

\[\adjustbox{max width=\linewidth}{$\displaystyle 
\langle u_{k,r},u_{k,s}\rangle=
\delta_{rs}\frac{64-2k}{64}\|Q_+^kF\|^2.
$}\]

Die Normen lauten 1, 465, 412200, 521164800, 833259571200. Die daraus
gebildete fünfdimensionale Variationsmatrix gibt 64 unabhängige
Richtungen unter \(-1.095812\Delta\). Das gesamte N=63-Komplement des
Nullbosonraums liegt über \(-3\Delta/4\). Minimax begrenzt den niedrigen
Raum auf genau 64 Dimensionen; seine injektive symmetrieverträgliche
Projektion auf die irreduzible duale 64 erzwingt ein einziges
Energieniveau.

Das beweist Existenz und Isolation. Seine Sichtbarkeit folgt aus den
positiven Entnahmemomenten. Mit

\[\adjustbox{max width=\linewidth}{$\displaystyle 
d=0.007737\Delta,\quad c=0.379636\Delta,
\quad a_{\max}=\frac{1.245656+1.158089}{64}\Delta
$}\]

gilt

\[\adjustbox{max width=\linewidth}{$\displaystyle 
Z_{\rm low}>\frac{c(1-1.245656/32)-a_{\max}}{c-d}
=\frac{40912436089}{46487375000}
=0.8800762806891118\ldots.
$}\]

Außerdem ist der niedrige Pol die kleinste Entnahmeenergie. Daher
\(a\ge\epsilon_{\rm low}Z_-\), also bereits ohne vollständige
Polauswertung

\[\adjustbox{max width=\linewidth}{$\displaystyle 
\frac{\epsilon_{\rm low}}\Delta
<\frac{1.245656+1.158089}{64-2(1.245656)}
=\frac{2403745}{61508688}=0.039079763821332\ldots.
$}\]

Diese Verschärfungen entstehen aus dem \textbf{Zusammenführen
kompatibler Beweise}, nicht aus einem neuen gewählten Parameter. Das
Restgewicht des gesamten CAR-Maßes ist somit kleiner als
0.119923719311\ldots{} .

\subsection{7. Einfache Organisation, aber keine falsche
Zwei-Linien-Lösung}\label{einfache-organisation-aber-keine-falsche-zwei-linien-luxf6sung}

Die ersten zwei orthonormalen Operatoren sind f und \(D/\sqrt S\). Die
erste Resolventenreduktion hat deshalb die exakte Form

\[\adjustbox{max width=\linewidth}{$\displaystyle 
\boxed{G(z)=\frac1{z-\displaystyle\frac{g^2S}{z-a_1-\Sigma_2(z)}}},
\qquad a_1=\Delta+\frac{7a}{S}.
$}\]

\(\Sigma_2\) ist die Resolvente des verbleibenden nativen Operatorraums,
gekoppelt an den dazu orthogonalen Rest von \(\mathcal L(D/\sqrt S)\).
Es wurde kein äußeres Bad hinzugefügt. Dies ist eine genaue Ordnung der
Rechnung, \textbf{keine abgeschlossene Berechnung von \(\Sigma_2\)} und
keine Vereinigung sämtlicher TOE-Aufgaben in einer einzigen bereits
gelösten Funktion.

\subsubsection{7.1 Warum man die Rückwirkung nicht exakt weglassen
darf}\label{warum-man-die-ruxfcckwirkung-nicht-exakt-weglassen-darf}

Angenommen, es gäbe nur je eine Entnahme- und Additionslinie mit
Energien \(\epsilon_-,\epsilon_+\). Symmetrie macht diese für alle r
gleich. Dann liefern die Bewegungsgleichungen auf demselben Grundzustand

\[\adjustbox{max width=\linewidth}{$\displaystyle 
D_r\Omega=-\frac{\epsilon_-}{g}f_r\Omega,
\qquad D_r^\dagger\Omega=\frac{\epsilon_+}{g}f_r^\dagger\Omega.
$}\]

Summe nach Multiplikation mit \(f_r^\dagger\) beziehungsweise \(f_r\)
und N=64 ergeben

\[\adjustbox{max width=\linewidth}{$\displaystyle 
H\Omega=\left[-32\epsilon_-+
(\Delta-\epsilon_++\epsilon_-)N_b\right]\Omega.
$}\]

Ein Eigenzustand mit negativer Energie kann keine feste Bosonenzahl
haben: Dann wäre \(\langle Q_++Q_-\rangle=0\) und seine Energie
\(\Delta\langle N_b\rangle\ge0\). Daher sind \(\Omega\) und
\(N_b\Omega\) unabhängig. Die angenommene Zweilinienform erzwingt
\(\epsilon_+-\epsilon_-=\Delta\).

Andererseits geben \(m_0=1,m_1=0\) für ein Zweilinienmaß
\(m_3/m_2=\epsilon_+-\epsilon_-\). Die exakt bestimmte Formel lautet
jedoch

\[\adjustbox{max width=\linewidth}{$\displaystyle 
\frac{m_3}{m_2}=\Delta+7a/S>\Delta.
$}\]

Widerspruch. Damit ist \(\Sigma_2\not\equiv0\) bewiesen. Mindestens ein
Kanal besitzt mehr als eine Energielinie. Der dominante isolierte
Entnahmepol und diese unvermeidliche Reststruktur widersprechen einander
nicht. Eine Zweilinienform könnte höchstens eine zu zertifizierende
Näherung sein.

\subsection{8. Was der Operationssatz jetzt tatsächlich
hergibt}\label{was-der-operationssatz-jetzt-tatsuxe4chlich-hergibt}

{\def\LTcaptype{none} % do not increment counter
\begin{longtable}[]{@{}
  >{\raggedright\arraybackslash}p{(\linewidth - 4\tabcolsep) * \real{0.3333}}
  >{\raggedright\arraybackslash}p{(\linewidth - 4\tabcolsep) * \real{0.3333}}
  >{\raggedright\arraybackslash}p{(\linewidth - 4\tabcolsep) * \real{0.3333}}@{}}
\toprule\noalign{}
\begin{minipage}[b]{\linewidth}\raggedright
Vertrag
\end{minipage} & \begin{minipage}[b]{\linewidth}\raggedright
Mathematisch abgesichert
\end{minipage} & \begin{minipage}[b]{\linewidth}\raggedright
Nicht dadurch verfügbar
\end{minipage} \\
\midrule\noalign{}
\endhead
\bottomrule\noalign{}
\endlastfoot
Markierter endlicher Matrixcompiler & Matrizen, Ordnungen, konkrete
endliche Syntheseidentitäten & Vollständiges Fockinstrument,
Präparation, physischer Zeitgenerator \\
H allein & Modellzeitentwicklung & Unabhängiges Schalten von X und
\(N_b\) \\
X und \(N_b\), wenn als Kontrollen gewährt & N=3-Kontrollalgebra der
Dimension 14 & Beliebige Modenoperationen \\
Zusätzlich dokumentierter vorzeichenrichtiger Clock-Lift &
N=3-Kontrollalgebra der Dimension \textbf{84} & Vollständige
Zustandsunterscheidung oder Ladung-eins-Instrument \\
\end{longtable}
}

Die neue Clock-Rechnung wurde übernommen und reproduziert. Die fünf
Multiplizitätszeilen für die sechs Clockphasen sind:

{\def\LTcaptype{none} % do not increment counter
\begin{longtable}[]{@{}lrrrrrr@{}}
\toprule\noalign{}
Teilraum & 0 & 1 & 2 & 3 & 4 & 5 \\
\midrule\noalign{}
\endhead
\bottomrule\noalign{}
\endlastfoot
Dunkle F3 & 6384 & 6280 & 6280 & 6384 & 6280 & 6280 \\
Dunkle BF & 16 & 8 & 8 & 16 & 8 & 8 \\
Aktiv 7 & 560 & 440 & 440 & 560 & 440 & 440 \\
Aktiv 10 & 112 & 88 & 88 & 112 & 88 & 88 \\
Aktiv 12 & 80 & 40 & 40 & 80 & 40 & 40 \\
\end{longtable}
}

Die aktiven Zeilen tragen zusätzlich einen Zweiniveaufaktor. Alle 45504
N=3-Zustände sind erfasst; der Kommutant hat weiterhin Dimension
240742144. Die alte Zahl 14 beschreibt den engeren Zweikontrollvertrag,
nicht den jetzt zusätzlich geprüften Clockvertrag.

\subsubsection{8.1 Konkrete
Präparationsgrenze}\label{konkrete-pruxe4parationsgrenze}

Alle genannten Fockkontrollen erhalten N. Ein Wort aus ihnen und jeder
einzeln zahlenerhaltende ausgewählte Messzweig bleiben im
Ausgangssektor. Sie können aus dem leeren N=0-Zustand \textbf{nicht} den
N=64-Grundzustand erzeugen. Diese Aussage betrifft zahlenerhaltende
Krauszweige; eine bloß U(1)-kovariante offene Dynamik kann sehr wohl
geladene Krausoperatoren haben.

Auch aus dem voll besetzten F führt bloßes \(e^{-itH}\) nicht zur
Konvergenz auf \(\Omega\): Der Grundzustandsüberlapp bleibt konstant.
Normiertes \(e^{-\tau H}F\) konvergiert mathematisch mit der bewiesenen
Lücke, benötigt aber eine gesonderte Instrument- und
Ressourcenherleitung.

\subsubsection{8.2 Der kleinste konkrete
Ladungstest}\label{der-kleinste-konkrete-ladungstest}

Ein Kandidat ist eine explizite Referenzmode c mit
\(T=\lambda(c^\dagger f_r+f_r^\dagger c)\). Sie erhält die gemeinsame
Ladung, verändert aber die Ladung der ursprünglichen Bank. Der endliche
Austausch- und Detuningtest ist exakt; \textbf{dieses T wurde nicht als
neuer nativer Hamiltonterm eingeführt}.

Für dieselbe N=64-Präparation unter \(H+\mu N\) gilt
\(G_\mu(z)=G_0(z-\mu)\), insbesondere \(m_1=\mu\). Das misst nur gegen
eine festgelegte Referenz: Ein gemeinsamer Zusatz
\(\mu(N_{\rm Bank}+N_{\rm Referenz})\) bleibt unsichtbar. Die Spektral-
schranken dieses Berichts gelten als Grundzustandsschranken nur für μ=0.

\subsection{9. Relativistisches Wörterbuch: Ausschluss und zwei präzise
Alternativen}\label{relativistisches-wuxf6rterbuch-ausschluss-und-zwei-pruxe4zise-alternativen}

Die Indexfrage darf nicht nachträglich durch einen Namen beantwortet
werden. Die Konventionen für Zweikomponentenspinoren wurden mit
\href{https://arxiv.org/abs/0812.1594}{Dreiner, Haber und Martin}
abgeglichen; die folgenden speziellen W-Tests sind eigene exakte
Tensorrechnungen.

\subsubsection{9.1 Was nicht funktioniert}\label{was-nicht-funktioniert}

Für 64 gleichhändige Weylfelder mit innerem Index I und einen skalaren
Vermittler ist \(M_A\otimes\varepsilon_{\rm Lorentz}\) symmetrisch. Die
Grassmann-Antisymmetrisierung verschwindet daher identisch, für alle 60
W-Zeilen. Außerdem lässt die volle irreduzible innere Darstellung
\(16\otimes4\) auf denselben 64 Komponenten nach Schur keine zusätzliche
kommutierende nichttriviale Lorentzspinorwirkung zu.

\subsubsection{9.2 Bereits bekannter nichtverschwindender
Typ}\label{bereits-bekannter-nichtverschwindender-typ}

Ein symmetrischer Lorentzspinortensor koppelt an das antisymmetrische M.
Der Vermittler hat dann einen passenden dualen (1,0)-Typ samt
adjungiertem Typ. Diese Variante besteht den algebraischen
Nichtnulltest, aber noch keinen vollständigen Kinetik-, Positivitäts-,
Constraints- oder Herkunftstest.

\subsubsection{9.3 Neue minimale skalare Alternative - ausdrücklich
zusätzlich}\label{neue-minimale-skalare-alternative---ausdruxfccklich-zusuxe4tzlich}

Will man zugleich W, gleiche Weylhand, lokale Bilinearität und einen
skalaren Vermittler behalten, kann man einen \textbf{unabhängigen}
Hilfsindex a=1,2 mit alternierender Form einführen:

\[\adjustbox{max width=\linewidth}{$\displaystyle 
b_A^\dagger(M_A)_{IJ}\varepsilon_{ab}
\varepsilon_{\alpha\beta}\psi_{I a\alpha}\psi_{J b\beta}
+\mathrm{h.c.}
$}\]

M und \(\varepsilon_{ab}\) sind jeweils antisymmetrisch; ihr Produkt ist
als innerer Kopplungstensor symmetrisch. Zusammen mit der Lorentz-
Epsilonform ist der Gesamttensor antisymmetrisch und nicht null. Alle 60
Kanäle bestehen den Test. Eine nichtverschwindende alternierende Form
existiert nicht in Dimension eins; in Dimension zwei ist sie bis auf
Normierung eindeutig. \textbf{In dieser eng benannten Klasse von
Tensorfaktor-Reparaturen ist die binäre Ergänzung minimal.}

Das ist noch keine gefundene native Lösung: Sie verdoppelt die inneren
Weylkomponenten von 64 auf 128, zusätzlich zu deren Lorentzspinorindex.
Ein Double-Cover-Minuszeichen stellt nicht automatisch zwei unabhängige
Felder bereit. Auch Nambu-Umbenennung \((f,f^\dagger)\) genügt nicht:
Das gemischte Produkt hat Ladung null statt −2, sodass derselbe
\(b^\dagger ff\)-Ladungsvertrag nicht erhalten bleibt.

Der entscheidende Quellenauftrag ist daher eng: Gibt es diese zweite
gleichgeladene, CAR-unabhängige Komponente bereits im tatsächlichen
Compilerprozess? Und liefert ihre Projektion genau das bisherige H und
die geprüfte Antwort? Ohne beides bleibt die skalare Alternative ein
zusätzliches Modell, nicht eine Erklärung des ursprünglichen.

\subsection{10. Nächste Schritte mit eindeutiger
Erfolgskontrolle}\label{nuxe4chste-schritte-mit-eindeutiger-erfolgskontrolle}

\begin{enumerate}
\def\labelenumi{\arabic{enumi}.}
\tightlist
\item
  \textbf{Native Quelle des Austauschoperators bestimmen.} Ein
  tatsächliches Operationswort mit Anfangszustand, Detektor,
  Adjungiertem, Ladungsbilanz und Record angeben. Ein weiteres neutrales
  Wort oder bloßer Algebraabschluss schließt diese Aufgabe nicht.
\item
  \textbf{Den Rest der nativen Antwort kontrollieren.} Das nächste Ziel
  ist \(\Sigma_2\) beziehungsweise sein erster Norm- und
  Momentkoeffizient, gemeinsam in N=63,64,65. Die vorliegenden
  Summenregeln und Polschranken sind zwingende Akzeptanztests. Den Rest
  auf null zu setzen ist exakt ausgeschlossen; eine Näherung braucht
  eine Restfehlergrenze.
\item
  \textbf{Die beiden Feldtypen an der Quelle entscheiden.} Entweder der
  symmetrische Vermittler mit korrektem Adjungierten und Kinetik, oder
  der skalare Typ mit wirklich nachgewiesener zweiter Komponente. Erst
  Nichtnullkopplung, Symmetrie, Ladung, CAR und positive Kinetik
  gemeinsam zählen als Feldadapter.
\item
  \textbf{Erst danach zwei operational bestimmte Teile verbinden.} Für
  einen tatsächlich hergeleiteten ungeraden Austausch wäre das
  projizierte Ein-Loch-Transfermatrixelement proportional zum jetzt
  eingeschlossenen Residuum. Ein solches formales Matrixelement beweist
  weder Verfügbarkeit des Austauschs noch einen isolierten
  Zweibank-Gesamtraum oder höhere Störungsordnungen. Eine räumliche
  Skalierung wurde hier nicht vorgezogen.
\end{enumerate}

{\def\LTcaptype{none} % do not increment counter
\begin{longtable}[]{@{}
  >{\raggedright\arraybackslash}p{(\linewidth - 4\tabcolsep) * \real{0.3333}}
  >{\raggedright\arraybackslash}p{(\linewidth - 4\tabcolsep) * \real{0.3333}}
  >{\raggedright\arraybackslash}p{(\linewidth - 4\tabcolsep) * \real{0.3333}}@{}}
\toprule\noalign{}
\begin{minipage}[b]{\linewidth}\raggedright
Tor
\end{minipage} & \begin{minipage}[b]{\linewidth}\raggedright
Nutzen dieser Revision
\end{minipage} & \begin{minipage}[b]{\linewidth}\raggedright
Entscheidender verbleibender Nachweis
\end{minipage} \\
\midrule\noalign{}
\endhead
\bottomrule\noalign{}
\endlastfoot
T1 & Exakter Clock-Kontrollabschluss und engere Ressourcenfrage &
Ursprüngliches vollständiges Operations- und Rahmenwörterbuch \\
T2 & Geladene native Antwort, sichtbarer Pol, drittes Moment &
Quelleneinbettung und renormiertes Half-Charge-Feld mit
Energie/Adjungiertem \\
T3 & Präzise Anforderung an ungeraden Austausch & Gemeinsamer
operationaler räumlicher, schließlich 3+1D-Träger \\
T4 & Falscher Skalartyp ausgeschlossen, minimale Alternativen & Chirales
Maß, Anomalien, Spiegelkontrolle, vollständige Feldkinetik \\
T5 & Pole und Reststruktur lokal kontrolliert & Gemeinsamer
wechselwirkender Grenzwert, Clusterstruktur, Streuung \\
T6 & Strengere interne/Lorentz-Indexbilanz & Familien, Massen und
Kopplungen auf demselben physikalischen Träger \\
T7 & Keine neue Spin-2-Konstruktion & Dynamischer Spin 2, Helizitäten
und universelle Kopplung \\
T8 & Eindeutiger Modellgrundzustand; Präparationslücke konkret &
Primitive Zustandswahl, Ressourcen, Records und Instrumente \\
\end{longtable}
}

\subsection{11. Erratum und
Versionsdisziplin}\label{erratum-und-versionsdisziplin}

Im eigenen v1.6.2-Text fehlte in Abschnitt 5.2 zwischen Entnahme- und
Additionsresolvente ein Pluszeichen. Richtig ist auf der dortigen
N=4-Referenz \(G=Z_h/(z+\epsilon_h)+a^\dagger(z+E_--H_5)^{-1}a\). Die
früheren Prüfrechnungen verwendeten die additive CAR-Gewichtsregel; der
Darstellungsfehler wird hier ausdrücklich korrigiert. Die archivierte
v1.6.2 wird nicht stillschweigend umgeschrieben.

Die zugelieferte v1.6.3 bleibt ebenfalls unverändert. Ihre Pol- und
Clockbefunde sind als übernommene und frisch reproduzierte Ergebnisse
kenntlich. Die engeren Schranken, das dritte Moment, der Ausschluss der
exakten Zweilinienantwort und die minimale skalare Hilfsindexalternative
sind die eigenen zusätzlichen Ergebnisse dieser Konsolidierung. Kein
literaturweiter Neuheitsanspruch, keine experimentelle Bestätigung und
kein Gesamtabschluss werden behauptet.

\begin{center}\rule{0.5\linewidth}{0.5pt}\end{center}

\section{Nutzerquelle A: vollständige Runde
v1.6.3}\label{nutzerquelle-a-vollstuxe4ndige-runde-v1.6.3}

Unveränderte Quelle. Nicht alle Aussagen darin werden übernommen.
Entscheidend sind die Quellenprüfung und Korrekturen in v1.6.7.

A. Rahmen der Runde v1.6.3 Ausgangspunkt ist v1.6.2 §14 mit den vier
Prioritäten in deiner Reihenfolge: (1) Operationssatz und Kommutant, (2)
nativer Grundzustand und geladene Antwort darauf, (3) Feldwörterbuch am
Tensor, (4) räumliche Skalierung erst danach. (4) ist ausdrücklich nicht
Gegenstand. Modell unverändert: H = ΔN\_b + gΣ\_A(b\_A†P\_A +
P\_A†b\_A), P\_A = Σ\_\{i\textless j\}W\_\{A,ij\}f\_jf\_i, 64
Fermionmoden, 60 Bosonmoden, N = N\_f + 2N\_b. Kein neuer Term.

Fundament common.py (PASS, 6 Prüfgruppen): Tensor W repo-lokal, SHA
3f00a089\ldots; WWᵀ = 8·I₆₀; 60 Kanäle × 8 disjunkte Paare auf 16 Moden,
jede Mode in genau 15 Kanälen; Gewichtserhaltung q\_i+q\_j = q\_A auf
allen 480 Paaren; Casimir-Identität 8WᵀW + 𝒞\_S + 𝒞\_C − 120·I = 0 auf
allen 2016 Zweiteilchenzuständen; induzierte Bosongeneratoren X\_B =
WΛ²(X)Wᵀ/8 gaußganzzahlig mit exakter Kovarianz; gepinnter Clock (SHA
9bf99de7\ldots) mit Periode 6, Slot-Permutation {[}2,0,1,4,3{]},
Vorzeichen −1, WΛ²G\_F = G\_B·W exakt.

B. Operationssatz und Kommutant (operations\_commutant.py, PASS, 162
Guards, 0 Bareiss-Fallbacks) Frage: Welche „unsichtbare`` Freiheit
bleibt, wenn nur bestimmte Operationen verfügbar sind? Gemessen als
Dimension des Kommutanten der von der Stufe erzeugten Algebra auf dem
N-Sektor. Kleiner Kommutant = mehr unterscheidbar.

Spektralzertifikat N=3 (exakt): S = C₃C₃ᵀ annulliert durch ∏(S−r) für r
∈ \{0,7,10,12\} (Residuum 0 nnz), Multiplizitäten 64/2880/576/320 aus
rationalen Projektorspuren, tr S = 29760, tr S² = 244800.

Irreduzibel-Zerlegung (Höchstgewichtsvektoren, exakt; Gewichte in
verdoppelten Koordinaten):

Sektor Block Dim Zerlegung unter so(10)⊕\allowbreak{}su(4) Casimir (𝒞\_S, 𝒞\_C) N=2
helle Paare↔Bosonen 60 (10, 6) ×1 (36, 20) N=2 dunkle Paare 1956 (120,
10) ⊕\allowbreak{} (126, 6), je ×1 (84,36), (100,20) N=3 dunkle Tripel = ker C₃ 37888
(560, 20″) ⊕\allowbreak{} (1200, 20) ⊕\allowbreak{} (672, 4̄), je ×1 (117,63), (141,39), (165,15)
N=3 ker S (der χ-Raum) 64 (16, 4̄) ×1 (45, 15) N=3 λ=12 320 (16, 20) ×1
(45, 39) N=3 λ=10 576 (144, 4̄) ×1 (85, 15) N=3 λ=7 2880 (144, 20) ×1
(85, 39) Alle Multiplizitäten sind 1; Summenprobe über
Weyl-Dimensionsformel bestanden. Meine vorab formulierte Hypothese
stimmt in allen Blöcken.

Kommutanten-Leiter:

Stufe (verfügbare Operationen) N=2 N=3 A: nur \{X, N\_b\} (die zwei
bisher geprüften Kontrollen) 3 829 536 1 444 233 216 B: A + Clock
(Ordnung 6) 641 760 240 742 144 Cc: A + 15 su(4)-Generatoren 30 376 2
247 168 Cs: A + 45 so(10)-Generatoren 172 1 648 C: A + alle 60
Generatoren 3 7 D: B + C 3 7 volle Algebra B(H\_N) 1 1 Lesart: Die zwei
Kontrollen lassen \textasciitilde1,4·10⁹ unsichtbare Freiheitsgrade in
N=3; der Clock allein reduziert um einen Faktor 6 (seine
Eigenwertmultiplizitäten sind fast gleichverteilt, z.B.
6384/6280/6280/6384/6280/6280 auf dem Dunkelraum); erst die volle innere
Symmetrie als Operationssatz kollabiert alles auf 7 Skalare - je einen
pro (Energieblock, Irrep). so(10) ist dabei die weit wirksamere Hälfte
(1648 vs.~2,2 Mio). Das ist genau die in v1.6.2 §10.3 geforderte
Abgrenzung: Mehrfachwiederholung der zwei Kontrollen beseitigt nichts;
die Existenz der Symmetrieoperationen als Operationen ist damit nicht
bewiesen, aber ihr Effekt ist jetzt exakt beziffert.

Clock-Befund: Der Clock wirkt auf keinem irreduziblen Block als Skalar
(er verschiebt Gewichte, Zeuge z.B. {[}−2,−2,0,−2,0{]}); er liegt also
nicht in der zusammenhängenden Symmetriegruppe, sondern ist ein äußeres
Element (Slot-Permutation). D = C gilt trotzdem, weil alle
Multiplizitäten 1 sind und der Clock jeden Block in sich abbildet.

C. Nativer Grundzustand N=64 (native\_ground\_state.py, PASS, 128
Guards: 119 exakt, 9 numerisch) C.1 Lochbild (Operatoridentität,
geguardet). Mit h\_i = f\_i† ist \textbar F⟩ =
f\_0†\ldots f\_63†\textbar0⟩ das Lochvakuum und P\_A = −Q\_A†, Q\_A† =
Σ\_\{i\textless j\}W\_\{A,ij\}h\_i†h\_j†. Unter (−1)\^{}\{N\_b\}: H ≅
H\_hole = ΔN\_b + gΣ(b\_A†Q\_A† + Q\_Ab\_A). Die Teilchen-Loch-Unitäre
trägt das Vorzeichen s(μ) = (−1)\^{}\{i+j\} pro Zweilochzustand;
Negativkontrolle: ohne s(μ) ist die Komplement-Abbildung keine
Identität.

C.2 Exakte Krylov-Kette v\_n = (T†)ⁿ\textbar F⟩, T† = Σb\_A†Q\_A†:

n Einträge in v\_n ν\_n = ‖v\_n‖² ν\_n/ν\_\{n−1\} 0 1 1 - 1 480 480 480
2 108 240 439 680 916 = 4·229 3 15 252 960 575 078 400 299520/229 ≈
1307,95 4 (nicht speicherbar) 952 296 652 800 (Spurnetzwerk) ≈ 1656,0 ν₃
wurde zweifach unabhängig erhalten: durch die
15-Millionen-Einträge-Vektorrechnung und durch die
Onishi/Wick-Spurnetzwerkformel (Abschnitt E). Exakte Übereinstimmung.

C.3 Exakte Guards: T v₁ = 480 v₀; T v₂ = 916 v₁ (Eindeutigkeit des
Ein-Boson-Singuletts (60⊗Λ²64̄)\^{}G, zuvor gruppentheoretisch
vorhergesagt); Adjungiertheit ⟨v₂\textbar Tv₃⟩ = ν₃; v₁, v₂ werden von 3
so(10)- und 3 su(4)-Generatoren annulliert (Singuletts);
⟨v₁\textbar Σ\_A Q\_AQ\_A†\textbar v₁⟩ = 458·480 = 219 840 - bestätigt
die korrigierte Formel Σ\_A Q\_AQ\_A† = 480 − (15N\_h +
C\_std\^{}\{Loch\})/2 mit dem Loch-Casimir 14 der (10,6) (meine erste
Annahme „480 − 15k`` war falsch und ist durch die Rechnung ersetzt).

C.4 Schließungstest. T v₃ = α v₂ + w₂ mit α = 299520/229 und w₂ ⊥ v₂,
‖w₂‖² = 5001523200/229 ≈ 2,184·10⁷ (exakt, Pythagoras-Identität ‖w₂‖² =
‖Tv₃‖² − ν₃²/ν₂ als Guard). Defekt ‖w₂‖²/‖Tv₃‖² = 2,90·10⁻⁵. Folge: Der
Zwei-Boson-Singulettraum ist \textgreater{} 1-dimensional, die Kette ist
exakt nicht geschlossen und das Modell nicht allein durch die ν\_n
lösbar; numerisch ist die Korrektur vernachlässigbar (Kopplung
g·√(‖w₂‖²/ν₃) = 0,195g, Energieänderung in der 6. Stelle).

C.5 Exakte Ritz-Matrix (Δ=1, g herausfaktorisiert, Basis v₀\ldots v₃,
w₂): Diagonale 0, 1, 2, 3, (2 für w₂); Nebendiagonalen 4√30, 2√229,
48√29770/229; w₂ koppelt nur an v₃ mit 3√598377/11908.

C.6 Ritz-Obergrenzen E\_K(64) (rigoros, Rayleigh-Ritz):

g/Δ K=1 K=2 K=3 (+w₂) Störungstheorie 2. Ordnung −480g² 1/20 −0,704159
−0,985178 −1,094232 −1,2 1/40 −0,241620 −0,288824 −0,295357 −0,3 1/100
−0,045894 −0,047851 −0,047894 −0,048 1/400 −0,002991 −0,003000 −0,003000
−0,003 Bei g/Δ = 1/20 ist die Kette nicht konvergiert (Kopplungen ≈ 1,1Δ
\ldots{} 1,6Δ je Stufe); dafür laufen ν₅ (und mit ν₄ die Verlängerung
auf K=5). Bei ≤ 1/40 ist die Konvergenz praktisch erreicht.

C.7 Observablen bei g/Δ = 1/20 auf dem K=3-Ritz-Zustand Ω (Näherung):
⟨N\_b⟩ = 0,8421 (Hellmann-Feynman ∂E/∂Δ = 0,8421 ✓); Überlapp
\textbar⟨F\textbar Ω⟩\textbar² = 0,397; Bosonzahlverteilung p₀\ldots p₃
= 0,397 / 0,397 / 0,172 / 0,034; Shannon-Untergrenze der
Boson-Loch-Verschränkung 1,66 bit (rigorose Untergrenze, da die
reduzierte Bosondichte in N\_b blockdiagonal ist). Worker-Behauptungen
„Überlapp \textgreater{} 1/4`` und „\textgreater{} 0,811 bit``:
konsistent, aber nicht konvergiert.

C.8 Geladene Ein-Teilchen-Antwort auf Ω (Guard auf r = 0, 5, 63,
Transitivität): Entnahme f\_r\textbar Ω⟩ (nach N=63): Z\_h = 1 −
⟨N\_b⟩/32 = 0,97368; Addition f\_r†\textbar Ω⟩ (nach N=65): Z\_add =
⟨N\_b⟩/32 = 0,02632; Z\_h + Z\_add = 1 exakt. Mittlere Anregungsenergie
der Entnahme ε\_h = 0,0311Δ, der Addition ε\_add = 1,1497Δ. Kompositnorm
Z\_χ = 23⟨N\_b⟩/480 = 0,0404 (v1.6.2-N=4-Referenz 0,0301 - anderer
Zustand, nur zur Einordnung).

C.9 Sektorauswahl und Worker-Behauptung. Untergrenzen aus H ≥ (1−θ)ΔN\_b
− (15g²/2θΔ)N\_f (Casimir-Identität + quadratische Ergänzung), korrekt
als max\_θ min\_\{N\_b\}:

g/Δ Ritz E(64) schärfster Konkurrent N=64 bewiesen global? 1/400
−0,0029996 N=63: −0,0029531 ja 1/100 −0,047894 N=63: −0,04725 ja 1/40
−0,295357 N=63: −0,2953125 ja (Abstand 4,5·10⁻⁵) 1/20 −1,094232 N=66:
−1,2; N=62: −1,1625 nein Bewiesener Bereich mit diesen Schranken: g/Δ ≤
0,0251 (Bisektion). Die Worker-Behauptung „eindeutiger
N=64-Singulett-Grundzustand für 0 \textless{} \textbar g\textbar/Δ ≤
1/20`` ist damit nicht reproduziert und nicht widerlegt: Die verwendeten
Untergrenzen sind zu grob; bei 1/20 müsste E(64) \textless{} −1,2 (gegen
N ≥ 66) und \textless{} −1,1625 (gegen N=62) gezeigt werden. Ob die
konvergierte Kette dorthin kommt, entscheidet der laufende ν₅-Schritt.
Chemisches Potential: μ* = −E/64 = 0,0171Δ \textless{} Δ/50, die frühere
Aussage „μ = Δ/50 wählt den leeren Zustand`` bleibt konsistent.

D. Feldwörterbuch am Tensor (field\_dictionary.py, PASS, 29 Guards) D.1
Allgemeine Regel (bewiesen, auf allen 60 Zeilen explizit nachgerechnet):
Für Grassmann-Felder ψ\_\{I,a\} koppelt der antisymmetrische Tensor
W\_\{IJ\} genau an den symmetrischen Teil des Spinorkerns K\_\{ab\}.

Zuordnung Spinorkern K Symmetrie Kopplung an W gleichhändige Weyl, Boson
Skalar (0,0) ε\_\{αβ\} antisym. 0 (60/60 Zeilen) gleichhändige Weyl,
Boson (1,0) σ\^{}\{μν\}ε sym. ≠ 0 (Koeffizientenrang 32 je Komponente)
Dirac/Majorana, Skalar C antisym. 0 Dirac/Majorana, Pseudoskalar Cγ⁵
antisym. 0 Dirac/Majorana, Axialvektor Cγ\^{}μγ⁵ antisym. 0
Dirac/Majorana, Vektor Cγ\^{}μ sym. ≠ 0 Dirac/Majorana, Tensor
Cσ\^{}\{μν\} sym. ≠ 0 Der Vermittler kann also nur ein Vektor- oder
antisymmetrisches Tensorfeld sein, niemals
Skalar/Pseudoskalar/Axialvektor. Er trägt U(1)\_N-Ladung 2 (komplexes
Feld); Vorbehalt (kein Satz): ein masseloser geladener Vektor ist kein
konsistentes freies Eichfeld.

D.2 Chiralitätsgradierungen. Trägergraph: 64 Moden, 480 Kanten, Grad 15,
dreiecksfrei, nicht bipartit - ungerader Zyklus {[}16, 45, 0, 30, 37{]}
(Länge 5, alle Kanten geprüft; unabhängig per BFS bestätigt). Folge: Es
gibt keine globale Chiralitätsaufteilung, unter der alle 60 Bosonen
Vektoren wären. Für jede der 44 getesteten G-kovarianten Gradierungen
ist jeder Kanal rein (0 gemischte Kanäle), aber der Bosonmultiplett
zerfällt in Typen:

Gradierung Vektor-Kanäle Tensor-Kanäle Stabilisator-Dim. Kommutant der
64 Identität (keine Aufspaltung) 0 60 60 1 (Schur; direkt gegengeprüft)
Pati-Salam-Spinorsplit Γ\_pq (10 Varianten) 24 36 36 = so(6)⊕\allowbreak{}so(4)⊕\allowbreak{}su(4)
2 Farbsplit (3 Varianten) 40 20 52 = so(10)⊕\allowbreak{}(su(2)²⊕\allowbreak{}u(1)) 2 gemischt (30
Varianten) 32 28 28 4 Ein einheitlicher Lorentztyp für alle 60
Vermittler erzwingt die gleichhändige Zuordnung → nur (1,0)⊕\allowbreak{}(0,1), also
ein antisymmetrischer Tensor B\_\{μν\}. Das ist kein symmetrisches
Spin-2-Feld. Kinetische Terme: einer pro irreduziblem Block des
Stabilisators (1, 2 oder 4), keine Familienaufspaltung 4 → 1+3.

E. Spurnetzwerk-Normen (norms\_by\_traces.py, läuft: n=5) Methode:
ν\_n/(n!)² = Σ\_\{\textbar m\textbar=n\} ‖Q†\^{}m F‖²/m! über bosonische
Kohärenzzustände, Onishi-Determinante det(1 − M̄M)\^{}\{1/2\} =
exp(−½Σ\_k tr((M̄M)\^{}k)/k) und Wick-Paarung → Summe über Zyklentypen
und Permutationen von Spurnetzwerken; alle Vorfaktoren exakte Fraction.
Auswertung auf dem Φ-Tensor Φ\_\{ij,kl\} = Σ\_A(M\_A)\{ij\}(M\_A)\{kl\}
(Einträge in \{−1,0,1\}, geguardet) mit optimalem Pfad; Exaktheit über
a-priori-Schranke \textless{} 2⁵³ pro Schritt, sonst Fallback modulo
fünf Primzahlen \textless{} 2¹³ mit CRT (Eindeutigkeit über
\textbar Netzwerk\textbar{} ≤ 60ⁿ·64\^{}\{\#Zyklen\}). Validierung: ν₁ =
480, ν₂ = 439 680 (brute force im Besetzungsbild), ν₃ = 575 078 400
(identisch mit der Vektorrechnung), CRT-Route allein reproduziert ν₃.
n=4 in 22 s (112 kanonische Netzwerke), ν₄ = 952 296 652 800. n=5 läuft
seit 17 CPU-Minuten.

F. Reparaturen an den Worker-Skripten (relevant für die Bewertung der
Zahlen) Vor dem Cursor-Absturz hinterlassene Skripte hatten folgende
Fehler, alle behoben und durch Guards abgesichert: Zyklus-Slots im
Spurnetzwerk (ergab ν₂ = 466560 statt 439680); Bosonengewicht ∏m\_A!
konstant 1; T wirkte bei wiederholtem Boson doppelt (Faktor 4 statt 2 →
T nicht adjungiert zu T†, ⟨v₁\textbar Tv₂⟩ = 453120 statt 439680);
Merge-Join mit inkonsistenter Ordnung für Loch in Mode 63;
Teilchen-Loch-Guard ohne s(μ); Guard ⟨v₁\textbar ΣQQ†\textbar v₁⟩ als
Norm der Summe statt Summe der Normen; Sektorschranke min-max statt
max-min (unzulässig optimistisch) und als Stub; w₂-Norm um q² = 229²
skaliert und an falsches Kettenglied gekoppelt; sympy-Eigenwerte für
≥5×5 Radikalmatrizen (Hänger); JSON-Tupelschlüssel; Schur-Kommutant über
instabile inkrementelle SVD; ungültiger „ungerader`` 4-Zyklus-Zeuge;
komplexe Generatoren mit vermischten Real-/Imaginärteilen im
Singulett-Guard.

Nebenbefund: Das replay\_manifest.json der Vorrunde v1.6.2 war um 10:06
durch einen Lauf ohne Zugriff auf den gesperrten Documents-Tensorpfad
auf FAIL gesetzt und normal.json geleert. Aus den intakten
optimized.json (PASS-Bedingung war Byteidentität) rekonstruiert; Summen
27143/27139/4 stimmen mit dem Delivery-Manifest.

G. Einordnung gegenüber T1-T8 Tor Beitrag dieser Runde Weiter offen T1
Exakte Kommutanten-Leiter A→B→Cc→Cs→C; Clock als äußeres Element
identifiziert Welche Symmetrieoperationen der Compiler tatsächlich
bereitstellt T2 Lochbild bewiesen; Z\_h + Z\_add = 1; ε\_h, ε\_add auf
dem Ritz-Zustand Konvergierter Zustand, Renormierung, Grenzraum T3 -
räumlicher Träger (ausgeklammert) T4 Vollständige Kern-Symmetrieregel;
Nichtbipartitheit → kein einheitliches Vektor-Wörterbuch chirales Maß,
Anomalien T5 Exakte Krylov-Kette, ν₁\ldots ν₄ zweifach unabhängig
Kontinuumslimes T7 (1,0)⊕\allowbreak{}(0,1) ist B\_\{μν\}, kein Spin 2 dynamischer
Spin 2 T8 Sektorauswahl bewiesen für g/Δ ≤ 0,0251; Worker-Bereich 1/20
nicht reproduziert; μ*-Fenster Zustandsfunktional, native Präparation
Kein Tor geschlossen. Keine RH-, Faktorisierungs- oder P-vs-NP-Aussage.

H. Was noch aussteht ν₅ (läuft) → Verlängerung der Ritz-Kette auf K=5,
Neubewertung von E(64) bei 1/20 und der Sektorfrage. replay.py über alle
fünf Prüfer normal + -OO mit Byteidentität, replay\_manifest.json.
RESULTS.md / EINFACH\_ERKLAERT.md, Kopie nach universal\_room/new/ als
v1.6.3, Eintrag in universal\_room/README.md.

\begin{center}\rule{0.5\linewidth}{0.5pt}\end{center}

\section{Nutzerquelle B: Korrektur und
Zwischenbericht}\label{nutzerquelle-b-korrektur-und-zwischenbericht}

Unveränderte Quelle. Nicht alle Aussagen darin werden übernommen.
Entscheidend sind die Quellenprüfung und Korrekturen in v1.6.7.

Korrektur verifiziert: S0 jetzt \textbf{3 829 536 / 1 444 233 216}
(deckt sich mit dem v1.6.2-Originalwert und der Parallelsession), 307
Guards, PASS in beiden Modi. Hier der ausführliche Zwischenbericht.

\begin{center}\rule{0.5\linewidth}{0.5pt}\end{center}

\section{Zwischenbericht: TFPT/Universalraum - fundamentale Lösung,
Stand 15.09.2026,
\textasciitilde11:45}\label{zwischenbericht-tfptuniversalraum---fundamentale-luxf6sung-stand-15.09.2026-1145}

\subsection{0. Gesamtlage: zwei Spuren}\label{gesamtlage-zwei-spuren}

Es existieren jetzt \textbf{zwei Contract-Ordner}, die sich unabhängig
auf denselben Kernpunkt zubewegt haben:

{\def\LTcaptype{none} % do not increment counter
\begin{longtable}[]{@{}
  >{\raggedright\arraybackslash}p{(\linewidth - 4\tabcolsep) * \real{0.3333}}
  >{\raggedright\arraybackslash}p{(\linewidth - 4\tabcolsep) * \real{0.3333}}
  >{\raggedright\arraybackslash}p{(\linewidth - 4\tabcolsep) * \real{0.3333}}@{}}
\toprule\noalign{}
\begin{minipage}[b]{\linewidth}\raggedright
Ordner
\end{minipage} & \begin{minipage}[b]{\linewidth}\raggedright
Ursprung
\end{minipage} & \begin{minipage}[b]{\linewidth}\raggedright
Stand
\end{minipage} \\
\midrule\noalign{}
\endhead
\bottomrule\noalign{}
\endlastfoot
\codeword{universalraum-operations-groundstate-20260915} & \textbf{diese
Session} (meine drei Arbeitspakete) & WP-A + WP-C fertig, WP-B-Sonde
läuft \\
\codeword{universalraum-native-ground-response-20260915} &
\textbf{Parallelsession} (andere Cursor-Sitzung, v1.6.4) & Grundzustand
+ geladene Antwort + Polsatz bereits konsolidiert \\
\end{longtable}
}

Die Parallelsession hat dabei einen extern eingegangenen \textbf{Polsatz
v1.6.3} reproduziert (317 Prüfbedingungen, byteidentisch) und eigene
Verschärfungen beigetragen. Beide Spuren validieren sich gegenseitig -
mit einer unten dokumentierten Abweichung, die ich inzwischen korrigiert
habe.

\textbf{Modellvertrag (unverändert in beiden Spuren):}
\(H=\Delta N_b+g\sum_A(b_A^\dagger P_A+P_A^\dagger b_A)\),
\(P_A=\sum_{i<j}W_{A,ij}f_jf_i\), \(N=N_f+2N_b\), 64 Fermionmoden, 60
Bosonmoden, \(WW^\dagger=8I_{60}\), 480 Vertizes, Prüfpunkt
\(g/\Delta=1/20\), kein \(\mu N\).

\begin{center}\rule{0.5\linewidth}{0.5pt}\end{center}

\subsection{\texorpdfstring{1. Fundament meiner Session:
\codeword{native\_source.py} (neu,
npz-frei)}{1. Fundament meiner Session:  (neu, npz-frei)}}\label{fundament-meiner-session-native_source.py-neu-npz-frei}

Problem: der gepinnte Tensor (\codeword{spinor\_tensors.npz}, SHA-256
\codeword{3f00a089…}) liegt unter \codeword{\textasciitilde{}/Documents}
- \textbf{für die Agenten-Shell macOS-gesperrt} (TCC). Die bisherigen
Verträge sind von dieser Shell nicht reproduzierbar.

Lösung: vollständige in-repo-Rekonstruktion aller Objekte aus der
Clifford/Außenalgebra-Konstruktion (die Rekonstruktion war in früheren
Runden bereits byte-exakt gegen das npz geprüft worden). \textbf{17
Guards, PASS}, darunter die volle Zwei-Teilchen-Casimiridentität
\(8W^\dagger W+4C_{\text{Spin}(10)}+4C_{\text{SU}(4)}=120\,I\) auf allen
2016 Paarzuständen. Exportiert: \codeword{W}, \codeword{J},
\codeword{C3}, Gewichte \codeword{FW/BW}, \codeword{BAR/ETA}, die 7
diskreten Symmetriegeneratoren, 60 Lie-Erzeuger, alle CAR-Helfer. Alle
drei Checker bauen nur darauf.

\begin{center}\rule{0.5\linewidth}{0.5pt}\end{center}

\subsection{2. WP-A: Der verfügbare Operationssatz (fertig, 307 exakte
Guards)}\label{wp-a-der-verfuxfcgbare-operationssatz-fertig-307-exakte-guards}

\textbf{Frage (v1.6.2 §10.3/§14):} Wie groß ist die unsichtbare Freiheit
(Kommutant) unter dem tatsächlich verfügbaren Operationssatz?

\subsubsection{Die Kommutanten-Leiter (N=3-Sektor, 45
504-dim)}\label{die-kommutanten-leiter-n3-sektor-45-504-dim}

{\def\LTcaptype{none} % do not increment counter
\begin{longtable}[]{@{}lrr@{}}
\toprule\noalign{}
Operationssatz & erzeugte Algebra & Kommutante \\
\midrule\noalign{}
\endhead
\bottomrule\noalign{}
\endlastfoot
S0: \{X, N\_b\} & 14 & \textbf{1 444 233 216} \\
S1: + Spin(10)×SU(4)-Rahmen & 743 583 744 & \textbf{7} \\
S2: + Clock-Lift & unverändert & 7 \\
S3: + Modenbesetzungen \{n\_r, m\_A\} & M₄₅₅₀₄ (2 070 614 016) &
\textbf{1} \\
S4: + geladene f-Instrumente & M₄₇₅₈₀ (N=2∪N=3) & 1 \\
\end{longtable}
}

N=2 analog: 3 829 536 → 3 → 3 → 1.

\textbf{Korrektur durch mich (Judge):} der Subagent hatte S0 als
\codeword{Σ 2m²} statt \codeword{Σ m²} gerechnet (1 452 961 792 statt 1
444 233 216) - der Guard war eine Tautologie. Jetzt korrigiert und gegen
den v1.6.2-Originalwert (\codeword{minimal\_interfaces.py}) sowie die
Parallelsession abgesichert; alle drei Quellen stimmen: \textbf{1 444
233 216}.

\subsubsection{Die Isotypie-Zerlegung von N=3 (beide Spuren unabhängig,
identisch)}\label{die-isotypie-zerlegung-von-n3-beide-spuren-unabhuxe4ngig-identisch}

{\def\LTcaptype{none} % do not increment counter
\begin{longtable}[]{@{}lrrrll@{}}
\toprule\noalign{}
Block & Dim & Mult & S-Eigenwert & (C\_S, C\_U) & Rolle \\
\midrule\noalign{}
\endhead
\bottomrule\noalign{}
\endlastfoot
(144, 20) & 2880 & 2 & \textbf{7} & (85, 39) & hell, X-gekoppelt \\
(144, 4̄) & 576 & 2 & \textbf{10} & (85, 15) & hell, X-gekoppelt \\
(16′, 20) & 320 & 2 & \textbf{12} & (45, 39) & hell, X-gekoppelt \\
(16′, 4̄) & 64 & 1 & 0 (ker S) & (45, 15) & χ-Kompositblock \\
(672, 4̄) & 2688 & 1 & - & (165, 15) & dunkel \\
(1200, 20) & 24 000 & 1 & - & (141, 39) & dunkel \\
(560, 20′) & 11 200 & 1 & - & (117, 63) & dunkel \\
\end{longtable}
}

Schur-Funktor-Zerlegung exakt geschält: Sym³(16) = 144⊕\allowbreak{}672, S₂₁(16) =
16′⊕\allowbreak{}144⊕\allowbreak{}1200, Λ³(16) = 560. In Λ³(16,4) multiplizitätsfrei; die drei
hellen Typen kommen je zweimal vor (eine Kopie pro Seite, X-verknüpft).
\textbf{Bemerkenswerter Befund (Parallelsession, von meiner Zerlegung
bestätigt): alle drei dunklen Blöcke tragen denselben Gesamtcasimir 45}
- die eine erhaltene Summeninvariante kann sie prinzipiell nicht
trennen; getrennt gelesen (141/4, 117/4, 165/4 bzw. 39/4, 63/4, 15/4)
trennen sie sofort.

\subsubsection{\texorpdfstring{Das Verfügbarkeits-Theorem
(Parallelsession,
\codeword{symmetry\_availability.py})}{Das Verfügbarkeits-Theorem (Parallelsession, )}}\label{das-verfuxfcgbarkeits-theorem-parallelsession-symmetry_availability.py}

Die gewährte Dynamik (Wörter aus X, N\_b) liegt in der
Verflechteralgebra (Dim Σmᵢ² = \textbf{16}); die 60 Lie-Erzeuger liegen
\textbf{nicht} darin - keine Komposition gewährter Dynamik erzeugt sie.
Die Symmetrie ist ein eigenständiges Primitiv: \emph{gewährt oder nicht
gewährt}. Der gesamte verbleibende symmetrieverträgliche Spielraum ist
\textbf{exakt 2 Dimensionen breit} = die Fähigkeit, die drei dunklen
Blöcke getrennt anzusprechen (= beide Casimire getrennt statt nur ihre
Summe). Die Clock sitzt im Normalisator, nicht im Zentralisator (mit X,
N\_b allein: Algebra 84, Kommutant 240 742 144; über voller Symmetrie
fügt sie nichts hinzu).

\textbf{Antwort auf die §10.3-Frage:} die operative Mehrdeutigkeit ist
auf die binäre Alternative „kontinuierliche Symmetrie verfügbar
ja/nein'' kollabiert - eingeklammert zwischen Kommutant \textbf{7} und
\textbf{1 444 233 216}.

\begin{center}\rule{0.5\linewidth}{0.5pt}\end{center}

\subsection{3. WP-C: Relativistisches Feldwörterbuch (fertig, 2161
exakte Guards,
-OO-stabil)}\label{wp-c-relativistisches-feldwuxf6rterbuch-fertig-2161-exakte-guards--oo-stabil}

Am tatsächlichen Tensor, vor jeder Kontinuumsrechnung:

{\def\LTcaptype{none} % do not increment counter
\begin{longtable}[]{@{}
  >{\raggedright\arraybackslash}p{(\linewidth - 4\tabcolsep) * \real{0.3000}}
  >{\raggedleft\arraybackslash}p{(\linewidth - 4\tabcolsep) * \real{0.4000}}
  >{\raggedright\arraybackslash}p{(\linewidth - 4\tabcolsep) * \real{0.3000}}@{}}
\toprule\noalign{}
\begin{minipage}[b]{\linewidth}\raggedright
Kanal
\end{minipage} & \begin{minipage}[b]{\linewidth}\raggedleft
Norm² pro Zeile
\end{minipage} & \begin{minipage}[b]{\linewidth}\raggedright
Status
\end{minipage} \\
\midrule\noalign{}
\endhead
\bottomrule\noalign{}
\endlastfoot
Skalar, gleichhändig ε\_\{αβ\}ψψ & \textbf{0} (alle 60 Zeilen, exakt) &
\textbf{verboten} \\
Skalar, gepunktet (0,0) & 0 & verboten \\
(1,0) symmetrisch \{I, σˣ, σᶻ\} & je 64, Summe \textbf{192} & erlaubt -
der Vertex ist \textbf{rein (1,0)+h.c.} \\
(0,1) konjugiert & je 64 & erlaubt \\
\end{longtable}
}

Korrektur meines Spec-Fehlers durch den Agenten: σʸ = −iε ist
\emph{anti}symmetrisch; die symmetrische Basis ist \{I, σˣ, σᶻ\}.
Dimensionsidentität C(128,2) = 2080 + 3·2016 = 8128 exakt. Kompositfeld
χ = Vermittler × Weyl: zerfällt (3/2,0)⊕\allowbreak{}(1/2,0), Projektoren Rang 4/2
exakt idempotent, flacher Spin-1/2-Anteil 1/3 (als Konvention
deklariert). Bilineare: 4096 = 1 (Singulett) + 45 + 15 (Adjungierte) +
4035 (Rest, ehrlich unbestimmt). Schur-Folge: kinetischer Operator auf
dem irreduziblen 64er-Multiplett = I₆₄ ⊗ d(p).

\textbf{Konsequenz:} jede relativistische Vervollständigung muss den 60
Vermittlern eine selbstduale Tensor-/feldstärkeartige Lorentzstruktur
((1,0), 3 Komponenten pro internem Label) geben - \textbf{der skalare
Kopplungskanal ist auf dem tatsächlichen Tensor exakt null und darf
nicht wieder verwendet werden}. Die Parallelsession ergänzte eine
minimale skalare Alternative mit extra alternierendem Hilfsindex (128
Komponenten) - ausdrücklich als \emph{zusätzliches} Modell markiert,
nicht als native Lösung.

\begin{center}\rule{0.5\linewidth}{0.5pt}\end{center}

\subsection{4. WP-B: Nativer Grundzustand (Parallelsession v1.6.4 +
meine laufende
Sonde)}\label{wp-b-nativer-grundzustand-parallelsession-v1.6.4-meine-laufende-sonde}

\textbf{Reproduziert (Parallelsession, frische Enumeration statt
Statusübernahme):}

{\def\LTcaptype{none} % do not increment counter
\begin{longtable}[]{@{}rrr@{}}
\toprule\noalign{}
k & Bosonkonfigurationen vollständig enumeriert & ‖Q₊ᵏF‖² \\
\midrule\noalign{}
\endhead
\bottomrule\noalign{}
\endlastfoot
0 & analytisch & 1 \\
1 & 60 Kanäle / 480 Paare & 480 \\
2 & 1 830 & 439 680 \\
3 & 37 820 & 575 078 400 \\
4 & 595 665 & 952 296 652 800 \\
\end{longtable}
}

\textbf{Modellinterner Satz:} für \(0<|g|/\Delta\le 1/20\) ist der
globale Grundzustand Ω \textbf{eindeutig, N=64,
Spin(10)×SU(4)-invariant}; am Prüfpunkt
\(-1{,}158089\,\Delta < E_0 < -1{,}129636\,\Delta\), Lücke
\(> 0{,}007737\,\Delta\). (Der volle Grundzustandsvektor ist damit nicht
ausgerechnet - ehrlich markiert.)

\textbf{Meine Session beigesteuert (Sonde, Teil 1):} die
\textbf{Z4-Zentrumsregel} - Singuletts existieren nur in Sektoren
\(N\equiv 0 \bmod 4\) (SU(4)-Zentrumsphase i\^{}N, niveauunabhängig,
exakt an Basiszuständen geprüft). Das erklärt, warum die Sektorwahl nur
unter N = \ldots, 56, 60, 64, 68, \ldots{} stattfindet; alle anderen
Sektoren haben Casimir-Floors.

\textbf{Noch laufend} (\href{aa1ce128-e016-4c52-ae8f-04d625b04c23}{WP-B
Sonde}): exakte Singulett-Multiplizitäten mult₁ (=1, der R-Zustand) und
mult₂ (Gewicht-0-Unterraum + Casimir-Kern), die exakten
Lanczos-Koeffizienten β₂², β₃² (Erwartung aus der Normsequenz: β₁²=480,
β₂²=439680/480=\textbf{916}, β₃²=299520/229 ≈ 1307,95 - diese
Verhältnisform gilt genau dann, wenn die Singuletts pro Niveau
eindimensional bleiben, also ist mult₂ der Zwickel), Casimir-Floor
c\_min=56 auf Λ², Sektorfloors N=56\ldots72, Gruppen-Census. Damit prüfe
ich die Kernzahlen der Parallelsession unabhängig nach.

\begin{center}\rule{0.5\linewidth}{0.5pt}\end{center}

\subsection{5. WP-D: Geladene Antwort auf dem Grundzustand
(Parallelsession
v1.6.4)}\label{wp-d-geladene-antwort-auf-dem-grundzustand-parallelsession-v1.6.4}

Genau die geforderte „dieselbe geladene Antwort auf dem abgesicherten
Zustand'':

\begin{itemize}
\tightlist
\item
  \textbf{Bewegungsgleichung (Operatoridentität, endlicher
  Teilchenkern):} \([H,f_r]=-gD_r\),
  \(D_r=\sum_{A,j}(M_A)_{rj}b_Af_j^\dagger\), Ladung −1 (nicht das
  frühere χ† mit Ladung +3). Kompositnorm:
  \(\langle\{D_r,D_s^\dagger\}\rangle=\delta_{rs}S\),
  \(S=15-\tfrac{7}{32}\bar b\).
\item
  \textbf{Exakte Spektralmomente:} \(Z_+=\bar b/32\),
  \(Z_-=1-\bar b/32\), \(a=(\Delta\bar b-E_0)/64\);
  \(m_0=1,\ m_1=0,\ m_2=g^2S,\ m_3=g^2(\Delta S+7a)\) (m₃ ist der eigene
  neue Schritt; unabhängig an der exakt lösbaren N=4/N=5-Antwort
  kontrolliert).
\item
  \textbf{Verschärfte rationale Schranken} (Polynom \(P(b)>0\),
  Wurzelisolation): \(0{,}842846<\bar b<1{,}245656\).
\item
  \textbf{Polsatz (v1.6.3, reproduziert):} isolierte Entnahmelinie,
  64-fach entartet (duale 64), \(Z_{\text{low}}>0{,}880076280689\ldots\)
  - \textbf{mehr als 88 \% des Spektralgewichts pro Mode} -, Polfenster
  \(0{,}007737<\epsilon/\Delta<0{,}039079764\), übrige Entnahmeenergien
  \(>0{,}379636\,\Delta\), Additionen \(>0{,}329636\,\Delta\).
\item
  \textbf{Zwei-Linien-Ausschluss:} eine exakte Zwei-Linien-Antwort
  erzwingt \(\epsilon_+-\epsilon_-=\Delta\), aber
  \(m_3/m_2=\Delta+7a/S>\Delta\) - Widerspruch. \(\Sigma_2\not\equiv0\)
  ist bewiesen; die Reststruktur ist unvermeidlich.
\end{itemize}

\begin{center}\rule{0.5\linewidth}{0.5pt}\end{center}

\subsection{6. Kreuzvalidierung der beiden
Spuren}\label{kreuzvalidierung-der-beiden-spuren}

{\def\LTcaptype{none} % do not increment counter
\begin{longtable}[]{@{}
  >{\raggedright\arraybackslash}p{(\linewidth - 6\tabcolsep) * \real{0.2500}}
  >{\raggedright\arraybackslash}p{(\linewidth - 6\tabcolsep) * \real{0.2500}}
  >{\raggedright\arraybackslash}p{(\linewidth - 6\tabcolsep) * \real{0.2500}}
  >{\raggedright\arraybackslash}p{(\linewidth - 6\tabcolsep) * \real{0.2500}}@{}}
\toprule\noalign{}
\begin{minipage}[b]{\linewidth}\raggedright
Befund
\end{minipage} & \begin{minipage}[b]{\linewidth}\raggedright
Meine Session
\end{minipage} & \begin{minipage}[b]{\linewidth}\raggedright
Parallelsession
\end{minipage} & \begin{minipage}[b]{\linewidth}\raggedright
Status
\end{minipage} \\
\midrule\noalign{}
\endhead
\bottomrule\noalign{}
\endlastfoot
N=3-Zerlegung (7 Blöcke, Mults 2,2,2,1,1,1,1) & ✓ & ✓ &
\textbf{identisch} \\
Casimirwerte (je Block) & (45,15), (85,15), \ldots{} & C = (C\_S+C\_U)/4
= 15, 25, 45, 21, 31, 45, 45 & \textbf{identisch} (zwei Normierungen) \\
Kommutant S1 (volle Symmetrie) & 7 & 7 & \textbf{identisch} \\
Kommutant S0 (X, N\_b) & \st{1 452 961 792} → \textbf{1 444 233 216} & 1
444 233 216 & \textbf{Diskrepanz gefunden + von mir korrigiert} \\
Skalarkanal exakt null & ✓ (60 Zeilen, beide Chiralitäten) & ✓ &
\textbf{identisch} \\
Grundzustandsnormen & Sonde läuft (β₁²=480 ✓) & 1, 480, 439 680, 575 078
400, 952 296 652 800 & Abgleich läuft \\
\end{longtable}
}

\begin{center}\rule{0.5\linewidth}{0.5pt}\end{center}

\subsection{7. Offen / nächste Schritte (in der Reihenfolge der
Direktive)}\label{offen-nuxe4chste-schritte-in-der-reihenfolge-der-direktive}

\begin{enumerate}
\def\labelenumi{\arabic{enumi}.}
\tightlist
\item
  \textbf{Sonde abschließen} (läuft): mult₂-Zertifikat, β-Kette,
  Sektorfloors → unabhängige Reproduktion der v1.6.4-Kernzahlen.
\item
  \textbf{Volle 33-Niveau-Kette:} falls mult\_k=1 weiter trägt, ist die
  Singulett-Kette durch die Normsequenz ‖Q₊ᵏF‖² (k≤32) bestimmt - dafür
  ist die Symmetriespur-Resummation nötig (die „vier Symmetriespuren''
  des Workers), nicht rohe Enumeration.
\item
  \textbf{Σ₂ kontrollieren} (erster Norm-/Momentkoeffizient, N=63/64/65
  gemeinsam) - die v1.6.4-Summenregeln sind die Akzeptanztests.
\item
  \textbf{Feldtyp an der Quelle entscheiden:} (1,0)-Vermittler mit
  Kinetik \emph{oder} nachgewiesene zweite Komponente für den skalaren
  Typ.
\item
  \textbf{Präparationslücke:} zahlenerhaltende Kontrollen erreichen N=64
  aus dem leeren Zustand nicht (Parallelsession, §8.3);
  imaginärzeitliches \(e^{-\tau H}F\) konvergiert mathematisch, braucht
  aber Instrumentenherleitung.
\item
  \textbf{Räumliche Skalierung: bewusst NICHT begonnen} - gemäß deiner
  Direktive trägt sie erst nach Zustand + Feldwörterbuch.
\item
  \textbf{Abschluss meines Contracts:} \codeword{replay.py} +
  \codeword{RESULTS.md} + README, wenn die Sonde zurück ist (inkl. des
  Erratums zum v1.6.2-Pluszeichen, das die Parallelsession dokumentiert
  hat).
\end{enumerate}

Kein T1-T8-Tor geschlossen, keine Promotion nach
\codeword{verification/}, kein Ledger/Paper-Eingriff - alles bleibt im
\codeword{experiments/}-Firewall.

\begin{center}\rule{0.5\linewidth}{0.5pt}\end{center}

\section{Nutzerquelle C: Multiplizitaeten und
Z4}\label{nutzerquelle-c-multiplizitaeten-und-z4}

Unveraenderte Nutzerquelle, als Vorschlag und nicht als Beweis
uebernommen. Die Pruefung und eigene Fortsetzung stehen in Teil B.

Ja. Der neue Stand verändert für mich die Richtung ziemlich fundamental.
Ich glaube inzwischen, dass wir an einer entscheidenden Stelle noch die
falschen Dinge miteinander identifizieren. Die stärkste Vereinfachung
ist aus meiner Sicht nicht ein weiterer cleverer Transferoperator. Sie
ist: Interne Identität, Ort und Zeit sind drei verschiedene Ebenen. TFPT
beschreibt die interne Identität bereits erstaunlich gut. Der
Universalraum könnte genau der bislang als „Multiplizität``,
„Kommutant`` oder „ungelöste Freiheitsgrade`` auftauchende relationale
Teil sein.

Wenn das stimmt, haben wir nicht zehn offene Probleme. Wir haben im
Wesentlichen einen Kategorienfehler. 1. Der wichtigste Hinweis steckt
direkt im neuen v1.6.6 Stand Der Clock ist jetzt tatsächlich als
konkretes Element von Spin(10) identifiziert. Das ist stark.
Gleichzeitig ist ausdrücklich gezeigt, dass daraus weder Hamiltonzeit
noch Raumrichtungen folgen. Und noch wichtiger: Der native 64
dimensionale Lochraum besitzt genau eine Energie. Wenn wir darin nur
verschiedene „Orte`` durch neue Charts definieren, gilt zwangsläufig
\[\adjustbox{max width=\linewidth}{$\displaystyle 
K=E_h S.
$}\]Dann sehen zwei Beschreibungen unterschiedlich aus, aber dynamisch
passiert exakt nichts. Die vorgeschlagene richtige Größe ist deshalb
\(G_{AB}(t)\), wobei man die statische Überlappung bei \(t=0\) von der
echten zeitlichen Änderung trennt. Das sagt für mich etwas sehr
Fundamentales: Die 64 Richtungen des Polraums sind keine 64 Orte. Sie
sind interne Freiheitsgrade. Der Vierzustandsversuch bestätigt das fast
schon mit einem Leuchtschild. Der Transfer 57 nach 58 funktioniert erst
mit einem zusätzlichen Mischer, aber dieser verändert eine SU(4)
Cartanladung. Das ist genau das Verhalten, das man erwarten würde, wenn
57 und 58 verschiedene interne Zustände sind und nicht „derselbe Zustand
links und rechts``. Wir haben also teilweise versucht, räumliche
Bewegung zu erzeugen, indem wir die Farbe des Teilchens verändern. Die
Algebra hat höflicherweise Nein gesagt. 2. Der Zwei Banken Versuch
verrät uns bereits, wie echte Bewegung aussehen muss Das Interessante an
v1.6.5 ist weniger der zusätzliche Link selbst. Es ist die Form des
resultierenden Niedrigenergieraums:
\[\adjustbox{max width=\linewidth}{$\displaystyle 
\mathcal H_{\rm low}
=
\mathbb C^2_{\rm links/rechts}
\otimes
\mathbb C^{64}_{\rm intern},
$}\]und der effektive Hamiltonoperator hat exakt die Form
\[\adjustbox{max width=\linewidth}{$\displaystyle 
H_{\rm eff}
=
h_{\rm links/rechts}
\otimes I_{64}.
$}\]Das \(I_{64}\) steht ausdrücklich in der Rechnung. Die interne
Überlagerung bleibt unverändert, während nur der Zwei Komponenten Faktor
bewegt wird. Das ist meiner Meinung nach der wichtigste Fingerzeig der
gesamten bisherigen Forschung. Nicht weil wir zwei Banken brauchen.
Sondern weil er zeigt, was „Raum`` representationstheoretisch sein muss:
\[\adjustbox{max width=\linewidth}{$\displaystyle 
\boxed{
\text{Materie} = \text{interne Darstellung},
\qquad
\text{Ort} = \text{Multiplizität derselben Darstellung}.
}
$}\]Das ist eine viel fundamentalere Trennung. 3. Ich würde den
Universalraum deshalb neu definieren Nehmen wir die innere Symmetrie
\[\adjustbox{max width=\linewidth}{$\displaystyle 
G=\mathrm{Spin}(10)\times SU(4).
$}\]Dann zerfällt ein globaler Hilbertraum ganz allgemein in
\[\adjustbox{max width=\linewidth}{$\displaystyle 
\boxed{
\mathcal H
=
\bigoplus_\lambda
V_\lambda\otimes M_\lambda
}
$}\]mit \(V_\lambda\): welche Art von Anregung es ist, also Ladungen,
interne Quantenzahlen und Darstellung, \(M_\lambda\): wie oft dieselbe
Art vorkommt, also der Multiplizitätsraum. Für jeden Hamiltonoperator,
der die innere Symmetrie erhält, gilt auf diesen Blöcken schematisch
\[\adjustbox{max width=\linewidth}{$\displaystyle 
H
=
\bigoplus_\lambda
I_{V_\lambda}\otimes h_\lambda.
$}\]Und plötzlich fällt vieles an seinen Platz. \(V_\lambda\) ist das
„Was``. \(M_\lambda\) ist der natürliche Kandidat für das „Wo``. Das ist
keine reine Erfindung. Bereits im nativen N=3 Raum treten dieselben
irreduziblen Typen teilweise mit Multiplizität zwei auf. Genau auf
solchen Multiplizitätsräumen kann die symmetrieerhaltende Dynamik
nichttrivial wirken. Vielleicht haben wir also einen Teil des Raumes die
ganze Zeit als „großen unerklärten Kommutanten`` betrachtet. Nicht der
gesamte Kommutant ist automatisch Raum. Aber wenn Raum emergiert, muss
er sehr wahrscheinlich dort leben und nicht innerhalb der internen 64.
Das ist für mich derzeit die eleganteste Interpretation des
Universalraums. 4. Damit wird auch klar, woher die fehlenden
Translationen kommen müssten In v1.6 existieren bereits die inneren Weyl
Matrizen \(A_\epsilon\). Für echte Bewegung fehlt aber gerade
\[\adjustbox{max width=\linewidth}{$\displaystyle 
U=\sum_\epsilon T_\epsilon\otimes A_\epsilon,
$}\]weil die \(T_\epsilon\) noch als räumliche Verschiebungen von außen
eingesetzt werden. Ohne sie summieren sich die inneren Matrizen einfach
zur Identität. TFPT\_Universalraum\_Update\_2026-09-15\_v1.6.pdfPDF
Bislang lautete die Frage: Woher bekommen wir \(T_\epsilon\)? Meine
Antwort wäre jetzt: Aus \(M_\lambda\). Nicht als neu hinzugeschriebene
Shiftoperatoren, sondern als Wirkung des ursprünglichen globalen
Hamiltonoperators auf verschiedene Vorkommen derselben internen
Darstellung. Dann wäre
\[\adjustbox{max width=\linewidth}{$\displaystyle 
T_\epsilon
$}\]kein zusätzliches Stück Raum. Es wäre die Darstellung der nativen
Dynamik auf dem Multiplizitätsraum. Das wäre ein echter Rückwärtsschluss
vom Universalraum auf TFPT und zugleich von TFPT auf den Universalraum.
5. Und jetzt kommt vielleicht der noch größere Punkt: unser heutiges
\(N\) könnte gar nicht fundamental sein Hier glaube ich, steckt eine
zweite massive Vereinfachung. Die volle E8 Struktur besitzt die bekannte
\(\mathbb Z_4\) Graduierung
\[\adjustbox{max width=\linewidth}{$\displaystyle 
\mathfrak g
=
\mathfrak g_0\oplus
\mathfrak g_1\oplus
\mathfrak g_2\oplus
\mathfrak g_3
$}\]mit unter anderem \[\adjustbox{max width=\linewidth}{$\displaystyle 
\mathfrak g_1=(16,4),
\qquad
\mathfrak g_2=(10,6),
\qquad
\mathfrak g_3=(\overline{16},\overline4),
$}\]und \[\adjustbox{max width=\linewidth}{$\displaystyle 
[\mathfrak g_a,\mathfrak g_b]
\subset
\mathfrak g_{a+b\;{\rm mod}\;4}.
$}\]Diese Struktur ist im Hauptdokument explizit.
TFPT\_Universalraum\_Hauptdokument\_2026-09-15\_v1.6.pdfPDF Unser
heutiger nativer Hamiltonoperator realisiert im Wesentlichen den Kanal
\[\adjustbox{max width=\linewidth}{$\displaystyle 
1+1\leftrightarrow2,
$}\]also zwei Fermionen zu einem Boson und zurück. Genau deshalb besitzt
er die schöne ganzzahlige Erhaltung
\[\adjustbox{max width=\linewidth}{$\displaystyle 
N=N_f+2N_b.
$}\]Aber jetzt die interessante Beobachtung: Die vollständige
\(\mathbb Z_4\) Struktur kann gar nicht konsistent zu einer ganzzahligen
Graduierung angehoben werden, wenn alle vier Gradkanäle physisch
realisiert werden. Denn setze \(q_1=1\). Aus
\[\adjustbox{max width=\linewidth}{$\displaystyle 
1+1\to2
$}\]folgt \(q_2=2\). Aus
\[\adjustbox{max width=\linewidth}{$\displaystyle 
1+3\to0
$}\]mit \(q_0=0\) folgt \(q_3=-1\). Aber aus
\[\adjustbox{max width=\linewidth}{$\displaystyle 
3+3\to2
$}\]würde dann \[\adjustbox{max width=\linewidth}{$\displaystyle 
q_2=-2
$}\]folgen. Also gleichzeitig \(q_2=2\) und \(q_2=-2\). Das funktioniert
exakt modulo vier, aber nicht als gewöhnliche ganze Ladung. Das ist eine
neue Schlussfolgerung aus den vorhandenen Strukturen, noch kein
bewiesener physischer TFPT Satz. Aber sie ist für mich extrem
interessant: \[\adjustbox{max width=\linewidth}{$\displaystyle 
\boxed{
\text{Die exakte }U(1)\text{ Erhaltung von }N
\text{ könnte nur eine Symmetrie unserer heutigen Teiltheorie sein.}
}
$}\]Die fundamentalere Quelle könnte lediglich die \(\mathbb Z_4\)
Graduierung erhalten. 6. Das könnte gleich mehrere heutige „Probleme``
verschwinden lassen Die Zustandswahl ist momentan unterbestimmt, weil
\[\adjustbox{max width=\linewidth}{$\displaystyle 
H\quad\text{und}\quad H+\mu N
$}\]dieselben inneren Symmetrien besitzen, aber völlig verschiedene
Vakuua auswählen können. Das ist inzwischen exakt demonstriert.
TFPT\_Universalraum\_Ergebnisse\_2026-09-15\_v1.6.mdMD Wenn aber die
vollständige Dynamik \(N\) gar nicht exakt erhält, sondern nur den Grad
modulo vier, ist \[\adjustbox{max width=\linewidth}{$\displaystyle 
\mu N
$}\]kein unsichtbarer beliebiger Zusatz mehr. Dann könnte dieselbe
vollständige Quelle auch das Vakuum auswählen, statt dass wir einen
Sektor und ein \(\mu\) nachträglich festlegen. Und plötzlich bekommt
eine ältere Merkwürdigkeit ebenfalls eine andere Bedeutung: Der
zusammengesetzte 64er Sektor hatte Ladung \(+3\), während ein Loch
Ladung \(-1\) hat. Die Dokumente mahnen völlig zu Recht, dass man das
nicht einfach umbenennen darf.
\codeword{TFPT\_Universalraum\_Minimale\_Fortsetzung\_Einfach\_2026-09-15\_v1.6.1.mdMD}
Aber modulo vier gilt natürlich
\[\adjustbox{max width=\linewidth}{$\displaystyle 
+3\equiv-1.
$}\]Und in v1.6.2 wurde sogar bereits ein konkreter Hintergrund
gefunden, auf dem genau
\[\adjustbox{max width=\linewidth}{$\displaystyle 
4-1=3
$}\]den Lochanschluss realisiert.
\codeword{TFPT\_Universalraum\_Konsolidierte\_Fortsetzung\_Einfach\_2026-09-15\_v1.6.2.mdMD}
Das riecht erheblich weniger nach Zufall als zuvor. Meine stärkste neue
Hypothese wäre daher: Wir haben bislang eine \(\mathbb Z_4\) graduierte
fundamentale Dynamik in einem Teilraum untersucht, in dem sie
versehentlich zu einer U(1) Erhaltung aufgeblasen wird.

Wenn das stimmt, erklärt es gleichzeitig die Präparationsmauer, die
\(\mu N\) Mehrdeutigkeit, den Unterschied \(+3\) gegen \(-1\) und
möglicherweise einen Teil der fehlenden Rückkopplung. 7. Dadurch wird
auch das Feldproblem wesentlich sauberer Der neue No Go Satz zeigt, dass
man auf demselben 64 dimensionalen internen Träger keine geeignete
Händigkeit einfach durch einen Basiswechsel erzeugen kann. Ein
zusätzlicher unabhängiger Zweierfaktor oder andere Strukturen bleiben
dagegen möglich. Ich würde diesen Zweierfaktor nicht mehr im internen E8
Raum suchen. Ich würde erwarten, dass er aus \(M_\lambda\) kommt. Also
etwa als zwei lokale Niedrigenergiekanäle eines Weyl Knotens, zwei
Orientierungszweige oder ein echtes geometrisches Dublett. Genau diese
Möglichkeit wird im neuen Text bereits als ernstzunehmend bezeichnet:
Eine geometrisch hergeleitete Zweifachheit könnte den zusätzlichen Index
erklären, müsste aber als echte unabhängige Freiheitsgrade konstruiert
werden. Damit würde auch diese Trennung elegant:
\[\adjustbox{max width=\linewidth}{$\displaystyle 
\boxed{
\text{Spin(10), SU(4), E8}=\text{innere Identität}
}
$}\]\[\adjustbox{max width=\linewidth}{$\displaystyle 
\boxed{
\text{Weyl Spin, Ort, Impuls}=\text{Niedrigenergiedynamik des Multiplizitätsraums}
}
$}\]Wir müssten keinen Lorentzspin mehr gewaltsam in die vorhandenen 64
internen Labels hineinpressen. 8. Und der Clock muss gar nicht zur Zeit
werden Auch hier glaube ich, dass wir zu viel vereinigen wollten. Der
neue Beweis hat bereits erklärt, was der Clock ist: ein bestimmtes
Spin(10) Element. Sehr gut. Fertig. Die physische Zeit ist dagegen die
Einparameterdynamik \[\adjustbox{max width=\linewidth}{$\displaystyle 
\alpha_t(A)=e^{itH}Ae^{-itH}.
$}\]Eine tatsächliche Uhr ist ein Teilsystem, dessen Zustand mit dieser
Dynamik korreliert und ausgelesen werden kann. Ein stationärer
Grundzustand selbst liefert keinen Zeitpfeil. Genau das stellt v1.6.6
klar.
\codeword{TFPT\_Universalraum\_Clock\_Gemeinsame\_Quelle\_Konsolidierung\_2026-09-15\_v1.6.6.mdMD}
Die elegante Lösung lautet deshalb nicht
\[\adjustbox{max width=\linewidth}{$\displaystyle 
\text{Clock}=\text{Zeit}.
$}\]Sondern \[\adjustbox{max width=\linewidth}{$\displaystyle 
\boxed{
\text{Clock}=\text{innere diskrete Symmetrie},
\qquad
\text{Zeit}=\text{Ordnung der realen Dynamik}.
}
$}\]Manchmal ist die schönste Vereinigung das Aufhören, zwei
verschiedene Dinge zwanghaft gleichzusetzen. 9. Das Gesamtbild wird
dadurch überraschend kompakt Ich würde den fundamentalen Kandidaten
heute ungefähr so schreiben:
\[\adjustbox{max width=\linewidth}{$\displaystyle 
\boxed{
\text{eine einzige globale, markierte, }\mathbb Z_4
\text{ graduierte Prozessalgebra}
}
$}\]mit einer physisch ausgewählten selbstadjungierten Dynamik
\(H_{\rm full}\). Ihre Darstellung zerfällt als
\[\adjustbox{max width=\linewidth}{$\displaystyle 
\boxed{
\mathcal H
=
\bigoplus_\lambda
V_\lambda^{\rm intern}\otimes M_\lambda^{\rm relational}.
}
$}\]Dann wäre: Erscheinung Fundamentale Herkunft Teilchenart und
Ladungen \(V_\lambda\) E8, Spin(10), SU(4) interne Algebra Clock innerer
Automorphismus Ort Multiplizitätsstruktur \(M_\lambda\) Bewegung
\(h_\lambda\) auf \(M_\lambda\) drei Raumdimensionen dreidimensionaler
Niedrigenergiegrenzwert von \(h_\lambda\) Weyl oder Dirac Spin lokale
Niedrigenergieentartung von \(M_\lambda\) Vakuum eindeutiger niedrigster
Zustand von \(H_{\rm full}\) Zeit Dynamik von \(H_{\rm full}\) Geometrie
Erreichbarkeit und Antwortkernel auf \(M_\lambda\) Gravitation
dynamischer Spin zwei Modus derselben relationalen Struktur

Das wäre tatsächlich eine einzige Architektur und kein Patchwork. 10.
Was ich jetzt nicht mehr machen würde Ich würde nicht weiter versuchen,
den Bosonmischer 0 nach 1 irgendwie doch aus mehr Kommutatoren
herauszuquetschen. Der Cartan Satz sagt uns ziemlich deutlich, dass
dieser konkrete Weg die falsche Art von Bewegung beschreibt. Ich würde
auch nicht mehr versuchen, die 64 Moden, die fünf Cliffordachsen oder
irgendwelche E8 Wurzelnummern direkt zu Raumkoordinaten zu erklären. Und
ich würde den großen Kommutanten nicht primär als Problem behandeln, das
man „wegkontrollieren`` muss. Vielleicht ist gerade die Multiplizität
das Zeug, aus dem Raum besteht. 11. Vier wirklich fundamentale Abnahmen
statt hundert weiterer Mikrotests 1. Vollständige Gradquelle: Die vier
E8 Grade müssen in einer einzigen adjungierten Prozessdarstellung
realisiert werden, mit den bereits geprüften W Vorzeichen als exakter
Teilstruktur. Der heutige Hamiltonoperator muss als Kompression daraus
folgen. Entscheidend ist, ob die vollständige Dynamik tatsächlich nur
\(\mathbb Z_4\) und nicht das künstlich stärkere \(N\) erhält. 2.
Multiplizitätsraum als Geometrie: Für eine feste Materiedarstellung,
idealerweise die native niedrige 64, muss die vollständige Quelle einen
wachsenden Multiplizitätsraum \(M_{64}\) erzeugen. Die native Dynamik
darauf muss nichttrivial sein. Keine externen Banken, kein eingesetztes
Hopping. Die Zwei Banken Form \(h\otimes I_{64}\) ist dabei der exakte
Solltyp. 3. Ein gemeinsamer Kontinuumssatz: Aus demselben \(h_{64}\)
müssen drei räumliche Richtungen, ein linearer Kegel, ein echter
Zweikomponenten Spinor und die chirale Struktur entstehen. Die heute
fehlenden \(T_\epsilon\) müssen dabei aus \(M_{64}\) kommen. Gelingt
das, fallen T2, T3, große Teile von T4 und T5 nicht einzeln, sondern
gemeinsam. 4. Vakuum und Gravitation: \(H_{\rm full}\) muss ohne frei
wählbares \(\mu\) einen eindeutigen physikalischen Zustand auswählen. Im
neutralen relationalen Sektor muss anschließend ein gaploser Spin zwei
Modus mit zwei Helizitäten und universeller Kopplung erscheinen. Wenn
dieser letzte Satz nicht kommt, haben wir eine schöne Materietheorie,
aber keine TOE. Mein derzeit stärkster Verdacht Die Fundamentalstruktur
ist nicht \[\adjustbox{max width=\linewidth}{$\displaystyle 
\text{Zelle}+\text{Link}+\text{Clock}+\text{Raum}+\text{Feld}.
$}\]Sondern eher \[\adjustbox{max width=\linewidth}{$\displaystyle 
\boxed{
\text{eine einzige }\mathbb Z_4\text{ graduierte Algebra zulässiger Veränderungen}
}
$}\]und alles andere sind unterschiedliche Zerlegungen ihrer
Darstellung: \[\adjustbox{max width=\linewidth}{$\displaystyle 
\boxed{
\text{interne Darstellung}
\;\times\;
\text{relationale Multiplizität}.
}
$}\]TFPT scheint bereits erstaunlich viel von der ersten Hälfte zu
kennen. Was wir Universalraum genannt haben, könnte im Kern die zweite
Hälfte sein. Und dann ist die wirklich spannende Umkehrung: Der große
bisher „unsichtbare`` Kommutant ist vielleicht nicht das, was wir noch
eliminieren müssen. Er könnte genau der Raum sein, den wir die ganze
Zeit gesucht haben.

Das ist aus meiner Sicht eine fundamental andere Forschungsrichtung als
v1.6.6 einfach weiterzuschreiben. Und sie erklärt auf einmal, warum
Clock, Cartan, Transport, der fehlende Zweierfaktor, die \(\mu N\)
Freiheit und die fehlenden Weyl Translationen alle genau dort klemmen,
wo sie klemmen.

\begin{center}\rule{0.5\linewidth}{0.5pt}\end{center}

\section{Nutzerquelle D: fundamentale
Reduktion}\label{nutzerquelle-d-fundamentale-reduktion}

Unveraenderte Nutzerquelle, als Vorschlag und nicht als Beweis
uebernommen. Die Pruefung und eigene Fortsetzung stehen in Teil B.

Ich würde den Ansatz jetzt tatsächlich auf eine andere Ebene heben. Der
Engpass ist inzwischen ziemlich klar: Wir verstehen einen inneren
Baustein und seine lokale Antwort immer besser. Was fehlt, ist die
Regel, aus der viele unterscheidbare Vorkommen dieses Bausteins, ihre
gemeinsame Bewegung und ihr physischer Zustand entstehen. Eine noch
genauere Lösung der einzelnen Bank kann diese fehlende Regel nicht
ersetzen. Das lässt sich aus eurem neuen Stand sogar mathematisch
begründen. Ich habe die Entwicklung bis v1.6.6 zusammengeführt, zentrale
Folgerungen herausgearbeitet und den ursprünglichen Paartensor
zusätzlich direkt geprüft. Die ausführliche Herleitung steht in der
\href{outputs/TFPT_Fundamentale_Reduktion_2026-09-15.md}{Forschungsnotiz
zur fundamentalen Reduktion}. 1. Was wir als tragenden Kern behalten
sollten Drei Dinge passen mittlerweile wirklich zusammen: - Ein innerer
algebraischer Baustein: die markierte Matrix-/E8-Struktur und der
konkrete Paartensor. - Eine funktionierende lokale Dynamik:
ursprüngliche Paarumwandlung, Rekopplung und - laut den aktuellen
Beweisen - ein eindeutiger Grundzustand des festgelegten
Hamiltonoperators. - Eine klar erkennbare Anregung: die isolierte
Entnahmelinie mit mehr als 88 Prozent des gesamten
Fermionspektralgewichts pro Mode. Der neue Clock-Anschluss vereinfacht
diesen Kern: Der dokumentierte Clock ist ein bestimmtes Element der
vorhandenen inneren Spin(10)-Symmetrie. Dafür braucht man keinen
getrennten mathematischen Mechanismus mehr. Das ist echte Vereinfachung.
Die Frage nach Raum, Bewegung und Zustandsauswahl bleibt jedoch
bestehen. Schon die v1.6-Grundlage trennt innere Walk-Amplituden von
tatsächlichen Ortsverschiebungen. 2. Die wichtigste Konsequenz: Auch das
perfekte lokale Spektrum liefert noch keinen Raum Diese Aussage würde
ich jetzt ins Zentrum stellen. Unter den dokumentierten Voraussetzungen
- innere Symmetrie, invarianter Grundzustand und ein irreduzibles
64er-Multiplett ursprünglicher Fermionoperatoren - gilt für die gesamte
Entnahmeantwort: \[\adjustbox{max width=\linewidth}{$\displaystyle 
C_{rs}(t)
=
\langle f_r\Omega,\,
e^{-it(H-E_0)}f_s\Omega\rangle
=
\delta_{rs}\,c(t).
$}\]Die Symmetrie erzwingt also: Alle inneren Richtungen haben dieselbe
skalare Antwort. In \(c(t)\) stecken sämtliche Nebenlinien und
Rückwirkungen, nicht nur der dominante Pol. Die Diagonalität steht
bereits im v1.6.4-Anhang. Ihre entscheidende Konsequenz lautet: Selbst
wenn wir die vollständige lokale Antwort exakt kennen, entsteht dadurch
kein Ortsindex zwischen diesen 64 inneren Richtungen.

Für zwei überlappende lineare Ansichten ergibt sich entsprechend nur
\[\adjustbox{max width=\linewidth}{$\displaystyle 
C_{uv}(t)=\langle u,v\rangle\,c(t).
$}\]Die Überlappung liefert den gemeinsamen Anteil; die Zeitentwicklung
liefert dieselbe lokale Antwort. Das wird durch einen anderen Namen für
die Ansichten nicht zur Ausbreitung. Diese Aussage verbietet weder
andere Vielteilchenprozesse noch Bewegung in einer größeren
Konstruktion. Sie zeigt aber, warum weitere lokale
Spektralverfeinerungen die grundlegende Herkunftsfrage nicht beantworten
können. 3. Wo Bewegung stattdessen mathematisch Platz hat Bei erhaltener
innerer Symmetrie lässt sich ein entsprechender Zustandsbereich
schreiben als \[\adjustbox{max width=\linewidth}{$\displaystyle 
\mathcal H
=
R_{\mathrm{intern}}\otimes\mathcal M,
\qquad
H=I_{\mathrm{intern}}\otimes h.
$}\]Dabei bedeutet: Bestandteil Bedeutung \(R_{\mathrm{intern}}\) Welche
innere Art von Anregung vorliegt \(\mathcal M\) Welche unabhängigen
Vorkommen derselben Art existieren \(h\) Wie sich diese Vorkommen
dynamisch verbinden

Der Fachbegriff für \(\mathcal M\) ist Multiplizitätsraum. Für den
isolierten einzelnen 64er-Pol ist dieser Raum eindimensional: Es gibt
dort nur eine Energie. Im gesamten Fockraum existieren weitere
Multiplizitäten; sie sind deshalb aber noch keine räumlichen Orte. Der
fundamentale nächste Gedanke ist daher: Raum müsste aus unabhängig
unterscheidbaren Vorkommen derselben inneren Struktur hervorgehen. Eine
mögliche spätere Beschreibung wäre \(\mathcal M\simeq\ell^2(X)\). Aber
dann müssen die Quelle und ihre Zusammensetzung erklären, was \(X\) ist,
welche Zugriffe lokal sind und wodurch Übergänge entstehen. Ein
nachträglich gewähltes dreidimensionales Gitter beantwortet das nicht.
Der Clock wirkt durch seine neue Identifikation innerhalb der inneren
Struktur. Er erzeugt allein noch keine solche Vielheit. 4. Die
Paritätsblockade sagt etwas über die Zusammensetzung - nicht über jede
gemeinsame Quelle Hier habe ich den ursprünglichen Tensor frisch
geprüft. Seine Paarungen verbinden bereits alle 64 Fermionmoden zu einer
einzigen zusammenhängenden Struktur. Für eine Teilmengenparität
\[\adjustbox{max width=\linewidth}{$\displaystyle 
\Pi_s=(-1)^{\sum_i s_i n_i}
$}\]muss an jeder ursprünglichen Paarung gelten:
\[\adjustbox{max width=\linewidth}{$\displaystyle 
s_i+s_j=0\pmod 2.
$}\]Weil der Paargraph zusammenhängend ist, bleiben - wenn die Bosonen
unverändert bleiben - nur Identität und globale Fermionparität. Es gibt
innerhalb dieser einen Bank keine echte Teilmenge ursprünglicher Moden
mit separat erhaltener Fermionparität. Das erklärt einen wichtigen
Unterschied: - Die eine native Bank erlaubt bereits innere Rekopplung. -
Separat kopierte Banken mit ausschließlich lokal geraden
Fermionoperationen bekommen zusätzliche lokale Erhaltungsgrößen, die den
einzelnen Lochtransfer blockieren. Die Art, wie wir Bausteine
zusammensetzen, kann also genau die Blockade erzeugen, die wir
anschließend mit einem neuen Mischer reparieren wollen. Das ist ein
Grund, die globale Zusammensetzung selbst zu untersuchen. Es beweist
noch keinen räumlichen Transport und hebt die Cartan-Schranke des
konkreten Mischers aus v1.6.6 nicht auf. 5. Auch die Zustandsfrage lässt
sich grundlegender und einfacher stellen Aus \(W\) und den angegebenen
Symmetrien allein folgt noch keine eindeutige Energiewahl. Euer eigener
Vergleich \[\adjustbox{max width=\linewidth}{$\displaystyle 
H_\mu=H+\mu N
$}\]zeigt das bereits: Dieselben inneren Symmetrien sind mit
unterschiedlichen globalen Grundzuständen vereinbar. Bei gleichem
präpariertem Eingang unterscheiden ausschließlich neutrale Beobachtungen
diese Entwicklungen trotzdem nicht. Daraus ziehe ich zwei Konsequenzen:
Erstens: Wenn eine sektorübergreifende Zustandsauswahl physisch relevant
ist, braucht sie zusätzliche Information - eine Quellregel, eine
Randbedingung oder ein ausdrücklich gesetztes Prinzip. „Symmetrisch``
und „einfach`` entscheiden diese Gegenbeispiele noch nicht. Zweitens:
Bei festem Gesamt-\(N\) ist \(\mu N\) nur eine Konstante. Eine
fundamental elegante Theorie muss solche unbeobachtbaren Unterschiede
nicht künstlich bestimmen. Sie sollte eindeutig sein bis auf Gleichheit
aller zugelassenen beobachtbaren Prozesse. Und noch eine Entlastung:
Eine fundamentale Theorie braucht keinen Experimentator außerhalb des
Universums, der ihr Vakuum aus dem leeren Zustand herstellt. Sie braucht
eine begründete Zustandsregel. Interne Präparationen und Detektoren sind
anschließend als Prozesse dieses Systems zu erklären. 6. Was für mich
jetzt als fundamentale Lösung zählen würde Ich würde die Arbeit um eine
explizite Quell- und Kompositionsregel organisieren:
\[\adjustbox{max width=\linewidth}{$\displaystyle 
\boxed{
\text{Quellregel}
\;\Longrightarrow\;
\text{gemeinsame Observablen, Dynamik, Zustand und Instrumente}
}
$}\]Die rechte Seite darf nicht aus unabhängig passend gewählten Teilen
bestehen. Dieselbe Regel müsste insbesondere erklären: 1. Warum mehrere
unabhängige Vorkommen derselben Anregung existieren. 2. Warum und wie
sie sich beeinflussen können. 3. Welcher physische Zustand und welche
beobachtbaren Energieunterschiede gelten. Der entscheidende
Bewegungsnachweis wäre dann sehr anschaulich: Auf derselben globalen
Quelle führe ich in Bereich A eine zulässige Operation aus oder lasse
sie aus. Ändert sich dadurch später die unbedingte Messstatistik in
Bereich B?

Damit unterscheidet man tatsächliche Wirkung von gemeinsamer
Anfangskorrelation, bloßer Überlappung und nachträglicher Auswahl
erfolgreicher Ausgänge. Solche Interventionswirkungen lassen sich über
zeitabhängige Kommutatoren untersuchen. Unter geeigneten lokalen
Wechselwirkungen führen entsprechende Schranken zu einem effektiven
Lichtkegel; diese Voraussetzungen müssen hier aber erst entstehen.
Bravyi, Hastings und Verstraete Eine solche gemeinsame Konstruktion
würde mehrere Herkunftsfragen zugleich beantworten. Danach wären
Raumdimension, relativistischer Feldtyp, Chiralität und Gravitation
Eigenschaften desselben Modells. Ihre Nachweise bleiben eigenständig -
etwa verschwindet das konkrete Grassmann-Vorzeichenproblem nicht
dadurch, dass die Quelle „global`` heißt. Meine Einschätzung Ich würde
den abgesicherten lokalen Grundzustand und Pol jetzt als Referenz
verwenden. Weitere Dezimalstellen sind sinnvoll, sobald eine konkrete
globale Konstruktion sie benötigt. Die Hauptarbeit sollte sich auf diese
Frage konzentrieren: Welche möglichst einfache Zusammensetzungsregel
macht aus dem bereits verstandenen inneren Baustein eine wachsende,
lokal unterscheidbare und gemeinsam dynamische Welt? Eine vollständige
solche Regel habe ich aus den vorliegenden Quellen nicht herleiten
können. Eine vollständige Lösung zu behaupten wäre deshalb falsch. Aber
der neue Stand erlaubt mehr als eine weitere Liste kleiner Aufgaben: Er
erklärt strukturell, warum die lokale Perfektionierung, der Clock und
bloße Überlappung die fehlende räumliche Ausführung nicht ersetzen
können - und welche zusätzliche Information eine wirkliche Lösung
enthalten muss. Die
\href{outputs/TFPT_Fundamentale_Reduktion_2026-09-15.md}{Forschungsnotiz}
enthält die Beweise, Voraussetzungen und Ergebnisse der frischen
Tensorprüfung. Die großen Grundzustands- und Transportzertifikate habe
ich dabei nicht erneut vollständig ausgeführt. - Eine globale Quellregel
konstruieren - Die verbleibende Unterbestimmtheit entscheiden - Das
Gesamtbild auf wenige Grundsätze reduzieren

12:00

\begin{center}\rule{0.5\linewidth}{0.5pt}\end{center}

\section{Nutzerquelle E: gemeinsamer ausfuehrbarer
Prozess}\label{nutzerquelle-e-gemeinsamer-ausfuehrbarer-prozess}

Unveraenderte Nutzerquelle, als Vorschlag und nicht als Beweis
uebernommen. Die Pruefung und eigene Fortsetzung stehen in Teil B.

Der wichtigste Satz zuerst: Nach dem heutigen Stand würde ich
Universalraum und TFPT nicht mehr als zwei konkurrierende Theorien
sehen. Der Universalraum ist der Kandidat für das tiefere Ding. TFPT ist
die außergewöhnlich starre mathematische Grammatik, die wir sehen, wenn
wir dieses Ding aus einer bestimmten Perspektive beschreiben. Oder in
deiner Welt gesprochen: Der Universalraum wäre die Runtime. TFPT wäre
Instruction Set, Typsystem und Compilerstruktur. Unsere beobachtete
Realität wäre ein ganz bestimmter laufender Prozess darauf. Das
Entscheidende ist: Wir haben inzwischen einige ziemlich harte Hinweise,
dass mehrere Dinge, die ursprünglich separat aussahen, tatsächlich
dieselbe innere Struktur sind. Aber wir haben noch nicht bewiesen, dass
diese Struktur die Natur beschreibt. Das ist die wichtigste Trennlinie.
Was ist der Universalraum überhaupt? Nicht ein noch größerer Raum, in
dem unser Universum schwimmt. Die radikalere Idee lautet: Ganz unten
gibt es zunächst überhaupt keinen Raum im üblichen Sinn. Es gibt
Zustände, erlaubte Veränderungen, Regeln zum Zusammensetzen dieser
Veränderungen, Aufzeichnungen vergangener Vorgänge und Regeln dafür, was
ausgelesen werden kann. Erst ein stabiles Muster dieser Beziehungen
könnte später als Entfernung, Zeit, Materie und Kausalität erscheinen.
Genau so wird der Universalraum im Hauptdokument definiert. Raum als
Netzwerk erreichbarer Beziehungen, Zeit als lesbare Ordnung von
Veränderungen und Materie als stabile oder wandernde Anregung sind dabei
bislang Forschungsziel, nicht bewiesene Ableitung. Bildlich: Stell dir
vor, du kennst keine Landkarte. Du weißt nur: Von A kann ich mit
Operation X nach B kommen, von B nach C, von A nach D nicht. Manche Wege
beeinflussen einander, manche hinterlassen Spuren. Wenn dieses Netz groß
genug und stabil genug ist, kannst du irgendwann feststellen: „A und B
sind Nachbarn``, „C ist weiter weg``, „Signale benötigen mindestens so
viele Schritte``. Die Landkarte entsteht dann aus Erreichbarkeit. Sie
war nicht vorher da. TFPT ist nun das Auffällige an dieser Geschichte:
Wir haben nicht einfach irgendein beliebiges Beziehungsnetz. Wir finden
eine extrem strukturierte algebraische Grammatik mit Spinoren, E8,
Spin(10), SU(4), komplexen Phasen, Fermionen, Bosonen und sehr
bestimmten signierten Kopplungen. Und genau hier wird es interessant.
Der mechanische Kern von TFPT ist erstaunlich einfach. Im aktuell
untersuchten nativen Modell gibt es 64 fermionische Moden und 60
bosonische Vermittler. Die Regel lautet im Wesentlichen:
\[\adjustbox{max width=\linewidth}{$\displaystyle 
\text{zwei Fermionen}\;\leftrightarrow\;\text{ein Boson}
$}\]Welche Fermionpaare mit welchem Boson verbunden sind, bestimmt ein
einziger signierter Tensor \(W\). Er enthält 480 Einträge mit
Vorzeichen. Die Dynamik ist
\[\adjustbox{max width=\linewidth}{$\displaystyle 
H=\Delta N_b+g\sum_A(b_A^\dagger P_A+P_A^\dagger b_A).
$}\]Dabei ist \(P_A\) schlicht die jeweils erlaubte Kombination aus zwei
Fermionen. Die Gesamtgröße \(N=N_f+2N_b\) bleibt erhalten. Spin(10) mal
SU(4) ist eine exakte innere Symmetrie dieses Modells. Noch ist das
ausdrücklich keine Lorentzsymmetrie unserer Raumzeit. Man kann sich das
als ein riesiges Lego System vorstellen. Es gibt nur sehr wenige Arten
von Steinen und eine sehr einfache Steckregel. Die Komplexität entsteht
durch Belegung und Zusammensetzung, nicht dadurch, dass für jedes
Phänomen eine neue Kraft erfunden wird. Genau das zeigte bereits die
Minimalrunde: Andere Belegungen erzeugen mit derselben Wechselwirkung
völlig neue Übergänge.
\codeword{TFPT\_Universalraum\_Minimale\_Fortsetzung\_2026-09-15\_v1.6.1.mdMD}
Die mögliche Entstehung unserer Realität würde ich momentan so sehen: 1.
Ganz unten steht eine globale Quelle. Noch keine Raumzeit, sondern eine
graduierte Algebra von Zuständen und erlaubten Operationen samt
positiven Wahrscheinlichkeiten und Aufzeichnungen. 2. TFPT ist ihre
lokale Grammatik. Aus der Quelle erscheinen Spin(10), SU(4), E8,
Ladungen, Fermion und Bosonkanäle sowie die zulässigen signierten
Paarumwandlungen. 3. Ein Zustand wird ausgewählt. Erst dadurch gibt es
so etwas wie ein Vakuum. Das ist wichtig, weil die Symmetrie allein den
Grundzustand nachweislich nicht auswählt. 4. Störungen dieses Zustands
bilden stabile Anregungen. Entfernst du beispielsweise ein Fermion,
entsteht im heutigen Modell überwiegend eine einzige klar isolierte
niedrige Anregung. 5. Wenn solche Anregungen zwischen lokal
unterscheidbaren Teilen tatsächlich wandern können, entsteht eine
Erreichbarkeitsstruktur. Aus ihr könnte Entfernung entstehen. 6. Wenn es
eine maximale Geschwindigkeit dieser Antwort gibt, entsteht ein
Kausalkegel. Erst dann wäre eine ernsthafte Brücke zur relativistischen
Raumzeit da. 7. Die stabilen niederenergetischen Anregungen dieses
Netzes erscheinen als Teilchen und Felder. Lokale Wechsel der
Beschreibung könnten sich als Eichsymmetrien zeigen. 8. Wenn die
Erreichbarkeitsstruktur selbst vom Zustand und seiner Energie abhängt,
könnte Geometrie dynamisch werden. Das wäre der Einstieg in Gravitation.
Der Clou wäre also: Raum, Zeit, Materie und vielleicht Gravitation wären
nicht vier zusätzliche Bausteine. Sie wären vier verschiedene
makroskopische Eigenschaften desselben fundamentalen Prozesses. Was wir
davon tatsächlich bewiesen haben Hier muss man ziemlich streng sein. Es
gibt bislang keinen Beweis, dass der Universalraum unsere physische
Realität ist. Es gibt aber mehrere Resultate, die weit über „lustige
gleiche Zahlen`` hinausgehen. Das erste ist E8. Die markierte maximale
Ordnung bestimmter komplexer 2 mal 2 Matrizen bildet tatsächlich ein
positives gerades unimodulares Gitter vom Rang acht mit exakt den 240
Norm zwei Vektoren von E8. Entscheidender als die Zahl 240 ist: Es
existiert eine konkrete markierungstreue Isometrie zum tatsächlichen
Compiler, die zusätzliche Struktur erhält. Das wurde reproduziert. Das
bedeutet: E8 ist hier keine Numerologie mehr. Innerhalb der Mathematik
ist die Verbindung eine Identität. Ob die Natur gerade diese Mathematik
benutzt, ist eine andere Frage. Noch stärker finde ich inzwischen den
Clock Befund. Der bislang separat wirkende endliche Clock lässt sich
explizit als Element derselben Spin(10) Struktur schreiben. Seine
Wirkung stimmt gleichzeitig auf Fermionen, Bosonen und dem
Kopplungstensor \(W\). Clock und innere Symmetrie sind also nicht zwei
zufällig kompatible Getriebe. Der Clock ist eine konkrete Stellung
desselben Getriebes. Das halte ich für einen der stärksten strukturellen
Befunde überhaupt. Aber interessant ist gerade die negative Seite:
Dieser Clock ist damit fast sicher nicht einfach „die Zeit``. Er sieht
viel eher nach einer internen Frame oder Orientierungsoperation aus. Die
physische Zeit muss zusätzlich aus echter Dynamik, Zustandsänderung und
lesbaren Records entstehen. Dann kommt der Zustand. Für den festgelegten
nativen Hamiltonoperator ist mathematisch bewiesen, dass im untersuchten
Kopplungsbereich ein eindeutiger globaler Grundzustand im Sektor
\(N=64\) existiert. Entfernt man ein Fermion, gibt es eine isolierte 64
fach entartete niedrige Linie, die mehr als 88,0076 Prozent des gesamten
normierten Fermionspektralgewichts trägt. Bildlich: Wir bauen aus der
sehr einfachen Paarregel ein riesiges Instrument. Ohne dass wir einen
Teilchenterm hineinschreiben, besitzt es einen bevorzugten Ruhezustand.
Zupfst du an einer der 64 Fermionsaiten, hörst du nicht chaotisches
Rauschen. Du hörst überwiegend einen sauber isolierten Ton. Das ist
interessant, weil genau so in Vielteilchenphysik ein emergentes
Quasiteilchen aussieht. Aber es ist noch kein Nachweis, dass dieser Ton
ein Elektron oder irgendein Standardmodellteilchen ist. Der Bericht sagt
das selbst ausdrücklich: relativistische Teilchendeutung, Masse und
räumliche Ausbreitung fehlen noch. Und dann kommt der vielleicht
wichtigste aktuelle Befund zur Bewegung. Wir haben inzwischen einen
echten kleinen Mechanismus gefunden:
\[\adjustbox{max width=\linewidth}{$\displaystyle 
\text{Fermionpaar}
\rightarrow
\text{Boson 0}
\rightarrow
\text{Boson 1}
\rightarrow
\text{anderes Fermionpaar}.
$}\]Mit einer zusätzlichen Bosonmischung überträgt dieser vollständige
native Vierzustandskanal eine Fermionmarke mit mehr als 99,3 Prozent
garantierter Wahrscheinlichkeit. Das ist kein künstlich abgeschnittener
Vierzustandsraum, sondern ein exakt invarianter Teil des vollständigen
nativen Wechselwirkungssystems. Aber, und das ist der entscheidende
Haken: Die Bosonmischung selbst haben wir noch nicht aus der Quelle
gewonnen. Sie verändert sogar eine SU(4) Cartanladung, welche der
bislang bekannte Operationssatz erhält. Deshalb können wir sie nicht
durch noch längeres Kombinieren derselben vorhandenen Operationen
herzaubern. Das ist für mich trotzdem enorm wertvoll. Wir suchen nicht
mehr nach „irgendeiner geheimnisvollen Raumoperation``. Wir wissen sehr
konkret, welche Art Ressource fehlt. Und was kann nun wirklich kein
Zufall mehr sein? Ich würde vier Ebenen unterscheiden. Dass irgendwo 64,
240 oder 248 auftauchen, kann Zufall oder Konstruktion sein. Das ist
schwache Evidenz. Dass eine komplette Gitterstruktur samt innerem
Produkt, Ganzzahligkeit, Vorzeichen und markierter Compilerabbildung
exakt E8 ergibt, ist mathematisch kein Zufall mehr. Es ist eine
Strukturidentität. Dass derselbe festgelegte Tensor \(W\) gleichzeitig
die Symmetrie, Paarwechselwirkungen, die geladene Anregungsstruktur und
die Clock Kovarianz trägt, ist nochmals stärker. Besonders die explizite
Identifikation des Clocks als Spin(10) Element reduziert die Zahl
unabhängiger Bausteine tatsächlich. Dass daraus schließlich ohne
erneutes Fitten beobachtete Naturkonstanten, drei Raumdimensionen,
Teilchenmassen, Mischungswinkel oder kosmologische Größen richtig
vorhergesagt werden, wäre physikalische Evidenz. Und genau Stufe vier
haben wir noch nicht. Deshalb würde ich nicht sagen: „Das kann unmöglich
Zufall sein, also ist die TOE bewiesen.`` Ich würde sagen: Die internen
Übereinstimmungen sind inzwischen zu strukturell, um sie als reine
Zahlenspielerei abzutun. Aber ob diese Mathematik die Natur auswählt,
ist noch völlig offen. Das ist ein großer Unterschied. Der neueste
Erkenntnissprung: Überlappung ist noch kein Raum Vor ein paar Runden war
meine favorisierte Idee: Vielleicht sind zwei „Banken`` gar nicht
unabhängig, sondern nur zwei lokale Ansichten desselben globalen
Universalraums. Das halte ich weiterhin für wahrscheinlich. Der neue
Audit hat aber etwas sehr Wichtiges klargemacht. Wenn zwei lokale
Beschreibungen sich überlappen, können sie bereits bei \(t=0\) dasselbe
Signal enthalten. Das sieht in einer Matrix wie eine Verbindung aus,
obwohl überhaupt nichts gewandert ist. Für den bereits bekannten 64
dimensionalen niedrigen Pol ist das besonders brutal: Alle Richtungen
haben dieselbe Energie. Verwendest du nur zwei verschiedene Basen
desselben Raums, gilt schlicht
\[\adjustbox{max width=\linewidth}{$\displaystyle 
K=E_h S.
$}\]Das Offdiagonale ist dann nur die Überlappung \(S\), kein Hop. Der
richtige Test ist deshalb die zeitabhängige Kreuzantwort
\[\adjustbox{max width=\linewidth}{$\displaystyle 
G_{AB}(t)
=
\langle\Omega|
f_B^\dagger
e^{-it(H-E_0)}
f_A
|\Omega\rangle.
$}\]\(G_{AB}(0)\) sagt: „Wie viel war ohnehin schon
dasselbe?{\kern0pt}`` Erst die Veränderung danach sagt: „Was ist
tatsächlich von A nach B passiert?{\kern0pt}`` Das ist meiner Meinung
nach ein fundamentaler Fortschritt, weil es uns den richtigen Begriff
von emergentem Raum liefert. Meine derzeit stärkste These zum
Universalraum Ich glaube inzwischen nicht, dass wir acht völlig
unabhängige T1 bis T8 Probleme lösen müssen. Ich glaube, dass wir drei
tieferliegende Lücken haben und dass viele der acht Probleme
gleichzeitig verschwinden könnten, wenn diese drei sauber geschlossen
werden. Die erste Lücke ist die globale Quelle. Wir brauchen ein
einziges mathematisches Objekt, das Algebra, Zustand, erlaubte
Operationen und Records gleichzeitig festlegt. Nicht erst eine Algebra
und später einen bequem ausgesuchten Hamiltonoperator. Dass Symmetrie
allein nicht genügt, ist sogar bewiesen: Man kann dem vorhandenen \(H\)
einen Term \(\mu N\) hinzufügen, dieselben inneren Symmetrien behalten
und trotzdem einen anderen Grundzustand auswählen.
TFPT\_Universalraum\_Ergebnisse\_2026-09-15\_v1.6.mdMD Meine Vermutung
ist deshalb: Der fundamentale Gegenstand ist kein Hamiltonoperator
allein. Er ist eher ein Tupel aus Operatoralgebra, positivem Zustand,
Dynamik und Recordstruktur. Der Universalraum ist genau dieses
Gesamtobjekt. TFPT ist seine algebraische Projektion. Die zweite Lücke
ist der native Transport. Mein stärkster Kandidat ist nicht ein
primitiver Term \(f_y^\dagger f_x\), also kein von Gott gegebener
Hoppingbefehl. Ich vermute, dass echter Transport immer eine
zusammengesetzte Operation ist:
\[\adjustbox{max width=\linewidth}{$\displaystyle 
\text{lokale Anregung}
\rightarrow
\text{gemeinsame Ressource}
\rightarrow
\text{lokale Anregung}.
$}\]Genau diese Architektur sehen wir im Vierzustandszeugen bereits
funktionieren. Das fehlende Teil könnte eine global erhaltene Ressource
sein, die lokal eine Cartanladung verändert und sie anderswo
kompensiert. Dann wäre die scheinbare lokale Symmetrieverletzung auf
globaler Ebene gar keine Verletzung. Das würde sehr elegant erklären,
warum wir den Mixer innerhalb einer einzelnen Bank nicht finden: Wir
suchen lokal nach einer Operation, deren Erhaltungsgesetz erst global
geschlossen wird. Das ist aktuell meine interessanteste konkrete
Hypothese. Die dritte Lücke ist die Übersetzung von Antwort zu
Geometrie. Ich würde Raum nicht aus den Labels 0 bis 63 und auch nicht
unmittelbar aus A3 oder drei Pauli Matrizen ableiten. Das Hauptdokument
warnt inzwischen selbst ausdrücklich davor, innere Dreidimensionalität
mit physischem Raum gleichzusetzen. Stattdessen würde ich physische
Entfernung aus dem dynamischen Antwortkernel definieren. Zwei Regionen
sind nah, wenn eine lokale Änderung mit wenigen nativen Operationen und
kurzer dynamischer Antwort die andere beeinflusst. Sie sind weit
entfernt, wenn viele Schritte nötig sind. Wenn die Anzahl erreichbarer
Regionen mit Radius \(r\) wie \(r^3\) wächst, bekommen wir einen ersten
Hinweis auf drei Dimensionen. Wenn gleichzeitig die Front der Antwort
eine universelle maximale Geschwindigkeit besitzt, bekommen wir einen
Lichtkegel. Dann wäre Raum wirklich emergent und nicht hineingemalt.
Dadurch bekommt auch Zeit eine viel sauberere Bedeutung Die aktuelle
Clock Entdeckung bringt mich zu einer ziemlich klaren These: Der TFPT
Clock ist vermutlich nicht Zeit. Er ist eine interne diskrete Frame
Operation. Physische Zeit entsteht erst, wenn drei Dinge zusammenkommen:
ein Hamiltonoperator, ein nichtstationärer Zustand und Records, anhand
derer sich „vorher`` und „nachher`` unterscheiden lassen. Das passt
auffällig gut dazu, dass ein stationärer Grundzustand unter \(H\) nur
eine globale Phase bekommt. Daraus entsteht kein Zeitpfeil. Ebenso
erzeugen bloße Basiswechsel keine echte Eichkrümmung. Diese Grenzen sind
im neuesten Audit ausdrücklich festgehalten. Also vielleicht: Clock =
Orientierung. Hamiltondynamik = Veränderung. Record = beobachtbare Zeit.
Das wäre wesentlich sauberer als zu versuchen, aus einer einzelnen
zyklischen Symmetrie sofort Sekunden zu bauen. Materie würde dann
ebenfalls ziemlich natürlich entstehen Das N64 Ergebnis zeigt bereits
den Prototyp. Das Vakuum ist nicht „nichts``. Es ist ein stark
korrelierter Grundzustand \(\Omega\). Ein Teilchen ist dann nicht
notwendig ein primitives Objekt. Es kann ein stabiler Pol in der Antwort
des Vakuums sein. Ein Elektron wäre in dieser Lesart ungefähr so
fundamental wie eine Schallwelle in einem Kristall: real, stabil und
messbar, aber nicht notwendigerweise der Grundbaustein des
darunterliegenden Systems. Der Unterschied zum banalen Quasiteilchenbild
wäre: Wenn derselbe Quellprozess zugleich Lorentzstruktur, Ladungen,
Statistik und Wechselwirkungen erzwingt, könnte das Quasiteilchen bei
niedriger Energie exakt wie ein fundamentales relativistisches Teilchen
aussehen. Noch haben wir diesen relativistischen Adapter nicht. Der
einfache Versuch, die 64 Labels einfach als gleichhändige Weylfelder zu
deklarieren, scheitert sogar basisunabhängig. Zusätzliche
Dirackomponenten, ein unabhängiger Zweierfaktor oder Ableitungsterme
sind dagegen nicht ausgeschlossen. Das sehe ich positiv: Die Mathematik
zwingt uns gerade dazu, den Lorentzspinor aus der globalen Dynamik
entstehen zu lassen, statt ihn nachträglich auf die internen Labels zu
kleben. Meine These zu Eichfeldern und Chiralität Eichfelder würde ich
ebenfalls nicht als neue Substanz hineinbauen. Wenn verschiedene lokale
TFPT Ansichten \(J_x\) in dieselbe globale Algebra eingebettet sind,
muss man festlegen, wie interne Frames beim Übergang von \(x\) zu \(y\)
miteinander verglichen werden. Diese Vergleichsoperation ist im Grunde
eine Connection. Aber: reine Basiswechsel teleskopieren um eine
geschlossene Schleife wieder zur Identität. Das wurde ausdrücklich
klargestellt. Für echte Krümmung braucht man eine dynamisch ausgewählte
Transportstruktur, nicht nur Koordinatenwechsel. Meine Vermutung wäre
deshalb: Eichfelder sind die interne Phase des realen
Transportoperators. Raumgeometrie ist seine Erreichbarkeitsstruktur.
Dasselbe Objekt hätte dann einen äußeren Teil, „wohin kann ich
gelangen?{\kern0pt}``, und einen inneren Teil, „wie dreht sich mein
Zustand dabei?{\kern0pt}``. Das wäre sehr elegant. Chiralität würde ich
anschließend nicht durch eine manuelle „links`` Markierung lösen,
sondern als Index oder spektralen Fluss dieses globalen
Transportoperators. Damit könnte aus einer global nahezu symmetrischen
Quelle bei niedriger Energie ein asymmetrischer chiraler Sektor
entstehen. Und Gravitation? Hier wäre meine Arbeitsthese: Wenn Raum
selbst aus dem Antwortoperator entsteht und der Antwortoperator vom
Zustand abhängt, dann verändert Energie automatisch die effektive
Geometrie. Dann ist \[\adjustbox{max width=\linewidth}{$\displaystyle 
\text{Materie}
\rightarrow
\text{verändert Antwortkernel}
\rightarrow
\text{verändert effektive Abstände und Kausalstruktur}.
$}\]Das klingt schon verdächtig nach Gravitation. Aber das allein reicht
nicht. Eine echte Lösung muss im Niedrigenergielimes einen masselosen
Spin zwei Sektor mit genau zwei Helizitäten, positiver Energie und
universeller Kopplung liefern. Genau diese Anforderungen bleiben
explizit offen.
\codeword{TFPT\_Universalraum\_Clock\_Gemeinsame\_Quelle\_Konsolidierung\_2026-09-15\_v1.6.6.mdMD}
Ich würde Gravitation deshalb ganz bewusst als letzten Test und nicht
als nächsten Trick behandeln. Wenn wir native Geometrie, Kegel und
chirale Materie haben, kann man fragen, ob kleine kollektive
Schwankungen dieser Geometrie automatisch Spin zwei ergeben. Wenn ja,
wird es richtig spannend. Und RH, Primzahlen, Faktorisierung und P gegen
NP? Hier würde ich momentan bremsen. Die gemeinsame Universalraumidee
kann erklären, warum Arithmetik als Schatten derselben
Kompositionsstruktur auftauchen könnte. Primitive Operationen könnten in
einem kommutativen Schatten wie Primfaktoren wirken. Das ist ein
reizvoller Gedanke. Das Hauptdokument formuliert genau diese engere
Hypothese. TFPT\_Universalraum\_Hauptdokument\_2026-09-15\_v1.6.pdfPDF
Aber primitive Schleifen eines allgemeinen Netzes sind nicht automatisch
Primzahlen. Gemischte primitive Zyklen entstehen ebenfalls. Für RH gibt
es immerhin ein messerscharfes Endziel: Würden wir aus derselben Quelle
einen positiven selbstadjungierten Operator \(A\) konstruieren und exakt
\[\adjustbox{max width=\linewidth}{$\displaystyle 
\frac{\Xi(z)}{\Xi(0)}
=
\det(I-z^2A^{-2})
$}\]beweisen, wäre das tatsächlich hinreichend. Diesen Operator und
diese Identität haben wir nicht. Ebenso existieren aktuell weder ein
neuer effizienter Faktorisierer noch eine P gegen NP Lösung.
\codeword{TFPT\_Universalraum\_Zwei\_Banken\_Konsolidierung\_2026-09-15\_v1.6.5.mdMD}
Ich würde deshalb Physik und Arithmetik vorerst parallel aus derselben
Quelle ableiten wollen, statt eine Seite gewaltsam aus der anderen zu
erklären. Treffen beide später auf exakt denselben Operator, wäre das
ein sehr starker Befund. Wenn ich alles auf eine fundamentale Vermutung
verdichten müsste Dann wäre es diese: Die fundamentale Realität besteht
nicht aus Teilchen in einem Raum. Sie besteht aus einer kleinen Algebra
erlaubter, phasenbehafteter Veränderungen und einem globalen Zustand, in
dem diese Veränderungen stattfinden. TFPT ist vermutlich eine
erstaunlich starre algebraische Darstellung dieser Regeln. Der
Universalraum ist die vollständige Ausführung dieser Regeln,
einschließlich Zustand, echter Operationsverfügbarkeit, Gedächtnis und
Auslesung. Raum ist das großskalige Muster dynamischer Erreichbarkeit.
Zeit ist die relationale Ordnung tatsächlich stattfindender und
aufgezeichneter Veränderungen. Materie sind stabile Pole und topologisch
geschützte Anregungen. Eichfelder beschreiben, wie interne Frames bei
realem Transport gegeneinander gedreht werden. Gravitation entsteht,
wenn der Zustand selbst diese Transportgeometrie verändert. Und der
nächste fundamentale Durchbruch wäre deshalb meiner Ansicht nach nicht
„noch eine größere E8 Rechnung``. Er wäre eine einzige globale
Quellenkonstruktion, in der wir zwei lokale TFPT Ansichten \(A\) und
\(B\) definieren und aus genau derselben primitiven Regel gleichzeitig
\[\adjustbox{max width=\linewidth}{$\displaystyle 
G_{AB}(0)
\qquad\text{und}\qquad
G_{AB}(t)
$}\]berechnen können, wobei \(G_{AB}(t)\) eine echte neue Antwort
enthält und nicht nur die ursprüngliche Überlappung umbenennt. Genau
diesen Test benennt auch die aktuelle v1.6.6 Konsolidierung als nächsten
entscheidenden Schritt.
\codeword{TFPT\_Universalraum\_Clock\_Gemeinsame\_Quelle\_Update\_2026-09-15\_v1.6.6.mdMD}
Wenn dieser eine Mechanismus nativ gelingt, würde ich unsere
Einschätzung deutlich hochstufen. Denn dann hätten wir erstmals in
derselben Quelle: innere Symmetrie, Clock, ausgewählten Zustand, stabile
Materieanregung und echte Bewegung. Dann fehlt nicht mehr „irgendwie
Raumzeit``. Dann hätten wir zum ersten Mal etwas, aus dem man Raumzeit
tatsächlich ableiten kann.

\begin{center}\rule{0.5\linewidth}{0.5pt}\end{center}

\section{Nutzerquelle F: Operationssatz, Grundzustand und
Feldwoerterbuch
v1.6.3}\label{nutzerquelle-f-operationssatz-grundzustand-und-feldwoerterbuch-v1.6.3}

Unveraenderte Nutzerquelle, als Vorschlag und nicht als Beweis
uebernommen. Die Pruefung und eigene Fortsetzung stehen in Teil B.

\section{TFPT / Universalraum: Operationssatz, nativer Grundzustand und
Feldwörterbuch}\label{tfpt-universalraum-operationssatz-nativer-grundzustand-und-feldwuxf6rterbuch}

\textbf{Forschungsfortsetzung v1.6.3 · 15. September 2026}

Diese Runde bearbeitet die vier Prioritäten aus v1.6.2 §14 in der dort
festgelegten Reihenfolge: den tatsächlich verfügbaren Operationssatz
exakt abgrenzen, den berichteten nativen Grundzustand mit seinem
ursprünglichen Hamilton- und Sektorvertrag reproduzieren und darauf die
geladene Antwort berechnen, das relativistische Feldwörterbuch am
tatsächlichen Tensor festhalten. Eine räumliche Skalierungsrechnung ist
ausdrücklich \textbf{nicht} Gegenstand; sie trägt erst, wenn Zustand,
Operationen und Feldtyp stehen. Haupt- und Update-PDF v1.6 bleiben
unverändert. Kein vollständiges T1-T8-Tor wird geschlossen.

\subsection{1. Ergebnis in einem Absatz}\label{ergebnis-in-einem-absatz}

Der N=64-Sektor des nativen Modells ist exakt als Loch-Paarerzeugung auf
dem gefüllten Zustand realisiert; seine Krylov-Kette ist bis zu 15 252
960 Einträgen exakt berechnet, ihre Normen ν₁\ldots ν₄ sind auf zwei
völlig unabhängigen Wegen (Besetzungsbild und Onishi/Wick-Spurnetzwerke)
identisch. Die Kette schließt bei zwei Bosonen \textbf{exakt nicht}
(Defekt 2,9·10⁻⁵), das Modell ist also nicht allein durch seine Normen
lösbar, numerisch aber fast. Rigorose Ritz-Obergrenzen, Bosonbesetzung,
Überlapp, Verschränkungsuntergrenze und die vollständige
Ein-Teilchen-Antwort (Z\_h + Z\_add = 1) liegen vor. Die
Worker-Behauptung eines eindeutigen N=64-Grundzustands bis g/Δ = 1/20
ist mit Casimir-Schranken nur bis g/Δ ≤ 0,0251 beweisbar und bei 1/20
\textbf{nicht reproduziert} (nicht widerlegt). Die Kommutanten-Leiter
zeigt exakt, wie viel unsichtbare Freiheit jede Operationsstufe
beseitigt: zwei Kontrollen lassen 1,44·10⁹ Freiheitsgrade in N=3, die
volle innere Symmetrie sieben. Das Feldwörterbuch ist am Tensor
entschieden: der Vermittler kann nur Vektor oder antisymmetrischer
Tensor sein; ein einheitliches Vektor-Wörterbuch existiert nicht
(Trägergraph nicht bipartit, 5-Zyklus {[}16,45,0,30,37{]}); der einzige
einheitliche Typ ist (1,0)⊕\allowbreak{}(0,1), ein B\_\{μν\}, kein Spin 2.

\subsection{2. Unveränderte Quelle und
Konventionen}\label{unveruxe4nderte-quelle-und-konventionen}

\[\adjustbox{max width=\linewidth}{$\displaystyle 
H=\Delta N_b+g\sum_{A=1}^{60}(b_A^\dagger P_A+P_A^\dagger b_A),\qquad
P_A=\sum_{i<j}W_{A,ij}f_jf_i,\qquad N=N_f+2N_b,\qquad WW^{\mathsf T}=8I_{60}.
$}\]

Der Tensor wird repo-lokal gelesen (SHA-256
\codeword{3f00a0892157a691e4ef26920c3702772c7bf76a250fad98ff04a079bb64b763}),
der Clock-Konstruktor ist gepinnt (\codeword{9bf99de7…}).
\codeword{common.py} prüft: 60 Kanäle × 8 disjunkte Paare, jede Mode in
15 Kanälen; Gewichtserhaltung auf allen 480 Paaren; die
Casimir-Identität 8WᵀW + 𝒞\_S + 𝒞\_C − 120 I = 0 auf allen 2016
Zweiteilchenzuständen (𝒞 = −ΣL², d.~h. viermal der Standard-Casimir);
induzierte Bosongeneratoren X\_B = WΛ²(X)Wᵀ/8 gaußganzzahlig und exakt
kovariant; Clock-Lift Periode sechs, Slot-Permutation (2,0,1,4,3),
Vorzeichen −1, WΛ²G\_F = G\_B W.

\subsection{3. Operationssatz und
Kommutant}\label{operationssatz-und-kommutant}

Gemessen wird die Dimension des Kommutanten der von einer
Operationsstufe erzeugten Algebra auf dem Sektor: kleiner Kommutant
bedeutet weniger unsichtbare Freiheit. Das Spektralzertifikat für N=3
ist reproduziert (S = C₃C₃ᵀ wird von ∏(S−r), r ∈ \{0,7,10,12\}, exakt
annulliert; Multiplizitäten 64, 2880, 576, 320 aus rationalen
Projektorspuren).

\subsubsection{3.1 Exakte Irreduzibel-Zerlegung
(Höchstgewichtsvektoren)}\label{exakte-irreduzibel-zerlegung-huxf6chstgewichtsvektoren}

{\def\LTcaptype{none} % do not increment counter
\begin{longtable}[]{@{}
  >{\raggedright\arraybackslash}p{(\linewidth - 8\tabcolsep) * \real{0.1875}}
  >{\raggedright\arraybackslash}p{(\linewidth - 8\tabcolsep) * \real{0.1875}}
  >{\raggedleft\arraybackslash}p{(\linewidth - 8\tabcolsep) * \real{0.2500}}
  >{\raggedright\arraybackslash}p{(\linewidth - 8\tabcolsep) * \real{0.1875}}
  >{\raggedright\arraybackslash}p{(\linewidth - 8\tabcolsep) * \real{0.1875}}@{}}
\toprule\noalign{}
\begin{minipage}[b]{\linewidth}\raggedright
Sektor
\end{minipage} & \begin{minipage}[b]{\linewidth}\raggedright
Block
\end{minipage} & \begin{minipage}[b]{\linewidth}\raggedleft
Dim
\end{minipage} & \begin{minipage}[b]{\linewidth}\raggedright
Zerlegung unter so(10)⊕\allowbreak{}su(4)
\end{minipage} & \begin{minipage}[b]{\linewidth}\raggedright
(𝒞\_S, 𝒞\_C)
\end{minipage} \\
\midrule\noalign{}
\endhead
\bottomrule\noalign{}
\endlastfoot
N=2 & helle Paare ↔ Bosonen & 60 & (10, 6) & (36, 20) \\
N=2 & dunkle Paare & 1956 & (120, 10) ⊕\allowbreak{} (126, 6) & (84,36), (100,20) \\
N=3 & dunkle Tripel = ker C₃ & 37888 & (560, 20″) ⊕\allowbreak{} (1200, 20) ⊕\allowbreak{} (672,
4̄) & (117,63), (141,39), (165,15) \\
N=3 & ker S & 64 & (16, 4̄) & (45, 15) \\
N=3 & λ = 12 & 320 & (16, 20) & (45, 39) \\
N=3 & λ = 10 & 576 & (144, 4̄) & (85, 15) \\
N=3 & λ = 7 & 2880 & (144, 20) & (85, 39) \\
\end{longtable}
}

Alle Multiplizitäten sind eins; Σ m·dim stimmt in jedem Block mit der
Weyl-Dimensionsformel. Höchstgewichte in verdoppelten Koordinaten, z. B.
(560,20″) ↔ ({[}3,3,1,1,−1{]},{[}3,3,3{]}), (672,4̄) ↔
({[}3,3,3,3,3{]},{[}1,1,−1{]}).

\subsubsection{3.2 Kommutanten-Leiter}\label{kommutanten-leiter}

{\def\LTcaptype{none} % do not increment counter
\begin{longtable}[]{@{}
  >{\raggedright\arraybackslash}p{(\linewidth - 4\tabcolsep) * \real{0.2727}}
  >{\raggedleft\arraybackslash}p{(\linewidth - 4\tabcolsep) * \real{0.3636}}
  >{\raggedleft\arraybackslash}p{(\linewidth - 4\tabcolsep) * \real{0.3636}}@{}}
\toprule\noalign{}
\begin{minipage}[b]{\linewidth}\raggedright
Verfügbare Operationen
\end{minipage} & \begin{minipage}[b]{\linewidth}\raggedleft
N=2
\end{minipage} & \begin{minipage}[b]{\linewidth}\raggedleft
N=3
\end{minipage} \\
\midrule\noalign{}
\endhead
\bottomrule\noalign{}
\endlastfoot
A: \{X, N\_b\} (die zwei bisher geprüften Kontrollen) & 3 829 536 & 1
444 233 216 \\
B: A + Clock (Ordnung 6) & 641 760 & 240 742 144 \\
Cc: A + 15 su(4)-Generatoren & 30 376 & 2 247 168 \\
Cs: A + 45 so(10)-Generatoren & 172 & 1 648 \\
C: A + alle 60 Generatoren & 3 & 7 \\
D: B + C & 3 & 7 \\
B(H\_N) & 1 & 1 \\
\end{longtable}
}

Für A ist die Blockstruktur M₃₇₈₈₈ ⊕\allowbreak{} M₆₄ ⊕\allowbreak{} (I₂⊗M₂₈₈₀) ⊕\allowbreak{} (I₂⊗M₅₇₆) ⊕\allowbreak{}
(I₂⊗M₃₂₀); für B folgen die Zahlen aus exakten Spuren der Clock-Potenzen
je Block (z. B. Multiplizitäten 6384/6280/6280/6384/6280/6280 auf dem
Dunkelraum); für die Lie-Stufen aus Σ m\_i². Der Clock ist auf
\textbf{keinem} irreduziblen Block ein Skalar (er verschiebt Gewichte;
Zeugen im JSON), also ein äußeres Element der Symmetrie, nicht in ihrer
zusammenhängenden Gruppe. D = C gilt trotzdem, weil alle Multiplizitäten
eins sind und der Clock jeden Block in sich abbildet.

\textbf{Folgerung.} Mehrfachwiederholung der zwei Kontrollen beseitigt
nichts; der Clock allein reduziert um einen Faktor sechs; erst die
innere Symmetrie als Operationssatz kollabiert die Freiheit auf je einen
Skalar pro (Energieblock, Irrep). so(10) ist dabei die wirksamere
Hälfte. Ob der Compiler diese Symmetrieoperationen tatsächlich
bereitstellt, ist damit nicht bewiesen; ihr Effekt ist exakt beziffert.

\subsection{4. Der native
N=64-Grundzustand}\label{der-native-n64-grundzustand}

\subsubsection{4.1 Lochbild}\label{lochbild}

Mit h\_i = f\_i† ist \textbar F⟩ = f₀†\ldots f₆₃†\textbar0⟩ das
Lochvakuum und als Operatoridentität P\_A = −Q\_A† mit Q\_A† =
Σ\_\{i\textless j\}W\_\{A,ij\}h\_i†h\_j†. Unter (−1)\^{}\{N\_b\} ist H
äquivalent zu

\[\adjustbox{max width=\linewidth}{$\displaystyle 
H_{\rm hole}=\Delta N_b+g\sum_A(b_A^\dagger Q_A^\dagger+Q_Ab_A),\qquad N_h=2N_b.
$}\]

Die Teilchen-Loch-Unitäre trägt das Vorzeichen s(μ) = (−1)\^{}\{i+j\} je
Zweilochzustand; die naive Komplement-Abbildung ohne s(μ) ist keine
Identität (Negativkontrolle bestanden).

\subsubsection{4.2 Exakte Krylov-Kette}\label{exakte-krylov-kette}

v\_n = (T†)ⁿ\textbar F⟩ mit T† = Σ\_A b\_A†Q\_A†:

{\def\LTcaptype{none} % do not increment counter
\begin{longtable}[]{@{}rrrl@{}}
\toprule\noalign{}
n & Einträge in v\_n & ν\_n = ‖v\_n‖² & ν\_n/ν\_\{n−1\} \\
\midrule\noalign{}
\endhead
\bottomrule\noalign{}
\endlastfoot
0 & 1 & 1 & - \\
1 & 480 & 480 & 480 \\
2 & 108 240 & 439 680 & 916 = 4·229 \\
3 & 15 252 960 & 575 078 400 & 299520/229 \\
4 & (nur Spurnetzwerk) & 952 296 652 800 & ≈ 1656,0 \\
5 & (nur Spurnetzwerk) & siehe \codeword{norms\_by\_traces.json} & \\
\end{longtable}
}

ν₁\ldots ν₃ sind auf zwei unabhängigen Wegen identisch (Besetzungsbild
mit exakten Ganzzahlamplituden; Spurnetzwerke, §6). Exakte Guards: T v₁
= 480 v₀; T v₂ = 916 v₁ (das Ein-Boson-Singulett (60⊗Λ²64̄)\^{}G ist
eindimensional); ⟨v₂\textbar T v₃⟩ = ν₃ (Adjungiertheit); v₁, v₂ sind
Singuletts unter je drei so(10)- und su(4)-Generatoren; ⟨v₁\textbar Σ\_A
Q\_AQ\_A†\textbar v₁⟩ = 458·480, was die Formel Σ\_A Q\_AQ\_A† = 480 −
(15N\_h + C\_std\^{}\{Loch\})/2 mit dem Loch-Casimir 14 der (10,6)
bestätigt - der Lochanteil von v₁ ist kein Singulett, nur der
Gesamtzustand.

\subsubsection{4.3 Schließung}\label{schlieuxdfung}

T v₃ = α v₂ + w₂ mit α = 299520/229, w₂ ⊥ v₂ und exakt

\[\adjustbox{max width=\linewidth}{$\displaystyle 
\|w_2\|^2=\|Tv_3\|^2-\frac{\nu_3^2}{\nu_2}=\frac{5001523200}{229}\approx 2{,}18\cdot10^7,
\qquad \frac{\|w_2\|^2}{\|Tv_3\|^2}=2{,}90\cdot10^{-5}.
$}\]

Der Zwei-Boson-Singulettraum ist also mehr als eindimensional; die Kette
ist \textbf{exakt nicht geschlossen}, und das Modell ist nicht allein
durch die ν\_n lösbar. Numerisch ist die Korrektur vernachlässigbar: w₂
koppelt nur an v₃ mit g·√(‖w₂‖²/ν₃) = 0,195 g.

\subsubsection{4.4 Ritz-Obergrenzen}\label{ritz-obergrenzen}

In der orthonormierten Basis \{v₀\ldots v\_K, w₂\} ist H (Δ = 1)
tridiagonal plus w₂-Zweig: Diagonale n bzw. 2; Nebendiagonalen g·4√30,
g·2√229, g·48√29770/229, \ldots; w₂-Kopplung g·3√598377/11908.

{\def\LTcaptype{none} % do not increment counter
\begin{longtable}[]{@{}lllll@{}}
\toprule\noalign{}
g/Δ & K=1 & K=2 & K=3 (+w₂) & −480g² \\
\midrule\noalign{}
\endhead
\bottomrule\noalign{}
\endlastfoot
1/20 & −0,704159 & −0,985178 & −1,094232 & −1,2 \\
1/40 & −0,241620 & −0,288824 & −0,295357 & −0,3 \\
1/100 & −0,045894 & −0,047851 & −0,047894 & −0,048 \\
1/400 & −0,002991 & −0,003000 & −0,003000 & −0,003 \\
\end{longtable}
}

Bei g/Δ = 1/20 ist die Kette bei K=3 nicht konvergiert (die
Stufenkopplungen liegen bei 1,1Δ\ldots1,6Δ); die Verlängerung mit ν₄, ν₅
steht in \codeword{native\_ground\_state.json} (Schlüssel
\codeword{ritz}). Bei g/Δ ≤ 1/40 ist die Konvergenz praktisch erreicht.

\subsubsection{4.5 Observablen auf dem Ritz-Zustand (g/Δ = 1/20,
K=3)}\label{observablen-auf-dem-ritz-zustand-gux3b4-120-k3}

⟨N\_b⟩ = 0,8421 (Hellmann-Feynman ∂E/∂Δ = 0,8421);
\textbar⟨F\textbar Ω⟩\textbar² = 0,397; Bosonzahlverteilung 0,397 /
0,397 / 0,172 / 0,034; Shannon-Untergrenze der Boson-Loch-Verschränkung
1,66 bit (rigoros, weil die reduzierte Bosondichte in N\_b blockdiagonal
ist). Die Worker-Angaben „Überlapp \textgreater{} 1/4`` und
„Verschränkung \textgreater{} 0,811 bit`` sind damit konsistent, aber
auf einem nichtkonvergierten Zustand.

\subsubsection{4.6 Geladene
Ein-Teilchen-Antwort}\label{geladene-ein-teilchen-antwort}

Für jede Mode r (Guard r ∈ \{0, 5, 63\}; Transitivität der
Quellsymmetrie):

\[\adjustbox{max width=\linewidth}{$\displaystyle 
Z_h=\|f_r\Omega\|^2=1-\frac{\langle N_b\rangle}{32}=0{,}97368,\qquad
Z_{\rm add}=\|f_r^\dagger\Omega\|^2=\frac{\langle N_b\rangle}{32}=0{,}02632,\qquad
Z_h+Z_{\rm add}=1.
$}\]

Mittlere Anregungsenergien: Entnahme ε\_h = 0,0311 Δ, Addition ε\_add =
1,1497 Δ. Kompositnorm Z\_χ = 23⟨N\_b⟩/480 = 0,0404 (die N=4-Referenz
aus v1.6.2 ergab 0,0301; verschiedene Zustände).

\subsubsection{4.7 Sektorauswahl und die
Worker-Behauptung}\label{sektorauswahl-und-die-worker-behauptung}

Untergrenzen aus der Casimir-Identität und quadratischer Ergänzung, H ≥
(1−θ)ΔN\_b − (15g²/2θΔ)N\_f, ausgewertet als max\_θ min\_\{N\_b\} je
Sektor:

{\def\LTcaptype{none} % do not increment counter
\begin{longtable}[]{@{}llll@{}}
\toprule\noalign{}
g/Δ & Ritz E(64) & schärfster Konkurrent & N=64 bewiesen global? \\
\midrule\noalign{}
\endhead
\bottomrule\noalign{}
\endlastfoot
1/400 & −0,0029996 & N=63: −0,0029531 & ja \\
1/100 & −0,047894 & N=63: −0,04725 & ja \\
1/40 & −0,295357 & N=63: −0,2953125 & ja \\
1/20 & −1,094232 & N=66: −1,2; N=62: −1,1625 & \textbf{nein} \\
\end{longtable}
}

Bewiesener Bereich mit diesen Schranken: g/Δ ≤ 0,0251. Die Behauptung
„eindeutiger Spin(10)×SU(4)-Singulett-Grundzustand für 0 \textless{}
\textbar g\textbar/Δ ≤ 1/20`` ist damit \textbf{nicht reproduziert und
nicht widerlegt}; bei 1/20 müsste E(64) \textless{} −1,2 gegen N ≥ 66
und \textless{} −1,1625 gegen N = 62 gezeigt werden. μ* = −E/64 = 0,0171
Δ liegt unter Δ/50; die frühere Aussage „μ = Δ/50 wählt den leeren
Zustand`` bleibt konsistent.

\subsection{5. Feldwörterbuch am tatsächlichen
Tensor}\label{feldwuxf6rterbuch-am-tatsuxe4chlichen-tensor}

\subsubsection{5.1 Symmetrieregel}\label{symmetrieregel}

Für Grassmann-Felder ψ\_\{I,a\} ist Σ\_\{IJ\}Σ\_\{ab\}
W\_\{IJ\}K\_\{ab\}ψ\_\{I,a\}ψ\_\{J,b\} = ½Σ(T − Tᵀ) mit T\_\{(Ia),(Jb)\}
= W\_\{IJ\}K\_\{ab\}. Weil W antisymmetrisch ist, verschwindet die
Kopplung genau dann, wenn K antisymmetrisch ist; \textbf{nur der
symmetrische Teil des Spinorkerns koppelt}. Auf allen 60 Zeilen explizit
geprüft:

{\def\LTcaptype{none} % do not increment counter
\begin{longtable}[]{@{}
  >{\raggedright\arraybackslash}p{(\linewidth - 4\tabcolsep) * \real{0.3333}}
  >{\raggedright\arraybackslash}p{(\linewidth - 4\tabcolsep) * \real{0.3333}}
  >{\raggedright\arraybackslash}p{(\linewidth - 4\tabcolsep) * \real{0.3333}}@{}}
\toprule\noalign{}
\begin{minipage}[b]{\linewidth}\raggedright
Zuordnung
\end{minipage} & \begin{minipage}[b]{\linewidth}\raggedright
Kern K
\end{minipage} & \begin{minipage}[b]{\linewidth}\raggedright
Kopplung
\end{minipage} \\
\midrule\noalign{}
\endhead
\bottomrule\noalign{}
\endlastfoot
gleichhändige Weyl, Boson Skalar (0,0) & ε antisym. & 0 \\
gleichhändige Weyl, Boson (1,0) & σ\^{}\{μν\}ε sym. & ≠ 0 (Rang 32 je
Komponente) \\
Dirac/Majorana: Skalar C, Pseudoskalar Cγ⁵, Axialvektor Cγ\^{}μγ⁵ &
antisym. & 0 \\
Dirac/Majorana: Vektor Cγ\^{}μ, Tensor Cσ\^{}\{μν\} & sym. & ≠ 0 \\
\end{longtable}
}

Der Vermittler ist damit Vektor oder antisymmetrischer Tensor, niemals
Skalar, Pseudoskalar oder Axialvektor. Er trägt U(1)\_N-Ladung 2; ein
masseloser geladener Vektor ist kein konsistentes freies Eichfeld
(Vorbehalt, kein Satz).

\subsubsection{5.2
Chiralitätsgradierungen}\label{chiralituxe4tsgradierungen}

Der Trägergraph (64 Moden, 480 Kanten, Grad 15, dreiecksfrei) ist
\textbf{nicht bipartit}; ein ungerader Zyklus ist {[}16, 45, 0, 30,
37{]}. Es gibt daher keine globale Chiralitätsaufteilung, unter der alle
60 Bosonen Vektoren wären. Für jede der 44 G-kovarianten Gradierungen
ist jeder Kanal rein (nie gemischt), aber der Bosonmultiplett zerfällt:

{\def\LTcaptype{none} % do not increment counter
\begin{longtable}[]{@{}
  >{\raggedright\arraybackslash}p{(\linewidth - 8\tabcolsep) * \real{0.1667}}
  >{\raggedleft\arraybackslash}p{(\linewidth - 8\tabcolsep) * \real{0.2222}}
  >{\raggedleft\arraybackslash}p{(\linewidth - 8\tabcolsep) * \real{0.2222}}
  >{\raggedright\arraybackslash}p{(\linewidth - 8\tabcolsep) * \real{0.1667}}
  >{\raggedleft\arraybackslash}p{(\linewidth - 8\tabcolsep) * \real{0.2222}}@{}}
\toprule\noalign{}
\begin{minipage}[b]{\linewidth}\raggedright
Gradierung
\end{minipage} & \begin{minipage}[b]{\linewidth}\raggedleft
Vektor-Kanäle
\end{minipage} & \begin{minipage}[b]{\linewidth}\raggedleft
Tensor-Kanäle
\end{minipage} & \begin{minipage}[b]{\linewidth}\raggedright
Stabilisator
\end{minipage} & \begin{minipage}[b]{\linewidth}\raggedleft
Kommutant der 64
\end{minipage} \\
\midrule\noalign{}
\endhead
\bottomrule\noalign{}
\endlastfoot
Identität & 0 & 60 & 60 & 1 \\
Pati-Salam-Spinorsplit (10 Varianten) & 24 & 36 & 36 = so(6)⊕\allowbreak{}so(4)⊕\allowbreak{}su(4)
& 2 \\
Farbsplit (3) & 40 & 20 & 52 & 2 \\
gemischt (30) & 32 & 28 & 28 & 4 \\
\end{longtable}
}

Ein einheitlicher Lorentztyp für alle 60 Vermittler erzwingt die
gleichhändige Zuordnung und damit (1,0)⊕\allowbreak{}(0,1): ein antisymmetrischer
Tensor B\_\{μν\}, kein symmetrisches Spin-2-Feld. Kinetische Terme:
einer pro irreduziblem Block des Stabilisators, keine
Familienaufspaltung 4 → 1+3. Ein Modenindex ist kein Raumpunkt.

\subsection{6. Spurnetzwerk-Normen}\label{spurnetzwerk-normen}

\[\adjustbox{max width=\linewidth}{$\displaystyle 
\frac{\nu_n}{(n!)^2}=\sum_{|m|=n}\frac{\|Q^{\dagger m}F\|^2}{m!}
=\Big[\text{Koeffizient von }s^n\Big]\int d\mu(z)\,
\exp\Big(-\tfrac12\sum_{k\ge1}\frac{s^k}{k}\operatorname{tr}\big((\bar ZM)(ZM)\big)^k\Big),
$}\]

über bosonische Kohärenzzustände, die Onishi-Determinante det(1 −
M̄M)\^{}\{1/2\} und Wick-Paarung der z-Variablen. Jedes Netzwerk ist eine
Kontraktion von n Kopien des Tensors Φ\_\{ij,kl\} =
Σ\_A(M\_A)\emph{\{ij\}(M\_A)}\{kl\} (Einträge in \{−1,0,1\}). Netzwerke
werden als 4-Bein-Multigraphen unter Knotenumbenennung und Slot-Tausch
kanonisiert (n=4: 120 → 27 Klassen). Exaktheit: float64-Kontraktion mit
Absolutwert-Begleitkontraktion (\textbar Σa\_kb\_k\textbar{} ≤
Σ\textbar a\_k\textbar\textbar b\_k\textbar{} \textless{} 2⁵³ je
Schritt), sonst Auswertung modulo fünf Primzahlen \textless{} 2¹³ mit
CRT (Eindeutigkeit aus \textbar Netzwerk\textbar{} ≤
60ⁿ·64\^{}\{\#Zyklen\}). Validierung: ν₁, ν₂ per Brute Force, ν₃ gegen
die Vektorrechnung, CRT-Route allein reproduziert ν₃.

\subsection{7. T1-T8}\label{t1-t8}

{\def\LTcaptype{none} % do not increment counter
\begin{longtable}[]{@{}
  >{\raggedright\arraybackslash}p{(\linewidth - 4\tabcolsep) * \real{0.3333}}
  >{\raggedright\arraybackslash}p{(\linewidth - 4\tabcolsep) * \real{0.3333}}
  >{\raggedright\arraybackslash}p{(\linewidth - 4\tabcolsep) * \real{0.3333}}@{}}
\toprule\noalign{}
\begin{minipage}[b]{\linewidth}\raggedright
Tor
\end{minipage} & \begin{minipage}[b]{\linewidth}\raggedright
Beitrag dieser Runde
\end{minipage} & \begin{minipage}[b]{\linewidth}\raggedright
Weiter offen
\end{minipage} \\
\midrule\noalign{}
\endhead
\bottomrule\noalign{}
\endlastfoot
T1 & Exakte Kommutanten-Leiter A→B→Cc→Cs→C; Clock als äußeres
Symmetrieelement & Welche Symmetrieoperationen der Compiler tatsächlich
bereitstellt \\
T2 & Lochbild bewiesen; Z\_h + Z\_add = 1; ε\_h, ε\_add & Konvergierter
Zustand, Renormierung, Grenzraum \\
T3 & - (ausgeklammert) & Räumlicher Träger \\
T4 & Kern-Symmetrieregel; kein einheitliches Vektor-Wörterbuch &
Chirales Maß, Anomalien \\
T5 & Exakte Krylov-Kette; ν₁\ldots ν₄ zweifach unabhängig &
Kontinuumslimes \\
T6 & Feldtyp-Census je Gradierung & Kopplungs- und Familiendaten \\
T7 & (1,0)⊕\allowbreak{}(0,1) ist B\_\{μν\}, kein Spin 2 & Dynamischer masseloser
Spin 2 \\
T8 & Sektorauswahl bewiesen bis g/Δ ≤ 0,0251; μ*-Fenster &
Zustandsfunktional, native Präparation \\
\end{longtable}
}

\textbf{Kein Tor wird als geschlossen markiert.} Keine RH-,
Faktorisierungs- oder P-versus-NP-Aussage; keine abgeleitete
Hylæan-Fähigkeit.

\subsection{8. Die jetzt entscheidenden
Prüfungen}\label{die-jetzt-entscheidenden-pruxfcfungen}

\begin{enumerate}
\def\labelenumi{\arabic{enumi}.}
\tightlist
\item
  \textbf{Untergrenzen verschärfen.} Die Casimir-Schranke trennt N=64
  bei 1/20 nicht von N=62 und N=66. Eine Jacobi-Vergleichsschranke mit
  rigorosen Normschranken für TT† je Bosonzahl oder eine Temple-Schranke
  auf der Ritz-Kette würde den Beweisbereich anheben.
\item
  \textbf{Den Zwei-Boson-Singulettraum ausschreiben.} w₂ ist der erste
  Zeuge; die volle Lanczos-Kette von H (nicht nur von T†) ab Tiefe 4
  braucht w₂ und seine Nachfolger.
\item
  \textbf{Operationen aus der Quelle.} Die Leiter sagt exakt, was
  so(10)-Operationen leisten würden; zu klären bleibt, welche davon der
  Compiler als markierte Operationen liefert. Der Clock ist ein äußeres
  Element und ersetzt sie nicht.
\item
  \textbf{Feldtyp vor Geometrie.} Jede chirale Aufspaltung macht den
  Bosontyp kanalabhängig (24/36 oder 40/20); eine räumliche Rechnung
  muss diese Aufteilung oder den einheitlichen Tensortyp explizit
  wählen.
\end{enumerate}

\subsection{9. Verifikation und nicht übernommene
Behauptungen}\label{verifikation-und-nicht-uxfcbernommene-behauptungen}

\codeword{replay.py} führt fünf Prüfer normal und unter \codeword{-OO}
aus und verlangt byteidentische Ausgaben; Status, Hashes und Guardsummen
stehen in \codeword{replay\_manifest.json}. Guardzahlen zählen
Prüfbedingungen, keine unabhängigen Theoreme. Exakt heißt
ganzzahlig/rational; numerisch heißt float64 mit Toleranz-Guard
(Ritz-Eigenwerte ab 3×3, Observablen, Erste-Moment-Auswertungen).

Nicht übernommen: der Worker-Beweis des N=64-Grundzustands bis 1/20;
eine physische Herleitung des Zustandsfunktionals; ein räumlicher
Träger; ein dynamischer Spin 2; die Existenz der Symmetrieoperationen
als Compiler- Primitive. Während der Runde festgestellte und behobene
Fehler der zunächst gelieferten Prüfer (u. a. Zyklus-Slots im
Spurnetzwerk, Bosonengewicht, doppelte T-Wirkung bei Bosonwiederholung,
Teilchen-Loch-Vorzeichen, min-max-Sektorschranke, skalierte w₂-Norm,
ungültiger Zykluszeuge) sind durch bleibende Guards abgesichert; keine
mathematische Aussage wurde an eine Erwartung angepasst.

Das \codeword{replay\_manifest.json} der Vorrunde v1.6.2 war durch einen
Lauf ohne Zugriff auf den externen Tensorpfad auf FAIL gesetzt worden;
es wurde aus den intakten \codeword{optimized.json} (PASS-Bedingung war
Byteidentität) rekonstruiert, Summen 27143/27139/4 stimmen mit dem
Delivery-Manifest.

\end{document}
